% Useful variables
\newcommand{\DocMainTitle}{SmallFloats.jl}
\newcommand{\DocVersion}{0.2.0}
\newcommand{\DocAuthors}{Jeffrey Sarnoff}
\newcommand{\JuliaVersion}{1.12.6}

% ---- Insert preamble
\input{preamble.tex}


\part{Home}


\chapter{SmallFloats.jl}



\label{2801651161961906345}{}


\emph{A conforming, performance-oriented Julia implementation of the IEEE P3109 draft standard — arithmetic formats for machine learning at bitwidths 3–16.}




\begin{minted}{jlcon}
julia> using SmallFloats

julia> Binary8p4se(1.6) + Binary8p4se(0.25)
Binary8p4se(1.875 ≡ 0x47)

julia> Binary5p3sf(0x08) == Binary5p3sf(1.0)     # an Unsigned is a CODE POINT
true
\end{minted}



\section{Why this package exists}



\label{13660426995343544750}{}


Small floating-point formats are easy to implement \emph{approximately} and surprisingly hard to implement \emph{exactly}. An FP8 \texttt{Exp} differs from a bit-exact one only in a handful of code points near rounding boundaries — invisible in a demo, decisive in a conformance suite, and quietly corrosive in research that compares number formats. SmallFloats.jl takes the position that for value sets this small there is no excuse for approximation:



\begin{itemize}
\item \textbf{Bit-exact defined results on every default path.} Every operation{\textquotesingle}s result is the correct projection of the mathematically exact value.


\item \textbf{One projection engine.} \texttt{RoundToPrecision → Saturate → Encode} is the single write path into a code point. There is no second, {\textquotedbl}fast but slightly different{\textquotedbl} rounding routine anywhere in the package.


\item \textbf{Approximation is opt-in, named, and measured.} Faster-but-inexact kernels live behind an explicit registry whose deviation bound κ is \emph{measured by exhaustive enumeration} at registration time; understated declarations are rejected.


\item \textbf{Enumerated where enumeration is affordable, and labelled where it is not.} ≈ 35 million verified units across 13 gates and tiers, each labelled \emph{exhaustive} or \emph{sampled} in a roll-call printed at the end of every run.


\item \textbf{Fast where it matters.} Table-gather kernels at fractions of a nanosecond per element; the scalar path fully specialized and allocation-free.

\end{itemize}


\section{Two things to know before you start}



\label{11297628053855568725}{}


\begin{tcolorbox}[toptitle=-1mm,bottomtitle=1mm,colback=admonition-warning!50!white,colframe=admonition-warning,title=\textbf{`Binary16p11se` is not `Float16`}]
Nor is \texttt{Binary16p8se} a \texttt{BFloat16}. The exponent biases differ by one, so \emph{every code point denotes a different value}. Converting between them is a conversion, never a reinterpretation.

\end{tcolorbox}


\begin{tcolorbox}[toptitle=-1mm,bottomtitle=1mm,colback=admonition-warning!50!white,colframe=admonition-warning,title=\textbf{Most format names are opt-in}]
\texttt{using SmallFloats} exports the \textbf{120 aliases at K ≤ 8}. The other 384 arrive with \texttt{using SmallFloats.Formats}, and are reachable without it as \texttt{SmallFloats.Binary16p6se} or \texttt{format(\allowbreak{}16, 6, true, true)}.

\end{tcolorbox}


\section{Where to go}



\label{11480021252683821972}{}


\begin{itemize}
\item \textbf{Get started} — \href{installation.md}{Installation \& Setup}, then the \href{quickstart.md}{Quickstart}: a working session in ten minutes.


\item \textbf{Understand the ideas} — the \href{mentalmodel.md}{Mental Model} (five ideas that carry the whole package), or \href{fifteenminutes.md}{Introducing P3109} if you want the standard itself first.


\item \textbf{Look something up} — the Reference section: formats, projection specs, the operation catalog, and the one-page \href{cheatsheet.md}{Cheat Sheet}.

\end{itemize}


\part{Getting Started}


\chapter{Installation \& Setup}



\label{3684371824697668327}{}


A working install in about five minutes.



\section{Requirements}



\label{13919788711114092546}{}


\begin{itemize}
\item \textbf{Julia 1.12} or later.


\item \textbf{\texttt{Quadmath}} (libquadmath{\textquotesingle}s \texttt{Float128}) is a hard dependency. It is used only inside the oracle and exact-fallback paths — never where it could change a result.

\end{itemize}


\section{Install}



\label{17214494489884574285}{}


The package is not yet registered, so install by URL:




\nopagebreak[4]
\begin{minted}{text}
pkg> add https://github.com/JeffreySarnoff/SmallFloats.jl
\end{minted}



Append \texttt{\#main} (or a tag/commit) to pin a revision:




\nopagebreak[4]
\begin{minted}{text}
pkg> add https://github.com/JeffreySarnoff/SmallFloats.jl#main
\end{minted}



For an editable install, \texttt{pkg> dev https:/\allowbreak{}/\allowbreak{}github.com/\allowbreak{}JeffreySarnoff/\allowbreak{}SmallFloats.jl} clones into \texttt{{\textasciitilde}/.julia/dev}, or use an existing clone:




\nopagebreak[4]
\begin{minted}{text}
pkg> dev path/to/SmallFloats.jl
\end{minted}



\section{Smoke test}



\label{27876697785071805}{}


Thirty seconds to confirm everything is in place:




\nopagebreak[4]
\begin{minted}{jlcon}
julia> using SmallFloats

julia> Binary8p4se(1.6) + Binary8p4se(0.25)
Binary8p4se(1.875 ≡ 0x47)

julia> conformance() !== nothing
true
\end{minted}



The first line proves the projection engine is live (1.6 rounds to 1.625, plus 0.25, lands exactly on 1.875); the second proves the conformance machinery is.



\section{Platform note: disabling Float128}



\label{2993233528556244241}{}


On platforms where libquadmath misbehaves, set an environment variable \textbf{before} loading the package:




\nopagebreak[4]
\begin{minted}{julia}
ENV["SmallFloats_Float128"] = "disable"
using SmallFloats
\end{minted}



This selects the pure-MPFR configuration. Results are \textbf{bit-identical} — that equivalence is itself part of the test suite; only oracle and build speed change.



\section{see also}



\label{9115667341246479961}{}


\href{quickstart.md}{Quickstart}.



\chapter{Quickstart}



\label{10347256324259028405}{}


One continuous session: two values, two ways of adding them, an array through a nonlinearity, and a block of quantized values. Every line is explained as it happens; the \hyperlinkref{9940662133030407839}{Mental Model} page picks up the {\textquotedbl}why{\textquotedbl} where this leaves off.



Assumes a \href{installation.md}{working install}. Start Julia and:




\nopagebreak[4]
\begin{minted}{jlcon}
julia> using SmallFloats
\end{minted}



\section{Make two values}



\label{1112983855294762154}{}



\begin{minted}{jlcon}
julia> x = Binary8p4se(1.6)
Binary8p4se(1.625 ≡ 0x45)
\end{minted}



You just used the most common format in the examples: \texttt{Binary8p4se} — 8 bits, precision 4, \textbf{s}igned, \textbf{e}xtended (it has ±Inf). The display tells you both halves of the story: \texttt{1.625} is the \emph{datum} (the exact value the format represents), \texttt{0x45} is the \emph{code point} (the 8-bit name of that datum).



Note what happened to \texttt{1.6}: the format has no datum at 1.6, so construction \emph{projected} — rounded to the nearest representable value, which is 1.625. Every write into a code point in this package goes through exactly one such projection.




\begin{minted}{jlcon}
julia> y = Binary8p4se(0x46)
Binary8p4se(1.75 ≡ 0x46)
\end{minted}



An \texttt{Unsigned} argument means \textbf{code point}, not number — \texttt{0x46} is the bit pattern, and the value is whatever that pattern denotes (1.75).




\begin{minted}{jlcon}
julia> Binary8p4se(0x02), Binary8p4se(2)
(Binary8p4se(0.001953125 ≡ 0x02), Binary8p4se(2.0 ≡ 0x48))
\end{minted}



Same digit, two very different results: \texttt{0x02} is code point 2; \texttt{2} is the numeric value two, projected.



\begin{tcolorbox}[toptitle=-1mm,bottomtitle=1mm,colback=admonition-warning!50!white,colframe=admonition-warning,title=\textbf{`Unsigned` means code point — at every width}]
\texttt{T(0x02)} constructs code point \texttt{0x02}; \texttt{T(2)} projects the numeric value two. Signedness of the integer type is what distinguishes them, and the rule holds for \texttt{UInt8}, \texttt{UInt16} and any other \texttt{Unsigned}. Use \texttt{Convert} when intent should be unmistakable.

\end{tcolorbox}


\section{Add them, two ways}



\label{9087301203278813898}{}



\begin{minted}{jlcon}
julia> x + y
Binary8p4se(3.5 ≡ 0x4e)
\end{minted}



Ordinary Julia arithmetic works. Under the hood, \texttt{+} computed the exact sum (1.625 + 1.75 = 3.375) and projected it once to the nearest datum, 3.5.



The same operation, written out in full:




\nopagebreak[4]
\begin{minted}{jlcon}
julia> Add(Binary8p4se, RNE_SN, x, y)
Binary8p4se(3.5 ≡ 0x4e)
\end{minted}



Every operation in the package has this explicit shape:




\nopagebreak[4]
\begin{minted}{text}
Op(result_format, projection_spec, operands...)
\end{minted}



\texttt{RNE\_\allowbreak{}SN} is a projection specification: \textbf{R}ound to \textbf{N}earest, ties \textbf{E}ven, with the draft{\textquotesingle}s default \textbf{S}aturation rules (\texttt{SN} = SatNone — the saturation mode is named for what it \emph{doesn{\textquotesingle}t} do: it does not clamp everything to finite, and it does not blindly propagate infinities; it applies the draft{\textquotesingle}s domain- and direction-dependent rows). It{\textquotesingle}s the package-wide default, which is why \texttt{+} and \texttt{Add(…, RNE\_\allowbreak{}SN, …)} agree.



Why would you want the long form? Because the projection is yours to choose. Watch an overflow under two different specs (\texttt{Binary8p4se} tops out at 224):




\nopagebreak[4]
\begin{minted}{jlcon}
julia> w = Binary8p4se(200.0); two = Binary8p4se(2.0);

julia> Multiply(Binary8p4se, RNE_SN, w, two)     # nearest: overflow → +Inf
Binary8p4se(Inf ≡ 0x7f)

julia> Multiply(Binary8p4se, RNE_SF, w, two)     # SF = clamp to finite range
Binary8p4se(224.0 ≡ 0x7e)
\end{minted}



Same operands, same exact product (400), different landing — the projection spec is part of the operation, not an afterthought.



\section{Send an array through a nonlinearity}



\label{6398029487376140467}{}



\begin{minted}{jlcon}
julia> A = Binary8p4se.([-1.0, 0.5, 2.0])
3-element Vector{Binary8p4se}:
 Binary8p4se(-1.0 ≡ 0xc0)
 Binary8p4se(0.5 ≡ 0x38)
 Binary8p4se(2.0 ≡ 0x48)

julia> Exp(Binary8p4se, RNE_SN, A)
3-element Vector{Binary8p4se}:
 Binary8p4se(0.375 ≡ 0x34)
 Binary8p4se(1.625 ≡ 0x45)
 Binary8p4se(7.5 ≡ 0x57)
\end{minted}



Broadcasting \texttt{Binary8p4se.(…)} projects each element, as you{\textquotesingle}d expect. The array call to \texttt{Exp} is where the package{\textquotesingle}s speed comes from: an 8-bit unary operation has only 256 possible inputs, so the first call builds a 256-byte result table — computed through the exact scalar path — and every element after that is a single lookup. You can see the table:




\nopagebreak[4]
\begin{minted}{jlcon}
julia> table_bytes()
256
\end{minted}



One 256-byte table, cached. The Base spelling \texttt{exp.(A)} rides the same machinery under the default projection.



\section{Create a block of quantized values}



\label{9146025431152277157}{}


Small formats have small ranges. \emph{Block} formats (the MX-style scheme) fight that by sharing one scale across a group of narrow elements: each represented value is \texttt{scale × element}.




\begin{minted}{jlcon}
julia> values = (100.0, -12.0, 0.5, 3.0);

julia> b = ConvertToBlockMaxAbsFinite(Binary8p1uf, Binary8p4se, RNE_SN, RNE_SN,
                                      Binary8p4se.(values))
Block{4, Binary8p1uf, Binary8p4se}(Binary8p1uf(64.0 ≡ 0x86), (…))
\end{minted}



Read the arguments: scale format \texttt{Binary8p1uf} (unsigned — scales are non-negative — with the huge range precision 1 buys), element format \texttt{Binary8p4se}, one projection spec for the scale and one for the elements, then the values. The algorithm picked scale 64 from the largest-magnitude element (100) and projected each element against it.



\begin{tcolorbox}[toptitle=-1mm,bottomtitle=1mm,colback=admonition-warning!50!white,colframe=admonition-warning,title=\textbf{Stage wide, then block-quantize}]
\texttt{ConvertToBlockMaxAbsFinite} takes already-\texttt{Binary} elements. If you stage \texttt{Float64} data through the \emph{element} format first, clamping can happen \textbf{before} the block scale absorbs it, and blocks buy you nothing. Stage through a wide-range format instead. (The example above is safe because its values already fit in \texttt{Binary8p4se}.)

\end{tcolorbox}


Read the block back as plain values:




\nopagebreak[4]
\begin{minted}{jlcon}
julia> ConvertFromBlock(Binary8p3se, RNE_SN, b)
(Binary8p3se(96.0 ≡ 0x5a), Binary8p3se(-12.0 ≡ 0xce), Binary8p3se(0.5 ≡ 0x3c), Binary8p3se(3.0 ≡ 0x46))
\end{minted}



Each \texttt{scale × element} was decoded exactly and projected once into the format you asked for. And the dot product — the operation blocks exist for — keeps every lane product and the accumulation \emph{exact} until a single final rounding:




\nopagebreak[4]
\begin{minted}{jlcon}
julia> BlockDotProduct(Binary8p4se, RNE_SN, b, b)
Binary8p4se(Inf ≡ 0x7f)
\end{minted}



The exact accumulated value exceeds \texttt{Binary8p4se}{\textquotesingle}s finite range, so the one final projection produces \texttt{Inf} under \texttt{RNE\_\allowbreak{}SN}.



\section{What just happened}



\label{2185968854617873464}{}


Four ideas carried this whole session, and they{\textquotesingle}re the whole package:



\begin{itemize}
\item[1. ] \textbf{A format is a finite set of datums} — \texttt{Binary8p4se} is 256 code points naming 256 values, no more.


\item[2. ] \textbf{Every write is one projection} — construction, arithmetic, conversion: exact math, then round × saturate into a code point.


\item[3. ] \textbf{The projection is explicit} — \texttt{RNE\_\allowbreak{}SN} is the default, not the law; \texttt{Add(T, ρ, …)} takes any spec.


\item[4. ] \textbf{Bulk work is table lookup} — finite functions are precomputed, exactly.

\end{itemize}


The \hyperlinkref{9940662133030407839}{Mental Model} develops each of these properly. If you{\textquotesingle}d rather learn by doing, the Tutorials start where this page ends, and the How-To Guides are organized by task.



\chapter{Mental Model}



\label{9940662133030407839}{}


Five ideas. Everything else in SmallFloats.jl — every page of the User Guide, every tutorial, every recipe — is an elaboration of one of them. Read this page once, slowly, and the rest of the documentation becomes reference material you consult, not prose you wade through.



\section{Idea 1 — A format is a finite set of datums}



\label{7753368820005915228}{}


A SmallFloats \emph{format} is not a floating-point type in the IEEE-754 sense of a continuum sampled at machine precision. It is a \textbf{finite, enumerated set of exact values} — the \emph{datums} — each with a name, its \emph{code point}. The format \texttt{Binary4p2se} (4 bits, precision 2, signed, extended) has 16 code points and they are all of them:




\nopagebreak[4]
\begin{minted}{jlcon}
julia> decode.(sort(Binary4p2se.(0x00:0x0f)))
16-element Vector{Float64}:
 -Inf
  -2.0
  -1.5
  -1.0
  -0.75
  -0.5
  -0.25
   0.0
   0.25
   0.5
   0.75
   1.0
   1.5
   2.0
  Inf
 NaN
\end{minted}



That listing \emph{is} the format. There is nothing between the grid points, no hidden state, no approximation in the set itself — the spacing you see (subnormals near zero, binades doubling outward) is the whole structure.



Three consequences fall out immediately:



\begin{itemize}
\item \textbf{One NaN, no negative zero.} The last code point is the format{\textquotesingle}s single NaN; the value \texttt{0.0} appears exactly once. IEEE-754 spends thousands of code points on NaN payloads and two on zeros; P3109 spends them on datums.


\item \textbf{The set has edges.} \texttt{-Inf}, \texttt{Inf}, and \texttt{NaN} occupy three of the sixteen names, which is why the largest \emph{finite} value is 2.0 and not something larger. Finite formats (suffix \texttt{f} instead of \texttt{e}) spend those two infinity code points on more finite datums instead.


\item \textbf{Larger formats are the same idea with more points.} \texttt{Binary8p4se} is 256 names; \texttt{Binary16p11se} is 65 536. The 504 formats are just the legal choices of how many names (bitwidth \texttt{K}), how much of the budget goes to significand vs exponent (precision \texttt{P}), and whether the set includes negatives (\texttt{s}/\allowbreak\texttt{u}) and infinities (\texttt{e}/\allowbreak\texttt{f}).

\end{itemize}


\emph{Now you can read:} \href{explain\_\allowbreak{}formats.md}{Formats, Aliases \& Representations} for why the type \texttt{Binary8p4se} is a concrete \emph{representation} of the abstract format \texttt{Binary\{8,4,true,true\}} — and why \texttt{===} is \texttt{false} but \texttt{<:} is \texttt{true}.



\section{Idea 2 — A value is a wrapper around a code point; decode is exact}



\label{561648171242528358}{}


Construct a value and you hold the code point; ask what it means and you get the datum back, exactly:




\nopagebreak[4]
\begin{minted}{jlcon}
julia> x = Binary8p4se(1.6)
Binary8p4se(1.625 ≡ 0x45)

julia> codepoint(x)
0x45

julia> decode(x)
1.625
\end{minted}



The display \texttt{1.625 ≡ 0x45} is the idea in one line: the code point \texttt{0x45} \emph{means} the datum 1.625, and the \texttt{≡} is exact, not approximate. \texttt{decode} never rounds — it is a table lookup from name to meaning.



The asymmetry that matters: \texttt{decode} is exact, but \textbf{construction is a projection}. \texttt{Binary8p4se(1.6)} did not store 1.6; it stored the name of the nearest datum, which happens to be 1.625. The value you get back is the value the format can say, not the value you asked for.



This is why the \texttt{Unsigned}-means-code-point rule (from the Quickstart) exists: \texttt{codepoint} and the constructor are inverses \emph{by design} —




\begin{minted}{jlcon}
julia> Binary8p4se(codepoint(x)) === x
true
\end{minted}



— and that round-trip is the format{\textquotesingle}s identity. A value is its code point; everything else is a view.



\emph{Now you can read:} \href{tutorial1\_\allowbreak{}values.md}{Values, Code Points \& Conversion} for the four ways in and two ways out, and the carriers \texttt{decode} returns at wide formats.



\section{Idea 3 — Every computation is exact math, then one projection}



\label{15574776794132898440}{}


When you compute, SmallFloats does not simulate low-precision arithmetic step by step. It computes the \textbf{mathematically exact result} of the operation on the decoded datums, then applies \textbf{one projection} — one rounding plus one saturation decision — to land on a code point:




\nopagebreak[4]
\begin{minted}{text}
exact result  →  RoundToPrecision  →  Saturate  →  code point
\end{minted}



\texttt{x + y} from the Quickstart: exact sum 3.375, projected once to 3.5. \texttt{200 × 2}: exact product 400 (not representable — projection decides between \texttt{Inf}, \texttt{224}, or \texttt{NaN}, and the spec says which).



Two things follow from {\textquotedbl}one projection, at the end{\textquotedbl}:



\begin{itemize}
\item \textbf{There is no accumulated rounding error inside an operation.} A fused \texttt{FMA(x, y, z)} rounds once, not twice. A \texttt{BlockDotProduct} accumulates all 32 lane products exactly and projects once at the end. Intermediate roundings are not {\textquotedbl}small{\textquotedbl} — they are \emph{absent}.


\item \textbf{The default path is bit-exact.} Because the exact math is carried far enough (Float64, Float128, or MPFR enclosures, escalating automatically), the code point you get is \emph{the} defined answer, not a good approximation of it. The package can prove this because the value sets are finite — see Idea 5.

\end{itemize}


\emph{Now you can read:} \href{explain\_\allowbreak{}projection\_\allowbreak{}contract.md}{The Projection Contract} for what {\textquotedbl}exact{\textquotedbl} means at each carrier width, and the symbolic-sticky trick that lets directed modes land exactly on asymptotes.



\needspace{17\baselineskip}
\section{Idea 4 — Projection = rounding mode × saturation mode}



\label{11791101062988838207}{}


A projection specification is the pair of two independent choices:




\begin{table}[h]
\centering
\begin{tabulary}{\linewidth}{L L L}
\toprule
 & what it decides & examples \\
\toprule
\textbf{rounding mode} & which neighbor a between-grid value lands on & nearest-ties-even, toward zero, toward ±∞, stochastic \\
\textbf{saturation mode} & what an out-of-range result becomes & clamp to finite (\texttt{SF}), propagate infinities (\texttt{SP}), draft rows (\texttt{SN}) \\
\bottomrule
\end{tabulary}

\end{table}



Watch one overflow resolve three ways (\texttt{Binary8p4se}, max finite 224):




\nopagebreak[4]
\begin{minted}{jlcon}
julia> w, two = Binary8p4se(200.0), Binary8p4se(2.0);

julia> Multiply(Binary8p4se, RNE_SN, w, two)   # nearest + draft rows → Inf
Binary8p4se(Inf ≡ 0x7f)

julia> Multiply(Binary8p4se, RNE_SF, w, two)   # nearest + clamp → 224
Binary8p4se(224.0 ≡ 0x7e)

julia> Multiply(Binary8p4se, RTZ_SN, w, two)   # toward zero + draft rows → 224
Binary8p4se(224.0 ≡ 0x7e)
\end{minted}



The first two differ in \emph{saturation} (SN vs SF); the first and third differ in \emph{rounding} (RNE vs RTZ) — and the third row is the subtle one: under \texttt{SatNone}, a directed mode that points \emph{away} from the overflow clamps to the extremal finite value, so \texttt{RTZ\_\allowbreak{}SN} agrees with \texttt{RNE\_\allowbreak{}SF} here for different reasons.



The full deterministic grid is 6 rounding modes × 3 saturation modes = 18 predefined specs (\texttt{RNE\_\allowbreak{}SN}, \texttt{RNA\_\allowbreak{}SF}, \texttt{RTP\_\allowbreak{}SP}, …), plus three stochastic families parameterized by a random-bit budget. \texttt{RNE\_\allowbreak{}SN} is the package-wide default — which is exactly why \texttt{+}, \texttt{exp}, and \texttt{T(1.6)} all agree with \texttt{Add(T, RNE\_\allowbreak{}SN, …)}, \texttt{Exp(T, RNE\_\allowbreak{}SN, …)}, and \texttt{Convert(T, RNE\_\allowbreak{}SN, …)}.



\begin{tcolorbox}[toptitle=-1mm,bottomtitle=1mm,colback=admonition-warning!50!white,colframe=admonition-warning,title=\textbf{The default is a session default}]
\texttt{+} and friends follow \texttt{DefaultProjection()}, which is process-global and stays set. Convenience syntax is for exploratory code; anything that must be reproducible names its spec explicitly. The Tutorials teach the safe \texttt{with\_\allowbreak{}default\_\allowbreak{}projection} form before showing \texttt{DefaultProjection!}.

\end{tcolorbox}


\emph{Now you can read:} \href{tutorial2\_\allowbreak{}projection.md}{Projection: Rounding \& Saturation} for the experiments, and \href{ref\_\allowbreak{}projections.md}{Projection Specifications} for the complete grid.



\section{Idea 5 — Approximation is opt-in, named, and measured}



\label{2713166630189569727}{}


The package{\textquotesingle}s speed comes from the finiteness of Idea 1: an 8-bit unary operation is a function with 256 inputs, so its \emph{entire} behavior is a 256-byte table, built once through the exact scalar path and then gathered per element (0.27 ns). The exactness is inherited: table and scalar are the same function by construction.



The same finiteness makes approximation \emph{honest}. If you register a faster but inexact kernel — say a hardware-style hard-tanh for \texttt{Tanh} — the registry does not take your word for its accuracy. It \textbf{measures} the worst-case deviation, exhaustively over every input, as a distance in code points:




\nopagebreak[4]
\begin{minted}{jlcon}
julia> κ, exhaustive = measure_kappa(
           x -> Clamp(Binary8p4se, RNE_SN, x,
                      Negate(one(Binary8p4se)), one(Binary8p4se)),
           :Tanh, Binary8p4se, (Binary8p4se,), RNE_SN)
(4.0, true)
\end{minted}



κ = 4 code points, verified over all 256 inputs — that is a measured fact, not a claim. Declare κ smaller than the measurement and registration is \emph{rejected}. Nothing approximate is reachable from the default API; an approximation exists only because you named it, and its κ is part of the conformance declaration the package exports.



This is the package{\textquotesingle}s answer to the problem stated in the Introduction: small formats are easy to implement approximately and hard to implement exactly. SmallFloats implements them exactly, and when you choose otherwise, it makes the choice explicit and quantified.



\emph{Now you can read:} \href{howto\_\allowbreak{}register\_\allowbreak{}approx.md}{Register \& Measure an Approximation}, and \href{ref\_\allowbreak{}conformance.md}{Conformance \& Approximation Registry}.



\section{The five ideas, as one picture}



\label{10474009627779219316}{}



\begin{minted}{text}
            Idea 1                    Idea 2                 Idea 3
   ┌──────────────────────┐   name   ┌──────────┐   exact    ┌────────────┐
   │  finite set of       │  ⇄       │  value = │   math     │  project   │
   │  datums (code points)│─────────→│ code pt. │───────────→│  once      │
   └──────────────────────┘  decode  └──────────┘            └────────────┘
                                                                  │
                          Idea 4                    Idea 5        ▼
                  projection = rounding × saturation      exact by default;
                                                          κ-measured if not
\end{minted}



Read left to right: a value is a name in a finite set (1–2); computing is exact math followed by one projection (3); the projection is your choice of rounding and saturation (4); and the result is exact unless you explicitly register otherwise (5).



\textbf{Where to go from here.} To build the ideas into skills: Values, Code Points \& Conversion. To look something up: the Reference pages. To understand \emph{why} the package is built this way: The Projection Contract and the rest of the Explanation section.



\chapter{Introducing P3109}



\label{1155828994189583929}{}


The IEEE P3109 draft — \emph{Arithmetic Formats for Machine Learning} — as this package sees it. Fifteen minutes here buys you the vocabulary every other page assumes. Claims are tagged with the draft clause they come from; the full transliteration is in the repository at \texttt{docs/other/IEEE\_\allowbreak{}D1.md} for when you want the normative text itself.



\section{The one law everything hangs on}



\label{17993345610855520848}{}


Every numeric operation in the draft is three steps (§4.3.2):




\nopagebreak[4]
\begin{minted}{text}
decode the operands  →  compute exactly over the extended reals  →  project once
\end{minted}



There is no arithmetic {\textquotedbl}in the format.{\textquotedbl} Operands are decoded to exact values, the operation is performed exactly, and a single projection — rounding plus saturation — writes the result back to a code point. The comparisons and predicates stop after decoding (they return Booleans); everything else follows the law. SmallFloats.jl is built as a direct implementation of this sentence.



\needspace{18\baselineskip}
\section{Formats: four parameters, 504 instances}



\label{12063105756466521911}{}


A format is \texttt{Binary\{K, P, Σ, Δ\}} (§3.1):




\begin{table}[h]
\centering
\begin{tabulary}{\linewidth}{L L L}
\toprule
parameter & meaning & range \\
\toprule
\texttt{K} & total bitwidth & 3 – 16 \\
\texttt{P} & significand precision, implicit bit included & signed: \texttt{0 < P < K}; unsigned: \texttt{0 < P ≤ K} \\
\texttt{Σ} & signedness: signed or unsigned &  \\
\texttt{Δ} & domain: \textbf{extended} (has ±Inf) or \textbf{finite} &  \\
\bottomrule
\end{tabulary}

\end{table}



The short name reads off the parameters: \texttt{Binary8p4se} = K 8, p 4, \textbf{s}igned, \textbf{e}xtended. The draft{\textquotesingle}s exponent bias is \emph{derived}, not chosen (§3.1, §4.14):




\nopagebreak[4]
\begin{minted}{text}
B = 2^(K−P−1)   (signed)        B = 2^(K−P)   (unsigned)
\end{minted}



This symmetric-about-the-midpoint bias is the single most consequential formula in the draft — see {\textquotedbl}Differences from IEEE 754{\textquotedbl} below.



\section{Datums and code points}



\label{15176353406766861258}{}


A format{\textquotesingle}s \emph{datum set} is a finite subset of the closed extended reals (ℝ plus ±Inf and NaN), in bijection with the code points \texttt{0 … 2{\textasciicircum}K − 1} (§3, Annex B). Every finite datum has a canonical form




\begin{minted}{text}
(−1)^sign · S · 2^(1−P) · 2^E        E > −B
\end{minted}



with significand \texttt{S} classifying it as normal (\texttt{2{\textasciicircum}(P−1) ≤ S < 2{\textasciicircum}P}) or subnormal (\texttt{0 < S < 2{\textasciicircum}(P−1)}).



The specials are where the draft parts company with IEEE-754 (Annex B, Table 3):



\begin{itemize}
\item \textbf{Exactly one zero.} There is no \texttt{−0}; a signed format has an odd number of finite datums.


\item \textbf{Exactly one NaN}, sitting at the code point IEEE-754 gives to \texttt{−0} in signed formats. No payloads, no quiet/signalling distinction — the thousands of code points IEEE spends on NaN encodings are datums here.


\item \textbf{±Inf only in extended formats}, adjacent to the finite extremes. A finite format spends those code points on more finite values.

\end{itemize}


\needspace{24\baselineskip}
\section{Projection: rounding × saturation}



\label{13445344687072648265}{}


A projection specification ρ is the pair \emph{(rounding mode, saturation mode)} (§4.2, §4.7.3–4.7.5).



Six deterministic rounding modes — \texttt{NearestTiesToEven}, \texttt{NearestTiesToAway}, \texttt{TowardPositive}, \texttt{TowardNegative}, \texttt{TowardZero}, \texttt{ToOdd} — plus three stochastic families \texttt{StochasticA/B/C}, which consume \texttt{N} random bits as an unsigned draw \texttt{R ∈ 0 … 2{\textasciicircum}N − 1} (§4.7.4; the draft deliberately leaves the bit source unspecified).



Three saturation modes decide what out-of-range results become (§4.7.5):




\begin{table}[h]
\centering
\begin{tabulary}{\linewidth}{L L L}
\toprule
mode & finite overflow & ±Inf input \\
\toprule
\texttt{SatFinite} & clamp to the finite extremes & clamp \\
\texttt{SatPropagate} & clamp & preserve representable infinities, else clamp \\
\texttt{SatNone} & the draft{\textquotesingle}s rows: to ±Inf, finite extreme, or NaN depending on rounding direction, signedness, and domain & preserve, else NaN \\
\bottomrule
\end{tabulary}

\end{table}



One ordering fact the whole clause turns on (§4.7.3 NOTE 1): \textbf{rounding precedes saturation}, but the rounding mode is passed \emph{into} saturation so a directed mode can force a finite result. Saturation itself never rounds.



\section{The operation catalog}



\label{4723431904892488085}{}


§4.9–§4.16 define the catalog: \texttt{Convert}; sign and basic arithmetic (\texttt{Abs} through \texttt{Divide}, fused \texttt{FMA} and \texttt{FAA}); the elementary functions (roots, exp/log families, trig, hyperbolic, π-scaled trig, \texttt{Softplus}, \texttt{Hypot}, \texttt{ArcTan2}); the extrema families (\texttt{Minimum}/\allowbreak\texttt{Maximum} and their NaN-quieting, magnitude, and finite variants, plus \texttt{Clamp}); comparisons, predicates, and \texttt{TotalOrder}; format-level queries (\texttt{BitwidthOf}, \texttt{MaxFiniteOf}, …); and the neighbour operations \texttt{NextGreaterThan}/\allowbreak\texttt{NextLessThan}.



Some results are visible only in the pattern tables, and they surprise:



\begin{itemize}
\item \texttt{Divide(x, 0) = NaN} for every \texttt{x} — including ±Inf (§4.10.5 NOTE 2). There is no IEEE-style ±Inf quotient.


\item \texttt{Recip(±Inf) = 0} and \texttt{Recip(0) = NaN} (§4.10.8).


\item \texttt{TotalOrder} puts NaN at an end — deterministic tie-breaking for search and sort (§4.12.1).


\item \texttt{NextGreaterThan(\allowbreak{}+Inf) = NaN}, unlike IEEE-754 \texttt{nextUp} (§4.16 NOTE 2).


\item \texttt{FAA} is associative by construction: \texttt{Add(\allowbreak{}Add(\allowbreak{}X,Y),Z) = Add(\allowbreak{}X,Add(\allowbreak{}Y,Z))} computed exactly (§4.10.7).

\end{itemize}


\section{Blocks and scaled operations}



\label{17754733778172498420}{}


Clause 5 adds the MX-style scheme: a \emph{block} is a scale factor in one format plus \texttt{B} elements in another, representing \texttt{scale × element} per lane. Block operations decode lanes exactly, apply the operation, and project against a result scale; the reductions (\texttt{BlockReduceAdd}, \texttt{BlockDotProduct}) accumulate exactly and project once. \texttt{ScaledOp} is the block-size-1 form. This is how the draft tracks dynamic range that a bare element format cannot hold.



\section{Conformance and κ}



\label{11385503936379929322}{}


The draft{\textquotesingle}s conformance machinery (§4.4–4.6) has two levels. A \emph{conforming implementation} produces the defined result for every operation, format, and mode. A \emph{κ-approximate implementation} may deviate, but must declare its maximum code-point distance κ from the defined result — measured along the total order, over inputs with finite defined results. Approximation is thus quantified and declared, never silent; SmallFloats implements both levels and measures κ by exhaustive enumeration at registration.



\needspace{18\baselineskip}
\section{Differences from IEEE 754}



\label{153811621810957742}{}


\texttt{Binary16p11se} looks like \texttt{Float16}; \texttt{Binary16p8se} looks like \texttt{BFloat16}. They are not (Annex F compares them explicitly):




\begin{table}[h]
\centering
\begin{tabulary}{\linewidth}{L L L L L}
\toprule
 & \texttt{Binary16p11se} & \texttt{Float16} & \texttt{Binary16p8se} & \texttt{BFloat16} \\
\toprule
exponent bias & \textbf{16} & \textbf{15} & \textbf{128} & \textbf{127} \\
largest finite & \textbf{65472} & \textbf{65504} & \textbf{3.3762e38} & \textbf{3.3895e38} \\
NaN encodings & \textbf{1} & \textbf{2046} & \textbf{1} & \textbf{2046} \\
negative zero & \textbf{no} & yes & \textbf{no} & yes \\
\bottomrule
\end{tabulary}

\end{table}



P3109{\textquotesingle}s bias formula makes the exponent range symmetric about the midpoint; IEEE-754{\textquotesingle}s is \texttt{2{\textasciicircum}(E−1) − 1}. The two differ by one, so \textbf{every code point denotes a value a factor of two away from its IEEE look-alike}. The same bits read as the other type silently halve or double every value. Converting between them is a conversion, never a reinterpretation.



\section{see also}



\label{9115667341246479961-dup1}{}


\href{mentalmodel.md}{Mental Model}, \href{papers\_\allowbreak{}index.md}{The Standard \& Design Papers}, \href{quickstart.md}{Quickstart}.



\part{Using SmallFloats}


\chapter{Tutorials}


\section{Values, Code Points \& Conversion}



\label{6242962761125980180}{}


This tutorial assumes the Quickstart: you{\textquotesingle}ve seen \texttt{Binary8p4se(1.6)} project onto \texttt{1.625 ≡ 0x45}, you{\textquotesingle}ve seen an \texttt{Unsigned} argument mean code point rather than number, and you{\textquotesingle}ve seen \texttt{Add(T, ρ, x, y)} alongside plain \texttt{+}. Here we slow down and go through every way a value gets made, read, stepped, and classified.



\subsection{Four ways in, two ways out}



\label{15348036207912288565}{}


A \texttt{Binary} value is an immutable wrapper around a code point. There are four constructors and two accessors:




\nopagebreak[4]
\begin{minted}{jlcon}
julia> x = Binary8p4se(1.6)              # construct from a Real: projects (rounds)
Binary8p4se(1.625 ≡ 0x45)

julia> Binary8p4se(0x45)                 # construct from a UInt8 CODE POINT (validated)
Binary8p4se(1.625 ≡ 0x45)

julia> rawvalue(Binary8p4se, 0x45)       # same, unchecked — the kernel-internal route
Binary8p4se(1.625 ≡ 0x45)

julia> Convert(Binary8p4se, RNE_SN, 3)   # explicit conversion, any mode
Binary8p4se(3.0 ≡ 0x4c)

julia> decode(x)                          # the exact datum, on the format's carrier
1.625

julia> codepoint(x)                       # the code point (extends Base.codepoint)
0x45
\end{minted}



\texttt{decode} never rounds — it is a lookup from name to meaning. Construction from a \texttt{Real} is the only place rounding happens on the way in.



\subsection{Important: \texttt{Unsigned} means code point, at every width}



\label{6585316590758660594}{}


The Quickstart already showed you \texttt{Binary8p4se(0x02)} and \texttt{Binary8p4se(2)} landing on different values. The rule generalizes to every integer width and every format, and it is worth stating as a standalone hazard because it never raises an error — it just silently gives you the wrong datum if you meant a number:




\nopagebreak[4]
\begin{minted}{jlcon}
julia> Binary8p4se(0x02), Binary8p4se(2)
(Binary8p4se(0.001953125 ≡ 0x02), Binary8p4se(2.0 ≡ 0x48))
\end{minted}



\begin{tcolorbox}[toptitle=-1mm,bottomtitle=1mm,colback=admonition-warning!50!white,colframe=admonition-warning,title=\textbf{`Unsigned` means code point — at every width}]
\texttt{T(0x02)} constructs code point \texttt{0x02}; \texttt{T(2)} projects the numeric value two. Signedness of the integer type is what distinguishes them, so the rule holds for \texttt{UInt8}, \texttt{UInt16} and any other \texttt{Unsigned}: \texttt{Binary16p6se(0x0002)} is code point 2, not the value 2.0. Use \texttt{Convert} when intent should be unmistakable.

Out-of-range codes throw for K < 8 (\texttt{Binary5p3sf(0x20)} is an error); the range check costs nothing measurable — 2.1 ns, identical to unchecked \texttt{rawvalue}. Round-tripping is \texttt{T(codepoint(x)) === x}.

\end{tcolorbox}


\subsection{\texttt{Convert} versus \texttt{convert}}



\label{4082966728771053146}{}


\texttt{Convert} (capitalized, from the spec-named register) is the explicit projection you saw above: \texttt{Convert(T, ρ, value)}, numeric for every integer. Base{\textquotesingle}s lowercase \texttt{convert} is a different, narrower operation that Julia{\textquotesingle}s promotion machinery calls, and it is value-preserving on \texttt{Unsigned} input rather than code-point reading:




\nopagebreak[4]
\begin{minted}{jlcon}
julia> Binary8p4se(0x02)
Binary8p4se(0.001953125 ≡ 0x02)
\end{minted}



Keep the two apart: \texttt{T(x::Unsigned)} reads a code point; \texttt{Convert(T, ρ, x)} always projects a numeric value, whatever its type. When you want projection and the argument might be an \texttt{Unsigned}, reach for \texttt{Convert}, not the constructor.



\subsection{The decode / codepoint round-trip}



\label{10860920436673360354}{}


\texttt{decode} and \texttt{codepoint} are inverses by construction — that round-trip is a format{\textquotesingle}s identity:




\nopagebreak[4]
\begin{minted}{jlcon}
julia> x = Binary8p4se(1.6)
Binary8p4se(1.625 ≡ 0x45)

julia> codepoint(x)
0x45

julia> decode(x)
1.625
\end{minted}



\texttt{decode} is \textbf{always exact}, but the type it returns is a property of the format, not of the package: \texttt{Float64} for the 432 formats whose exponent range it holds, \texttt{Float128} or an exact dyadic carrier for the other 72. \texttt{Float64(x)} always returns a \texttt{Float64} and rounds where the datum does not fit.



\subsection{Enumerating a whole format}



\label{15296297375297823522}{}


Small formats are small enough to look at in full. \texttt{Binary4p2se} has 16 code points, sorted here by the total order (single NaN last):




\nopagebreak[4]
\begin{minipage}{\linewidth}
\begin{minted}{julia}
decode.(sort(Binary4p2se.(0x00:0x0f)))      # broadcast the code-point constructor
\end{minted}




\begin{minted}{text}
[-Inf, -2.0, -1.5, -1.0, -0.75, -0.5, -0.25, 0.0,
  0.25, 0.5, 0.75, 1.0, 1.5, 2.0, Inf, NaN]
\end{minted}
\end{minipage}




That listing \emph{is} the format: one NaN, no negative zero, subnormal spacing near zero, binades doubling outward. This is a good way to build intuition for a format before committing to it.



\subsection{Stepping and classification}



\label{15274975594442253874}{}


Stepping moves to the adjacent code point; classification names what kind of datum you{\textquotesingle}re holding:




\nopagebreak[4]
\begin{minted}{jlcon}
julia> NextGreaterThan(Binary8p4se(1.0))
Binary8p4se(1.125 ≡ 0x41)

julia> Class(Binary8p4se(0.01))           # draft classification
ClassPosNormal::FPClass = 0x06
\end{minted}



\texttt{NextGreaterThan}/\allowbreak\texttt{NextLessThan} are the draft{\textquotesingle}s Next operations; \texttt{nextfloat} and \texttt{prevfloat} alias them, so both spellings work. Standard predicates apply too: \texttt{isnan}, \texttt{isinf}, \texttt{isfinite}, \texttt{iszero}, \texttt{signbit}, \texttt{issubnormal}:




\nopagebreak[4]
\begin{minted}{jlcon}
julia> Binary8p4se(1e9)                   # overflow under the default spec → +Inf
Binary8p4se(Inf ≡ 0x7f)

julia> Binary8p4se(-0.0)                  # no −0: projects to the single zero
Binary8p4se(0.0 ≡ 0x00)

julia> isnan(Binary8p4se(NaN)), isfinite(Binary8p4se(2.0))
(true, true)
\end{minted}



\subsection{\texttt{zero}, \texttt{one}, \texttt{eps}, \texttt{typemax}}



\label{17547000303368754034}{}


The usual numeric-type vocabulary works on every format too — \texttt{zero}, \texttt{one}, \texttt{eps}, \texttt{typemin}, \texttt{typemax} — alongside the predicates and stepping operations above; they answer from the format{\textquotesingle}s own grid, not from any IEEE assumption about it.



\subsection{Try it}



\label{14108577375059315810}{}


Using only what this tutorial covered: what does \texttt{decode(\allowbreak{}Binary8p4se(\allowbreak{}0x01))} return, and how is it different from \texttt{decode(Binary8p4se(1))}? Work it out before checking.



<details> <summary>Answer</summary>



\texttt{Binary8p4se(0x01)} reads code point 1 — the smallest positive subnormal for this format, \texttt{0.0009765625} (shown as \texttt{MinPositiveOf(\allowbreak{}Binary8p4se)} in the User Guide). \texttt{Binary8p4se(1)} projects the numeric value one onto the nearest datum, \texttt{1.0}. Same digit, two different constructors, two very different answers — the pitfall from this tutorial in miniature.



</details>



\subsection{see also}



\label{9115667341246479961-dup2}{}


\href{tutorial2\_\allowbreak{}projection.md}{Projection: Rounding \& Saturation}, \href{ref\_\allowbreak{}formats.md}{Format Names \& Queries}.



\section{Projection: Rounding \& Saturation}



\label{5636167919932701518}{}


Values, Code Points \& Conversion covered how a value gets made — \texttt{T(x)} projects, \texttt{T(c::Unsigned)} reads a code point, \texttt{decode}/\allowbreak\texttt{codepoint} round-trip exactly. This tutorial is about the projection step itself: what the six rounding modes and three saturation modes do, and how to choose one without leaving a footgun for the rest of your session.



\subsection{One value, every rounding mode}



\label{14660386703135810651}{}


\texttt{2.30078125} sits between the \texttt{Binary8p4se} grid points 2.25 and 2.5 (ulp = 0.25 in \texttt{[2, 4)}). Run it through all six deterministic rounding modes, with \texttt{SatNone}:




\nopagebreak[4]
\begin{minipage}{\linewidth}
\begin{minted}{julia}
using SmallFloats

for μ in (NearestTiesToEven(), NearestTiesToAway(), TowardPositive(),
          TowardNegative(), TowardZero(), ToOdd())
    v = Convert(Binary8p4se, ProjSpec(μ, SatNone()), 2.30078125)
    println(rpad(string(typeof(μ)), 20), " → ", decode(v))
end
\end{minted}




\begin{minted}{text}
NearestTiesToEven    → 2.25
NearestTiesToAway    → 2.25
TowardPositive       → 2.5
TowardNegative       → 2.25
TowardZero           → 2.25
ToOdd                → 2.25
\end{minted}
\end{minipage}




(\texttt{ToOdd} keeps 2.25 because its significand is already odd; try 2.05 to see it move.)



That loop is the whole rounding half of a \texttt{ProjSpec} laid out at once: nearest (two tie-breaking rules), the two directed-toward-infinity modes, toward zero, and round-to-odd. Cross those six against the three saturation modes (\texttt{SatFinite}, \texttt{SatPropagate}, \texttt{SatNone}) and you have the full deterministic grid — 18 predefined constants, each named \texttt{R<mode>\_\allowbreak{}Sat<mode>} (\texttt{RNE\_\allowbreak{}SN}, \texttt{RTZ\_\allowbreak{}SF}, and so on), plus three stochastic rounding families layered on top. \texttt{RNE\_\allowbreak{}SN} is the package{\textquotesingle}s default, which is why \texttt{+} agreed with \texttt{Add(T, RNE\_\allowbreak{}SN, ...)} in the Quickstart.



\subsection{The overflow triptych}



\label{6238881727759123804}{}


Saturation only shows its hand when a result falls outside the format{\textquotesingle}s range. \texttt{Binary8p4se} tops out at \texttt{MaxFinite = 224}; multiply 200 by 2 (exact product



\begin{itemize}
\item[1. ] under three specs:

\end{itemize}



\begin{minted}{jlcon}
julia> w, two = Binary8p4se(200.0), Binary8p4se(2.0);

julia> Multiply(Binary8p4se, RNE_SN, w, two)       # overflow → ±Inf
Binary8p4se(Inf ≡ 0x7f)

julia> Multiply(Binary8p4se, RNE_SF, w, two)     # overflow → MaxFinite
Binary8p4se(224.0 ≡ 0x7e)

julia> Multiply(Binary8p4se, RTZ_SN, w, two)
Binary8p4se(224.0 ≡ 0x7e)   # SatNone + directed-toward-zero clamps finite
\end{minted}



Read the three rows as two separate axes moving one at a time. Row 1 to row 2 holds rounding fixed (\texttt{RNE}) and switches saturation (\texttt{SN} → \texttt{SF}): same overflow, but \texttt{SatFinite} clamps to the extremal finite value instead of going to infinity. Row 1 to row 3 holds saturation fixed (\texttt{SN}) and switches rounding (\texttt{RNE} → \texttt{RTZ}): a directed mode that points \emph{away} from the overflow lands on the finite extreme even under \texttt{SatNone}.



\begin{tcolorbox}[toptitle=-1mm,bottomtitle=1mm,colback=admonition-note!50!white,colframe=admonition-note,title=\textbf{Note}]
\texttt{SatNone} is not {\textquotedbl}no saturation handling{\textquotedbl}: it is the draft{\textquotesingle}s row set in which nearest modes overflow to ±Inf (or NaN for finite formats) while directed modes that point away from the overflow clamp to the extremal finite value.

\end{tcolorbox}


\subsection{Session defaults, used safely}



\label{11919694300743617532}{}


Every \texttt{ProjSpec} above was named explicitly as the second argument. But \texttt{+}, \texttt{exp}, and \texttt{T(x)} all read a \emph{session default} instead — \texttt{DefaultProjection()}, which starts at \texttt{RNE\_\allowbreak{}SN}. You will often want to explore under a different default without touching the rest of your session. Do that with \texttt{with\_\allowbreak{}default\_\allowbreak{}projection}, which sets the default for the duration of a block and restores it afterward, even if the block throws:




\nopagebreak[4]
\begin{minted}{julia}
with_default_projection(RTZ_SF) do
    Binary8p4se(1e9)          # clamps, inside this block only
end
\end{minted}



Outside that block, \texttt{DefaultProjection()} is exactly what it was before you called it. This is the form to reach for by default — \texttt{with\_\allowbreak{}default\_\allowbreak{}type} and \texttt{with\_\allowbreak{}default\_\allowbreak{}returntype} follow the same shape for the other two session defaults.



Sometimes you do want to change the default for the rest of a script or REPL session — for instance, to make every later construction saturate instead of overflowing to infinity. That{\textquotesingle}s what \texttt{DefaultProjection!} is for, and it comes with a hazard the safe form above doesn{\textquotesingle}t have:



\begin{tcolorbox}[toptitle=-1mm,bottomtitle=1mm,colback=admonition-warning!50!white,colframe=admonition-warning,title=\textbf{The session defaults are process-global and they stay set}]
\texttt{DefaultProjection!} and its siblings mutate state shared by every module in the process. Nothing scopes them for you: set one in a script, a notebook cell, or a REPL session and \emph{every} later \texttt{T(x::Real)}, \texttt{a + b} and \texttt{Exp(x)} follows it until something sets it back.

\end{tcolorbox}


If you do call it, restore the default explicitly when you{\textquotesingle}re done:




\nopagebreak[4]
\begin{minted}{jlcon}
julia> x, y = Binary8p4se(200.0), Binary8p4se(2.0);

julia> x * y                              # RNE_SN: overflow → +Inf
Binary8p4se(Inf ≡ 0x7f)

julia> DefaultProjection!(RTZ_SF);

julia> x * y                              # now clamps to MaxFinite
Binary8p4se(224.0 ≡ 0x7e)

julia> DefaultProjection!(RNE_SN)         # restore before moving on — see below
(NearestTiesToEven, SatNone)

julia> Multiply(Binary8p4se, RNE_SN, x, y)   # explicit ρ is unaffected
Binary8p4se(Inf ≡ 0x7f)
\end{minted}



The last line is the important contrast: naming \texttt{ρ} explicitly, as in every \texttt{Op(T, ρ, ...)} call so far, is never affected by the session default in either direction. The default only governs the convenience forms (\texttt{x * y}, bare \texttt{T(x)}, \texttt{Exp(x)}).



\subsection{First contact with stochastic rounding}



\label{398785361657144390}{}


Everything above has been deterministic: the same operands under the same \texttt{ProjSpec} always give the same code point. The rounding-mode list also includes stochastic families — \texttt{StochasticA\{N\}}, \texttt{StochasticB\{N\}}, \texttt{StochasticC\{N\}} — whose outcome depends on \texttt{N} random bits drawn per projection. You can pin that draw explicitly with the \texttt{R} keyword, which makes a single stochastic projection exactly reproducible:




\nopagebreak[4]
\begin{minted}{jlcon}
julia> σ = RSA_SN();                 # ≡ ProjSpec(StochasticA{8}(), SatNone())

julia> Add(Binary8p4se, σ, Binary8p4se(2.0), Binary8p4se(0.03125); rng = Xoshiro(1))
Binary8p4se(2.0 ≡ 0x48)

julia> Add(Binary8p4se, σ, Binary8p4se(2.0), Binary8p4se(0.03125); R = 0)
Binary8p4se(2.0 ≡ 0x48)      # smallest draw: rounds down

julia> Add(Binary8p4se, σ, Binary8p4se(2.0), Binary8p4se(0.03125); R = 255)
Binary8p4se(2.25 ≡ 0x49)     # largest draw: rounds up
\end{minted}



The exact fraction here is 1/8 of an ulp, so over all 256 draws exactly 32 round up — \texttt{StochasticA} is unbiased in expectation. \texttt{R} must lie in \texttt{0:2{\textasciicircum}N-1}. That{\textquotesingle}s as far as stochastic rounding goes in this tutorial; Tutorial 5 is entirely about it.



\subsection{Try it}



\label{14108577375059315810-dup1}{}


Using \texttt{with\_\allowbreak{}default\_\allowbreak{}projection}, compute \texttt{Binary8p4se(1e9)} under \texttt{RTP\_\allowbreak{}SF} (toward positive infinity, clamp to finite) without disturbing the session default, then confirm the default is still \texttt{RNE\_\allowbreak{}SN} afterward.



<details> <summary>Answer</summary>




\begin{minted}{julia}
with_default_projection(RTP_SF) do
    Binary8p4se(1e9)
end
\end{minted}



gives \texttt{Binary8p4se(\allowbreak{}224.0 ≡ 0x7e)} (clamped to \texttt{MaxFinite}, since \texttt{RTP} points toward the overflow and \texttt{SatFinite} clamps). Immediately after the block, \texttt{DefaultProjection()} is back to \texttt{RNE\_\allowbreak{}SN} — the whole point of the combinator over calling \texttt{DefaultProjection!} by hand.



</details>



\subsection{see also}



\label{9115667341246479961-dup3}{}


\href{tutorial3\_\allowbreak{}arrays.md}{Operations \& Arrays}, \href{ref\_\allowbreak{}projections.md}{Projection Specifications}.



\section{Operations \& Arrays}



\label{4519771177622430269}{}


Values, Code Points \& Conversion and Projection: Rounding \& Saturation worked entirely with scalars: \texttt{T(x)} to construct, \texttt{Convert}/\allowbreak\texttt{Op(T, ρ, ...)} to project explicitly, and \texttt{+}/\allowbreak\texttt{exp}/... as the same-format, default-projection convenience forms. This tutorial scales that up to arrays, and introduces the machinery — table gathers, \texttt{vmap}, the \texttt{Op} callable, counting sort — that makes bulk work fast.



\subsection{The two registers, again}



\label{12349619288888336319}{}


Every scalar operation used on those pages came from one of two registers, and the same split carries over to arrays:



\begin{itemize}
\item \textbf{The spec-named register} exposes every draft operation under its draft name with an explicit result format and spec: \texttt{Op(fr, ρ, operands...)}. This is \texttt{Add}, \texttt{Multiply}, \texttt{Exp}, and the rest of the catalog from Projection: Rounding \& Saturation, where you name \texttt{ρ} yourself.


\item \textbf{The Base register} makes each format an ordinary Julia number under its \emph{default} spec (\texttt{RNE\_\allowbreak{}SN} unless you changed it): \texttt{+ - * /}, \texttt{exp}, \texttt{sqrt}, \texttt{min}, \texttt{max}, comparisons — same-format operands only, no silent cross-format promotion.

\end{itemize}


Both registers exist because they serve different code: library code that must be insensitive to the session default names \texttt{ρ} explicitly through the spec-named register; exploratory code reads naturally through the Base register{\textquotesingle}s operators. Both extend to arrays without any new vocabulary.



\subsection{Array calls: \texttt{Add(T, ρ, A, B)}, \texttt{Exp(T, ρ, A)}}



\label{5567456132284147075}{}


Every registered operation has array methods with exactly the same argument shape as the scalar form, operands replaced by arrays:




\nopagebreak[4]
\begin{minted}{jlcon}
julia> A = Binary8p4se.(randn(Xoshiro(7), 4) .* 2)
4-element Vector{Binary8p4se}:
 Binary8p4se(-0.875 ≡ 0xbe)
 Binary8p4se(4.0 ≡ 0x50)
 Binary8p4se(-3.25 ≡ 0xcd)
 Binary8p4se(2.25 ≡ 0x49)

julia> Exp(Binary8p4se, RNE_SN, A)
4-element Vector{Binary8p4se}:
 Binary8p4se(0.40625 ≡ 0x35)
 Binary8p4se(56.0 ≡ 0x6e)
 Binary8p4se(0.0390625 ≡ 0x1a)
 Binary8p4se(9.0 ≡ 0x59)
\end{minted}



\texttt{Add(T, ρ, A, B)} reads the same way, elementwise, for a two-array call.



\subsection{Watching a table build}



\label{8296939794981940873}{}


For pure (non-stochastic) specs, unary and binary array calls run as \textbf{table gathers}: the first call builds and caches a result table for that exact \texttt{(op, formats, ρ)} specialization — 256 bytes for a unary 8-bit operation, 64 KiB for binary — and every later element costs a single lookup. You can watch this happen directly:




\nopagebreak[4]
\begin{minipage}{\linewidth}
\begin{minted}{julia}
using SmallFloats, Random
empty_tables!()
rng = Xoshiro(4)
pre = Binary8p4se.(randn(rng, 4096) .* 2.5)       # pre-activations, ±2.5σ
act = Tanh(Binary8p4se, RNE_SN, pre)          # table-gather kernel

(table_bytes(),
 round(count(v -> abs(decode(v)) == 1.0, act) / 4096; digits=3),
 length(unique(codepoint.(act))),
 decode(NextLessThan(one(Binary8p4se))))
\end{minted}




\begin{minted}{text}
(256, 0.384, 100, 0.9375)
# LUT size; fraction saturated to ±1; distinct output codes; last value below 1
\end{minted}
\end{minipage}




\texttt{empty\_\allowbreak{}tables!()} clears the cache; \texttt{table\_\allowbreak{}bytes()} reports its current footprint. A second unary activation over the same array grows the cache rather than replacing it:




\nopagebreak[4]
\begin{minipage}{\linewidth}
\begin{minted}{julia}
actS = Softplus(Binary8p4se, RNE_SN, pre)
(length(unique(codepoint.(actS))), table_bytes())
\end{minted}




\begin{minted}{text}
(61, 512)          # softplus keeps 61 distinct codes here; two LUTs now cached
\end{minted}
\end{minipage}




Every table entry is built through the exact scalar path, so it is bit-identical by construction to calling the scalar operation element by element — the table is a cache, not a different, faster-but-approximate implementation. Stochastic calls never use a table: they always run the scalar pipeline per element, because each element draws its own random bits.



\subsection{\texttt{vmap} and \texttt{vmap!}}



\label{5255729870829091668}{}


\texttt{vmap}/\allowbreak\texttt{vmap!} are the generic, table-aware entry points behind the named array calls above. \texttt{vmap!} is the in-place form and is what you want when you already have a destination array to fill rather than a fresh one to allocate:




\nopagebreak[4]
\begin{minted}{julia}
C = vmap(:Add, T, ρ, A, B)
vmap!(C, Val(:Add), T, ρ, A, B)
\end{minted}



Reach for \texttt{vmap!} in a loop that reuses a buffer; reach for the named form (\texttt{Add(T, ρ, A, B)}) everywhere else — they run the same kernel.



\subsection{The \texttt{Op} callable, with \texttt{map} and broadcast}



\label{3754941402779852898}{}


\texttt{Op(T, ρ, A)} above already reads like \texttt{map}. When you want Julia{\textquotesingle}s own array verbs — \texttt{map}, broadcast — rather than the named array call, bind the operation, result format, and projection into an \texttt{Op} value and use it as a function:




\nopagebreak[4]
\begin{minted}{jlcon}
julia> A = Binary8p4se.([1.5, 0.25]); B = Binary8p4se.([0.25, 1.5]);

julia> map(Op(:Add, Binary8p4se, RNE_SN), A, B)
2-element Vector{Binary8p4se}:
 Binary8p4se(1.75 ≡ 0x46)
 Binary8p4se(1.75 ≡ 0x46)

julia> Op(:Exp, Binary8p4se, RNE_SN).(A)
2-element Vector{Binary8p4se}:
 Binary8p4se(4.5 ≡ 0x51)
 Binary8p4se(1.25 ≡ 0x42)
\end{minted}



\texttt{Op} is a singleton type — the operation, format, and projection all live in its type parameters — so a call specializes exactly as \texttt{Add(\allowbreak{}Binary8p4se, RNE\_\allowbreak{}SN, x, y)} does and allocates nothing beyond the array \texttt{map} or broadcast itself builds.



\subsection{Sorting}



\label{12368189431511364825}{}


Sorting is special-cased for every \texttt{Binary} format: values compare through integer order keys, and vectors of \texttt{Binary} sort with an O(n) \textbf{counting sort} installed as the default algorithm. \texttt{sort(A)} just works, NaN sorts last (first under \texttt{rev=true}), matching Base{\textquotesingle}s conventions:




\nopagebreak[4]
\begin{minipage}{\linewidth}
\begin{minted}{julia}
decode.(sort(Binary4p2se.(0x00:0x0f)))      # broadcast the code-point constructor
\end{minted}




\begin{minted}{text}
[-Inf, -2.0, -1.5, -1.0, -0.75, -0.5, -0.25, 0.0,
  0.25, 0.5, 0.75, 1.0, 1.5, 2.0, Inf, NaN]
\end{minted}
\end{minipage}




\texttt{TotalOrder(x, y)} exposes the draft{\textquotesingle}s total order directly if you need the comparison itself rather than a sorted array — single NaN largest.



\subsection{Try it}



\label{14108577375059315810-dup2}{}


Given \texttt{A = Binary8p4se.(\allowbreak{}[1.5, 0.25])} and \texttt{B = Binary8p4se.(\allowbreak{}[0.25, 1.5])} from above, write the \texttt{Op}-callable version of \texttt{Add(\allowbreak{}Binary8p4se, RNE\_\allowbreak{}SN, A, B)} using broadcast rather than \texttt{map}.



<details> <summary>Answer</summary>




\begin{minted}{julia}
Op(:Add, Binary8p4se, RNE_SN).(A, B)
\end{minted}



This broadcasts the same bound callable used with \texttt{map} above and returns the same two-element vector, \texttt{[1.75, 1.75]} in \texttt{Binary8p4se}.



</details>



\subsection{see also}



\label{9115667341246479961-dup4}{}


\href{tutorial4\_\allowbreak{}blocks.md}{Blocks \& MX-Style Quantization}, \href{ref\_\allowbreak{}operations.md}{Operation Catalog}, \href{ref\_\allowbreak{}arrays\_\allowbreak{}blocks.md}{Arrays, Blocks \& Packed Storage}.



\section{Blocks \& MX-Style Quantization}



\label{5821215397688818864}{}


Operations \& Arrays covered arrays of a single format. This tutorial covers \texttt{Block} values, which pair one shared scale with a group of narrow elements — the MX-style scheme for buying dynamic range that a single narrow format cannot hold on its own.



\subsection{Scale × element}



\label{9198442399466125325}{}


Small formats have small ranges. A \texttt{Block\{B,FS,FE\}} fights that by sharing one scale, in format \texttt{FS}, across \texttt{B} elements in format \texttt{FE}; the value each element \emph{represents} is \texttt{scale × element}, not the element alone. Construct one directly from a scale and its elements:




\nopagebreak[4]
\begin{minted}{jlcon}
julia> b = Block(Binary8p1uf(4.0), Binary8p4se(1.5), Binary8p4se(-0.75),
                 Binary8p4se(2.0), Binary8p4se(0.5))
Block{4, Binary8p1uf, Binary8p4se}(Binary8p1uf(4.0 ≡ 0x82), (…))
\end{minted}



Read the type parameters: block size 4, scale format \texttt{Binary8p1uf} (unsigned — scales are non-negative — with the huge range that precision 1 buys), element format \texttt{Binary8p4se}. The represented values are \texttt{4.0 × 1.5 = 6.0}, \texttt{4.0 × -0.75 = -3.0}, and so on.



\subsection{\texttt{ConvertToBlockMaxAbsFinite}}



\label{2152059938679189190}{}


Building a block by hand from a scale you already chose is one path; the more common one is to quantize a tuple of values and let the draft{\textquotesingle}s algorithm pick the scale for you, from the maximum finite \texttt{|element|}, then project each element against it:




\nopagebreak[4]
\begin{minted}{jlcon}
julia> ConvertToBlockMaxAbsFinite(Binary8p1uf, Binary8p4se, RNE_SN, RNE_SN,
           (Binary8p4se(100.0), Binary8p4se(-12.0), Binary8p4se(0.5), Binary8p4se(3.0)))
Block{4, Binary8p1uf, Binary8p4se}(Binary8p1uf(64.0 ≡ 0x86), (…))
\end{minted}



Read the arguments in order: scale format, element format, a projection spec for the scale, a projection spec for the elements, then the values. The algorithm picked scale 64 from the largest-magnitude element (100) and projected each of the four elements against that scale.



\subsection{\texttt{ConvertFromBlock}}



\label{15940086769284149764}{}


Read a block back out as plain values in a format of your choosing: each \texttt{scale × element} is decoded exactly and projected once.




\begin{minted}{jlcon}
julia> ConvertFromBlock(Binary8p3se, RNE_SN, b)   # decode scale×elem, project
(Binary8p3se(6.0 ≡ 0x4a), Binary8p3se(-3.0 ≡ 0xc6), Binary8p3se(8.0 ≡ 0x4c), Binary8p3se(2.0 ≡ 0x44))
\end{minted}



\subsection{The stage-wide pitfall}



\label{13625156544789482334}{}


\texttt{ConvertToBlockMaxAbsFinite} takes already-\texttt{Binary} elements as input. If you stage your raw \texttt{Float64} data through the narrow \emph{element} format before quantizing into a block, \texttt{SatFinite} clamps the large values \textbf{before} the block scale ever gets a chance to absorb them — and the block buys you nothing. This is the single most consequential mistake with this API, so watch it happen on real data before you write code that makes it.



Each row of \texttt{W} here lives at a different scale, spanning roughly \texttt{2{\textasciicircum}16} across the matrix — exactly the shape MX-style blocking exists for:




\nopagebreak[4]
\begin{minipage}{\linewidth}
\begin{minted}{julia}
using SmallFloats, Random, Statistics

rng = Xoshiro(3)
W = randn(rng, 64, 64) .* 2 .* (2.0 .^ (collect(0:63) ./ 4))   # per-row scales

function mx_rows(W, ::Type{FST}, ::Type{FS}, ::Type{FE}, ::Val{B}) where {FST,FS,FE,B}
    n, m = size(W)
    [begin
        # stage through a WIDE-RANGE format so nothing clamps before scaling
        seg = ntuple(k -> Convert(FST, RNE_SF, W[i, (j-1)*B + k]), Val(B))
        ConvertToBlockMaxAbsFinite(FS, FE, RNE_SN, RNE_SN, seg)
    end for i in 1:n, j in 1:m ÷ B]
end

blocks = mx_rows(W, Binary8p2se, Binary8p1uf, Binary8p4se, Val(32))
recon = [let b = blocks[i, (j-1) ÷ 32 + 1]
             decode(b.s) * decode(b.x[(j-1) % 32 + 1])
         end for i in 1:64, j in 1:64]
plain = [decode(Convert(Binary8p4se, RNE_SF, w)) for w in W]

relerr(A) = sqrt(mean(((W .- A) ./ max.(abs.(W), 1e-9)).^2))
relerr(plain), relerr(recon)
\end{minted}




\begin{minted}{text}
(0.616, 0.109)      # plain 8p4se clamps the large rows; MX tracks them
\end{minted}
\end{minipage}




The plain, unblocked \texttt{Binary8p4se} conversion clamps every row whose scale exceeds what \texttt{Binary8p4se} alone can hold — relative error 0.616. Staging through the wide-range \texttt{Binary8p2se} first, then letting the per-row block scale absorb the range, brings that down to 0.109 — the residual is now just the staging format{\textquotesingle}s own precision.



\begin{tcolorbox}[toptitle=-1mm,bottomtitle=1mm,colback=admonition-warning!50!white,colframe=admonition-warning,title=\textbf{Stage wide, then block-quantize}]
\texttt{ConvertToBlockMaxAbsFinite} takes already-\texttt{Binary} elements. If you stage your \texttt{Float64} data through the \emph{element} format first, \texttt{SatFinite} clamps the large values \textbf{before} the block scale can absorb them, and blocks buy you nothing — we measured exactly that (0.617 vs 0.616) with \texttt{Binary8p4se} staging. Stage through a wide-range format (\texttt{Binary8p2se} above); the residual 0.109 here is the staging format{\textquotesingle}s own precision, which bounds what any downstream scheme can keep.

\end{tcolorbox}


At \texttt{B = 32} with an 8-bit scale the storage cost is 8.25 bits/value.



\subsection{\texttt{BlockDotProduct}: one rounding, at the end}



\label{17638104271940170310}{}


\texttt{BlockDotProduct} computes every lane product and the accumulation \emph{exactly}, in a wide carrier, and projects only once, at the very end — there is no hidden intermediate rounding lane by lane:




\nopagebreak[4]
\begin{minipage}{\linewidth}
\begin{minted}{julia}
rng = Xoshiro(9)
a64, b64 = randn(rng, 32), randn(rng, 32)
qb(v) = ConvertToBlockMaxAbsFinite(Binary8p1uf, Binary8p4se, RNE_SN, RNE_SN,
            ntuple(i -> Convert(Binary8p4se, RNE_SF, v[i]), Val(32)))
dq = BlockDotProduct(Binary8p4se, RNE_SN, qb(a64), qb(b64))
(a64'b64, decode(dq))
\end{minted}




\begin{minted}{text}
(5.0634, 4.5)       # difference is input quantization only, never accumulation error
\end{minted}
\end{minipage}




The gap between \texttt{5.0634} and \texttt{4.5} here is entirely the up-front quantization of \texttt{a64} and \texttt{b64} into 4-bit-precision elements — accumulation itself contributes nothing, because the lane products and their sum are carried exactly until the single final projection. Reductions besides the dot product follow the same shape: \texttt{BlockReduceAdd}, \texttt{BlockReduceMultiply}.



\subsection{\texttt{BlockVector}}



\label{13615094720265409582}{}


For many blocks of the same shape, \texttt{BlockVector} stores them in a structure-of-arrays layout rather than as an array of individually boxed \texttt{Block} values — a memory-bound storage form for blocks, the same way packed code-point storage exists for bare values.



\subsection{Try it}



\label{14108577375059315810-dup3}{}


You have raw \texttt{Float64} data with widely varying per-row magnitude, and you want to quantize it into \texttt{Block\{32,Binary8p1uf,Binary8p4se\}}. Name the staging mistake to avoid, and the one-line fix.



<details> <summary>Answer</summary>



The mistake is converting straight to \texttt{Binary8p4se}, the narrow element format.



Passing those narrowed elements to \texttt{ConvertToBlockMaxAbsFinite} lets \texttt{SatFinite} clamp large rows before the block scale can absorb them. Stage through a wide-range format first, such as \texttt{Convert(\allowbreak{}Binary8p2se, RNE\_\allowbreak{}SF, w)}, and pass those staged elements onward—the \texttt{seg = ntuple(\allowbreak{}k -> Convert(\allowbreak{}FST, RNE\_\allowbreak{}SF, ...), ...)} line in the demo above.



</details>



\subsection{see also}



\label{9115667341246479961-dup5}{}


\href{tutorial5\_\allowbreak{}stochastic.md}{Stochastic Rounding}, \href{ref\_\allowbreak{}arrays\_\allowbreak{}blocks.md}{Arrays, Blocks \& Packed Storage}.



\section{Stochastic Rounding}



\label{4094877813010619789}{}


Projection: Rounding \& Saturation introduced the stochastic families — \texttt{StochasticA\{N\}}, \texttt{StochasticB\{N\}}, \texttt{StochasticC\{N\}} — and showed one projection pinned exactly with \texttt{R = 0} and \texttt{R = 255}. This tutorial builds the full picture: why stochastic rounding is unbiased in expectation, how to audit that exhaustively over every draw, how to make a stochastic stream reproducible, and the payoff — an accumulator that keeps absorbing small updates where nearest rounding stalls.



\subsection{Unbiasedness: nearest is not, stochastic is (on average)}



\label{10570852926174133617}{}


Nearest rounding maps \texttt{0.30078125} to \texttt{0.3125} in \texttt{Binary8p4se} — every single time. That{\textquotesingle}s a systematic bias of about \texttt{+0.0117}. Stochastic rounding gets it right \emph{on average} instead of every time:




\nopagebreak[4]
\begin{minipage}{\linewidth}
\begin{minted}{julia}
target = 0.30078125                       # ν = 3/16 of an ulp above 0.296875
σ = RSA_SN(16)                       # StochasticA{16} with SatNone
rng = Xoshiro(1)
m = mean(decode(Convert(Binary8p4se, σ, target; rng)) for _ in 1:200_000)
(decode(Binary8p4se(target)), m)
\end{minted}




\begin{minted}{text}
(0.3125, 0.300765)                        # RNE result vs stochastic mean ≈ target
\end{minted}
\end{minipage}




Over 200,000 draws, the stochastic mean lands within noise of the true target, \texttt{0.30078125} — while the deterministic result is stuck 0.0117 high on every single call. This is the property that makes stochastic rounding matter for accumulating many small contributions (gradients, activation statistics) in a low-precision format.



\subsection{The R-sweep: auditing the distribution exactly}



\label{1146114789827069351}{}


Every stochastic projection is a deterministic function of the draw \texttt{R}, so its \emph{distribution} is checkable exactly rather than just approximately by sampling. For \texttt{x = 2 + 3/64} in \texttt{Binary8p4se} the fraction above the lower grid point is \texttt{ν = 3/16} of an ulp, so \texttt{StochasticA\{4\}} — a 4-bit budget, 16 possible draws — must round up for exactly 3 of the 16 draws:




\nopagebreak[4]
\begin{minipage}{\linewidth}
\begin{minted}{julia}
σ4 = RSA_SN(4)                  # ≡ ProjSpec(StochasticA{4}(), SatNone())
x = 2.0 + 3/64
count(decode(SmallFloats.project(Binary8p4se, σ4, x; R)) == 2.25 for R in 0:15)
\end{minted}




\begin{minted}{text}
3
\end{minted}
\end{minipage}




Sweeping \texttt{R} over its whole range turns {\textquotedbl}is my stochastic pipeline unbiased?{\textquotedbl} from a statistical question into an exhaustive one — this is the same pattern the shipped test suite uses to verify every stochastic mode.



\subsection{Reproducibility}



\label{2141891870398561726}{}


A stochastic projection needs random bits from somewhere, and you control the source the same three ways as any other random draw in the package: implicitly from the task-local RNG, from an explicit \texttt{rng} you pass in, or by fixing the draw \texttt{R} itself (as in the sweep above and in Projection: Rounding \& Saturation). For a reproducible \emph{stream} rather than one pinned draw, pass a seeded \texttt{rng}:




\nopagebreak[4]
\begin{minted}{jlcon}
julia> σ = RSA_SN();                 # ≡ ProjSpec(StochasticA{8}(), SatNone())

julia> Add(Binary8p4se, σ, Binary8p4se(2.0), Binary8p4se(0.03125); rng = Xoshiro(1))
Binary8p4se(2.0 ≡ 0x48)
\end{minted}



Calling this again with a freshly-seeded \texttt{Xoshiro(1)} reproduces exactly the same result, the same way \texttt{Random.seed!} and explicit \texttt{AbstractRNG} arguments work for \texttt{rand}/\allowbreak\texttt{randn} elsewhere in the package. Prefer a seeded \texttt{rng} in tests and anywhere else you need to replay a run; reach for explicit \texttt{R} only when you need to force one specific draw, as in the sweep above.



The three controls compose the way you{\textquotesingle}d expect: no \texttt{rng} and no \texttt{R} draws from the task-local generator (least reproducible, fine for exploration); an explicit \texttt{rng} pins the whole stream (reproducible across a run, and the right default for anything you{\textquotesingle}ll re-run); explicit \texttt{R} pins one draw exactly (the right tool for a unit test that checks one specific rounding outcome). All three reach the same stochastic kernel — only the source of randomness changes.



\subsection{The payoff: swamping and stall}



\label{15557402075891912974}{}


Accumulate 400 gradient-style steps of \texttt{0.011} into a \texttt{Binary8p3se} accumulator. Under nearest rounding, once the accumulator{\textquotesingle}s magnitude dwarfs the increment, every add rounds straight back to the accumulator — it \textbf{stalls}, permanently. Stochastic rounding keeps absorbing the increments in expectation instead:




\nopagebreak[4]
\begin{minipage}{\linewidth}
\begin{minted}{julia}
function accumulate_demo(nsteps, g)
    σ = RSA_SN(16)
    rng = Xoshiro(11)
    acc_rne = Binary8p3se(0.0); acc_sto = Binary8p3se(0.0)
    for _ in 1:nsteps
        acc_rne = Add(Binary8p3se, RNE_SN, acc_rne, Convert(Binary8p3se, RNE_SN, g))
        acc_sto = Add(Binary8p3se, σ, acc_sto, Convert(Binary8p3se, σ, g; rng); rng)
    end
    (nsteps * g, decode(acc_rne), decode(acc_sto))
end
accumulate_demo(400, 0.011)
\end{minted}




\begin{minted}{text}
(4.4, 0.125, 5.0)   # exact sum, RNE accumulator (stalled!), stochastic accumulator
\end{minted}
\end{minipage}




The exact sum after 400 steps is 4.4. The nearest-rounding accumulator gets stuck at 0.125 — once the accumulator reaches a binade where the ulp exceeds \texttt{0.011}, every subsequent \texttt{+0.011} rounds back to the same code point, and no number of further steps will move it. The stochastic accumulator lands on 5.0 on this seed: noisy, but unbiased, and it never permanently stops moving the way the nearest accumulator does. In practice you keep a wider accumulator format when you can, and reach for stochastic rounding when you can{\textquotesingle}t.



\subsection{Try it}



\label{14108577375059315810-dup4}{}


Without running any code: for \texttt{StochasticA\{N\}} with a budget of \texttt{N} bits, about what fraction of an ulp{\textquotesingle}s worth of draws would you expect to round \emph{down} if the true value sits exactly at \texttt{ν = 1/4} of an ulp above the lower grid point? Use the R-sweep logic above to reason it out, then check against the pattern in that section.



<details> <summary>Answer</summary>



By the same exhaustive logic as the \texttt{ν = 3/16} sweep above, a value sitting at \texttt{ν = 1/4} of an ulp rounds up on the fraction of draws proportional to \texttt{ν} itself — one quarter of the \texttt{2{\textasciicircum}N} possible draws round up, and the remaining three quarters round down. \texttt{StochasticA} is unbiased precisely because that fraction always equals \texttt{ν}, whatever binade or format you{\textquotesingle}re in.



</details>



\subsection{see also}



\label{9115667341246479961-dup6}{}


\href{howto\_\allowbreak{}choose\_\allowbreak{}format.md}{How-To Guides}, \href{ref\_\allowbreak{}formats.md}{Format Names \& Queries}, \href{ref\_\allowbreak{}projections.md}{Projection Specifications}.



\chapter{How-To Guides}


\section{Choose a Format}



\label{4793595142642214899}{}


Pick a \texttt{Binary\{K,P,SGN,EXT\}} format that fits a specific tensor, signal, or quantity.



\subsection{Ingredients}



\label{13776634711169142479}{}


\begin{itemize}
\item A sample of the data (or a plausible generator for it) so you can measure damage before committing.


\item \texttt{formatname}, \texttt{MaxFiniteOf}, \texttt{decode} for inspecting candidates.


\item \texttt{RNE\_\allowbreak{}SF} (or another \texttt{ProjSpec}) to project a sample through each candidate.

\end{itemize}


\subsection{Recipe}



\label{111379208343275940}{}


\textbf{1. Enumerate the format before committing.} Small formats are small enough to look at in full. Sorting \texttt{Binary4p2se}{\textquotesingle}s 16 code points shows the subnormal spacing, the binade structure, and the dynamic range in one line:




\nopagebreak[4]
\begin{minipage}{\linewidth}
\begin{minted}{julia}
decode.(sort(Binary4p2se.(0x00:0x0f)))
\end{minted}




\begin{minted}{text}
[-Inf, -2.0, -1.5, -1.0, -0.75, -0.5, -0.25, 0.0,
  0.25, 0.5, 0.75, 1.0, 1.5, 2.0, Inf, NaN]
\end{minted}
\end{minipage}




\textbf{2. At fixed bitwidth, precision trades against range.} For a unit-scale Gaussian tensor, RMSE roughly doubles per lost significand bit while \texttt{MaxFinite} grows explosively — spend bits on precision unless the data{\textquotesingle}s magnitude actually needs the range:




\nopagebreak[4]
\begin{minipage}{\linewidth}
\begin{minted}{julia}
rng = Xoshiro(7); x = randn(rng, 50_000)
for F in (Binary8p5se, Binary8p4se, Binary8p3se, Binary8p2se)
    back = [decode(Convert(F, RNE_SF, xi)) for xi in x]
    println(rpad(formatname(F), 12), " rmse = ", round(sqrt(mean((x .- back).^2)); sigdigits=3),
            "   maxfinite = ", decode(MaxFiniteOf(F)))
end
\end{minted}




\begin{minted}{text}
Binary8p5se  rmse = 0.0133   maxfinite = 15.0
Binary8p4se  rmse = 0.0266   maxfinite = 224.0
Binary8p3se  rmse = 0.0528   maxfinite = 49152.0
Binary8p2se  rmse = 0.103    maxfinite = 2.147483648e9
\end{minted}
\end{minipage}




If your data{\textquotesingle}s dynamic range genuinely exceeds one binade{\textquotesingle}s reach at your precision budget, don{\textquotesingle}t buy that range with a wider element — get it from a block scale instead.



\textbf{3. Bounded-in-[0,1] quantities want unsigned.} A membership degree, a probability, a normalized score — anything that is never negative — is a job for an unsigned finite format, not a signed one spending a code on a sign it never uses. \texttt{Binary8p6uf} gives a [0, 15.5] range at 1/32 resolution near 1:




\nopagebreak[4]
\begin{minipage}{\linewidth}
\begin{minted}{julia}
F = Binary8p6uf
membership(x, a, b, c, d) = clamp(min((x - a) / (b - a), 1.0, (d - x) / (d - c)), 0.0, 1.0)
warm = F(membership(26.0, 15, 20, 24, 28))   # membership of 26 °C in "warm"
hot  = F(membership(26.0, 24, 30, 100, 101)) # … and in "hot"

(decode(warm), decode(hot),
 decode(Minimum(F, RNE_SN, warm, hot)),      # warm AND hot
 decode(Maximum(F, RNE_SN, warm, hot)),      # warm OR hot
 decode(Multiply(F, RNE_SF, warm, hot)))   # product t-norm
\end{minted}




\begin{minted}{text}
(0.5, 0.3359375, 0.3359375, 0.5, 0.16796875)
\end{minted}
\end{minipage}




\textbf{4. Pick finite (\texttt{f}) when \texttt{Inf} would be meaningless.} If nothing in your domain should ever legitimately produce or consume an infinity — a scale factor, a bounded score, a probability-like quantity — a finite-domain format spends the code points IEEE-style formats give to ±Inf on two more finite datums instead. Reach for \texttt{e} (extended) only when overflow-to-infinity is a real outcome you want to detect downstream.



\begin{tcolorbox}[toptitle=-1mm,bottomtitle=1mm,colback=admonition-warning!50!white,colframe=admonition-warning,title=\textbf{An unsigned format is not free precision}]
Going unsigned doubles your positive-side resolution only if the quantity truly never needs a sign. Do not use an unsigned format for anything that can legitimately go negative (gradients, residuals, log-odds) — projecting a negative \texttt{Real} into an unsigned format is not a sign flip, it is a domain violation and produces a different failure than you likely intend. Reach for the fuzzy-inference example above only when the quantity is already bounded in \texttt{[0, ∞)} by construction.

\end{tcolorbox}


\subsection{What can go wrong}



\label{7385351218440078181}{}


\begin{tcolorbox}[toptitle=-1mm,bottomtitle=1mm,colback=admonition-warning!50!white,colframe=admonition-warning,title=\textbf{Enumeration intuition does not scale past small K}]
\texttt{decode.(\allowbreak{}sort(\allowbreak{}F.(\allowbreak{}0x00:...)))} is a genuine format-choice tool for \texttt{K} small enough to print — beyond roughly \texttt{K = 8} the code-point space is too large to eyeball, so switch to the RMSE/MaxFinite sweep instead.

\end{tcolorbox}


\begin{tcolorbox}[toptitle=-1mm,bottomtitle=1mm,colback=admonition-note!50!white,colframe=admonition-note,title=\textbf{RMSE alone hides overflow}]
A format can have a low RMSE on the bulk of your data and still overflow its tails to \texttt{Inf}. Always check \texttt{count(\allowbreak{}isinf, back) /\allowbreak{} length(\allowbreak{}back)} alongside RMSE and max error — see Quantize a Tensor or Model for the full three-number check.

\end{tcolorbox}


\subsection{see also}



\label{9115667341246479961-dup7}{}


\href{howto\_\allowbreak{}quantize\_\allowbreak{}tensor.md}{Quantize a Tensor or Model}, \href{howto\_\allowbreak{}blocks\_\allowbreak{}dynamic\_\allowbreak{}range.md}{Use Blocks to Track Dynamic Range}.



\section{Quantize a Tensor or Model}



\label{4573358416020164411}{}


Project a full weight tensor into a small format, measure the damage, and pack the result for storage.



\subsection{Ingredients}



\label{13776634711169142479-dup1}{}


\begin{itemize}
\item \texttt{Binary8p4se.(...)} (or any format, broadcast) to project every element.


\item \texttt{decode.(...)} to bring the quantized values back to \texttt{Float64} for measurement.


\item \texttt{mean}, \texttt{maximum}, \texttt{count}/\allowbreak\texttt{isinf} from \texttt{Statistics}/\allowbreak\texttt{Base} for MSE, max error, and overflow fraction.


\item \texttt{rawvalue}, \texttt{PackedVector} to pack the quantized codes for storage.

\end{itemize}


\subsection{Recipe}



\label{111379208343275940-dup1}{}


\textbf{1. Project the tensor.} Broadcasting the format constructor projects every weight independently, exactly as a single value would:




\nopagebreak[4]
\begin{minipage}{\linewidth}
\begin{minted}{julia}
using SmallFloats, Random, Statistics

rng = Xoshiro(42)
w = randn(rng, 10_000) .* 0.25          # typical trained-weight scale
q = Binary8p4se.(w)                     # project every weight
back = decode.(q)

mean((w .- back).^2), maximum(abs.(w .- back)), count(isinf, back) / length(back)
\end{minted}




\begin{minted}{text}
(4.41e-5, 0.0312, 0.0)                  # MSE, max |error|, overflow fraction
\end{minted}
\end{minipage}




Report all three numbers together, not just MSE: the max error is exactly half an ulp of the largest binade the data reached (the worst case nearest rounding permits), and the overflow fraction catches tail values that clamped or blew up to \texttt{Inf} — a failure mode MSE alone can hide if it affects only a few elements.



\textbf{2. Pack the quantized model for storage.} Once you have code points, a \texttt{PackedVector} stores them at \texttt{bitwidth(F)} bits per element instead of one byte (or two) per element:




\nopagebreak[4]
\begin{minipage}{\linewidth}
\begin{minted}{julia}
n = 1_000_000
model = [rawvalue(Binary5p2se, UInt8(rand(0:31))) for _ in 1:n]
pv = PackedVector(model)
(sizeof(model), sizeof(pv.data))
\end{minted}




\begin{minted}{text}
(1000000, 625000)   # 8 bits → 5 bits per value; indexing and vmap work directly on pv
\end{minted}
\end{minipage}




\texttt{Binary5p2se} is 5 bits wide, so 1,000,000 values pack into 625,000 bytes instead of 1,000,000 — indexing, \texttt{collect}, and \texttt{vmap} all work directly on \texttt{pv} without a separate unpacking step.



\textbf{3. Verify the quantizer, don{\textquotesingle}t just trust it.} Formats are small enough to check exhaustively rather than spot-check. Enumerate every in-range code point of the source format, project it, and compare against an independent brute-force nearest search:




\nopagebreak[4]
\begin{minipage}{\linewidth}
\begin{minted}{julia}
using SmallFloats

function ref_nearest_distance(::Type{T}, x) where {T}
    fins = [v for v in (rawvalue(T, UInt8(c)) for c in 0:255) if isfinite(decode(v))]
    minimum(abs(setprecision(() -> BigFloat(decode(v)) - BigFloat(x), BigFloat, 256))
            for v in fins)
end

function verify_quantizer()
    ok = true
    for c in 0x00:0xff
        x = decode(rawvalue(Binary8p3se, c))
        (isfinite(x) && abs(x) <= decode(MaxFiniteOf(Binary8p4se))) || continue
        got = Convert(Binary8p4se, RNE_SN, x)
        ok &= abs(decode(got) - x) == Float64(ref_nearest_distance(Binary8p4se, x))
    end
    ok
end
verify_quantizer()
\end{minted}




\begin{minted}{text}
true
\end{minted}
\end{minipage}




This pattern — enumerate the inputs, compare against an independently written reference — costs a few hundred comparisons and is available to your own quantization code at trivial cost.



\subsection{What can go wrong}



\label{7385351218440078181-dup1}{}


\begin{tcolorbox}[toptitle=-1mm,bottomtitle=1mm,colback=admonition-warning!50!white,colframe=admonition-warning,title=\textbf{MSE looks fine while overflow silently eats the tail}]
A tensor with a few large outliers can show excellent MSE while those outliers all clamp to \texttt{Inf} (or \texttt{MaxFinite} under \texttt{SatFinite}). Always check the overflow fraction, and if it is nonzero, either widen the format{\textquotesingle}s range or move to a block-scaled format instead of a bare element format.

\end{tcolorbox}


\begin{tcolorbox}[toptitle=-1mm,bottomtitle=1mm,colback=admonition-default!50!white,colframe=admonition-default,title=\textbf{Packing is for storage, not compute}]
\texttt{PackedVector} is deliberately not in-place-arithmetic-capable: the rule is store packed, compute unpacked. Don{\textquotesingle}t try to quantize \emph{into} a \texttt{PackedVector} element by element — build the plain \texttt{Vector\{F\}} first (as above), then wrap it with \texttt{PackedVector}.

\end{tcolorbox}


\subsection{see also}



\label{9115667341246479961-dup8}{}


\href{howto\_\allowbreak{}choose\_\allowbreak{}format.md}{Choose a Format}, \href{howto\_\allowbreak{}packed\_\allowbreak{}storage.md}{Work with Packed Storage}, \href{howto\_\allowbreak{}blocks\_\allowbreak{}dynamic\_\allowbreak{}range.md}{Use Blocks to Track Dynamic Range}.



\section{Use Blocks to Track Dynamic Range}



\label{7223351510969591640}{}


Keep a bare element format from clamping data whose dynamic range exceeds one binade, by sharing a per-group scale (the MX-style scheme).



\subsection{Ingredients}



\label{13776634711169142479-dup2}{}


\begin{itemize}
\item A wide-range staging format (e.g. \texttt{Binary8p2se}) to hold raw values before scaling.


\item \texttt{ConvertToBlockMaxAbsFinite(\allowbreak{}FS, FE, scale\_\allowbreak{}ρ, element\_\allowbreak{}ρ, values\_\allowbreak{}tuple)} to build one block per group, picking the scale from the group{\textquotesingle}s max |element|.


\item \texttt{Val(B)} for the block size.


\item \texttt{relerr} (or your own metric) to compare plain quantization against blocked quantization.

\end{itemize}


\subsection{Recipe}



\label{111379208343275940-dup2}{}


\textbf{1. Build data whose scale varies by row.} Here each row of a 64×64 matrix lives at a different scale, spanning roughly \texttt{2{\textasciicircum}16} across the whole matrix — exactly the situation a single bare format cannot track:




\nopagebreak[4]
\begin{minted}{julia}
using SmallFloats, Random, Statistics

rng = Xoshiro(3)
W = randn(rng, 64, 64) .* 2 .* (2.0 .^ (collect(0:63) ./ 4))   # per-row scales
\end{minted}



\textbf{2. Stage wide, then block-quantize per row.} Convert each row{\textquotesingle}s raw values through a \emph{wide-range} format first (so nothing clamps before the block scale sees it), then hand fixed-size chunks to \texttt{ConvertToBlockMaxAbsFinite}:




\nopagebreak[4]
\begin{minipage}{\linewidth}
\begin{minted}{julia}
function mx_rows(W, ::Type{FST}, ::Type{FS}, ::Type{FE}, ::Val{B}) where {FST,FS,FE,B}
    n, m = size(W)
    [begin
        # stage through a WIDE-RANGE format so nothing clamps before scaling
        seg = ntuple(k -> Convert(FST, RNE_SF, W[i, (j-1)*B + k]), Val(B))
        ConvertToBlockMaxAbsFinite(FS, FE, RNE_SN, RNE_SN, seg)
    end for i in 1:n, j in 1:m ÷ B]
end

blocks = mx_rows(W, Binary8p2se, Binary8p1uf, Binary8p4se, Val(32))
recon = [let b = blocks[i, (j-1) ÷ 32 + 1]
             decode(b.s) * decode(b.x[(j-1) % 32 + 1])
         end for i in 1:64, j in 1:64]
plain = [decode(Convert(Binary8p4se, RNE_SF, w)) for w in W]

relerr(A) = sqrt(mean(((W .- A) ./ max.(abs.(W), 1e-9)).^2))
relerr(plain), relerr(recon)
\end{minted}




\begin{minted}{text}
(0.616, 0.109)      # plain 8p4se clamps the large rows; MX tracks them
\end{minted}
\end{minipage}




Plain \texttt{Binary8p4se} quantization clamps the large rows and posts a relative error of 0.616; the block scheme tracks each row{\textquotesingle}s own scale and cuts that to 0.109 — the residual there is the staging format{\textquotesingle}s own precision, the floor any downstream scheme can reach.



At \texttt{B = 32} with an 8-bit scale, the storage cost is 8.25 bits/value — slightly more than the bare element format{\textquotesingle}s 8 bits, in exchange for tracking dynamic range the bare format cannot.



\subsection{What can go wrong}



\label{7385351218440078181-dup2}{}


\begin{tcolorbox}[toptitle=-1mm,bottomtitle=1mm,colback=admonition-warning!50!white,colframe=admonition-warning,title=\textbf{Stage wide, then block-quantize}]
\texttt{ConvertToBlockMaxAbsFinite} takes already-\texttt{Binary} elements. If you stage your \texttt{Float64} data through the \emph{element} format first, \texttt{SatFinite} clamps the large values \textbf{before} the block scale can absorb them, and blocks buy you nothing — we measured exactly that (0.617 vs 0.616) with \texttt{Binary8p4se} staging. Stage through a wide-range format (\texttt{Binary8p2se} above); the residual 0.109 here is the staging format{\textquotesingle}s own precision, which bounds what any downstream scheme can keep.

\end{tcolorbox}


\begin{tcolorbox}[toptitle=-1mm,bottomtitle=1mm,colback=admonition-note!50!white,colframe=admonition-note,title=\textbf{The scale format should be unsigned}]
Scales are non-negative by construction, so the scale format (\texttt{FS} above) should be unsigned with as much range as you can spare — \texttt{Binary8p1uf} is the canonical MX scale shape: precision 1, unsigned, extended for maximal dynamic range at minimal bits.

\end{tcolorbox}


\subsection{see also}



\label{9115667341246479961-dup9}{}


\href{howto\_\allowbreak{}choose\_\allowbreak{}format.md}{Choose a Format}, \href{howto\_\allowbreak{}quantize\_\allowbreak{}tensor.md}{Quantize a Tensor or Model}.



\section{Get Reproducible Stochastic Results}



\label{8071451965980986143}{}


Make stochastic-rounding code reproducible — across runs, and inside test suites — by fixing the entropy source.



\subsection{Ingredients}



\label{13776634711169142479-dup3}{}


\begin{itemize}
\item An explicit \texttt{AbstractRNG} (e.g. \texttt{Xoshiro(seed)}) passed as \texttt{rng}.


\item An explicit \texttt{R} (the raw stochastic draw) for tests that must be bit-exact.


\item \texttt{DefaultRNG!} if you want the implicit (no-\texttt{rng}-argument) forms to be reproducible too.


\item The \texttt{projection} keyword on scalar \texttt{rand}/\allowbreak\texttt{randn} forms.

\end{itemize}


\subsection{Recipe}



\label{111379208343275940-dup3}{}


\textbf{1. Prefer an explicit rng stream over the task-local default.} Passing any \texttt{AbstractRNG} first gives a stream independent of global state:




\nopagebreak[4]
\begin{minted}{jlcon}
julia> rand(Xoshiro(1), Binary8p4se)
Binary8p4se(0.0703125 ≡ 0x21)

julia> rand(Xoshiro(8), Binary8p4se, 4)
4-element Vector{Binary8p4se}:
 Binary8p4se(0.40625 ≡ 0x35)
 Binary8p4se(0.021484375 ≡ 0x13)
 Binary8p4se(0.75 ≡ 0x3c)
 Binary8p4se(0.28125 ≡ 0x31)
\end{minted}



If you rely on the implicit task-local generator instead, \texttt{Random.seed!} controls it exactly as for any other type:




\nopagebreak[4]
\begin{minted}{jlcon}
julia> Random.seed!(1234); rand(Binary8p4se)
Binary8p4se(0.3125 ≡ 0x32)

julia> Random.seed!(1234); rand(Binary8p4se)      # same seed, same draw
Binary8p4se(0.3125 ≡ 0x32)
\end{minted}



\textbf{2. A stochastic projection draws from the same rng as the value draw.} One underlying \texttt{Float64} draw, landed on the grid four different ways — the seed fixes the draw, the projection decides the landing:




\nopagebreak[4]
\begin{minted}{jlcon}
julia> rand(Xoshiro(2), Float64)                # the underlying uniform draw
0.0022505868897625403

julia> rand(Xoshiro(2), Binary8p4se)            # default: floor (stays in [0,1))
Binary8p4se(0.001953125 ≡ 0x02)

julia> rand(Xoshiro(2), Binary8p4se; projection = RTP_SN)   # ceiling
Binary8p4se(0.0029296875 ≡ 0x03)

julia> rand(Xoshiro(2), Binary8p4se; projection = RNE_SN)   # nearest
Binary8p4se(0.001953125 ≡ 0x02)

julia> rand(Xoshiro(2), Binary8p4se; projection = RSA_SN(8))  # stochastic
Binary8p4se(0.001953125 ≡ 0x02)
\end{minted}



Because the stochastic projection draws its random bits from the \emph{same} rng as the uniform draw, seeded streams stay reproducible end to end — you do not need a second seed for the rounding decision.



\textbf{3. For a test, fix the raw draw \texttt{R} instead of an rng.} This makes a single projection exactly reproducible without depending on an RNG{\textquotesingle}s internal algorithm at all — the right tool when you want to test both edges of a stochastic rounding decision deterministically:




\nopagebreak[4]
\begin{minted}{jlcon}
julia> σ = RSA_SN();                 # ≡ ProjSpec(StochasticA{8}(), SatNone())

julia> Add(Binary8p4se, σ, Binary8p4se(2.0), Binary8p4se(0.03125); rng = Xoshiro(1))
Binary8p4se(2.0 ≡ 0x48)

julia> Add(Binary8p4se, σ, Binary8p4se(2.0), Binary8p4se(0.03125); R = 0)
Binary8p4se(2.0 ≡ 0x48)      # smallest draw: rounds down

julia> Add(Binary8p4se, σ, Binary8p4se(2.0), Binary8p4se(0.03125); R = 255)
Binary8p4se(2.25 ≡ 0x49)     # largest draw: rounds up
\end{minted}



\textbf{4. Set \texttt{DefaultRNG!} if you want implicit calls to be reproducible too.} \texttt{DefaultRNG} is one of the six session defaults (initial value: the \texttt{Xoshiro} type); set it once at the top of a script or test file if you want every call that doesn{\textquotesingle}t name an rng to still be controlled by your seed.



\subsection{What can go wrong}



\label{7385351218440078181-dup3}{}


\begin{tcolorbox}[toptitle=-1mm,bottomtitle=1mm,colback=admonition-warning!50!white,colframe=admonition-warning,title=\textbf{Array and `!` forms always use the defaults}]
\texttt{rand(T, dims)}, \texttt{randn!(A)}, and friends always use the default projection (floor for \texttt{rand}, nearest + \texttt{SatFinite} for \texttt{randn}) — there is no \texttt{projection} keyword on the array forms. For arrays under another projection, draw scalars explicitly: \texttt{[rand(\allowbreak{}rng, T; projection = ρ) for \_\allowbreak{} in 1:n]}.

\end{tcolorbox}


\begin{tcolorbox}[toptitle=-1mm,bottomtitle=1mm,colback=admonition-note!50!white,colframe=admonition-note,title=\textbf{R must lie in `0:2{\textasciicircum}N-1`}]
For an \texttt{N}-bit stochastic mode, an out-of-range explicit \texttt{R} is a contract violation, not a silently wrapped value — pick \texttt{R} from the mode{\textquotesingle}s actual bit budget (\texttt{RSA\_\allowbreak{}SN(8)} takes \texttt{R ∈ 0:255}).

\end{tcolorbox}


\begin{tcolorbox}[toptitle=-1mm,bottomtitle=1mm,colback=admonition-warning!50!white,colframe=admonition-warning,title=\textbf{Nearest/ceiling can return exactly 1.0 from `rand`}]
The default floor projection is the contract-keeper for \texttt{rand}: it guarantees results stay in \texttt{[0, 1)}. Opting into \texttt{RTP\_\allowbreak{}SN} or \texttt{RNE\_\allowbreak{}SN} can return exactly \texttt{1.0} (mass near the top rounds up) — only opt out of the default if your code can tolerate that.

\end{tcolorbox}


\subsection{see also}



\label{9115667341246479961-dup10}{}


\href{howto\_\allowbreak{}choose\_\allowbreak{}format.md}{Choose a Format}, \href{howto\_\allowbreak{}quantize\_\allowbreak{}tensor.md}{Quantize a Tensor or Model}.



\section{Work with Packed Storage}



\label{11408754518297066540}{}


Store a vector of small-format values at its true bit width instead of one byte (or two) per element.



\subsection{Ingredients}



\label{13776634711169142479-dup4}{}


\begin{itemize}
\item \texttt{PackedVector(v)} built from an existing \texttt{Vector\{F\}}.


\item Standard indexing, \texttt{collect}, iteration — \texttt{PackedVector} is a full \texttt{AbstractVector\{F\}}.


\item \texttt{vmap} for computing directly on packed data (it unpacks tiles internally).


\item The sizing rule: \texttt{bitwidth(F)} bits per element, rounded up to whole storage words.

\end{itemize}


\subsection{Recipe}



\label{111379208343275940-dup4}{}


\textbf{1. Build a packed vector from ordinary values.} \texttt{PackedVector\{F\}} stores each code point in exactly \texttt{bitwidth(F)} bits:




\nopagebreak[4]
\begin{minted}{jlcon}
julia> v = Binary5p2se.(rand(Xoshiro(3), 6) .* 4);

julia> pv = PackedVector(v); sizeof(pv.data)
8                     # 6 × 5 bits = 30 bits → one 64-bit word

julia> pv[3] == v[3]
true
\end{minted}



Six values at 5 bits each is 30 bits, which rounds up to one 64-bit backing word — \texttt{sizeof(pv.data)} reports 8 bytes.



\textbf{2. Index, collect, and iterate like any other vector.} \texttt{PackedVector} is a full \texttt{AbstractVector\{F\}}:




\nopagebreak[4]
\begin{minted}{julia}
pv[2]
pv[2] = Binary5p2se(0.75)
collect(pv)
\end{minted}



\textbf{3. Compute with \texttt{vmap} directly on the packed vector.} \texttt{vmap} accepts a \texttt{PackedVector} argument directly, unpacking cache-friendly tiles internally — you do not need to \texttt{collect} first:




\nopagebreak[4]
\begin{minted}{julia}
out = vmap(:Exp, Binary5p2se, ρ, pv)
\end{minted}



\textbf{4. Size larger vectors the same way.} The packing ratio is exactly \texttt{bitwidth(F) / 8} against a plain byte-per-element vector (for \texttt{K ≤ 8} formats). At a million elements of a 5-bit format:




\nopagebreak[4]
\begin{minipage}{\linewidth}
\begin{minted}{julia}
n = 1_000_000
model = [rawvalue(Binary5p2se, UInt8(rand(0:31))) for _ in 1:n]
pv = PackedVector(model)
(sizeof(model), sizeof(pv.data))
\end{minted}




\begin{minted}{text}
(1000000, 625000)   # 8 bits → 5 bits per value; indexing and vmap work directly on pv
\end{minted}
\end{minipage}




1,000,000 bytes of plain storage becomes 625,000 bytes packed — \texttt{n × K / 8} bytes, in general, rounded to whole words.



\subsection{What can go wrong}



\label{7385351218440078181-dup4}{}


\begin{tcolorbox}[toptitle=-1mm,bottomtitle=1mm,colback=admonition-warning!50!white,colframe=admonition-warning,title=\textbf{There is deliberately no in-place packed arithmetic}]
The rule is \emph{store packed, compute unpacked}. Don{\textquotesingle}t reach for an in-place packed \texttt{+=}-style operation — it doesn{\textquotesingle}t exist. Compute with \texttt{vmap} (which unpacks internally and returns a fresh result) or unpack with \texttt{collect} first, compute, then re-pack with \texttt{PackedVector} if you want to keep the result compressed at rest.

\end{tcolorbox}


\begin{tcolorbox}[toptitle=-1mm,bottomtitle=1mm,colback=admonition-default!50!white,colframe=admonition-default,title=\textbf{Packing helps storage bandwidth, not per-element compute speed}]
Reach for \texttt{PackedVector} when memory footprint or bandwidth is the bottleneck (e.g. a large quantized model kept resident). If you are only computing, not storing, an ordinary \texttt{Vector\{F\}} (one byte per element at \texttt{K ≤ 8}) is simpler and does not pay any unpacking cost.

\end{tcolorbox}


\subsection{see also}



\label{9115667341246479961-dup11}{}


\href{howto\_\allowbreak{}quantize\_\allowbreak{}tensor.md}{Quantize a Tensor or Model}, \href{howto\_\allowbreak{}choose\_\allowbreak{}format.md}{Choose a Format}.



\section{Register \& Measure an Approximation (κ)}



\label{2701220102296233845}{}


Register a faster, inexact kernel in place of the bit-exact default, with an honest, measured error bound.



\subsection{Ingredients}



\label{13776634711169142479-dup5}{}


\begin{itemize}
\item \texttt{measure\_\allowbreak{}kappa(\allowbreak{}fn, opname, resultformat, operandformats, ρ)} to measure the worst-case code-point deviation exhaustively.


\item \texttt{register\_\allowbreak{}approx!(\allowbreak{}name, opname, resultformat, operandformats, ρ, fn; κ)} to register the implementation under a declared bound.


\item \texttt{approx}, \texttt{list\_\allowbreak{}approx}, \texttt{kappa}, \texttt{kappa\_\allowbreak{}measured} to retrieve and inspect registered approximations.


\item \texttt{unregister\_\allowbreak{}approx!} to remove one.

\end{itemize}


\subsection{Recipe}



\label{111379208343275940-dup5}{}


\textbf{1. Measure κ before declaring anything.} \texttt{measure\_\allowbreak{}kappa} runs your approximate function against the defined operation over every input and returns the worst code-point distance, exhaustively verified:




\nopagebreak[4]
\begin{minipage}{\linewidth}
\begin{minted}{julia}
step2(x) = (r = Exp(Binary8p4se, RNE_SN, x);
            isfinite(decode(r)) ? NextGreaterThan(NextGreaterThan(r)) : r)

measure_kappa(step2, :Exp, Binary8p4se, (Binary8p4se,), RNE_SN)
\end{minted}




\begin{minted}{text}
(2.0, true)                # κ = 2 code points, verified over all 256 inputs
\end{minted}
\end{minipage}




\textbf{2. Register with the measured (or a larger) κ.} \texttt{register\_\allowbreak{}approx!} takes the same signature plus the implementation and a declared \texttt{κ}:




\nopagebreak[4]
\begin{minipage}{\linewidth}
\begin{minted}{julia}
register_approx!(:cheater, :Exp, Binary8p4se, (Binary8p4se,), RNE_SN, step2; κ=1)
\end{minted}




\begin{minted}{text}
ERROR: ArgumentError: declared κ = 1.0 understates measured κ = 2.0 — registration rejected
\end{minted}
\end{minipage}




Declaring \texttt{κ = 1} when the measured deviation is 2 is rejected outright — the registry measures your honesty at registration time; you cannot understate the bound.



\textbf{3. Register with a truthful κ and retrieve it.} Using \texttt{κ=1} failed above; declaring the true value (or higher) succeeds. The same shape applied to a hardware-style approximation (\texttt{hardtanh} standing in for \texttt{Tanh}) shows the full round trip:




\nopagebreak[4]
\begin{minipage}{\linewidth}
\begin{minted}{julia}
one4 = Binary8p4se(1.0)
hardtanh(x) = Clamp(Binary8p4se, RNE_SN, x, Negate(one4), one4)

κ, exhaustive = measure_kappa(hardtanh, :Tanh, Binary8p4se, (Binary8p4se,), RNE_SN)
(κ, exhaustive)
\end{minted}




\begin{minted}{text}
(4.0, true)                # worst deviation: 4 code points, verified on all 256 inputs
\end{minted}




\begin{minted}{julia}
impl = register_approx!(:hardtanh_act, :Tanh, Binary8p4se, (Binary8p4se,),
                        RNE_SN, hardtanh; κ=4)
(kappa(:hardtanh_act), kappa_measured(impl), :hardtanh_act in list_approx())
\end{minted}




\begin{minted}{text}
(4.0, 4.0, true)           # declared = measured; conformance_report() now lists it
\end{minted}
\end{minipage}




\textbf{4. Remove a registration when you{\textquotesingle}re done with it.}




\begin{minted}{julia}
unregister_approx!(:hardtanh_act)
\end{minted}



\subsection{What can go wrong}



\label{7385351218440078181-dup5}{}


\begin{tcolorbox}[toptitle=-1mm,bottomtitle=1mm,colback=admonition-warning!50!white,colframe=admonition-warning,title=\textbf{Understating κ is not silently clamped — it throws}]
You cannot register \texttt{:cheater} with \texttt{κ=1} {\textquotedbl}just to see{\textquotedbl} and get a warning; \texttt{register\_\allowbreak{}approx!} throws \texttt{ArgumentError} and refuses the registration entirely. There is no partial-registration state to clean up.

\end{tcolorbox}


\begin{tcolorbox}[toptitle=-1mm,bottomtitle=1mm,colback=admonition-note!50!white,colframe=admonition-note,title=\textbf{NaN-mismatching implementations need κ = NaN, not a number}]
If your approximate kernel disagrees with the defined operation on whether a NaN is produced, you must acknowledge this explicitly by declaring \texttt{κ = NaN} — a numeric κ asserts NaN behavior matches.

\end{tcolorbox}


\begin{tcolorbox}[toptitle=-1mm,bottomtitle=1mm,colback=admonition-note!50!white,colframe=admonition-note,title=\textbf{Approximate implementations never enter the default API}]
Registering \texttt{:my\_\allowbreak{}fast\_\allowbreak{}exp} or \texttt{:hardtanh\_\allowbreak{}act} does not change what \texttt{Exp} or \texttt{Tanh} return by default — the bit-exact path is always the one \texttt{Exp}, \texttt{+}, and friends use. Retrieve an approximation explicitly with \texttt{approx(:name)} and call it yourself; \texttt{list\_\allowbreak{}approx()} enumerates what{\textquotesingle}s registered, and \texttt{conformance\_\allowbreak{}report()} lists them alongside the exact catalog.

\end{tcolorbox}


\subsection{see also}



\label{9115667341246479961-dup12}{}


\href{howto\_\allowbreak{}conformance.md}{Read \& Export the Conformance Declaration}.



\section{Convert To/From Float16 \& BFloat16 Safely}



\label{6879013464833569958}{}


Move data between SmallFloats formats and IEEE \texttt{Float16}/\allowbreak\texttt{BFloat16} without silently changing its value by a factor of two.



\subsection{Ingredients}



\label{13776634711169142479-dup6}{}


\begin{itemize}
\item \texttt{Convert(T, ρ, x)} — the only safe path between formats with different exponent bias.


\item \texttt{datumsexact} to check whether the two datum sets permit an exact round trip.


\item Never \texttt{reinterpret} across formats with different bias — see the warning below.

\end{itemize}


\needspace{22\baselineskip}
\subsection{Recipe}



\label{111379208343275940-dup6}{}


\textbf{1. Understand the mismatch before you touch either format.} The K ≤ 16 grid contains two formats that look like \texttt{Float16} and \texttt{BFloat16} but are not — the difference is the kind that produces correct-looking numbers rather than errors:




\begin{table}[h]
\centering
\begin{tabulary}{\linewidth}{L L L L L}
\toprule
 & \texttt{Binary16p11se} & \texttt{Float16} & \texttt{Binary16p8se} & \texttt{BFloat16} \\
\toprule
precision \texttt{P} & 11 & 11 & 8 & 8 \\
exponent bias & \textbf{16} & \textbf{15} & \textbf{128} & \textbf{127} \\
largest finite & \textbf{65472} & \textbf{65504} & \textbf{3.3762e38} & \textbf{3.3895e38} \\
NaN encodings & \textbf{1} & \textbf{2046} & \textbf{1} & \textbf{2046} \\
negative zero & \textbf{no} & yes & \textbf{no} & yes \\
domain is a parameter & \textbf{yes} (\texttt{se}/\allowbreak\texttt{sf}) & no & \textbf{yes} & no \\
\bottomrule
\end{tabulary}

\end{table}



P3109 defines the bias so the exponent range is symmetric about the format{\textquotesingle}s midpoint; IEEE-754 defines it as \texttt{2{\textasciicircum}(E-1) − 1}. The two differ by one, so \textbf{every datum in the format is a factor of two away from its IEEE counterpart} — \texttt{Binary16p11se} and \texttt{Float16} share a significand layout and disagree on the value of every code point.



\textbf{2. Always convert — never reinterpret.} \texttt{Convert} is a real conversion: \texttt{Float16(1.5)} becomes \texttt{Binary16p11se(1.5)}, unchanged in value:




\nopagebreak[4]
\begin{minted}{jlcon}
julia> using SmallFloats.Formats

julia> Convert(Binary16p11se, RNE_SF, Float16(1.5))    # a conversion: 1.5 ↦ 1.5
Binary16p11se(1.5 ≡ 0x4200)

julia> reinterpret(Binary16p11se, Float16(1.5))        # the SAME BITS: 1.5 ↦ 0.75
Binary16p11se(0.75 ≡ 0x3e00)
\end{minted}



\texttt{Float16(1.5)} has the bit pattern \texttt{0x3e00}; read as a \texttt{Binary16p11se}, those same bits denote \textbf{0.75}, because the exponent bias differs by one. Nothing errors, nothing warns, and the answer is off by a factor of two.



\textbf{3. Check \texttt{datumsexact} when you need to know if the round trip is lossless.} The package converts to and from both \texttt{Float16} and \texttt{BFloat16} exactly, wherever the datum sets permit — \texttt{datumsexact} tells you where that holds so you don{\textquotesingle}t have to assume.



\textbf{4. Remember the other differences are draft decisions, not bugs.} A single NaN encoding means there is no payload to propagate and no quiet/signalling distinction — the 2045 code points IEEE spends on NaN payloads are ordinary datums here. No negative zero means \texttt{-0.0} and \texttt{0.0} are the same code point. The domain parameter (\texttt{e} for extended with ±Inf, \texttt{f} for finite) has no IEEE analogue: \texttt{Binary16p11sf} spends the two infinity code points on finite values instead. None of these need {\textquotedbl}fixing{\textquotedbl} before converting — \texttt{Convert} already accounts for all of them.



\subsection{What can go wrong}



\label{7385351218440078181-dup6}{}


\begin{tcolorbox}[toptitle=-1mm,bottomtitle=1mm,colback=admonition-warning!50!white,colframe=admonition-warning,title=\textbf{`reinterpret` silently returns the wrong value, off by 2×}]
\texttt{reinterpret(\allowbreak{}Binary16p11se, Float16(\allowbreak{}1.5))} gives \texttt{0.75}, not \texttt{1.5} — the same bit pattern, read under a different exponent bias. This is the single most dangerous operation in cross-format interop: it never errors, so a pipeline that accidentally reinterprets instead of converting produces plausible-looking numbers that are systematically wrong by a factor of two. Always use \texttt{Convert}, never \texttt{reinterpret}, when moving between a SmallFloats format and \texttt{Float16}/\allowbreak\texttt{BFloat16}.

\end{tcolorbox}


\begin{tcolorbox}[toptitle=-1mm,bottomtitle=1mm,colback=admonition-note!50!white,colframe=admonition-note,title=\textbf{If you want IEEE semantics, use Float16/BFloat16 directly}]
\texttt{Binary16p11se} and \texttt{Binary16p8se} are not drop-in replacements for \texttt{Float16}/\allowbreak\texttt{BFloat16} — they exist for the P3109 draft{\textquotesingle}s own reasons (symmetric exponent range, single NaN, no negative zero, parametric domain). If your downstream consumer specifically needs IEEE bit patterns, keep the data in \texttt{Float16}/\allowbreak\texttt{BFloat16} and convert only at the SmallFloats boundary.

\end{tcolorbox}


\subsection{see also}



\label{9115667341246479961-dup13}{}


\href{howto\_\allowbreak{}choose\_\allowbreak{}format.md}{Choose a Format}, \href{howto\_\allowbreak{}conformance.md}{Read \& Export the Conformance Declaration}.



\section{Read \& Export the Conformance Declaration}



\label{2397441964831923039}{}


Produce a machine-readable statement of exactly which formats, operations, and approximations a session used, and attach it to an experiment{\textquotesingle}s artifacts.



\subsection{Ingredients}



\label{13776634711169142479-dup7}{}


\begin{itemize}
\item \texttt{conformance()} — the live declaration as a Julia value.


\item \texttt{conformance\_\allowbreak{}report()} — a printed, human-readable form.


\item \texttt{conformance\_\allowbreak{}dict()} — a plain nested \texttt{Dict\{String,Any\}} for JSON/TOML serialization.


\item \texttt{draft\_\allowbreak{}revision()} — the P3109 draft revision the package implements.


\item Any JSON/TOML writer your project already depends on.

\end{itemize}


\subsection{Recipe}



\label{111379208343275940-dup7}{}


\textbf{1. Inspect the live declaration.} \texttt{conformance()} returns the formats, operations, mode vocabulary, the table specializations actually instantiated this session, and every registered approximation. \texttt{conformance\_\allowbreak{}report()} prints the same information for a human to read at the REPL.



\textbf{2. Export it as a plain \texttt{Dict} for serialization.}




\begin{minted}{julia}
d = conformance_dict()          # plain nested Dict{String,Any}
(d["package"], length(d["formats"]), length(d["operations"]),
 [s["op"] for s in d["cached_specializations"]])
\end{minted}



Serialize \texttt{d} with any JSON/TOML writer to attach a machine-readable conformance statement to experiment artifacts — for example, alongside a saved model checkpoint or a benchmark result, so a later reader can confirm exactly which formats and approximations produced it without re-running the session.



\textbf{3. Record the draft revision alongside it.} \texttt{draft\_\allowbreak{}revision()} reports which revision of the P3109 draft this package implements — include it next to the conformance dict so an artifact remains interpretable even if the draft itself changes in a later package version.



\textbf{4. Confirm registered approximations show up.} Any approximation registered with \texttt{register\_\allowbreak{}approx!} appears in the conformance declaration automatically — \texttt{conformance\_\allowbreak{}report()} lists it alongside the exact operation catalog, so the declaration always reflects the true provenance of a session{\textquotesingle}s results, exact and approximate alike.



\subsection{What can go wrong}



\label{7385351218440078181-dup7}{}


\begin{tcolorbox}[toptitle=-1mm,bottomtitle=1mm,colback=admonition-note!50!white,colframe=admonition-note,title=\textbf{The declaration is live, not static}]
\texttt{conformance()} and \texttt{conformance\_\allowbreak{}dict()} reflect the \emph{current} session — including which table specializations have actually been built so far. Call them at the point you want to snapshot (typically right before saving results), not at the top of a script, or the declaration may under-report what the run actually exercised.

\end{tcolorbox}


\begin{tcolorbox}[toptitle=-1mm,bottomtitle=1mm,colback=admonition-warning!50!white,colframe=admonition-warning,title=\textbf{A conformance dict does not travel with unregistered approximations}]
If you registered an approximation in one session and load results from it in a fresh session, \texttt{conformance\_\allowbreak{}dict()} in the new session will not know about that approximation unless you re-register it. Save the exported dict itself with the artifact — do not rely on regenerating it later from a differently configured session.

\end{tcolorbox}


\begin{tcolorbox}[toptitle=-1mm,bottomtitle=1mm,colback=admonition-note!50!white,colframe=admonition-note,title=\textbf{JSON/TOML serialization is your responsibility, not the package{\textquotesingle}s}]
\texttt{conformance\_\allowbreak{}dict()} returns a plain nested \texttt{Dict}; the package does not ship a JSON or TOML writer. Use whatever serialization library your project already depends on — the dict is deliberately unopinionated about the output format.

\end{tcolorbox}


\subsection{see also}



\label{9115667341246479961-dup14}{}


\href{howto\_\allowbreak{}register\_\allowbreak{}approx.md}{Register \& Measure an Approximation (κ)}, \href{howto\_\allowbreak{}choose\_\allowbreak{}format.md}{Choose a Format}.



\chapter{Explanation}


\section{Formats, Aliases \& Representations}



\label{3847673756616657637}{}


Every value in SmallFloats.jl belongs to a \emph{format}, and a format is a Julia type. This page explains what that type actually is, why the package draws a hard line between the abstract idea of a format and the concrete type your code uses, and why that line matters for both correctness and speed.



\subsection{The parametric type}



\label{13610640477998426652}{}


Every format is an instance of the same four-parameter type:




\nopagebreak[4]
\begin{minted}{julia}
Binary{K,P,SGN,EXT}
\end{minted}



\texttt{K} is the bitwidth, \texttt{3 ≤ K ≤ 16}. \texttt{P} is the precision — significand bits, including the implicit bit. \texttt{SGN} is a \texttt{Bool}: signed or unsigned. \texttt{EXT} is a \texttt{Bool}: extended (the domain includes ±Inf) or finite (it does not). Every legal combination of these four parameters is a distinct format, and the package ships all of them: \textbf{504 formats} in total.



That is the whole parameter space. There is no fifth axis, and — this is the detail that trips people up when they start comparing formats — \texttt{K} is not an independent budget. It is spent entirely by the other three: \texttt{K = P + E + SGN} where \texttt{E} is the exponent width implied by \texttt{K}, \texttt{P}, and \texttt{SGN}. Raising precision by one bit necessarily lowers exponent range, because there is nowhere else in the word for the extra bit to come from.



\subsection{\texttt{<:}, not \texttt{===}}



\label{3394570003399755675}{}


\texttt{Binary\{K,P,SGN,EXT\}} is an abstract type. It is not what a value{\textquotesingle}s type actually is, and it is not what you should write as an array element type or a constructor target. The concrete type is a \emph{representation} — a draft-named alias like \texttt{Binary8p4se} — related to the abstract format by subtyping, not equality:




\nopagebreak[4]
\begin{minted}{jlcon}
julia> Binary8p4se <: Binary{8,4,true,true}      # `<:`, not `===` — see below
true

julia> formatname(Binary{6,3,true,false})
:Binary6p3sf

julia> SmallFloats.Binary16p6se                  # always reachable, unexported
Binary16p6se

julia> format(16, 6, true, true)                 # programmatic, runtime params
Binary16p6se

julia> using SmallFloats.Formats                 # …or bring all 504 into scope

julia> Binary16p6se
Binary16p6se
\end{minted}



\texttt{Binary8p4se} is a concrete subtype of \texttt{Binary\{8,4,true,true\}}, but the two type objects are distinct. The distinction is not an implementation accident; it comes from the draft{\textquotesingle}s own vocabulary. A \emph{format} is a datum set and an encoding — an abstract description of which values exist and how they are named. A \emph{representation} is a specific choice of machine word that carries a code point: how many bits, what storage unit, what bit layout. That choice sits below the level of description the standard cares about. \texttt{Binary} names the format; \texttt{Binary8p4se} is the representation you actually compute with.



The suffix on an alias reads directly: \texttt{p4} is precision 4, \texttt{s}/\allowbreak\texttt{u} is signed/unsigned, \texttt{e}/\allowbreak\texttt{f} is extended/finite. \texttt{Binary8p4se} is 8-bit, precision 4, signed, extended — read the name and you have the four parameters without consulting a table.



\subsection{Why 504 formats, and why only 120 are exported}



\label{9834957682784680318}{}


Every legal \texttt{(K, P, SGN, EXT)} combination at \texttt{K ∈ 3:16} gets an alias, because the whole point of the package is that formats are cheap to have — each one is 256 (or fewer) code points and a handful of type-parameter queries, not a hand-written implementation. But 504 identifiers is more than any one script needs in scope at once, so \texttt{using SmallFloats} exports only the \textbf{120 aliases at K ≤ 8} — the sizes people actually reach for most: sub-byte and byte-sized formats for quantization, embeddings, and packed storage.



The other 384 (\texttt{K} from 9 to 16) are one opt-in away. You never need to define anything to reach them — they already exist as bindings, just not exported ones:



\begin{itemize}
\item \texttt{SmallFloats.Binary16p6se} — fully qualified, always available regardless of what{\textquotesingle}s exported.


\item \texttt{using SmallFloats.Formats} — brings every one of the 504 aliases into scope at once.


\item \texttt{format(K, P, SGN, EXT)} — the programmatic route, for when the parameters are runtime values rather than something you can spell as a literal type name.

\end{itemize}


\subsection{Format introspection}



\label{7504779454545473321}{}


Every format answers a fixed set of queries about its own structure, drawn from the draft{\textquotesingle}s Group M, in both a Julia-flavored spelling and the draft-named form:




\nopagebreak[4]
\begin{minted}{jlcon}
julia> bitwidth(Binary8p4se), precision(Binary8p4se), expbias(Binary8p4se)
(8, 4, 8)

julia> MaxFiniteOf(Binary8p4se)
Binary8p4se(224.0 ≡ 0x7e)

julia> MinPositiveOf(Binary8p4se)
Binary8p4se(0.0009765625 ≡ 0x01)
\end{minted}



\texttt{MinFiniteOf}, \texttt{MinNormalOf}, \texttt{MaxSubnormalOf}, \texttt{expbitwidth}, \texttt{trailingsigbits}, \texttt{issigned}, \texttt{isextended}, and their draft-named counterparts (\texttt{BitwidthOf}, \texttt{PrecisionOf}, \texttt{ExponentBiasOf}, …) round out the set.



The important property of every one of these queries is that it is a pure function of the type parameters — nothing about a particular value changes the answer — so it works identically whether you hand it a type or a value:




\nopagebreak[4]
\begin{minted}{jlcon}
julia> x = Binary8p4se(1.6);

julia> BitwidthOf(x), PrecisionOf(x), SignednessOf(x), DomainOf(x)
(8, 4, true, true)

julia> MaxFiniteOf(x)
Binary8p4se(224.0 ≡ 0x7e)
\end{minted}



\texttt{bitwidth(x)} and \texttt{bitwidth(typeof(x))} are the same query, and both fold to a compile-time literal — there is no runtime cost to asking a value what kind of format it lives in.



\begin{tcolorbox}[toptitle=-1mm,bottomtitle=1mm,colback=admonition-default!50!white,colframe=admonition-default,title=\textbf{The abstract format as an array element type}]
Because \texttt{Binary\{K,P,Σ,Δ\}} is abstract, declaring an array with it as the element type gives you a non-\texttt{isbits} element type, and every element is boxed:


\nopagebreak[4]
\begin{minted}{jlcon}
julia> isbitstype(eltype(Vector{Binary{8,4,true,true}}(undef, 3)))
false

julia> isbitstype(eltype(Vector{Binary8p4se}(undef, 3)))
true
\end{minted}

Nothing about this errors. Every operation still returns the right answer, because the package normalizes the abstract format at every constructor and entry point it can see. What you silently lose is the entire performance story: table gathers, zero-allocation scalar calls, specialized dispatch — all of it depends on the element type being concrete. This is the sharper edge of the type distinction above: comparing the concrete alias with the abstract format makes their different roles visible, while the abstract-array mistake costs you real performance and says nothing at all.

\texttt{similar(a, T)} normalizes the request for you, but \texttt{Vector\{T\}(undef, n)} cannot be intercepted at the call site. Always spell the alias, or use \texttt{format(K, P, Σ, Δ)} when the parameters are runtime values.

\end{tcolorbox}


\subsection{Go deeper}



\label{4168117981347147797}{}


\begin{quote}
Continue with \textbf{Values, Code Points \& Conversion} (Insights track) for how a \texttt{Binary} value stores a code point, how \texttt{decode} recovers the exact datum, and the four ways to construct one.

\end{quote}


\section{The Projection Contract}



\label{15998783979828160471}{}


Every number that enters a \texttt{Binary} format goes through the same pipeline, regardless of whether it arrives via a constructor, an arithmetic operator, or an explicit conversion. This page explains what that pipeline promises, what {\textquotedbl}exact{\textquotedbl} means at each stage, and the handful of cases where the contract refuses rather than bends.



\subsection{One write path}



\label{6607324653425258239}{}


There is exactly one way a code point gets produced in SmallFloats.jl:




\nopagebreak[4]
\begin{minted}{text}
exact result  →  RoundToPrecision  →  Saturate  →  code point
\end{minted}



The operation is computed \emph{exactly} first — as exact real arithmetic, not as a simulation of some intermediate low-precision format — and only then is the result rounded to the target precision and, if it falls outside the representable range, saturated according to the chosen saturation mode. This holds for a scalar \texttt{Add}, for a fused \texttt{FMA} (one rounding, not two), and for a full \texttt{BlockDotProduct} (every lane product and the accumulation exact, projected once at the end). There is no second path that skips a step or rounds twice; every code point the package ever produces came out of this same pipeline.



The consequence worth sitting with: there is no accumulated rounding error \emph{inside} an operation. Rounding error exists only at the single boundary where the exact result crosses into a code point — never before it.



\subsection{What {\textquotedbl}exact{\textquotedbl} means, per carrier}



\label{9135858451050429722}{}


{\textquotedbl}Exact math{\textquotedbl} has to happen somewhere concrete, and the carrier it happens on depends on the format. \texttt{decode(x)} always returns the format{\textquotesingle}s exact datum, never a rounded approximation of it — but the \emph{type} it returns is a property of the format, not a package-wide constant.



For 432 of the 504 formats, the datum{\textquotesingle}s exponent range fits inside a \texttt{Float64}, so \texttt{decode} returns a \texttt{Float64} and the exact arithmetic runs there. For the other 72 — formats whose exponent range or precision a \texttt{Float64} cannot hold without rounding — \texttt{decode} returns \texttt{Float128} or an exact dyadic carrier instead, so no precision is lost representing the datum in the first place. \texttt{datumcarrier(T)} names which carrier a given format uses, and it is this carrier, not \texttt{Float64} universally, that the exact-math half of the pipeline runs on.



This distinction matters because \texttt{Float64(x)} always returns a \texttt{Float64} regardless of the format{\textquotesingle}s own carrier, and rounds wherever the datum does not fit — it is a display/interop convenience, not the internal computation path. At the user level, \texttt{promotecarrier(T)} names the \emph{promotion} carrier — \texttt{Float64} for the 432 formats at the ordinary rung, \texttt{BigFloat} for the other 72 — which is the type an operation between a \texttt{Binary} value and an ordinary number lands in.



\subsection{Symbolic sticky: projecting {\textquotedbl}just below{\textquotedbl} a number}



\label{1121493031502271068}{}


Some values are not just numbers — they are the limit of a sequence approaching a number from one side, and a directed rounding mode needs to know which side to land the projection on. \texttt{tanh} approaching 1 as its argument grows is the canonical case: the exact mathematical limit is 1, but \texttt{tanh(x) < 1} for every finite \texttt{x}, so a correct projection under \texttt{TowardNegative} must round \emph{down from} 1, not treat 1 as an ordinary carrier value sitting exactly on a code point.



The engine handles this with a \texttt{sticky} argument, \texttt{sticky ∈ \{-1, 0, +1\}}, meaning {\textquotedbl}the true value is the carrier value plus an infinitesimal of this sign.{\textquotedbl} This lets \texttt{project} land exactly on the correct neighbor for enclosure endpoints and asymptotes without inventing a fake carrier value that isn{\textquotesingle}t quite 1:




\nopagebreak[4]
\begin{minipage}{\linewidth}
\begin{minted}{julia}
using SmallFloats: project
project(Binary8p4se, ProjSpec(TowardNegative(), SatNone()), 1.0; sticky=-1)
\end{minted}




\begin{minted}{text}
Binary8p4se(0.9375 ≡ 0x3f)     # == NextLessThan(1.0): the engine crossed the binade
\end{minted}
\end{minipage}




Read this carefully: the carrier value passed in is \texttt{1.0}, exactly representable in \texttt{Binary8p4se}. Without \texttt{sticky}, a directed-toward-negative projection of an exactly-representable value is the identity — it would return \texttt{1.0} unchanged. \texttt{sticky=-1} tells the engine that the true value sits an infinitesimal \emph{below} 1.0, so \texttt{TowardNegative} must move to the next lower code point instead of stopping at 1.0 itself — and it does, crossing the binade boundary down to \texttt{0.9375}, which is exactly \texttt{NextLessThan(1.0)}. This is how the package gets asymptotic behavior like \texttt{tanh → 1⁻} exactly right under every rounding mode, rather than approximately right under most of them.



\subsection{Why \texttt{Rational} inputs are refused, not double-rounded}



\label{16856155639614633457}{}


\texttt{Convert} accepts \texttt{Binary} values of any format, \texttt{Float16}/\allowbreak\texttt{Float32}/\texttt{Float64}, \texttt{Float128}, any \texttt{Integer}, and \texttt{BigFloat} — every type whose values the package can project \emph{exactly} onto the carrier it needs. \texttt{Float128} in particular projects directly, preserving all 113 significand bits, which matters when a value sits just above a rounding midpoint: staging through \texttt{Float64} first would round it onto the midpoint and then break the tie the wrong way, a genuine double rounding that changes the answer.



\texttt{Rational} is conspicuously missing from that list, and that is deliberate rather than an oversight. A \texttt{Rational\{Int\}} can represent a value that no supported carrier holds exactly (a denominator that is not a power of two, for instance), so converting it would require rounding once to reach a carrier and again to reach the format — silently, with no record of it having happened. That would put a second, invisible rounding step inside what looks like a single conversion, which is exactly the guarantee the one-write-path contract exists to prevent.



So \texttt{Convert} refuses \texttt{Rational} inputs outright, with an error that tells you what to do about it: convert explicitly to a supported carrier first (\texttt{Binary8p4se(Float64(π))} for irrational values, or \texttt{Binary8p4se(Float64(r))} for a \texttt{Rational} \texttt{r}) and thereby own the double rounding as a choice you made, not a default the package made for you.



\subsection{Why the contract is a single, named function}



\label{5104813108012279296}{}


All of this — exact math, rounding, saturation, and the sticky adjustment for asymptotic limits — lives behind one function, \texttt{project}, rather than being reimplemented at each call site that needs a code point. That centralization is what makes the {\textquotedbl}one write path{\textquotedbl} claim checkable rather than aspirational: every operation in the package, from a scalar \texttt{Add} to a table build to a block reduction, is required to produce its code point by calling through \texttt{project}, and the test suite is able to assert this because there is exactly one place the assertion needs to look. A second hand-rolled rounding step anywhere in the codebase would be a bug by definition, not a matter of style.



This is also why the table-gather performance strategy described elsewhere in the documentation costs nothing in correctness. A cached table entry and a fresh scalar computation are guaranteed identical precisely because both routes go through \texttt{project} — the table is not a separate, hopefully-equivalent fast path, it is the \emph{same} path, memoized. Nothing about caching results for speed introduces a second contract to keep in sync with the first.



The same reasoning extends to approximate kernels. When a registered approximation replaces the default path for one operation, it is measured against exactly this contract{\textquotesingle}s output — the exact-math-then-project result — rather than against some looser notion of {\textquotedbl}close enough.{\textquotedbl} The projection contract is not just how the default path stays correct; it is also the fixed reference every deliberate approximation in the package is honestly scored against.



That is what makes a declared κ a fact about the \emph{difference} between two functions, rather than a vague impression of {\textquotedbl}close enough{\textquotedbl} — both sides of the comparison ultimately trace back to the same contract this page describes.



\subsection{Go deeper}



\label{4168117981347147797-dup1}{}


\begin{quote}
Continue with \textbf{Projection: Rounding \& Saturation} (Insights track) for the hands-on experiments — every rounding mode and saturation mode combination worked through on real overflow and midpoint cases.

\end{quote}


\section{Session Defaults \& Their Scope}



\label{493033648811986665}{}


Six values control what {\textquotedbl}the default{\textquotedbl} means for any convenience call in SmallFloats.jl — the format \texttt{T(2.1)} produces, the projection \texttt{x + y} follows, the RNG \texttt{rand(T)} draws from. This page explains what those six defaults are, how they stay consistent with each other, why they are dangerous to set casually, and the safe pattern for consuming them.



\needspace{23\baselineskip}
\subsection{The six defaults}



\label{6312792895918690226}{}


Six session-wide defaults are readable as \texttt{DefaultX()} and settable as \texttt{DefaultX!(v)}:




\begin{table}[h]
\centering
\begin{tabulary}{\linewidth}{L L L}
\toprule
default & initial value & setter accepts \\
\toprule
\texttt{DefaultType} & \texttt{Binary8p2se} & any fully-parameterized \texttt{Binary} type \\
\texttt{DefaultReturnType} & \texttt{Binary8p2se} & any fully-parameterized \texttt{Binary} type \\
\texttt{DefaultRoundingMode} & \texttt{NearestTiesToEven()} & mode instance or type \\
\texttt{DefaultSaturationMode} & \texttt{SatNone()} & mode instance or type \\
\texttt{DefaultProjection} & \texttt{RNE\_\allowbreak{}SN} & a \texttt{ProjSpec}, or \texttt{(mode, sat)} \\
\texttt{DefaultRNG} & the \texttt{Xoshiro} type & RNG type or instance \\
\texttt{DefaultRbits} & \texttt{8} & \texttt{Int} in \texttt{1:60} \\
\bottomrule
\end{tabulary}

\end{table}



\subsection{Coherence in both directions}



\label{18340556388588567673}{}


\texttt{DefaultProjection} and its two component defaults, \texttt{DefaultRoundingMode} and \texttt{DefaultSaturationMode}, are kept coherent no matter which one you change. Setting \texttt{DefaultRoundingMode!} or \texttt{DefaultSaturationMode!} rebuilds \texttt{DefaultProjection} around the new component; setting \texttt{DefaultProjection!} directly decomposes it back into both components. The invariant \texttt{DefaultProjection(\allowbreak{}) === ProjSpec(\allowbreak{}DefaultRoundingMode(\allowbreak{}), DefaultSaturationMode(\allowbreak{}))} holds always, regardless of which of the three setters you used last:




\nopagebreak[4]
\begin{minted}{jlcon}
julia> DefaultRoundingMode!(TowardZero())
TowardZero()

julia> DefaultProjection()               # followed the component
(TowardZero, SatNone)

julia> DefaultProjection!(RNA_SF)
(NearestTiesToAway, SatFinite)

julia> DefaultRoundingMode(), DefaultSaturationMode()   # followed the projection
(NearestTiesToAway(), SatFinite())

julia> DefaultProjection!(RNE_SN)        # restore: these setters are PROCESS-wide
(NearestTiesToEven, SatNone)
\end{minted}



\subsection{The defaults are process-global}



\label{5299094382095427161}{}


\begin{tcolorbox}[toptitle=-1mm,bottomtitle=1mm,colback=admonition-warning!50!white,colframe=admonition-warning,title=\textbf{The session defaults are process-global and they stay set}]
\texttt{DefaultProjection!} and its siblings mutate state shared by every module in the process. Nothing scopes them for you: set one in a script, a notebook cell, or a REPL session and \emph{every} later \texttt{T(x::Real)}, \texttt{a + b}, and \texttt{Exp(x)} follows it until something sets it back.

Prefer \texttt{with\_\allowbreak{}default\_\allowbreak{}projection} (and \texttt{with\_\allowbreak{}default\_\allowbreak{}type}, \texttt{with\_\allowbreak{}default\_\allowbreak{}returntype}), which restore on exit even if the body throws:


\nopagebreak[4]
\begin{minted}{julia}
with_default_projection(RTZ_SF) do
    Binary8p4se(1e9)          # clamps, inside this block only
end
\end{minted}

Every example in the User Guide assumes the pristine \texttt{(\allowbreak{}NearestTiesToEven, SatNone)}, which is why each demonstration that changes a default restores it explicitly afterward.

\end{tcolorbox}


\texttt{DefaultProjection} is not merely advisory — \texttt{default\_\allowbreak{}projspec} reads it, so the same-format convenience methods (\texttt{a + b}, \texttt{Exp(x)}, …), the Base-register operators, and \texttt{T(x::Real)} construction all follow it:




\nopagebreak[4]
\begin{minted}{jlcon}
julia> x, y = Binary8p4se(200.0), Binary8p4se(2.0);

julia> x * y                              # RNE_SN: overflow → +Inf
Binary8p4se(Inf ≡ 0x7f)

julia> DefaultProjection!(RTZ_SF);

julia> x * y                              # now clamps to MaxFinite
Binary8p4se(224.0 ≡ 0x7e)

julia> DefaultProjection!(RNE_SN)         # restore before moving on — see below
(NearestTiesToEven, SatNone)

julia> Multiply(Binary8p4se, RNE_SN, x, y)   # explicit ρ is unaffected
Binary8p4se(Inf ≡ 0x7f)
\end{minted}



\begin{tcolorbox}[toptitle=-1mm,bottomtitle=1mm,colback=admonition-warning!50!white,colframe=admonition-warning,title=\textbf{Changing the default is a global semantic change}]
Every caller of the convenience forms — including code in other packages — sees it. Library code that needs a specific projection should name it explicitly rather than rely on the session default.

\end{tcolorbox}


\subsection{The \texttt{with\_\allowbreak{}default\_\allowbreak{}*} combinators}



\label{6052462280974839264}{}


To \emph{consume} a default in your own code without paying dynamic-dispatch costs, go through the \texttt{with\_\allowbreak{}default\_\allowbreak{}*} combinators — \texttt{with\_\allowbreak{}default\_\allowbreak{}type}, \texttt{with\_\allowbreak{}default\_\allowbreak{}returntype}, \texttt{with\_\allowbreak{}default\_\allowbreak{}projection} — which call \texttt{f(default, args...)}:




\nopagebreak[4]
\begin{minted}{jlcon}
julia> with_default_type((T, x) -> T(x), 1.5)
Binary8p2se(1.5 ≡ 0x41)
\end{minted}



\subsubsection{The speculation-guard cost model}



\label{6559940910047111362}{}


Following the default costs nothing while it holds its initial value: the convenience forms consume it through a speculation guard, the same mechanism the \texttt{with\_\allowbreak{}default\_\allowbreak{}*} combinators use, so they compile against the constant and stay allocation-free with concretely inferred results — this is pinned in the test suite, not an incidental optimization. Once you change the default, they cross a function barrier instead: one dynamic dispatch per call, with everything inside that barrier still fully specialized. The explicit forms, \texttt{Add(T, ρ, x, y)}, are unaffected either way — they never read a default in the first place.



\subsubsection{Allocation contract}



\label{8451993120491484882}{}


When \texttt{f}{\textquotesingle}s result type does not depend on the default — \texttt{with\_\allowbreak{}default\_\allowbreak{}projection} with the formats fixed by the caller is the normal shape — the call is zero-allocation with a concretely inferred result, pinned in the test suite. When the result{\textquotesingle}s type \emph{is} the default (\texttt{with\_\allowbreak{}default\_\allowbreak{}type} used as a constructor, where the returned value{\textquotesingle}s type is whatever \texttt{DefaultType} currently names), the value is still computed on the specialized path, but it boxes once where it escapes the function barrier — the irreducible cost of returning a runtime-chosen type.



\subsection{There is no accumulator default}



\label{4054721483846715076}{}


\begin{tcolorbox}[toptitle=-1mm,bottomtitle=1mm,colback=admonition-note!50!white,colframe=admonition-note,title=\textbf{There is no accumulator default}]
\texttt{DefaultAccumulatorType} existed through v0.1.0 and was removed. The wide-precision accumulator used inside a reduction is not a matter of preference: the span filter measures the operands and picks the cheapest carrier that is \emph{provably exact} for them, so the reduction returns the defined result regardless of which carrier that turns out to be. A knob overriding that choice would make \texttt{BlockDotProduct} inexact from the default API — a much larger cost than losing a configuration option. To trade speed for nothing else, set \texttt{ENV[{\textquotedbl}SmallFloats\_\allowbreak{}Float128{\textquotedbl}] = {\textquotedbl}disable{\textquotedbl}}, which is bit-identical by construction and only affects build/oracle speed.

\end{tcolorbox}


\subsection{Go deeper}



\label{4168117981347147797-dup2}{}


\begin{quote}
Continue with \textbf{Projection: Rounding \& Saturation} (Insights track) to see the default projection interact with explicit rounding and saturation choices across a full set of worked overflow examples.

\end{quote}


\section{Binary16p11se Is Not Float16}



\label{11052350688152029480}{}


\begin{tcolorbox}[toptitle=-1mm,bottomtitle=1mm,colback=admonition-warning!50!white,colframe=admonition-warning,title=\textbf{Same bit layout, different values}]
\texttt{Binary16p11se} shares a significand layout with IEEE-754 \texttt{Float16}, and \texttt{Binary16p8se} shares one with the machine-learning \texttt{BFloat16}. Neither pair is interchangeable. Reinterpreting the bits of one as the other produces a number that looks entirely plausible and is wrong by a factor of two.

\end{tcolorbox}


The K ≤ 16 grid contains two formats that resemble well-known IEEE and ML-community 16-bit types closely enough to be mistaken for them. This page explains exactly where they diverge, why the divergence is a single-bit bias difference rather than a structural one, and how to convert between them safely.



\needspace{19\baselineskip}
\subsection{The comparison}



\label{1819503745329534087}{}



\begin{table}[h]
\centering
\begin{tabulary}{\linewidth}{L L L L L}
\toprule
 & \texttt{Binary16p11se} & \texttt{Float16} & \texttt{Binary16p8se} & \texttt{BFloat16} \\
\toprule
precision \texttt{P} & 11 & 11 & 8 & 8 \\
exponent bias & \textbf{16} & \textbf{15} & \textbf{128} & \textbf{127} \\
largest finite & \textbf{65472} & \textbf{65504} & \textbf{3.3762e38} & \textbf{3.3895e38} \\
NaN encodings & \textbf{1} & \textbf{2046} & \textbf{1} & \textbf{2046} \\
negative zero & \textbf{no} & yes & \textbf{no} & yes \\
domain is a parameter & \textbf{yes} (\texttt{se}/\allowbreak\texttt{sf}) & no & \textbf{yes} & no \\
\bottomrule
\end{tabulary}

\end{table}



Every bolded pair looks close enough to be the same thing at a glance, and none of them are.



\subsection{The bias is the key difference}



\label{7300750043427387462}{}


The root cause of every row above is a single design choice about how the exponent bias is defined. P3109 sets the bias so that the exponent range is symmetric about the format{\textquotesingle}s midpoint. IEEE-754 defines it as \texttt{2{\textasciicircum}(E-1) − 1}. The two rules differ by exactly one for these bit layouts — bias 16 versus 15, bias 128 versus 127 — and that difference of one propagates through every code point in the format. \texttt{Binary16p11se} and \texttt{Float16} share a significand layout bit-for-bit and disagree on the \emph{value} of every single code point they hold. It is not that a few values differ at the edges; the entire mapping from bit pattern to number is shifted by a power of two.



\subsection{65472 versus 65504}



\label{13070067715832661291}{}


The largest-finite row is a direct consequence of that shift, plus one more detail specific to extended formats. \texttt{Binary16p11se}{\textquotesingle}s largest finite value is \texttt{(2{\textasciicircum}11 − 1) · 2{\textasciicircum}5 = 65472}. \texttt{Float16}{\textquotesingle}s is \texttt{65504}, one ulp larger. The gap is not the bias difference alone — it{\textquotesingle}s because \texttt{Binary16p11se} is an extended format (\texttt{se}), meaning its top code point in each sign is reserved for \texttt{Inf}, one step short of the largest value the bit pattern could otherwise reach. \texttt{Float16} reserves an entire exponent field for NaN/Inf encodings, which lands its largest finite value one ulp higher in the same layout. The two numbers are close enough that a bug producing \texttt{65472} instead of \texttt{65504} (or vice versa) would not look like an obviously wrong answer — it would look like rounding.



\subsection{Single NaN, no −0, domain as a parameter}



\label{2357685936865613328}{}


The remaining rows are draft design decisions, not incidental byproducts of the bias shift:



\begin{itemize}
\item \textbf{A single NaN encoding} means there is no payload to propagate and no quiet/signalling distinction. The 2045 extra code points IEEE spends on NaN payloads are ordinary datums in a P3109 format instead.


\item \textbf{No negative zero} means \texttt{-0.0} and \texttt{0.0} are the same code point, which is why a sign-symmetric format has an odd number of finite values rather than an even one.


\item \textbf{The domain parameter} — \texttt{e} for extended (with ±Inf), \texttt{f} for finite — has no IEEE analogue at all. \texttt{Binary16p11sf} spends the two code points \texttt{Binary16p11se} gives to infinity on two more finite values instead. There is no such choice available in \texttt{Float16}; the domain is fixed by the IEEE standard.

\end{itemize}


\subsection{Convert, never reinterpret}



\label{8897596886508283514}{}


Because the two formats disagree on the value of every code point, moving a value from one to the other must be a genuine numeric conversion — recompute the code point that represents the same number — never a bit-for-bit reinterpretation:




\nopagebreak[4]
\begin{minted}{jlcon}
julia> using SmallFloats.Formats

julia> Convert(Binary16p11se, RNE_SF, Float16(1.5))    # a conversion: 1.5 ↦ 1.5
Binary16p11se(1.5 ≡ 0x4200)

julia> reinterpret(Binary16p11se, Float16(1.5))        # the SAME BITS: 1.5 ↦ 0.75
Binary16p11se(0.75 ≡ 0x3e00)
\end{minted}



The second line demonstrates the hazard in one expression. \texttt{Float16(1.5)} has the bit pattern \texttt{0x3e00}. Read as a \texttt{Binary16p11se} code point, those same bits denote not 1.5 but \texttt{0.75}, because the exponent bias differs by one. Nothing errors, nothing warns, and the returned value is a plausible, finite, correctly-typed number that is off by a factor of exactly two. This is the single most important habit this page can leave you with: \texttt{Convert} between these formats, always; never \texttt{reinterpret}.



\subsection{Why P3109 chose a different bias rule at all}



\label{12056613435512522735}{}


It is worth asking why the draft did not simply reuse the IEEE bias convention and avoid this entire class of error. The answer is that P3109 formats span a much wider range of bitwidths — 3 through 16 bits — and a single bias rule has to behave sensibly at every one of them, including the very smallest formats where an off-by-one in the bias has a proportionally much larger effect on the usable exponent range. Defining the bias so the exponent range sits symmetrically about the format{\textquotesingle}s midpoint is the choice that generalizes cleanly across that whole span; the IEEE \texttt{2{\textasciicircum}(E-1) - 1} convention was fixed for a small, specific set of binary interchange formats and was never asked to generalize to a 3-bit format. The two rules agree in spirit — both center the exponent range as sensibly as the bit budget allows — and disagree by exactly one at every bitwidth where they overlap, which is precisely why \texttt{Binary16p11se} and \texttt{Float16}, sharing a significand layout, do not share a value mapping.



\subsection{Interop via \texttt{datumsexact}}



\label{12181598389657671511}{}


If IEEE semantics are what you actually want, use \texttt{Float16} and \texttt{BFloat16} directly — the package converts to and from both, exactly, wherever the two datum sets permit it. \texttt{datumsexact} is the query that tells you whether a given conversion between two formats (or a format and \texttt{Float16}/\allowbreak\texttt{BFloat16}) is exact for every value, so you can check before relying on a round trip rather than discover a rounding after the fact.



This is the same discipline the rest of the package applies to every conversion: rather than promising exactness everywhere and hoping, it exposes a query you can call before trusting a round trip. \texttt{datumsexact} between \texttt{Binary16p11se} and \texttt{Float16} reports false in the direction that matters here, and that answer is available before you write a single line of code that depends on the two formats agreeing.



\subsection{The general lesson}



\label{5601922061357718911}{}


The \texttt{Binary16p11se}/\allowbreak\texttt{Float16} pair is the sharpest illustration in the whole package of a rule worth generalizing: two types that look alike are not interchangeable just because their bit patterns are the same shape. Nothing about Julia{\textquotesingle}s type system stops you from calling \texttt{reinterpret} between any two \texttt{isbits} types of the same size, and nothing about the result looking like a normal number will tell you it{\textquotesingle}s wrong. The only defense is knowing, for any two formats you might be tempted to treat as equivalent, whether they actually share a value mapping — and \texttt{Convert} plus \texttt{datumsexact} are the tools that answer that question honestly instead of assuming it.



\subsection{Go deeper}



\label{4168117981347147797-dup3}{}


\begin{quote}
Continue with \textbf{Values, Code Points \& Conversion} (Insights track) for the general rule this page is a special case of: \texttt{UInt8} means code point, every other numeric type means value, and \texttt{Convert} is how you move a value between any two supported representations. That page also covers the carrier \texttt{decode} returns at these wider formats, which is the other half of getting \texttt{Binary16p11se} and \texttt{Binary16p8se} interop right.

\end{quote}


\section{Random Values: What the Draws Mean}



\label{14079953364687431431}{}


\texttt{rand} and \texttt{randn} work on every SmallFloats format, in every setup a Julia programmer already expects — implicit RNG, explicit RNG, arrays, in-place fills, and a choice of projection. This page explains what the draws actually mean statistically, since a floor-projected uniform and a nearest-projected normal are subtly different guarantees, and why that difference is the right default for each.



\subsection{What the draws mean}



\label{2715999854479597255}{}


\textbf{\texttt{rand}} draws the real uniform distribution on \texttt{[0, 1)} at \texttt{Float64} precision and \emph{floor}-projects it onto the format{\textquotesingle}s grid. Floor-projection means each code point receives exactly the real measure of its interval — the code point representing \texttt{0.25} gets drawn with probability proportional to the width of the half-open interval it floor-covers, not the width implied by rounding to nearest. This is what makes \texttt{rand} a correct discrete approximation of the continuous uniform distribution rather than an approximation biased toward the interval midpoints. The result is always in \texttt{[0, 1)}; it never reaches \texttt{1.0}, matching the half-open convention every other Julia numeric type uses for \texttt{rand}.



\textbf{\texttt{randn}} projects a standard-normal \texttt{Float64} draw round-to-nearest, and combines that with \texttt{SatFinite} saturation. Rounding to nearest is the natural choice for a normal draw, since there is no asymmetric floor/ceiling convention to preserve the way there is for a uniform distribution. The \texttt{SatFinite} half of the pairing matters more than it might look: tail draws that land beyond \texttt{MaxFiniteOf(T)} clamp to the extremal finite datum rather than escaping to \texttt{Inf} or \texttt{NaN}. \texttt{randn} never returns \texttt{±Inf} or \texttt{NaN} — a guarantee that matters most for the smallest formats, where the finite range is narrow enough that a naive Gaussian tail draw would routinely overflow. \texttt{Binary3p1se}, for instance, has \texttt{MaxFinite = 1.0}; without the clamp, most \texttt{randn} draws from a genuine standard normal would land outside the format entirely.



\texttt{randn} also requires a signed format, since a negative draw is a normal occurrence, not an edge case, and an unsigned format has nowhere to put one:




\nopagebreak[4]
\begin{minted}{jlcon}
julia> randn(Binary6p3ue)
ERROR: ArgumentError: randn requires a signed format; Binary6p3ue cannot represent negative draws
\end{minted}



\subsection{Setup 1 — quick and implicit}



\label{9816463656748245006}{}


With no RNG argument, draws come from Julia{\textquotesingle}s task-local default generator (a \texttt{Xoshiro}), and \texttt{Random.seed!} controls them exactly as it would for any other numeric type:




\nopagebreak[4]
\begin{minted}{jlcon}
julia> Random.seed!(1234); rand(Binary8p4se)
Binary8p4se(0.3125 ≡ 0x32)

julia> Random.seed!(1234); rand(Binary8p4se)      # same seed, same draw
Binary8p4se(0.3125 ≡ 0x32)
\end{minted}



\subsection{Setup 2 — an explicit RNG stream}



\label{10345610731072531391}{}


Pass any \texttt{AbstractRNG} as the first argument for a reproducible stream that is independent of global task-local state:




\nopagebreak[4]
\begin{minted}{jlcon}
julia> rand(Xoshiro(1), Binary8p4se)
Binary8p4se(0.0703125 ≡ 0x21)

julia> rand(Xoshiro(8), Binary8p4se, 4)
4-element Vector{Binary8p4se}:
 Binary8p4se(0.40625 ≡ 0x35)
 Binary8p4se(0.021484375 ≡ 0x13)
 Binary8p4se(0.75 ≡ 0x3c)
 Binary8p4se(0.28125 ≡ 0x31)
\end{minted}



\subsection{Setup 3 — arrays and in-place}



\label{4815011202830101406}{}


\texttt{rand(T, dims...)} and \texttt{randn(T, dims...)} build an \texttt{Array\{T\}} directly. \texttt{rand!(A)} and \texttt{randn!(A)} fill an existing array in place, which costs one byte per element at \texttt{K ≤ 8} and two above it — cheap enough to preallocate once and reuse across a hot loop:




\nopagebreak[4]
\begin{minted}{jlcon}
julia> A = Vector{Binary8p4se}(undef, 3); randn!(Xoshiro(2), A)
3-element Vector{Binary8p4se}:
 Binary8p4se(-0.005859375 ≡ 0x86)
 Binary8p4se(1.75 ≡ 0x46)
 Binary8p4se(-1.0 ≡ 0xc0)

julia> randn(Xoshiro(3), Binary8p4se, 2, 3) |> typeof
Matrix{Binary8p4se}
\end{minted}



Array draws feed directly into the storage layers when packing is the eventual goal: \texttt{PackedVector(\allowbreak{}rand(\allowbreak{}Binary4p2se, n))} builds the packed representation straight from a batch of draws.



\subsection{Setup 4 — choosing the projection}



\label{8806886239720002746}{}


The scalar \texttt{::Type} forms take a \texttt{projection} keyword to land the draw under any \texttt{ProjSpec} other than the defaults:




\nopagebreak[4]
\begin{minted}{jlcon}
julia> rand(Xoshiro(2), Binary8p4se; projection = RTP_SN)    # ceiling
Binary8p4se(0.0029296875 ≡ 0x03)

julia> randn(Xoshiro(6), Binary8p4se; projection = RTZ_SF) # toward zero
Binary8p4se(-1.875 ≡ 0xc7)

julia> rand(Xoshiro(4), Binary8p4se; projection = RSA_SN(8)) # stochastic
Binary8p4se(0.8125 ≡ 0x3d)
\end{minted}



A stochastic projection draws its random bits from the \emph{same} RNG the value itself came from, so a seeded stream stays fully reproducible even when the projection is stochastic. The defaults exist because they are the contract-keepers: floor for \texttt{rand} guarantees the correct measure per code point, and nearest-plus-\texttt{SatFinite} for \texttt{randn} guarantees no \texttt{±Inf}/\allowbreak\texttt{NaN}. Opting into a different projection can break either guarantee — a ceiling projection can produce exactly \texttt{1.0} from \texttt{rand}, and a \texttt{SatNone} or \texttt{SatPropagate} projection can let \texttt{randn} return \texttt{±Inf} or \texttt{NaN} from a tail draw. The array and \texttt{!} forms always use the defaults; if you need another projection applied across an array, draw scalars explicitly: \texttt{[rand(\allowbreak{}rng, T; projection = ρ) for \_\allowbreak{} in 1:n]}.



The tail-clamp guarantee is easiest to see on the smallest formats, where the finite range is narrow enough that the difference between {\textquotedbl}clamp{\textquotedbl} and {\textquotedbl}overflow{\textquotedbl} shows up on almost every draw — \texttt{Binary3p1se}, with \texttt{MaxFinite = 1.0}, is the canonical example, and is worked through with full transcripts in the Understanding-track random-values walkthrough.



\subsection{Why floor for \texttt{rand} but nearest for \texttt{randn}}



\label{1791254741157626894}{}


The asymmetry between the two default rounding choices is not an inconsistency — each distribution has a different notion of what a correct discrete sample means. A uniform distribution on \texttt{[0, 1)} assigns probability mass in direct proportion to interval width, so the only projection that reproduces the \emph{right} distribution over code points is one that maps each half-open sub-interval to the code point it floor-covers. Rounding a uniform draw to the \emph{nearest} code point instead would bias the result toward interval midpoints, silently distorting the distribution measure. A normal distribution has no such half-open convention to preserve — nearest is simply the natural, unbiased choice, and the only extra concern is the tails, which is where \texttt{SatFinite} earns its place in the pairing.



\subsection{Why this matters for the smallest formats}



\label{5428556917797162578}{}


The finite range shrinks fast as bitwidth drops, and the safety net \texttt{randn}{\textquotesingle}s \texttt{SatFinite} provides matters most exactly where the range is tightest. A \texttt{Binary8p4se} array of \texttt{randn} draws will very rarely reach \texttt{MaxFiniteOf} at all, because the finite range is wide relative to a standard normal{\textquotesingle}s typical spread. A \texttt{Binary3p1se} array is a different story — with \texttt{MaxFinite = 1.0}, a meaningful fraction of an honest standard normal{\textquotesingle}s mass lies beyond the representable range, and every one of those draws needs somewhere defined to land. \texttt{SatFinite} is what makes {\textquotedbl}somewhere defined{\textquotedbl} mean the extremal finite value rather than \texttt{Inf} or \texttt{NaN}, so code that consumes a \texttt{Binary3p1se} array of draws never has to check for non-finite values it did not ask for.



\subsection{Go deeper}



\label{4168117981347147797-dup4}{}


\begin{quote}
Continue with \textbf{Random Draws \& Reproducibility} (Insights track) for the worked K = 3 tail-clamp comparison and the full set of seeded-stream experiments.

\end{quote}


\section{Why No Cross-Format Promotion}



\label{8873299272814736834}{}


Mixing two \texttt{Binary} formats in an arithmetic expression fails. Mixing a \texttt{Binary} value with an ordinary Julia number succeeds, but the result is not a \texttt{Binary} value. Both behaviors are deliberate, and this page explains the reasoning behind each, along with why the obvious third option — a \texttt{promote\_\allowbreak{}rule} that produces a wider \texttt{Binary} format automatically — was considered and rejected.



\subsection{Mixed \texttt{Binary} formats: deliberate failure, not silent widening}



\label{12941061487779548011}{}


Adding two values from different formats does not silently pick a format wide enough to hold both and round into it. It fails, with an error that names the problem:




\nopagebreak[4]
\begin{minted}{jlcon}
julia> x = Binary8p4se(1.5);

julia> x + Convert(Binary8p4se, RNE_SN, Binary5p3sf(1.0))
Binary8p4se(2.5 ≡ 0x4a)
\end{minted}



The line above only works because the \texttt{Binary5p3sf} value was converted \emph{explicitly} first. Without that conversion, \texttt{x + y} for \texttt{x::Binary8p4se} and \texttt{y::Binary5p3sf} raises a \texttt{promotion ... failed to change} error rather than silently choosing some third format to compute in. This is not a gap in the promotion machinery — it{\textquotesingle}s the package{\textquotesingle}s stated position that mixing formats is a decision the caller makes explicitly, with \texttt{Convert}, never one the package makes for you by picking a resolution.



\subsection{\texttt{Binary} with ordinary numbers: promotes to the promotion carrier}



\label{3293283507695187362}{}


Combining a \texttt{Binary} value with an ordinary Julia number — an \texttt{Int}, a \texttt{Float64} — does not fail, but it does not produce a \texttt{Binary} value either. It promotes to the format{\textquotesingle}s \textbf{promotion carrier}, an ordinary float type, and the result is that ordinary float:




\nopagebreak[4]
\begin{minted}{jlcon}
julia> x + 2.0
3.5
\end{minted}



For 432 of the 504 formats — the ones whose exponent range and precision fit inside a \texttt{Float64} without rounding — that carrier is \texttt{Float64}, which is what the example above returns. For the other 72, the carrier is \texttt{BigFloat}: a \texttt{Float64} cannot hold a \texttt{Binary16p1uf} datum exactly, and promoting into \texttt{Float64} anyway would silently turn a perfectly finite value into \texttt{±Inf} with no P3109 operation anywhere in sight. \texttt{promotecarrier(T)} names the target carrier for any format, so the promotion is queryable rather than a fact you have to remember.



If you actually want the result back in the original format, project it back explicitly — the package will not do this step for you, for the same reason it won{\textquotesingle}t guess which format to promote two mixed \texttt{Binary} values into:




\nopagebreak[4]
\begin{minted}{julia}
Convert(Binary8p4se, RNE_SN, decode(x) + 2.0)
\end{minted}



\subsection{Why widen is not another format}



\label{6193834797985144004}{}


It{\textquotesingle}s tempting to think the fix for mixed-format arithmetic is obvious: define a \texttt{promote\_\allowbreak{}rule} between two \texttt{Binary} types that picks the smallest format containing both operands{\textquotesingle} datum sets, the way \texttt{promote\_\allowbreak{}type(\allowbreak{}Int, Float64)} picks \texttt{Float64}. A companion internal study worked this out in detail, and the answer is that this rule cannot be built soundly. Formats trade precision for range through the identity \texttt{K = P + E + σ} — raising precision by one bit necessarily lowers exponent range by the same amount — so two formats at the same bitwidth are usually \emph{incomparable}: neither contains the other{\textquotesingle}s datum set. Measured over the full 504-format grid, a {\textquotedbl}smallest format containing both{\textquotedbl} answer exists for only about 60\% of format pairs; the other 40\% have no common format at all within the K ≤ 16 grid, so any \texttt{Binary}-valued rule would have to fall back to something else for those pairs anyway.



Worse, even where the smallest containing format does exist, picking it pairwise is not associative — because Julia{\textquotesingle}s \texttt{promote\_\allowbreak{}type(a, b, c)} reduces pairwise, a non-associative rule would make the \emph{result format} depend on the order the operands happened to appear in a \texttt{+} chain, which is exactly the kind of instability the package refuses to ship. And even in the cases where the rule is well-defined, it is often surprising on its own terms: two formats that agree on everything except whether they carry infinities do not promote to either input{\textquotesingle}s own bitwidth — they need one bit more, because the wider format has to hold both the finite format{\textquotesingle}s largest value \emph{and} an infinity code point, and at the narrower width there is no code point left for both. A silent promotion landing in a format the user never named, at a width neither operand has, is the wrong kind of convenience for a package whose entire premise is that you choose your result format. The full argument, with the exact percentages and the associativity counterexample, is worked out in the promotion-rules study (docs/other/promoterules.md in the repository).



\subsection{Why this differs from \texttt{Int}/\allowbreak\texttt{Float64} promotion}



\label{7463766682595139252}{}


It helps to be precise about what makes format promotion a different problem from ordinary numeric promotion. \texttt{promote\_\allowbreak{}type(\allowbreak{}Int, Float64)} works because the target, \texttt{Float64}, is fixed in advance and holds every \texttt{Int} exactly (up to 2{\textasciicircum}53) — the rule never has to ask {\textquotedbl}which format contains both operands{\textquotedbl} because there is only one plausible answer and it never runs out of room. Two \texttt{Binary} formats have no such fixed universal target: the format space is a lattice of finite datum sets, not a single always-sufficient supertype, and which format contains which shifts with every combination of precision, range, signedness, and domain. The question \texttt{promote\_\allowbreak{}type} answers for ordinary numbers is trivial; the same question for two arbitrary \texttt{Binary} formats is a genuine combinatorial problem, and the promotion-rules study exists because that problem was worked out in full rather than assumed away.



\subsection{Explicit \texttt{Convert} as the answer}



\label{705282074153534982}{}


Given that no automatic rule is both sound and predictable, the package{\textquotesingle}s answer is to make the choice explicit every time: \texttt{Convert(T, ρ, x)} moves any value into the format \texttt{T} you name, under the projection \texttt{ρ} you choose. There is no ambiguity about which format the computation happens in, no order-dependence, and no case where the rule silently falls back to a carrier for some inputs and a \texttt{Binary} format for others. The lattice of which format contains which is real and occasionally useful to know explicitly, but it belongs in a query you call when you want the answer, not in an operator that decides for you every time two formats meet.



This is the same philosophy that governs every other write path in the package: one explicit projection, chosen by the caller, rather than a rule that infers intent from context. A \texttt{promote\_\allowbreak{}rule} between two \texttt{Binary} formats would be exactly the kind of implicit rounding decision the rest of the design goes out of its way to avoid — it would round somewhere other than inside \texttt{project}, which is the one invariant the whole package is built not to break. Treating cross-format arithmetic as a decision you make explicitly, rather than a convenience the library infers, is not a limitation bolted on around an otherwise-automatic feature; it is the same rule that governs rounding modes, saturation modes, and result formats applied consistently to one more place where an implicit choice could have hidden.



\subsection{What this buys you as a caller}



\label{18348306953458602459}{}


The practical upshot is that a \texttt{Binary} expression never surprises you about which format it computed in. If every operand in an expression shares a format, the result is that format, under whatever projection was in play. If you mix in an ordinary number, the result predictably drops to the promotion carrier, and you can query \texttt{promotecarrier(T)} in advance to know which carrier that will be. And if you try to mix two different \texttt{Binary} formats without converting, the error fires immediately rather than letting a silently-widened result propagate three functions deeper before something notices the format changed. All three outcomes are things you can predict from the types alone, before running anything — which is the actual goal a promotion rule is supposed to serve, achieved here by refusing the one case that could not serve it honestly.



\subsection{Go deeper}



\label{4168117981347147797-dup5}{}


\begin{quote}
Continue with \textbf{Values, Code Points \& Conversion} (Insights track) for the full set of \texttt{Convert} argument types and the constructor rules that make explicit conversion the answer to every cross-format situation.

\end{quote}


\section{Performance Model}



\label{11678570700420481952}{}


Performance advice for SmallFloats.jl is scattered across the User Guide, the cheat sheet, and the technical examples because it applies at every layer — construction, arithmetic, arrays, storage. This page pulls all of it into one spine: the handful of decisions that determine whether your code runs at scalar speed (tens of nanoseconds) or falls off a cliff into microsecond dynamic dispatch.



\subsection{Pass formats statically}



\label{14626146111254921222}{}


Every scalar entry point in the package — \texttt{Add}, \texttt{project}, a constructor — fully specializes when the format type is known at compile time: through a \texttt{const} binding, a type parameter, or an ordinary function argument. Under that condition, a scalar \texttt{Add} runs at roughly 18–26 ns and \texttt{project} at roughly 13 ns, with zero allocations.



The moment a format type is read from a \textbf{non-\texttt{const} global}, Julia can no longer resolve the method at compile time, and every call pays for dynamic dispatch instead — roughly 1 µs for a keyword call, a difference of nearly two orders of magnitude for the identical computation:



\begin{tcolorbox}[toptitle=-1mm,bottomtitle=1mm,colback=admonition-default!50!white,colframe=admonition-default,title=\textbf{The 60× benchmarking slowdown}]
The same benchmark, run with the format type \texttt{T} read from a non-\texttt{const} global instead of passed as a type parameter, measures around 1 µs — Julia{\textquotesingle}s dynamic keyword dispatch overhead, not this package{\textquotesingle}s arithmetic. Two real project post-mortems trace their {\textquotedbl}SmallFloats is slow{\textquotedbl} reports to exactly this mistake. The shipped \texttt{benchmarking/\allowbreak{}benchmarking.jl} asserts zero warm-path allocation before it believes any number it produces, and your own benchmarks should do the same:


\nopagebreak[4]
\begin{minipage}{\linewidth}
\begin{minted}{julia}
using Chairmarks, SmallFloats, Random
using Statistics: median

function bench_add(::Type{T}) where {T<:Binary}          # T: type parameter, not a global
    pool = [rawvalue(T, rand(UInt8)) for _ in 1:4096]
    @be (rand(pool), rand(pool)) (t -> Add(T, RNE_SN, t[1], t[2]))(_)
end

b = bench_add(Binary8p4se)
(round(median(b).time * 1e9; digits=1), median(b).allocs)
\end{minted}


\begin{minted}{text}
(16.3, 0.0)          # ns per full scalar Add, zero allocations
\end{minted}
\end{minipage}


\end{tcolorbox}


The fix, when a format genuinely has to come from a runtime value, is one function barrier: write \texttt{f(::Type\{T\}, …) where \{T\}} and call into it once the type is known, and everything inside specializes at full speed again.



Format operands entering \texttt{Chairmarks}-style benchmarks specifically must come from the untimed \texttt{setup} phase, and if you retrieve an operation function reflectively (\texttt{getfield(\allowbreak{}SmallFloats, op)}), pass it through an argument barrier so it specializes too — the same discipline that makes the benchmark trustworthy also makes ordinary calling code fast.



\subsection{Convenience forms are free at the initial default}



\label{4134116807364174795}{}


\texttt{x + y}, \texttt{Exp(x)}, and \texttt{T(2.1)} all read the session default projection through a speculation guard rather than a dynamic lookup. While the default holds its initial value, the guard lets these calls compile against a constant, so they are allocation-free with concretely inferred results — a property pinned in the test suite, not an incidental benchmark result. The moment you change the default, these same calls cross a function barrier instead: one dynamic dispatch per call, with everything inside that barrier still fully specialized. Code that must stay insensitive to whatever the session default happens to be should name its projection explicitly — \texttt{Add(T, ρ, x, y)} with a \texttt{const} \texttt{ρ} — rather than rely on the convenience spelling in a hot loop.



\subsection{Bulk work belongs in array calls}



\label{3141937120956456863}{}


The table-gather kernels that back array operations run roughly 50× faster than the scalar path per element, because after the first call for a given \texttt{(op, formats, ρ)} specialization, every later element is a single table lookup rather than a full evaluation. That first call pays a one-time table build: roughly 0.4 ms for a unary table, tens of milliseconds for an 8×8 binary table, up to a few milliseconds for a 2 MiB ternary table. Benchmark warm calls separately from the first cold call, or the build cost will dominate a measurement that is supposed to be about steady-state throughput. Once warm, measured throughput is around 0.27 ns/element for unary gathers and 0.5 ns/element for binary gathers.



\subsubsection{Ternary tiers scale with bitwidth}



\label{15486319803806300295}{}


Ternary operations (\texttt{FMA}, \texttt{FAA}, \texttt{Clamp}) ride the same table-gather mechanism whenever the combined operand bitwidth keeps the resulting table affordable, and the package picks the tier automatically — no action needed on your part:



\begin{itemize}
\item \textbf{Eager} (up to 256 KiB; every all-\texttt{K ≤ 6} signature): built on the first array call, exactly like a unary or binary table.


\item \textbf{Adaptive} (up to 2 MiB; the \texttt{K = 7} band): built only once a signature has processed enough elements to earn its build, held in a byte-bounded, LRU-evicted cache — so a signature used once in passing doesn{\textquotesingle}t pay for a table it will never reuse.


\item \textbf{Compute} (\texttt{K = 8}; a 16 MiB table stops being a cache win): the scalar pipeline runs per element instead, optionally threaded for long arrays (roughly 4× at 4 threads).

\end{itemize}


Every table entry, eager or adaptive, is built through the exact scalar path, so a table lookup and the equivalent scalar call are bit-identical by construction — the speed comes with no accuracy trade at all.



\subsection{Stochastic array calls are always per-element}



\label{7899356641446882989}{}


Deterministic array operations use cached tables; stochastic ones cannot, because each element needs its own independent random draw. Stochastic array calls always run the scalar pipeline once per element, consuming one draw per projection. Pass an explicit \texttt{rng} when you need the result reproducible — the array and \texttt{!} forms always use whatever RNG you supply, so this is the one lever available for controlling a stochastic array call{\textquotesingle}s output exactly.



\subsection{Memory: \texttt{PackedVector} and \texttt{BlockVector}}



\label{7347418221015137253}{}


When storage bandwidth matters more than direct byte-level access, \texttt{PackedVector\{F\}} stores code points at \texttt{bitwidth(F)} bits per element rather than rounding up to a whole byte or word — a \texttt{Binary5p2se} vector at 5 bits/element instead of 8. It behaves as a full \texttt{AbstractVector\{F\}} for indexing, \texttt{collect}, and iteration, and \texttt{vmap} accepts it directly, unpacking cache-friendly tiles internally rather than materializing the whole array. The governing rule is \emph{store packed, compute unpacked}: there is deliberately no in-place packed arithmetic, because computing directly on sub-byte-packed bits would give up more in kernel complexity than it could ever save in bandwidth. \texttt{BlockVector} is the analogous structure for many same-shape \texttt{Block}s, holding them in a structure-of-arrays layout instead of an array of individually-boxed blocks.



\texttt{table\_\allowbreak{}bytes()} reports the current cache footprint across unary, binary, and ternary tables alike; \texttt{empty\_\allowbreak{}tables!()} resets it, which is useful both for isolating a benchmark from a previous session{\textquotesingle}s warm tables and for freeing memory in a long-running process that has touched many format specializations.



\subsection{Sorting: O(n) counting sort}



\label{898878640439474703}{}


Sorting is special-cased rather than falling through to Base{\textquotesingle}s generic comparison sort. \texttt{Binary} values compare through integer order keys, and a vector of \texttt{Binary} values sorts with an \textbf{O(n) counting sort} installed as the default algorithm — roughly 8× the stock comparison sort at 64K elements, with \texttt{rev=true} supported at the same cost. \texttt{sort(A)} simply picks this path up automatically; NaN sorts last under the default order (first under \texttt{rev=true}), matching Base{\textquotesingle}s own conventions for float sorting so existing code does not need to special-case \texttt{Binary} vectors.



\subsection{Do not use \texttt{@fastmath}}



\label{1769973427897829437}{}


Every number in the performance model above depends on the one-write-path contract: exact math, then exactly one projection. \texttt{@fastmath} exists in Julia to relax IEEE semantics for speed — reordering operations, assuming no NaNs, dropping distinctions the standard is careful about — and every one of those relaxations is exactly the kind of silent rounding shortcut this package{\textquotesingle}s whole design refuses to take. Applying \texttt{@fastmath} to code using \texttt{Binary} values does not make the arithmetic faster in any meaningful way, because the actual bottleneck the sections above address is dispatch and table-gather behavior, not floating-point instruction scheduling — but it can silently invalidate the exactness guarantee that makes the table-gather speedup safe to rely on in the first place. If a computation is slow, the fix is one of the items above: static format types, a function barrier, array calls instead of scalar loops, or \texttt{PackedVector} for memory pressure — never \texttt{@fastmath}.



\subsection{Go deeper}



\label{4168117981347147797-dup6}{}


\begin{quote}
Continue with \textbf{Benchmarking Without Fooling Yourself} (Understanding track) for the full benchmark doctrine — Chairmarks setup discipline, argument barriers for reflectively-retrieved functions, and the warm/cold table-build measurement split.

\end{quote}


\chapter{Reference}


\section{Reference: Format Names \& Queries}



\label{10754142039234372555}{}


Dry listing of the format naming grammar, the alias/type distinction, and every format-query accessor. For narrative explanation, see the Mental Model page; for a worked introduction, see the Quickstart and Values, Code Points \& Conversion.



\needspace{41\baselineskip}
\subsection{Naming grammar}



\label{14025844087195034522}{}



\begin{minted}{text}
Binary K p P (s|u) (e|f)
       │   │  │     └─ extended (Inf) or finite domain
       │   │  └────── signed or unsigned
       │   └───────── significand precision, including the implicit bit
       └───────────── total bitwidth, 3 through 16
\end{minted}




\begin{table}[h]
\centering
\begin{tabulary}{\linewidth}{L L L}
\toprule
Field & Values & Meaning \\
\toprule
\texttt{K} & 3–16 & total bitwidth \\
\texttt{P} & 1–\texttt{K} & significand precision, including the implicit bit \\
\texttt{s}|\texttt{u} & signed / unsigned & sign field present or absent \\
\texttt{e}|\texttt{f} & extended / finite & domain includes ±Inf, or spends those code points on finite datums \\
\bottomrule
\end{tabulary}

\end{table}



There are 504 legal \texttt{(K,P,SGN,EXT)} combinations, each with a draft-named alias. 120 aliases at \texttt{K ≤ 8} are exported by \texttt{using SmallFloats}.




\begin{table}[h]
\centering
\begin{tabulary}{\linewidth}{L L L L L}
\toprule
Example alias & K & P & Signed & Domain \\
\toprule
\texttt{Binary8p4se} & 8 & 4 & signed & extended \\
\texttt{Binary8p4sf} & 8 & 4 & signed & finite \\
\texttt{Binary6p3ue} & 6 & 3 & unsigned & extended \\
\texttt{Binary5p5uf} & 5 & 5 & unsigned & finite \\
\texttt{Binary16p6se} & 16 & 6 & signed & extended (opt-in export) \\
\texttt{Binary16p1uf} & 16 & 1 & unsigned & finite \\
\bottomrule
\end{tabulary}

\end{table}



\subsection{\texttt{<:} versus \texttt{===}}



\label{7480538980344996601}{}


\begin{itemize}
\item \texttt{Binary} is abstract; \texttt{Binary8p4se} is its concrete representation.


\item \texttt{Binary8p4se <: Binary\{8,4,true,true\}} is \texttt{true}.


\item \texttt{Binary8p4se === Binary\{8,4,true,true\}} is \texttt{false}. Use the alias, or \texttt{format(K, P, Σ, Δ)} for runtime parameters. See Values, Code Points \& Conversion.

\end{itemize}


\needspace{19\baselineskip}
\subsection{\texttt{format()} and opt-in access}



\label{14939675529699646064}{}



\begin{table}[h]
\centering
\begin{tabulary}{\linewidth}{L L L}
\toprule
Form & Scope & Result \\
\toprule
\texttt{using SmallFloats} & exports 120 aliases, \texttt{K ≤ 8} & alias directly in scope \\
\texttt{using SmallFloats.Formats} & brings all 504 aliases into scope & alias directly in scope \\
\texttt{SmallFloats.Binary16p6se} & unexported, always reachable & alias by qualified path \\
\texttt{format(K, P, Σ, Δ)} & programmatic, runtime parameters & the alias for those parameters \\
\texttt{formatname(Binary\{K,P,Σ,Δ\})} & programmatic & \texttt{Symbol} of the alias name \\
\bottomrule
\end{tabulary}

\end{table}



\needspace{26\baselineskip}
\subsection{Format query table}



\label{3219646311892515237}{}


Every query accepts either a \texttt{Type} or a \texttt{value}; \texttt{bitwidth(\allowbreak{}x) ≡ bitwidth(\allowbreak{}typeof(\allowbreak{}x))}, and both fold to a literal constant.




\begin{table}[h]
\centering
\begin{tabulary}{\linewidth}{L L L L}
\toprule
Query & Returns & Draft-named form & Teaching page \\
\toprule
\texttt{bitwidth(T)} & \texttt{K} & \texttt{BitwidthOf(T)} & Cheat Sheet \\
\texttt{precision(T)} & \texttt{P} & \texttt{PrecisionOf(T)} & Cheat Sheet \\
\texttt{issigned(T)} & \texttt{Bool} & \texttt{SignednessOf(T)} & Cheat Sheet \\
\texttt{isextended(T)} & \texttt{Bool} & \texttt{DomainOf(T)} & Cheat Sheet \\
\texttt{expbias(T)} & exponent bias & \texttt{ExponentBiasOf(T)} & Cheat Sheet \\
\texttt{expbitwidth(T)} & exponent field width & — & Cheat Sheet \\
\texttt{trailingsigbits(T)} & trailing significand bits & — & Cheat Sheet \\
\texttt{MaxFiniteOf(T)} & largest finite value, typed & — & Cheat Sheet \\
\texttt{MinFiniteOf(T)} & smallest-magnitude finite value, typed & — & Cheat Sheet \\
\texttt{MinPositiveOf(T)} & smallest positive value, typed & — & Cheat Sheet \\
\texttt{MinNormalOf(T)} & smallest positive normal value, typed & — & Cheat Sheet \\
\texttt{MaxSubnormalOf(T)} & largest subnormal value, typed & — & Cheat Sheet \\
\bottomrule
\end{tabulary}

\end{table}



Draft-named forms exist for every query in Group M, not only the four listed above (\texttt{BitwidthOf}, \texttt{PrecisionOf}, \texttt{SignednessOf}, \texttt{DomainOf}, \texttt{ExponentBiasOf}, and their neighbors); Julia-style and draft-named forms are interchangeable.



\subsection{Value-and-type call forms}



\label{7040199918633646645}{}



\begin{minted}{julia}
bitwidth(Binary8p4se)        # type form
x = Binary8p4se(1.6)
bitwidth(x)                  # value form, same answer, folds to 8
MaxFiniteOf(Binary8p4se)     # Binary8p4se(224.0 ≡ 0x7e)
MaxFiniteOf(x)               # same result
\end{minted}



\subsection{see also}



\label{9115667341246479961-dup15}{}


\href{mentalmodel.md}{Mental Model}, \href{tutorial1\_\allowbreak{}values.md}{Values, Code Points \& Conversion}, \href{howto\_\allowbreak{}choose\_\allowbreak{}format.md}{Choose a Format}.



\section{Reference: Projection Specifications}



\label{6573491883575706896}{}


Dry listing of \texttt{ProjSpec}, the predefined deterministic grid, the stochastic constructor families, and their accessors. For narrative explanation, see Projection: Rounding \& Saturation.



\subsection{\texttt{ProjSpec} constructor}



\label{14281805247734099455}{}



\begin{minted}{julia}
ProjSpec(rounding_mode, saturation_mode)
\end{minted}



A \texttt{ProjSpec} is the pair \texttt{(\allowbreak{}rounding mode, saturation mode)}. Both components are zero-size singleton types; a \texttt{ProjSpec} costs nothing at runtime and fully specializes every call that reaches it.



\needspace{25\baselineskip}
\subsection{Deterministic predefined specs (6 × 3)}



\label{12417255764667262599}{}


Named \texttt{R<mode>\_\allowbreak{}Sat<mode>}. Every cell below is an exported constant.




\begin{table}[h]
\centering
\begin{tabulary}{\linewidth}{L L L L L}
\toprule
Rounding & \texttt{SatFinite} & \texttt{SatPropagate} & \texttt{SatNone} & Teaching page \\
\toprule
\texttt{NearestTiesToEven} & \texttt{RNE\_\allowbreak{}SF} & \texttt{RNE\_\allowbreak{}SP} & \texttt{RNE\_\allowbreak{}SN} & Projection: Rounding \& Saturation \\
\texttt{NearestTiesToAway} & \texttt{RNA\_\allowbreak{}SF} & \texttt{RNA\_\allowbreak{}SP} & \texttt{RNA\_\allowbreak{}SN} & Projection: Rounding \& Saturation \\
\texttt{TowardPositive} & \texttt{RTP\_\allowbreak{}SF} & \texttt{RTP\_\allowbreak{}SP} & \texttt{RTP\_\allowbreak{}SN} & Projection: Rounding \& Saturation \\
\texttt{TowardNegative} & \texttt{RTN\_\allowbreak{}SF} & \texttt{RTN\_\allowbreak{}SP} & \texttt{RTN\_\allowbreak{}SN} & Projection: Rounding \& Saturation \\
\texttt{TowardZero} & \texttt{RTZ\_\allowbreak{}SF} & \texttt{RTZ\_\allowbreak{}SP} & \texttt{RTZ\_\allowbreak{}SN} & Projection: Rounding \& Saturation \\
\texttt{ToOdd} & \texttt{RTO\_\allowbreak{}SF} & \texttt{RTO\_\allowbreak{}SP} & \texttt{RTO\_\allowbreak{}SN} & Projection: Rounding \& Saturation \\
\bottomrule
\end{tabulary}

\end{table}



\texttt{RNE\_\allowbreak{}SN} is the package-wide initial default (\texttt{DefaultProjection()}).



\needspace{21\baselineskip}
\subsection{Stochastic families}



\label{8787228723720205682}{}


Constructors, not constants — the random-bit budget \texttt{N} is a type parameter. \texttt{N ∈ 1:60}; omitting \texttt{N} uses the default \texttt{N = 8} (\texttt{SmallFloats.DEFAULT\_\allowbreak{}RBITS}).




\begin{table}[h]
\centering
\begin{tabulary}{\linewidth}{L L L L L}
\toprule
Variant & \texttt{SatFinite} & \texttt{SatPropagate} & \texttt{SatNone} & Teaching page \\
\toprule
\texttt{StochasticA} & \texttt{RSA\_\allowbreak{}SF(N)} & \texttt{RSA\_\allowbreak{}SP(N)} & \texttt{RSA\_\allowbreak{}SN(N)} & Projection: Rounding \& Saturation \\
\texttt{StochasticB} & \texttt{RSB\_\allowbreak{}SF(N)} & \texttt{RSB\_\allowbreak{}SP(N)} & \texttt{RSB\_\allowbreak{}SN(N)} & Projection: Rounding \& Saturation \\
\texttt{StochasticC} & \texttt{RSC\_\allowbreak{}SF(N)} & \texttt{RSC\_\allowbreak{}SP(N)} & \texttt{RSC\_\allowbreak{}SN(N)} & Projection: Rounding \& Saturation \\
\bottomrule
\end{tabulary}

\end{table}




\begin{minted}{julia}
RSA_SN()      # StochasticA, N = 8, SatNone — the default form
RSA_SN(16)    # StochasticA, N = 16, SatNone
RSC_SF(16) === ProjSpec(StochasticC{16}(), SatFinite())   # true
\end{minted}



\needspace{18\baselineskip}
\subsection{Accessors}



\label{11927386850142820610}{}



\begin{table}[h]
\centering
\begin{tabulary}{\linewidth}{L L L}
\toprule
Accessor & Returns & Teaching page \\
\toprule
\texttt{roundingmode(ρ)} & the rounding-mode instance & Projection: Rounding \& Saturation \\
\texttt{saturationmode(ρ)} & the saturation-mode instance & Projection: Rounding \& Saturation \\
\texttt{isstochastic(ρ)} & \texttt{Bool} & Projection: Rounding \& Saturation \\
\texttt{nrandbits(ρ)} & \texttt{N} for a stochastic spec & Projection: Rounding \& Saturation \\
\bottomrule
\end{tabulary}

\end{table}



\needspace{17\baselineskip}
\subsection{Saturation-mode meaning}



\label{8872045410401590770}{}



\begin{table}[h]
\centering
\begin{tabulary}{\linewidth}{L L}
\toprule
Mode & Out-of-range behavior \\
\toprule
\texttt{SatFinite} & clamp everything to the finite range \\
\texttt{SatPropagate} & preserve representable infinities; clamp other overflow \\
\texttt{SatNone} & apply the draft{\textquotesingle}s domain-, signedness-, and direction-dependent rows (nearest modes overflow to ±Inf or NaN for finite formats; directed modes pointing away from the overflow clamp to the extremal finite value) \\
\bottomrule
\end{tabulary}

\end{table}



\subsection{Explicit \texttt{R} range rule}



\label{11101438195765952950}{}


For an \texttt{N}-bit stochastic mode, an explicit draw \texttt{R} passed as a keyword must satisfy \texttt{0 ≤ R ≤ 2{\textasciicircum}N - 1}. Passing \texttt{R} outside that range is an error.




\begin{minted}{julia}
Add(T, σ, x, y; rng = Xoshiro(1))   # reproducible stream
Add(T, σ, x, y; R = 17)             # exact draw, ideal for tests
\end{minted}



\subsection{see also}



\label{9115667341246479961-dup16}{}


\href{tutorial2\_\allowbreak{}projection.md}{Projection: Rounding \& Saturation}, \href{ref\_\allowbreak{}defaults.md}{Session Defaults}.



\section{Reference: Operation Catalog}



\label{17973223157804438522}{}


The canonical operation table. 52 operations total: 30 unary, 18 binary, 3 ternary, 1 conversion. Source: the scalar operation catalog in the Cheat Sheet, cross-checked against Julia Compatibility. For semantics and worked examples, see the User Guide and Julia Compatibility pages.



\needspace{47\baselineskip}
\subsection{Unary (30)}



\label{2812026322430809105}{}



\begin{table}[h]
\centering
\begin{tabulary}{\linewidth}{L L L L L}
\toprule
Operation & Arity & Base spelling & Notes & Teaching page \\
\toprule
\texttt{Abs} & unary & \texttt{abs} &  & Cheat Sheet \\
\texttt{Negate} & unary & unary \texttt{-} &  & Cheat Sheet \\
\texttt{Sqrt} & unary & \texttt{sqrt} &  & Cheat Sheet \\
\texttt{RSqrt} & unary & — & no Base counterpart; call by draft name & Julia Compatibility \\
\texttt{Recip} & unary & \texttt{inv} & \texttt{Recip(±Inf) = 0} & Cheat Sheet \\
\texttt{Exp} & unary & \texttt{exp} &  & Cheat Sheet \\
\texttt{Exp2} & unary & \texttt{exp2} &  & Cheat Sheet \\
\texttt{ExpMinusOne} & unary & \texttt{expm1} &  & Cheat Sheet \\
\texttt{Log} & unary & \texttt{log} &  & Cheat Sheet \\
\texttt{Log2} & unary & \texttt{log2} &  & Cheat Sheet \\
\texttt{LogOnePlus} & unary & \texttt{log1p} &  & Cheat Sheet \\
\texttt{Softplus} & unary & — & no Base counterpart; call by draft name & Julia Compatibility \\
\texttt{Sin} & unary & \texttt{sin} &  & Cheat Sheet \\
\texttt{Cos} & unary & \texttt{cos} &  & Cheat Sheet \\
\texttt{Tan} & unary & \texttt{tan} &  & Cheat Sheet \\
\texttt{ArcSin} & unary & \texttt{asin} &  & Cheat Sheet \\
\texttt{ArcCos} & unary & \texttt{acos} &  & Cheat Sheet \\
\texttt{ArcTan} & unary & \texttt{atan} & single-argument form; see \texttt{ArcTan2} for two-argument & Cheat Sheet \\
\texttt{Sinh} & unary & \texttt{sinh} &  & Cheat Sheet \\
\texttt{Cosh} & unary & \texttt{cosh} &  & Cheat Sheet \\
\texttt{Tanh} & unary & \texttt{tanh} &  & Cheat Sheet \\
\texttt{ArcSinh} & unary & \texttt{asinh} &  & Julia Compatibility \\
\texttt{ArcCosh} & unary & \texttt{acosh} &  & Julia Compatibility \\
\texttt{ArcTanh} & unary & \texttt{atanh} &  & Julia Compatibility \\
\texttt{SinPi} & unary & \texttt{sinpi} & reduces exactly mod 2 first & Cheat Sheet \\
\texttt{CosPi} & unary & \texttt{cospi} & reduces exactly mod 2 first & Cheat Sheet \\
\texttt{TanPi} & unary & \texttt{tanpi} & reduces exactly mod 2 first & Cheat Sheet \\
\texttt{ArcSinPi} & unary & — & no Base counterpart; call by draft name & Julia Compatibility \\
\texttt{ArcCosPi} & unary & — & no Base counterpart; call by draft name & Julia Compatibility \\
\texttt{ArcTanPi} & unary & — & no Base counterpart; call by draft name & Julia Compatibility \\
\bottomrule
\end{tabulary}

\end{table}



Trig of ±Inf is NaN for all rows above.



\needspace{43\baselineskip}
\subsection{Binary (18)}



\label{9493958918153706001}{}



\begin{table}[h]
\centering
\begin{tabulary}{\linewidth}{L L L L L}
\toprule
Operation & Arity & Base spelling & Notes & Teaching page \\
\toprule
\texttt{CopySign} & binary & \texttt{copysign} &  & Julia Compatibility \\
\texttt{Add} & binary & \texttt{+} &  & Cheat Sheet \\
\texttt{Subtract} & binary & \texttt{-} &  & Cheat Sheet \\
\texttt{Multiply} & binary & \texttt{*} &  & Cheat Sheet \\
\texttt{Divide} & binary & \texttt{/} & \texttt{x / 0 → NaN} for every \texttt{x}, including ±Inf & Cheat Sheet \\
\texttt{Hypot} & binary & \texttt{hypot} &  & Julia Compatibility \\
\texttt{ArcTan2} & binary & \texttt{atan(y, x)} & Base{\textquotesingle}s \texttt{(y, x)} argument order & Cheat Sheet \\
\texttt{ArcTan2Pi} & binary & — & no Base counterpart; call by draft name & Julia Compatibility \\
\texttt{Maximum} & binary & \texttt{max} & NaN-propagating, exactly Base{\textquotesingle}s float semantics & Cheat Sheet \\
\texttt{Minimum} & binary & \texttt{min} & NaN-propagating, exactly Base{\textquotesingle}s float semantics & Cheat Sheet \\
\texttt{MaximumNumber} & binary & — & NaN-ignoring extremum; no Base counterpart & Julia Compatibility \\
\texttt{MinimumNumber} & binary & — & NaN-ignoring extremum; no Base counterpart & Julia Compatibility \\
\texttt{MaximumMagnitude} & binary & — & magnitude extremum; no Base counterpart & Julia Compatibility \\
\texttt{MinimumMagnitude} & binary & — & magnitude extremum; no Base counterpart & Julia Compatibility \\
\texttt{MaximumMagnitudeNumber} & binary & — & NaN-ignoring magnitude extremum; no Base counterpart & Julia Compatibility \\
\texttt{MinimumMagnitudeNumber} & binary & — & NaN-ignoring magnitude extremum; no Base counterpart & Julia Compatibility \\
\texttt{MinimumFinite} & binary & — & finite-preferring extremum; no Base counterpart & Julia Compatibility \\
\texttt{MaximumFinite} & binary & — & finite-preferring extremum; no Base counterpart & Julia Compatibility \\
\bottomrule
\end{tabulary}

\end{table}



\texttt{0 · ∞ → NaN}. A NaN operand generally propagates, except the \texttt{*Number}/\allowbreak\texttt{*Finite} extremum variants, which prefer a non-NaN (and, for the \texttt{Finite} variants, finite) operand.



\needspace{16\baselineskip}
\subsection{Ternary (3)}



\label{1836852574804533888}{}



\begin{table}[h]
\centering
\begin{tabulary}{\linewidth}{L L L L L}
\toprule
Operation & Arity & Base spelling & Notes & Teaching page \\
\toprule
\texttt{FMA} & ternary & \texttt{fma}, \texttt{muladd} & one rounding for both Base spellings & Cheat Sheet \\
\texttt{FAA} & ternary & — & fused add-add; no Base counterpart & Julia Compatibility \\
\texttt{Clamp} & ternary & \texttt{clamp} &  & Cheat Sheet \\
\bottomrule
\end{tabulary}

\end{table}



\needspace{15\baselineskip}
\subsection{Conversion}



\label{87850402300235514}{}



\begin{table}[h]
\centering
\begin{tabulary}{\linewidth}{L L L L L}
\toprule
Operation & Arity & Base spelling & Notes & Teaching page \\
\toprule
\texttt{Convert} & conversion & — & the one operation accepting non-\texttt{Binary} operands: \texttt{Binary} (any format), \texttt{Float16/32/64}, \texttt{Float128}, \texttt{Integer}, \texttt{BigFloat} & Values, Code Points \& Conversion \\
\bottomrule
\end{tabulary}

\end{table}



Catalog total: 30 unary + 18 binary + 3 ternary + 1 conversion = 52.



\needspace{18\baselineskip}
\subsection{Base spellings that are componentwise composites}



\label{5840106601692650242}{}


These Base functions are not single draft operations; each is a componentwise tuple of draft operations run under the session default.




\begin{table}[h]
\centering
\begin{tabulary}{\linewidth}{L L}
\toprule
Base spelling & Composite of \\
\toprule
\texttt{sincos} & \texttt{(Sin(x), Cos(x))} \\
\texttt{sincospi} & \texttt{(SinPi(x), CosPi(x))} \\
\texttt{minmax} & \texttt{(\allowbreak{}Minimum(\allowbreak{}x, y), Maximum(\allowbreak{}x, y))} \\
\bottomrule
\end{tabulary}

\end{table}



\needspace{27\baselineskip}
\subsection{Deliberately unmapped draft operations}



\label{14662348581302551618}{}


No Base spelling exists for these; call them by draft name. Base has no counterpart for any row in this list.




\begin{table}[h]
\centering
\begin{tabulary}{\linewidth}{L L}
\toprule
Draft operation & Reason unmapped \\
\toprule
\texttt{RSqrt} & Base has \texttt{inv(sqrt(x))} as two roundings, not one; no single-rounding Base spelling \\
\texttt{Softplus} & no Base equivalent \\
\texttt{ArcSinPi} & Base has only the \texttt{sinpi} family (no inverse) \\
\texttt{ArcCosPi} & Base has only the \texttt{sinpi} family (no inverse) \\
\texttt{ArcTanPi} & Base has only the \texttt{sinpi} family (no inverse) \\
\texttt{ArcTan2Pi} & Base has only the \texttt{sinpi} family (no inverse) \\
\texttt{MaximumNumber}, \texttt{MinimumNumber} & Base has no NaN-ignoring extremum pair \\
\texttt{MaximumMagnitude}, \texttt{MinimumMagnitude} & Base has no magnitude-extremum pair \\
\texttt{MaximumMagnitudeNumber}, \texttt{MinimumMagnitudeNumber} & Base has no NaN-ignoring magnitude-extremum pair \\
\texttt{MinimumFinite}, \texttt{MaximumFinite} & Base has no finite-preferring extremum pair \\
\texttt{FAA} & Base has no fused add-add \\
\bottomrule
\end{tabulary}

\end{table}



The mapping is a declarative partition over the full op list; the test suite asserts it is exhaustive — every operation is either mapped to a Base spelling or listed in \texttt{\_\allowbreak{}NO\_\allowbreak{}BASE\_\allowbreak{}COUNTERPART}.



\subsection{Call shapes}



\label{8684074389003448897}{}


Explicit form, every operation:




\nopagebreak[4]
\begin{minted}{julia}
Op(result_format, ρ, operands...; rng, R)
\end{minted}



\texttt{rng} and \texttt{R} apply only to stochastic \texttt{ρ}.



Convenience form, same-format operands, \texttt{DefaultProjection()}:




\nopagebreak[4]
\begin{minted}{julia}
Op(x...)
\end{minted}



\subsection{see also}



\label{9115667341246479961-dup17}{}


\href{cheatsheet.md}{Cheat Sheet}, \href{explain\_\allowbreak{}no\_\allowbreak{}promotion.md}{Why No Cross-Format Promotion}, \href{tutorial1\_\allowbreak{}values.md}{Values, Code Points \& Conversion}.



\section{Reference: Session Defaults}



\label{17183561231861118727}{}


Dry listing of the seven process-global session defaults, their coherence invariant, the \texttt{with\_\allowbreak{}default\_\allowbreak{}*} combinators, and the cost model. For the hazards of process-global state, see the User Guide{\textquotesingle}s session-defaults section and Julia Compatibility.



\needspace{24\baselineskip}
\subsection{Getter / setter table}



\label{17510285445625325441}{}


Read with \texttt{DefaultX()}, set with \texttt{DefaultX!(v)}.




\begin{table}[h]
\centering
\begin{tabulary}{\linewidth}{L L L L}
\toprule
Default & Initial value & Setter accepts & Teaching page \\
\toprule
\texttt{DefaultType} & \texttt{Binary8p2se} & any fully-parameterized \texttt{Binary} type & Cheat Sheet \\
\texttt{DefaultReturnType} & \texttt{Binary8p2se} & any fully-parameterized \texttt{Binary} type & Cheat Sheet \\
\texttt{DefaultRoundingMode} & \texttt{NearestTiesToEven()} & mode instance or type & Cheat Sheet \\
\texttt{DefaultSaturationMode} & \texttt{SatNone()} & mode instance or type & Cheat Sheet \\
\texttt{DefaultProjection} & \texttt{RNE\_\allowbreak{}SN} & a \texttt{ProjSpec}, or \texttt{(mode, sat)} & Cheat Sheet \\
\texttt{DefaultRNG} & the \texttt{Xoshiro} type & RNG type or instance & Cheat Sheet \\
\texttt{DefaultRbits} & \texttt{8} & \texttt{Int} in \texttt{1:60} & Cheat Sheet \\
\bottomrule
\end{tabulary}

\end{table}




\begin{minted}{julia}
DefaultType()            # Binary8p2se     DefaultType!(Binary8p4se)
DefaultReturnType()      # Binary8p2se     DefaultReturnType!(Binary8p3se)
DefaultRoundingMode()    # NearestTiesToEven()
DefaultSaturationMode()  # SatNone()
DefaultProjection()      # RNE_SN
DefaultRNG()             # Xoshiro         DefaultRNG!(Xoshiro(42))
DefaultRbits()           # 8               DefaultRbits!(16)
\end{minted}



\needspace{21\baselineskip}
\subsection{Coherence invariant}



\label{2615016547748953455}{}



\begin{minted}{julia}
DefaultProjection() === ProjSpec(DefaultRoundingMode(), DefaultSaturationMode())
\end{minted}



This holds at all times, in both directions:




\begin{table}[h]
\centering
\begin{tabulary}{\linewidth}{L L}
\toprule
Action & Effect \\
\toprule
\texttt{DefaultRoundingMode!(\allowbreak{}mode)} & rebuilds \texttt{DefaultProjection} around the new mode, holding saturation fixed \\
\texttt{DefaultSaturationMode!(\allowbreak{}mode)} & rebuilds \texttt{DefaultProjection} around the new mode, holding rounding fixed \\
\texttt{DefaultProjection!(ρ)} & decomposes \texttt{ρ} back into both components \\
\bottomrule
\end{tabulary}

\end{table}



There is no \texttt{DefaultAccumulatorType}: removed after v0.1.0. Reduction accumulator carriers are picked by the span filter (provably exact for the operands), not by a session preference.



\needspace{21\baselineskip}
\subsection{\texttt{with\_\allowbreak{}default\_\allowbreak{}*} combinators}



\label{5124425984651407666}{}



\begin{table}[h]
\centering
\begin{tabulary}{\linewidth}{L L L}
\toprule
Combinator & Signature & Teaching page \\
\toprule
\texttt{with\_\allowbreak{}default\_\allowbreak{}type} & \texttt{with\_\allowbreak{}default\_\allowbreak{}type(\allowbreak{}f, args...)} calls \texttt{f(\allowbreak{}DefaultType(\allowbreak{}), args...)} & Cheat Sheet \\
\texttt{with\_\allowbreak{}default\_\allowbreak{}returntype} & \texttt{with\_\allowbreak{}default\_\allowbreak{}returntype(\allowbreak{}f, args...)} calls \texttt{f(\allowbreak{}DefaultReturnType(\allowbreak{}), args...)} & Cheat Sheet \\
\texttt{with\_\allowbreak{}default\_\allowbreak{}projection} & \texttt{with\_\allowbreak{}default\_\allowbreak{}projection(\allowbreak{}f, args...)} calls \texttt{f(\allowbreak{}DefaultProjection(\allowbreak{}), args...)} & Cheat Sheet \\
\bottomrule
\end{tabulary}

\end{table}




\begin{minted}{julia}
with_default_type((T, x) -> T(x), 1.5)          # Binary8p2se(1.5 ≡ 0x41)
with_default_projection((ρ, x, y) -> Add(T, ρ, x, y), x, y)
with_default_returntype(f, args...)              # f(DefaultReturnType(), args...)
\end{minted}



Each combinator restores nothing itself — it only \emph{reads} the default at call time via \texttt{f(default, args...)}. For scoped mutate-and-restore, the same names also support a do-block form that restores the prior value on exit, even if the body throws.



\subsection{Cost model}



\label{14729244558823917650}{}


\begin{itemize}
\item \textbf{Speculation guard.} Convenience forms (\texttt{x + y}, \texttt{Exp(x)}, \texttt{T(2.1)}) and the \texttt{with\_\allowbreak{}default\_\allowbreak{}*} combinators read the default through a speculation guard: allocation-free and concretely inferred while the default holds its initial value.


\item \textbf{One barrier dispatch after a change.} Once a default is changed, each call crosses a function barrier: one dynamic dispatch at entry, everything inside still fully specialized.


\item \textbf{Explicit forms are unaffected.} \texttt{Add(T, ρ, x, y)} with a \texttt{const} or argument-bound \texttt{ρ} never consults a session default and never pays this cost either way.


\item \textbf{Boxing at escape.} When the combinator{\textquotesingle}s result type \emph{is} the default itself (\texttt{with\_\allowbreak{}default\_\allowbreak{}type} used as a constructor), the value computes on the specialized path but boxes once where it escapes the call — the irreducible cost of a runtime-chosen type. When the result type does not depend on the default, the call is zero-allocation with a concretely inferred result.

\end{itemize}


\subsection{see also}



\label{9115667341246479961-dup18}{}


\href{quickstart.md}{Quickstart}, \href{explain\_\allowbreak{}no\_\allowbreak{}promotion.md}{Why No Cross-Format Promotion}, \href{ref\_\allowbreak{}projections.md}{Projection Specifications}.



\section{Reference: Arrays, Blocks \& Packed Storage}



\label{789702907541812790}{}


Signatures only. For narrative treatment, see the User Guide{\textquotesingle}s Arrays, kernels, and sorting section, and the Blocks and scaled operations section.



\needspace{20\baselineskip}
\subsection{Array operation forms}



\label{1032371883808262493}{}



\begin{table}[h]
\centering
\begin{tabulary}{\linewidth}{L L L}
\toprule
Form & Signature & Teaching page \\
\toprule
Unary array & \texttt{Op(fr, ρ, A)} & Cheat Sheet \\
Binary array & \texttt{Op(fr, ρ, A, B)} & Cheat Sheet \\
Ternary array & \texttt{Op(fr, ρ, A, B, C)} & Cheat Sheet \\
Generic map & \texttt{vmap(\allowbreak{}:Op, fr, ρ, operands...)} & Cheat Sheet \\
Generic map, in-place & \texttt{vmap!(\allowbreak{}C, Val(\allowbreak{}:Op), fr, ρ, operands...)} & Cheat Sheet \\
Callable operation object & \texttt{Op(:Op, fr, ρ)} then \texttt{.(A)} or \texttt{map(\allowbreak{}Op(\allowbreak{}:Op, fr, ρ), A, B)} & User Guide \\
\bottomrule
\end{tabulary}

\end{table}



Operands and destination must have matching axes.



\needspace{16\baselineskip}
\subsection{Table cache introspection}



\label{14157903985024423279}{}



\begin{table}[h]
\centering
\begin{tabulary}{\linewidth}{L L L L}
\toprule
Function & Signature & Notes & Teaching page \\
\toprule
\texttt{table\_\allowbreak{}bytes} & \texttt{table\_\allowbreak{}bytes()} & current cache footprint, bytes & Cheat Sheet \\
\texttt{empty\_\allowbreak{}tables!} & \texttt{empty\_\allowbreak{}tables!()} & reset the table cache & Cheat Sheet \\
\bottomrule
\end{tabulary}

\end{table}



Deterministic unary/binary array calls use cached result tables (256 B unary, 64 KiB binary). Ternary tables tier by combined operand bitwidth: eager (≤ 18 bits total), adaptive with LRU eviction (≤ 21 bits total), compute per element at \texttt{K = 8}. Stochastic operations always compute per element and consume one draw per projection.



\needspace{25\baselineskip}
\subsection{Block functions}



\label{5661144322417720737}{}



\begin{table}[h]
\centering
\begin{tabulary}{\linewidth}{L L L}
\toprule
Function & Signature & Teaching page \\
\toprule
\texttt{blocksize} & \texttt{blocksize(b)} & Cheat Sheet \\
\texttt{scaleformat} & \texttt{scaleformat(b)} & Cheat Sheet \\
\texttt{elemformat} & \texttt{elemformat(b)} & Cheat Sheet \\
\texttt{ConvertFromBlock} & \texttt{ConvertFromBlock(\allowbreak{}fr, ρ, b)} & Cheat Sheet \\
\texttt{ConvertToBlock} & \texttt{ConvertToBlock(\allowbreak{}fs, fr, ρ, values\_\allowbreak{}tuple, scale)} & Cheat Sheet \\
\texttt{ConvertToBlockMaxAbsFinite} & \texttt{ConvertToBlockMaxAbsFinite(\allowbreak{}fs, fr, scale\_\allowbreak{}ρ, element\_\allowbreak{}ρ, values\_\allowbreak{}tuple)} & Cheat Sheet \\
\texttt{BlockReduceAdd} & \texttt{BlockReduceAdd(fr, ρ, b)} & Cheat Sheet \\
\texttt{BlockReduceMultiply} & \texttt{BlockReduceMultiply(\allowbreak{}fr, ρ, b)} & Cheat Sheet \\
\texttt{BlockDotProduct} & \texttt{BlockDotProduct(\allowbreak{}fr, ρ, bx, by)} & Cheat Sheet \\
\bottomrule
\end{tabulary}

\end{table}



\texttt{BlockReduceAdd}, \texttt{BlockReduceMultiply}, and \texttt{BlockDotProduct} perform their lane products and accumulation exactly, projecting once at the end — no hidden intermediate rounding.



\needspace{18\baselineskip}
\subsubsection{Generated \texttt{BlockOp} / \texttt{ScaledOp} pattern}



\label{13411109538340054149}{}


Every scalar operation except \texttt{Convert} has a generated block form and a generated scaled form.




\begin{table}[h]
\centering
\begin{tabulary}{\linewidth}{L L}
\toprule
Pattern & Signature \\
\toprule
\texttt{BlockOp} & \texttt{BlockOp(\allowbreak{}fr, ρ, b1, b2, result\_\allowbreak{}scale)} (arity follows the underlying \texttt{Op}) \\
\texttt{ScaledOp} & \texttt{ScaledOp(\allowbreak{}fr, ρ, scale1, x1, scale2, x2, ...)} (block size 1; arity follows the underlying \texttt{Op}) \\
\bottomrule
\end{tabulary}

\end{table}



Examples:




\nopagebreak[4]
\begin{minted}{julia}
BlockAdd(T, ρ, b1, b2, result_scale)
BlockExp(T, ρ, b, result_scale)

ScaledAdd(T, ρ, scale1, x1, scale2, x2)
ScaledExp(T, ρ, scale, x)
\end{minted}



\subsection{\texttt{BlockVector}}



\label{13615094720265409582-dup1}{}


Stores many equal-shape blocks in a structure-of-arrays layout. No additional scalar API beyond the \texttt{Block} operations above; construct from a collection of same-shape \texttt{Block}s.



\needspace{20\baselineskip}
\subsection{\texttt{PackedVector}}



\label{17791294993201593296}{}



\begin{table}[h]
\centering
\begin{tabulary}{\linewidth}{L L L}
\toprule
Operation & Signature & Notes \\
\toprule
Construct & \texttt{PackedVector(v)} & \texttt{v::AbstractVector\{F\}} \\
Index (read) & \texttt{pv[i]} & returns an \texttt{F} \\
Index (write) & \texttt{pv[i] = x} & \texttt{x::F} \\
Collect & \texttt{collect(pv)} & materializes a \texttt{Vector\{F\}} \\
\texttt{sizeof} rule & \texttt{sizeof(pv.data)} & \texttt{ceil(\allowbreak{}length(\allowbreak{}pv) * bitwidth(\allowbreak{}F) /\allowbreak{} 8)} bytes, rounded to the storage word \\
\texttt{vmap} acceptance & \texttt{vmap(:Op, F, ρ, pv)} & accepted directly; tiles unpack internally \\
\bottomrule
\end{tabulary}

\end{table}



\texttt{PackedVector} stores each code point in exactly \texttt{bitwidth(F)} bits. Computation unpacks tiles internally; packed arithmetic is deliberately not in-place — store packed, compute unpacked.



\subsection{see also}



\label{9115667341246479961-dup19}{}


\href{cheatsheet.md}{Cheat Sheet}, \href{tutorial3\_\allowbreak{}arrays.md}{Operations \& Arrays}, \href{howto\_\allowbreak{}blocks\_\allowbreak{}dynamic\_\allowbreak{}range.md}{Use Blocks to Track Dynamic Range}.



\section{Reference: Conformance \& Approximation Registry}



\label{6802684063430148207}{}


Dry listing of the conformance-declaration functions and the κ-approximation registry. For narrative explanation, see the User Guide{\textquotesingle}s Conformance and κ-approximate implementations section.



\needspace{21\baselineskip}
\subsection{Conformance functions}



\label{8429749652979037907}{}



\begin{table}[h]
\centering
\begin{tabulary}{\linewidth}{L L L L}
\toprule
Function & Signature & Returns & Teaching page \\
\toprule
\texttt{conformance} & \texttt{conformance()} & the live declaration: formats, operations, mode vocabulary, table specializations instantiated this session, registered approximations & User Guide \\
\texttt{conformance\_\allowbreak{}dict} & \texttt{conformance\_\allowbreak{}dict()} & a plain nested \texttt{Dict}, for JSON/TOML serialization & User Guide \\
\texttt{conformance\_\allowbreak{}report} & \texttt{conformance\_\allowbreak{}report()} & prints the declaration & User Guide \\
\texttt{draft\_\allowbreak{}revision} & \texttt{draft\_\allowbreak{}revision()} & the P3109 draft revision this package conforms to & User Guide \\
\bottomrule
\end{tabulary}

\end{table}



\needspace{26\baselineskip}
\subsection{κ registry}



\label{958257881295565054}{}



\begin{table}[h]
\centering
\begin{tabulary}{\linewidth}{L L L L}
\toprule
Function & Signature & Returns & Teaching page \\
\toprule
\texttt{measure\_\allowbreak{}kappa} & \texttt{measure\_\allowbreak{}kappa(\allowbreak{}fn, op, fr, (\allowbreak{}argformats...), ρ)} & \texttt{(κ, exhaustive)} — measured max code-point deviation and whether the measurement was exhaustive & User Guide \\
\texttt{register\_\allowbreak{}approx!} & \texttt{register\_\allowbreak{}approx!(\allowbreak{}name, op, fr, args, ρ, fn; κ)} & registers \texttt{fn} as approximation \texttt{name} for \texttt{op}, declaring bound \texttt{κ} & User Guide \\
\texttt{approx} & \texttt{approx(name)} & the registered implementation object & User Guide \\
\texttt{kappa} & \texttt{kappa(impl)} & the declared κ bound & User Guide \\
\texttt{kappa\_\allowbreak{}measured} & \texttt{kappa\_\allowbreak{}measured(impl)} & the measured κ from registration & User Guide \\
\texttt{list\_\allowbreak{}approx} & \texttt{list\_\allowbreak{}approx()} & all registered approximation names & User Guide \\
\texttt{unregister\_\allowbreak{}approx!} & \texttt{unregister\_\allowbreak{}approx!(\allowbreak{}name)} & removes a registered approximation & User Guide \\
\bottomrule
\end{tabulary}

\end{table}




\begin{minted}{julia}
κ, exhaustive = measure_kappa(fn, :Exp, T, (T,), ρ)
register_approx!(:fast_exp, :Exp, T, (T,), ρ, fn; κ)
impl = approx(:fast_exp)
kappa(impl)
kappa_measured(impl)
list_approx()
unregister_approx!(:fast_exp)
\end{minted}



\subsection{Understatement-rejection rule}



\label{10633883490468249847}{}


κ is measured exhaustively at registration time whenever the domain permits exhaustive enumeration. Declaring a κ smaller than the measured deviation throws — an implementation cannot register itself as more accurate than it measures.



\subsection{κ = NaN acknowledgment rule}



\label{17355671808123045861}{}


An implementation whose NaN behavior differs from the default API{\textquotesingle}s (a NaN mismatch against the exhaustive reference) must be registered with an explicit \texttt{κ = NaN}. This is an acknowledgment, not a numeric bound: it cannot be produced by measurement alone and must be supplied by the caller of \texttt{register\_\allowbreak{}approx!}.



Approximate implementations registered through this path are never substituted into the default API; the default API is bit-exact, always.



\subsection{see also}



\label{9115667341246479961-dup20}{}


\href{howto\_\allowbreak{}conformance.md}{Read \& Export the Conformance Declaration}, \href{under\_\allowbreak{}verification.md}{Verification \& Benchmark Doctrine}.



\section{External Reference}



\label{5597696625397351643}{}


Docstrings for the documented \textbf{public} surface — the exported names. The full export list is far larger (400+ names, including the 120 exported format aliases — the other 384 are opt-in via \texttt{using SmallFloats.Formats} — the predefined projection-spec grid, and the generated operation and block registers); names without docstrings are covered descriptively in the User Guide; unexported machinery is in the Reference: Internal (docstrings) page.



The index below groups the exported surface by family and points each family at the descriptive reference page that documents its full table. The \texttt{@autodocs} block that follows renders the docstrings themselves, alphabetically, as generated from source.



\needspace{31\baselineskip}
\subsection{Index}



\label{6663683553518785561}{}



\begin{table}[h]
\centering
\begin{tabulary}{\linewidth}{L L L}
\toprule
Family & Scope & Descriptive reference page \\
\toprule
Arrays, blocks, packed storage & \texttt{Op(fr,ρ,A,...)}, \texttt{vmap}/\allowbreak\texttt{vmap!}, \texttt{Block}, \texttt{BlockVector}, \texttt{PackedVector}, \texttt{table\_\allowbreak{}bytes}, \texttt{empty\_\allowbreak{}tables!} & Reference: Arrays, Blocks \& Packed Storage \\
Conformance and approximation registry & \texttt{conformance}, \texttt{conformance\_\allowbreak{}dict}, \texttt{conformance\_\allowbreak{}report}, \texttt{draft\_\allowbreak{}revision}, \texttt{measure\_\allowbreak{}kappa}, \texttt{register\_\allowbreak{}approx!}, \texttt{approx}, \texttt{kappa}, \texttt{kappa\_\allowbreak{}measured}, \texttt{list\_\allowbreak{}approx}, \texttt{unregister\_\allowbreak{}approx!} & Reference: Conformance \& Approximation Registry \\
Formats & 504 format aliases (120 exported at \texttt{K ≤ 8}), \texttt{format}, \texttt{formatname}, and the format-query accessors & Reference: Format Names \& Queries \\
Operations & the 52-operation catalog: 30 unary, 18 binary, 3 ternary, \texttt{Convert} & Reference: Operation Catalog \\
Projection specs & \texttt{ProjSpec}, the 6×3 deterministic grid, the \texttt{RSA}/\allowbreak\texttt{RSB}/\texttt{RSC} stochastic constructors, \texttt{roundingmode}, \texttt{saturationmode}, \texttt{isstochastic}, \texttt{nrandbits} & Reference: Projection Specifications \\
Session defaults & \texttt{DefaultType}, \texttt{DefaultReturnType}, \texttt{DefaultRoundingMode}, \texttt{DefaultSaturationMode}, \texttt{DefaultProjection}, \texttt{DefaultRNG}, \texttt{DefaultRbits}, and their setters and \texttt{with\_\allowbreak{}default\_\allowbreak{}*} combinators & Reference: Session Defaults \\
\bottomrule
\end{tabulary}

\end{table}


\par\medskip\noindent\begin{minipage}{\linewidth}
\hypertarget{1538304370509133722}{\texttt{SmallFloats.SmallFloats}}  -- {Module.}
\nopagebreak
\begin{adjustwidth}{2em}{0pt}


\begin{minted}{julia}
SmallFloats
\end{minted}

A conforming, performance-oriented Julia implementation of the IEEE P3109 draft standard for arithmetic formats for machine learning (bitwidths 3–16).

Bit-exact defined results on every default path; the projection engine (\texttt{RoundToPrecision → Saturate → Encode}) is the single write path into a code point; approximate fast paths exist only behind the explicit κ registry (\texttt{register\_\allowbreak{}approx!} / \texttt{approx}), never substituted silently. See \texttt{conformance()} for the live declaration and \texttt{SmallFloats.draft\_\allowbreak{}revision(\allowbreak{})} for the implemented draft.

\textbf{Performance note}

Pass format types through \texttt{const} bindings, type parameters, or function arguments. All exported entry points fully specialize when the format type is statically known (measured: a complete scalar \texttt{Add} ≈ 26 ns, \texttt{project} ≈ 13 ns, zero allocations). Calling through a \textbf{non-\texttt{const} global} format type forces dynamic dispatch on every call — measured at {\textasciitilde}1 µs per scalar keyword call — which is Julia semantics, not package overhead; a single function barrier (\texttt{f(::Type\{T\}, args...) where \{T\}}) restores full speed. The test suite pins the specialization properties (concrete return types, zero warm-path allocation) as deterministic regressions.



\end{adjustwidth}
\end{minipage}
\par\medskip

\par\medskip\noindent\begin{minipage}{\linewidth}
\hypertarget{1153421143103638867}{\texttt{SmallFloats.RNA\_\allowbreak{}SF}}  -- {Constant.}
\nopagebreak
\begin{adjustwidth}{2em}{0pt}

(NearestTiesToAway, SatFinite).



\end{adjustwidth}
\end{minipage}
\par\medskip

\par\medskip\noindent\begin{minipage}{\linewidth}
\hypertarget{1497742833584286063}{\texttt{SmallFloats.RNA\_\allowbreak{}SN}}  -- {Constant.}
\nopagebreak
\begin{adjustwidth}{2em}{0pt}

(NearestTiesToAway, SatNone).



\end{adjustwidth}
\end{minipage}
\par\medskip

\par\medskip\noindent\begin{minipage}{\linewidth}
\hypertarget{9868279548188878236}{\texttt{SmallFloats.RNA\_\allowbreak{}SP}}  -- {Constant.}
\nopagebreak
\begin{adjustwidth}{2em}{0pt}

(NearestTiesToAway, SatPropagate).



\end{adjustwidth}
\end{minipage}
\par\medskip

\par\medskip\noindent\begin{minipage}{\linewidth}
\hypertarget{3026972477457364217}{\texttt{SmallFloats.RNE\_\allowbreak{}SF}}  -- {Constant.}
\nopagebreak
\begin{adjustwidth}{2em}{0pt}

(NearestTiesToEven, SatFinite).



\end{adjustwidth}
\end{minipage}
\par\medskip

\par\medskip\noindent\begin{minipage}{\linewidth}
\hypertarget{10331688225281198727}{\texttt{SmallFloats.RNE\_\allowbreak{}SN}}  -- {Constant.}
\nopagebreak
\begin{adjustwidth}{2em}{0pt}

(NearestTiesToEven, SatNone) — the package-wide default ρ.



\end{adjustwidth}
\end{minipage}
\par\medskip

\par\medskip\noindent\begin{minipage}{\linewidth}
\hypertarget{4678886371087216428}{\texttt{SmallFloats.RNE\_\allowbreak{}SP}}  -- {Constant.}
\nopagebreak
\begin{adjustwidth}{2em}{0pt}

(NearestTiesToEven, SatPropagate).



\end{adjustwidth}
\end{minipage}
\par\medskip

\par\medskip\noindent\begin{minipage}{\linewidth}
\hypertarget{5549571836457789421}{\texttt{SmallFloats.RTN\_\allowbreak{}SF}}  -- {Constant.}
\nopagebreak
\begin{adjustwidth}{2em}{0pt}

(TowardNegative, SatFinite).



\end{adjustwidth}
\end{minipage}
\par\medskip

\par\medskip\noindent\begin{minipage}{\linewidth}
\hypertarget{16640126877541699580}{\texttt{SmallFloats.RTN\_\allowbreak{}SN}}  -- {Constant.}
\nopagebreak
\begin{adjustwidth}{2em}{0pt}

(TowardNegative, SatNone).



\end{adjustwidth}
\end{minipage}
\par\medskip

\par\medskip\noindent\begin{minipage}{\linewidth}
\hypertarget{4658287144677563280}{\texttt{SmallFloats.RTN\_\allowbreak{}SP}}  -- {Constant.}
\nopagebreak
\begin{adjustwidth}{2em}{0pt}

(TowardNegative, SatPropagate).



\end{adjustwidth}
\end{minipage}
\par\medskip

\par\medskip\noindent\begin{minipage}{\linewidth}
\hypertarget{17370994210103718270}{\texttt{SmallFloats.RTO\_\allowbreak{}SF}}  -- {Constant.}
\nopagebreak
\begin{adjustwidth}{2em}{0pt}

(ToOdd, SatFinite).



\end{adjustwidth}
\end{minipage}
\par\medskip

\par\medskip\noindent\begin{minipage}{\linewidth}
\hypertarget{8851151326547106986}{\texttt{SmallFloats.RTO\_\allowbreak{}SN}}  -- {Constant.}
\nopagebreak
\begin{adjustwidth}{2em}{0pt}

(ToOdd, SatNone).



\end{adjustwidth}
\end{minipage}
\par\medskip

\par\medskip\noindent\begin{minipage}{\linewidth}
\hypertarget{6365993733773253093}{\texttt{SmallFloats.RTO\_\allowbreak{}SP}}  -- {Constant.}
\nopagebreak
\begin{adjustwidth}{2em}{0pt}

(ToOdd, SatPropagate).



\end{adjustwidth}
\end{minipage}
\par\medskip

\par\medskip\noindent\begin{minipage}{\linewidth}
\hypertarget{10475668134315902188}{\texttt{SmallFloats.RTP\_\allowbreak{}SF}}  -- {Constant.}
\nopagebreak
\begin{adjustwidth}{2em}{0pt}

(TowardPositive, SatFinite).



\end{adjustwidth}
\end{minipage}
\par\medskip

\par\medskip\noindent\begin{minipage}{\linewidth}
\hypertarget{1691628013094057374}{\texttt{SmallFloats.RTP\_\allowbreak{}SN}}  -- {Constant.}
\nopagebreak
\begin{adjustwidth}{2em}{0pt}

(TowardPositive, SatNone).



\end{adjustwidth}
\end{minipage}
\par\medskip

\par\medskip\noindent\begin{minipage}{\linewidth}
\hypertarget{1355188992532382511}{\texttt{SmallFloats.RTP\_\allowbreak{}SP}}  -- {Constant.}
\nopagebreak
\begin{adjustwidth}{2em}{0pt}

(TowardPositive, SatPropagate).



\end{adjustwidth}
\end{minipage}
\par\medskip

\par\medskip\noindent\begin{minipage}{\linewidth}
\hypertarget{16398804727682363821}{\texttt{SmallFloats.RTZ\_\allowbreak{}SF}}  -- {Constant.}
\nopagebreak
\begin{adjustwidth}{2em}{0pt}

(TowardZero, SatFinite).



\end{adjustwidth}
\end{minipage}
\par\medskip

\par\medskip\noindent\begin{minipage}{\linewidth}
\hypertarget{18041415609987892616}{\texttt{SmallFloats.RTZ\_\allowbreak{}SN}}  -- {Constant.}
\nopagebreak
\begin{adjustwidth}{2em}{0pt}

(TowardZero, SatNone).



\end{adjustwidth}
\end{minipage}
\par\medskip

\par\medskip\noindent\begin{minipage}{\linewidth}
\hypertarget{7796365331314406520}{\texttt{SmallFloats.RTZ\_\allowbreak{}SP}}  -- {Constant.}
\nopagebreak
\begin{adjustwidth}{2em}{0pt}

(TowardZero, SatPropagate).



\end{adjustwidth}
\end{minipage}
\par\medskip

\par\medskip\noindent\begin{minipage}{\linewidth}
\hypertarget{17988065543319926718}{\texttt{SmallFloats.ApproxImpl}}  -- {Type.}
\nopagebreak
\begin{adjustwidth}{2em}{0pt}

A registered κ-approximate implementation of one operation specialization.



\end{adjustwidth}
\end{minipage}
\par\medskip

\par\medskip\noindent\begin{minipage}{\linewidth}
\hypertarget{10793991096487538164}{\texttt{SmallFloats.Binary}}  -- {Type.}
\nopagebreak
\begin{adjustwidth}{2em}{0pt}


\begin{minted}{julia}
Binary{K,P,SGN,EXT} <: AbstractFloat
\end{minted}

A P3109 \textbf{format}: bitwidth \texttt{K ∈ 3:16}, precision \texttt{P}, signedness \texttt{SGN::Bool} (\texttt{true} = Signed), domain \texttt{EXT::Bool} (\texttt{true} = Extended, i.e. the datum set includes infinities).

\texttt{Binary} is \textbf{abstract}, and that is the draft{\textquotesingle}s own distinction: a format is {\textquotedbl}a datum set and an encoding{\textquotedbl} (D1 §3.1), while the width of the machine word carrying a code point is a representation choice below the standard{\textquotesingle}s level of description. \texttt{Code8} is that representation. The parameterized name is the \emph{format}; the draft-named alias is its \emph{representation}:


\nopagebreak[4]
\begin{minted}{julia}
Binary{8,4,true,true}   # the format  — abstract
Binary8p4se             # Code8{8,4,true,true}, its representation — concrete
\end{minted}

Consequently \texttt{Binary8p4se === Binary\{8,4,true,true\}} is \textbf{\texttt{false}} and \texttt{Binary8p4se <: Binary\{8,4,true,true\}} is \textbf{\texttt{true}}. Use the alias, or \texttt{format(K, P, Σ, Δ)}, wherever a concrete type is needed — an array element type, a \texttt{similar}, a constructor target.

The code point occupies the low \texttt{K} bits of the storage unit; the high bits are maintained zero as a representation invariant. Construct raw code points with \texttt{T(c::Unsigned)} (validated code-point construction), \texttt{rawvalue(T, c)} (unchecked kernel route), or \texttt{Code8\{...\}(Val(:code), x)}; construct from numeric values with \texttt{T(x::Real)} (default projection spec) or \texttt{Convert}. \texttt{Unsigned} is the one argument-type class meaning \emph{code point}, at \textbf{every} width — \texttt{T(0x02)} is code point 2 whether \texttt{T}{\textquotesingle}s unit is a \texttt{UInt8} or a \texttt{UInt16}; all other \texttt{Real}s mean \emph{value}.

There are \texttt{4K − 2} formats at each \texttt{K} (\texttt{P ∈ 1:K} × signed/unsigned × finite/extended, less the two \texttt{P == K} signed combinations, which have no exponent bits), so \texttt{KMIN:KMAX} = 504 in total.



\end{adjustwidth}
\end{minipage}
\par\medskip

\par\medskip\noindent\begin{minipage}{\linewidth}
\hypertarget{16226061866935425277}{\texttt{SmallFloats.Block}}  -- {Type.}
\nopagebreak
\begin{adjustwidth}{2em}{0pt}


\begin{minted}{julia}
Block{B, FS<:Binary, FE<:Binary}
\end{minted}

A P3109 block \texttt{(s, [x₁ … x\_\allowbreak{}B])}: scale factor in format \texttt{FS}, \texttt{B ≥ 1} elements in format \texttt{FE} (draft §5). \texttt{isbits}; blocks live in registers/stack.



\end{adjustwidth}
\end{minipage}
\par\medskip

\par\medskip\noindent\begin{minipage}{\linewidth}
\hypertarget{2816898984109632562}{\texttt{SmallFloats.BlockVector}}  -- {Type.}
\nopagebreak
\begin{adjustwidth}{2em}{0pt}


\begin{minted}{julia}
BlockVector{B,FS,FE} <: AbstractVector{Block{B,FS,FE}}
\end{minted}

Structure-of-arrays storage: \texttt{scales::Vector\{FS\}}, \texttt{elems::Matrix\{FE\}} (B × n, column-major so each block{\textquotesingle}s elements are one contiguous column).



\end{adjustwidth}
\end{minipage}
\par\medskip

\par\medskip\noindent\begin{minipage}{\linewidth}
\hypertarget{5292295670398946918}{\texttt{SmallFloats.FPClass}}  -- {Type.}
\nopagebreak
\begin{adjustwidth}{2em}{0pt}

Eight-way datum classification returned by \texttt{Class} (draft §4.13.1).



\end{adjustwidth}
\end{minipage}
\par\medskip

\par\medskip\noindent\begin{minipage}{\linewidth}
\hypertarget{4095773949561946178}{\texttt{SmallFloats.NearestTiesToAway}}  -- {Type.}
\nopagebreak
\begin{adjustwidth}{2em}{0pt}

Round to nearest; ties away from zero (draft §4.7.1).



\end{adjustwidth}
\end{minipage}
\par\medskip

\par\medskip\noindent\begin{minipage}{\linewidth}
\hypertarget{13007042087745522853}{\texttt{SmallFloats.NearestTiesToEven}}  -- {Type.}
\nopagebreak
\begin{adjustwidth}{2em}{0pt}

Round to nearest; ties to the even code point (IEEE default; draft §4.7.1).



\end{adjustwidth}
\end{minipage}
\par\medskip

\par\medskip\noindent\begin{minipage}{\linewidth}
\hypertarget{9195041457639280132}{\texttt{SmallFloats.Op}}  -- {Type.}
\nopagebreak
\begin{adjustwidth}{2em}{0pt}


\begin{minted}{julia}
Op(name::Symbol, FR::Type{<:Binary}, ρ::ProjSpec) -> Op
\end{minted}

A registry operation bound to a result format and a projection, as a \textbf{callable} — so the idiomatic array verbs work directly:


\nopagebreak[4]
\begin{minted}{julia}
map(Op(:Add, Binary8p4se, RNE_SN), A, B)
Op(:Exp, Binary8p4se, RNE_SN).(A)
\end{minted}

\texttt{Op} is a singleton type: the operation, format and projection all live in the type parameters, so calls specialize exactly as \texttt{Add(\allowbreak{}Binary8p4se, RNE\_\allowbreak{}SN, x, y)} does. Use \hyperlinkref{5240195886209104356}{\texttt{vmap!}} when you have a destination array to fill — it is the in-place, table-aware kernel and this type is a thin front for it.



\end{adjustwidth}
\end{minipage}
\par\medskip

\par\medskip\noindent\begin{minipage}{\linewidth}
\hypertarget{4163091736558650519}{\texttt{SmallFloats.PackedVector}}  -- {Type.}
\nopagebreak
\begin{adjustwidth}{2em}{0pt}


\begin{minted}{julia}
PackedVector{F<:Binary} <: AbstractVector{F}
\end{minted}

Bit-packed storage of \texttt{F} code points at \texttt{bitwidth(F)} bits per element. Construct with \texttt{PackedVector(A::AbstractVector\{F\})}; recover with \texttt{Vector(pv)} or \texttt{collect(pv)}. Memory: \texttt{⌈n·K/64⌉} words vs \texttt{n·sizeof(\allowbreak{}codeunit\_\allowbreak{}type(\allowbreak{}F))} unpacked.

\textbf{It saves nothing when \texttt{K == 8·sizeof(\allowbreak{}codeunit\_\allowbreak{}type(\allowbreak{}F))}} — K = 8 and K = 16 — where the packed layout is bit-for-bit the unpacked one and every access pays the shift/mask/splice for no return. \texttt{packing\_\allowbreak{}saves(F)} says so; see its docstring for why that is a query rather than an error.



\end{adjustwidth}
\end{minipage}
\par\medskip

\par\medskip\noindent\begin{minipage}{\linewidth}
\hypertarget{14580192752130828046}{\texttt{SmallFloats.ProjSpec}}  -- {Type.}
\nopagebreak
\begin{adjustwidth}{2em}{0pt}


\begin{minted}{julia}
ProjSpec{R<:RoundingMode3109, S<:SaturationMode}
\end{minted}

A projection specification ρ = (rounding mode, saturation mode), draft §4.2. Zero-size; construct as \texttt{ProjSpec(\allowbreak{}NearestTiesToEven(\allowbreak{}), SatNone(\allowbreak{}))} or via the exported constants. Kernels specialize on its type.



\end{adjustwidth}
\end{minipage}
\par\medskip

\par\medskip\noindent\begin{minipage}{\linewidth}
\hypertarget{8332080681071361269}{\texttt{SmallFloats.StochasticA}}  -- {Type.}
\nopagebreak
\begin{adjustwidth}{2em}{0pt}

Stochastic rounding, variant A of draft §4.7.4: RoundAway ⟺ ⌊ν·2{\textasciicircum}N⌋ + R ≥ 2{\textasciicircum}N.



\end{adjustwidth}
\end{minipage}
\par\medskip

\par\medskip\noindent\begin{minipage}{\linewidth}
\hypertarget{7252621162593156278}{\texttt{SmallFloats.StochasticB}}  -- {Type.}
\nopagebreak
\begin{adjustwidth}{2em}{0pt}

Stochastic rounding, variant B: RoundAway ⟺ ⌊ν·2{\textasciicircum}(N+1)⌋ + (2R+1) ≥ 2{\textasciicircum}(N+1).



\end{adjustwidth}
\end{minipage}
\par\medskip

\par\medskip\noindent\begin{minipage}{\linewidth}
\hypertarget{11628007139636489123}{\texttt{SmallFloats.StochasticC}}  -- {Type.}
\nopagebreak
\begin{adjustwidth}{2em}{0pt}

Stochastic rounding, variant C: RoundAway ⟺ RNITE(ν·2{\textasciicircum}N) + R ≥ 2{\textasciicircum}N.



\end{adjustwidth}
\end{minipage}
\par\medskip

\par\medskip\noindent\begin{minipage}{\linewidth}
\hypertarget{18226116418678579012}{\texttt{SmallFloats.ToOdd}}  -- {Type.}
\nopagebreak
\begin{adjustwidth}{2em}{0pt}

Round inexact results to the nearest \emph{odd} code point (draft §4.7.3; sticky-friendly).



\end{adjustwidth}
\end{minipage}
\par\medskip

\par\medskip\noindent\begin{minipage}{\linewidth}
\hypertarget{11984278085027763854}{\texttt{SmallFloats.TowardNegative}}  -- {Type.}
\nopagebreak
\begin{adjustwidth}{2em}{0pt}

Directed rounding toward −∞ (draft §4.7.2).



\end{adjustwidth}
\end{minipage}
\par\medskip

\par\medskip\noindent\begin{minipage}{\linewidth}
\hypertarget{14592637289702736603}{\texttt{SmallFloats.TowardPositive}}  -- {Type.}
\nopagebreak
\begin{adjustwidth}{2em}{0pt}

Directed rounding toward +∞ (draft §4.7.2).



\end{adjustwidth}
\end{minipage}
\par\medskip

\par\medskip\noindent\begin{minipage}{\linewidth}
\hypertarget{14300163576002667725}{\texttt{SmallFloats.TowardZero}}  -- {Type.}
\nopagebreak
\begin{adjustwidth}{2em}{0pt}

Directed rounding toward zero — truncation (draft §4.7.2).



\end{adjustwidth}
\end{minipage}
\par\medskip

\par\medskip\noindent\begin{minipage}{\linewidth}
\hypertarget{14325508049649243681}{\texttt{SmallFloats.BlockDotProduct}}  -- {Method.}
\nopagebreak
\begin{adjustwidth}{2em}{0pt}

BlockDotProduct (draft §5.3.2). Lane products can carry 64 significant bits, so the products and their sum are formed in the exact accumulator, with the ∞/NaN fold algebra resolved on the Float64 classifications first.



\end{adjustwidth}
\end{minipage}
\par\medskip

\par\medskip\noindent\begin{minipage}{\linewidth}
\hypertarget{13767263915777893890}{\texttt{SmallFloats.BlockReduceAdd}}  -- {Method.}
\nopagebreak
\begin{adjustwidth}{2em}{0pt}

BlockReduceAdd (draft §5.3.1): project(reduce(ωAdd, [0, X…])).



\end{adjustwidth}
\end{minipage}
\par\medskip

\par\medskip\noindent\begin{minipage}{\linewidth}
\hypertarget{17748774098335709824}{\texttt{SmallFloats.BlockReduceMultiply}}  -- {Method.}
\nopagebreak
\begin{adjustwidth}{2em}{0pt}

BlockReduceMultiply (draft §5.3.1): project(reduce(ωMultiply, [1, X…])).



\end{adjustwidth}
\end{minipage}
\par\medskip

\par\medskip\noindent\begin{minipage}{\linewidth}
\hypertarget{17727979567356838264}{\texttt{SmallFloats.Class}}  -- {Method.}
\nopagebreak
\begin{adjustwidth}{2em}{0pt}

Class(v) -> FPClass (draft §4.13.1): classify \texttt{v} as NaN, ±Inf, ±normal, ±subnormal, or zero.



\end{adjustwidth}
\end{minipage}
\par\medskip

\par\medskip\noindent\begin{minipage}{\linewidth}
\hypertarget{829764723835457682}{\texttt{SmallFloats.Convert}}  -- {Method.}
\nopagebreak
\begin{adjustwidth}{2em}{0pt}


\begin{minted}{julia}
Convert(fr, ρ, A::AbstractArray{<:Union{Float16,Float32,Float64,BFloat16}}; rng) -> Array{fr}
\end{minted}

Bulk ingestion: exact widening per element, then projection under ρ. Unlike the Binary-array Convert (a Shape-A table gather), external inputs are not enumerable, so this is a Shape-B loop; stochastic ρ draws per element.



\end{adjustwidth}
\end{minipage}
\par\medskip

\par\medskip\noindent\begin{minipage}{\linewidth}
\hypertarget{2118065800926972834}{\texttt{SmallFloats.Convert}}  -- {Method.}
\nopagebreak
\begin{adjustwidth}{2em}{0pt}


\begin{minted}{julia}
Convert(fr, ρ, x) -> fr
\end{minted}

Draft §4.9 Convert⟨f\emph{x, f}r, ρ⟩. Accepts \texttt{Binary}, IEEE binary16/32/64 floats (widened exactly to the Float64 carrier), \texttt{Float128} (projected directly), \texttt{Integer} (exact via a sufficiently wide BigFloat), and \texttt{BigFloat} (projected directly; the caller warrants the value is exact).



\end{adjustwidth}
\end{minipage}
\par\medskip

\par\medskip\noindent\begin{minipage}{\linewidth}
\hypertarget{15995489740638035986}{\texttt{SmallFloats.ConvertFromBlock}}  -- {Method.}
\nopagebreak
\begin{adjustwidth}{2em}{0pt}

ConvertFromBlock (§5.2.1): blockdecode, then project each element (no scale division).



\end{adjustwidth}
\end{minipage}
\par\medskip

\par\medskip\noindent\begin{minipage}{\linewidth}
\hypertarget{11399645714333645993}{\texttt{SmallFloats.ConvertToBlock}}  -- {Method.}
\nopagebreak
\begin{adjustwidth}{2em}{0pt}

ConvertToBlock (§5.2.2): decode elements, BlockProject against the supplied scale.



\end{adjustwidth}
\end{minipage}
\par\medskip

\par\medskip\noindent\begin{minipage}{\linewidth}
\hypertarget{9309677842871381966}{\texttt{SmallFloats.ConvertToBlockMaxAbsFinite}}  -- {Method.}
\nopagebreak
\begin{adjustwidth}{2em}{0pt}

ConvertToBlockMaxAbsFinite (§5.2.3): S = reduce(ωMaximumFinite, [NaN, |X₁|…]), project S under ρs, then BlockProject the elements under ρ.



\end{adjustwidth}
\end{minipage}
\par\medskip

\par\medskip\noindent\begin{minipage}{\linewidth}
\hypertarget{17118068279635247347}{\texttt{SmallFloats.DefaultProjection}}  -- {Method.}
\nopagebreak
\begin{adjustwidth}{2em}{0pt}


\begin{minted}{julia}
DefaultProjection() -> ProjSpec
DefaultProjection!(ρ::ProjSpec) -> ρ
DefaultProjection!(m, s) -> ProjSpec
\end{minted}

The session{\textquotesingle}s default projection specification. Initialized to \texttt{(\allowbreak{}DefaultRoundingMode(\allowbreak{}), DefaultSaturationMode(\allowbreak{}))} = \texttt{RNE\_\allowbreak{}SN}. Setting it directly decomposes ρ into \hyperlinkref{2306866772383893895}{\texttt{DefaultRoundingMode}} and \hyperlinkref{1523060637620834303}{\texttt{DefaultSaturationMode}}, so the three stay coherent in both directions.



\end{adjustwidth}
\end{minipage}
\par\medskip

\par\medskip\noindent\begin{minipage}{\linewidth}
\hypertarget{2737063410370640279}{\texttt{SmallFloats.DefaultRNG}}  -- {Method.}
\nopagebreak
\begin{adjustwidth}{2em}{0pt}


\begin{minted}{julia}
DefaultRNG() -> Type{<:AbstractRNG} | AbstractRNG
DefaultRNG!(rng) -> rng
\end{minted}

The session{\textquotesingle}s default random-number generator for stochastic rounding. Initialized to the \texttt{Xoshiro} \emph{type} (a fresh generator per use); may be set to an \texttt{AbstractRNG} instance for a reproducible stream.



\end{adjustwidth}
\end{minipage}
\par\medskip

\par\medskip\noindent\begin{minipage}{\linewidth}
\hypertarget{2297064248264928244}{\texttt{SmallFloats.DefaultRbits}}  -- {Method.}
\nopagebreak
\begin{adjustwidth}{2em}{0pt}


\begin{minted}{julia}
DefaultRbits() -> Int
DefaultRbits!(n::Int) -> n
\end{minted}

The session{\textquotesingle}s default random-bit budget N for the stochastic rounding families (\texttt{StochasticA\{N\}}, \texttt{StochasticB\{N\}}, \texttt{StochasticC\{N\}}). Initialized to 8; must satisfy 1 ≤ N ≤ 60.



\end{adjustwidth}
\end{minipage}
\par\medskip

\par\medskip\noindent\begin{minipage}{\linewidth}
\hypertarget{2178409582565511335}{\texttt{SmallFloats.DefaultReturnType}}  -- {Method.}
\nopagebreak
\begin{adjustwidth}{2em}{0pt}


\begin{minted}{julia}
DefaultReturnType() -> Type{<:Binary}
DefaultReturnType!(T::Type{<:Binary}) -> T
\end{minted}

The session{\textquotesingle}s default \emph{result} format — the format an operation projects into when the caller does not name one. Initialized to \texttt{Binary8p2se}. Independent of \texttt{DefaultType}, which is the default \emph{operand} format.



\end{adjustwidth}
\end{minipage}
\par\medskip

\par\medskip\noindent\begin{minipage}{\linewidth}
\hypertarget{2306866772383893895}{\texttt{SmallFloats.DefaultRoundingMode}}  -- {Method.}
\nopagebreak
\begin{adjustwidth}{2em}{0pt}


\begin{minted}{julia}
DefaultRoundingMode() -> RoundingMode3109
DefaultRoundingMode!(m) -> m
\end{minted}

The session{\textquotesingle}s default rounding mode. Initialized to \texttt{NearestTiesToEven()}. Setting it rebuilds \hyperlinkref{17118068279635247347}{\texttt{DefaultProjection}} from the new mode and the current \hyperlinkref{1523060637620834303}{\texttt{DefaultSaturationMode}}. Accepts an instance or a fully-parameterized type (\texttt{NearestTiesToAway}, \texttt{StochasticA\{8\}}, …).



\end{adjustwidth}
\end{minipage}
\par\medskip

\par\medskip\noindent\begin{minipage}{\linewidth}
\hypertarget{1523060637620834303}{\texttt{SmallFloats.DefaultSaturationMode}}  -- {Method.}
\nopagebreak
\begin{adjustwidth}{2em}{0pt}


\begin{minted}{julia}
DefaultSaturationMode() -> SaturationMode
DefaultSaturationMode!(s) -> s
\end{minted}

The session{\textquotesingle}s default saturation mode. Initialized to \texttt{SatNone()}. Setting it rebuilds \hyperlinkref{17118068279635247347}{\texttt{DefaultProjection}} from the current \hyperlinkref{2306866772383893895}{\texttt{DefaultRoundingMode}} and the new mode. Accepts an instance or a type.



\end{adjustwidth}
\end{minipage}
\par\medskip

\par\medskip\noindent\begin{minipage}{\linewidth}
\hypertarget{17514657111736502113}{\texttt{SmallFloats.NextGreaterThan}}  -- {Method.}
\nopagebreak
\begin{adjustwidth}{2em}{0pt}

NextGreaterThan(v) (draft §4.16): the least datum greater than \texttt{v} in the total order — one step up the code lattice, with NaN → NaN and MaxFinite/+Inf → NaN at the top. \texttt{Base.nextfloat} on \texttt{Binary} is this operation.



\end{adjustwidth}
\end{minipage}
\par\medskip

\par\medskip\noindent\begin{minipage}{\linewidth}
\hypertarget{5658488599993001773}{\texttt{SmallFloats.NextLessThan}}  -- {Method.}
\nopagebreak
\begin{adjustwidth}{2em}{0pt}

NextLessThan(v) (draft §4.16): the greatest datum less than \texttt{v} in the total order — one step down the code lattice, with NaN → NaN and MinFinite/−Inf → NaN at the bottom. \texttt{Base.prevfloat} on \texttt{Binary} is this operation.



\end{adjustwidth}
\end{minipage}
\par\medskip

\par\medskip\noindent\begin{minipage}{\linewidth}
\hypertarget{3159928313134809876}{\texttt{SmallFloats.RoundOf}}  -- {Function.}
\nopagebreak
\begin{adjustwidth}{2em}{0pt}

Draft-named accessor: RoundOf(ρ) — the rounding mode of ρ (§4.2).



\end{adjustwidth}
\end{minipage}
\par\medskip

\par\medskip\noindent\begin{minipage}{\linewidth}
\hypertarget{16516058050886536028}{\texttt{SmallFloats.SatOf}}  -- {Function.}
\nopagebreak
\begin{adjustwidth}{2em}{0pt}

Draft-named accessor: SatOf(ρ) — the saturation mode of ρ (§4.2).



\end{adjustwidth}
\end{minipage}
\par\medskip

\par\medskip\noindent\begin{minipage}{\linewidth}
\hypertarget{8603138667834636985}{\texttt{SmallFloats.TotalOrder}}  -- {Method.}
\nopagebreak
\begin{adjustwidth}{2em}{0pt}

TotalOrder⟨fx,fy⟩ (draft §4.12.1): x ≤ y in the total order (single NaN largest). Same-format comparisons run on the integer order key (bitops plan K1) — proven equivalent to the decode order exhaustively over all pairs; cross-format keys are not comparable, so mixed formats keep the decode path.



\end{adjustwidth}
\end{minipage}
\par\medskip

\par\medskip\noindent\begin{minipage}{\linewidth}
\hypertarget{9294149665862826323}{\texttt{SmallFloats.approx}}  -- {Method.}
\nopagebreak
\begin{adjustwidth}{2em}{0pt}

Retrieve a registered approximate implementation (callable via \texttt{.fn}).



\end{adjustwidth}
\end{minipage}
\par\medskip

\par\medskip\noindent\begin{minipage}{\linewidth}
\hypertarget{15155075425824800578}{\texttt{SmallFloats.bitwidth}}  -- {Method.}
\nopagebreak
\begin{adjustwidth}{2em}{0pt}

Format bitwidth K (3–8).



\end{adjustwidth}
\end{minipage}
\par\medskip

\par\medskip\noindent\begin{minipage}{\linewidth}
\hypertarget{9130444423673118751}{\texttt{SmallFloats.codedistance}}  -- {Method.}
\nopagebreak
\begin{adjustwidth}{2em}{0pt}

Code-point distance along the total order between two values of one format.



\end{adjustwidth}
\end{minipage}
\par\medskip

\par\medskip\noindent\begin{minipage}{\linewidth}
\hypertarget{16705531398006417767}{\texttt{SmallFloats.codeunit\_\allowbreak{}type}}  -- {Method.}
\nopagebreak
\begin{adjustwidth}{2em}{0pt}

The storage unit of a format{\textquotesingle}s code point: \texttt{UInt8} for K ≤ 8, \texttt{UInt16} above.



\end{adjustwidth}
\end{minipage}
\par\medskip

\par\medskip\noindent\begin{minipage}{\linewidth}
\hypertarget{16371925287560310251}{\texttt{SmallFloats.conformance}}  -- {Method.}
\nopagebreak
\begin{adjustwidth}{2em}{0pt}


\begin{minted}{julia}
conformance() -> ConformanceDeclaration
\end{minted}

The package{\textquotesingle}s draft-§4.6-style conformance declaration, derived live from the operation registry, the table cache (the specializations actually instantiated), and the approximate-implementation registry. Serialize with \texttt{conformance\_\allowbreak{}dict} or render with \texttt{conformance\_\allowbreak{}report}.



\end{adjustwidth}
\end{minipage}
\par\medskip

\par\medskip\noindent\begin{minipage}{\linewidth}
\hypertarget{132348358124900907}{\texttt{SmallFloats.conformance\_\allowbreak{}dict}}  -- {Function.}
\nopagebreak
\begin{adjustwidth}{2em}{0pt}

Nested \texttt{Dict\{String,Any\}} form of the declaration — serializable by any JSON/TOML writer.



\end{adjustwidth}
\end{minipage}
\par\medskip

\par\medskip\noindent\begin{minipage}{\linewidth}
\hypertarget{7789835443308218455}{\texttt{SmallFloats.conformance\_\allowbreak{}report}}  -- {Function.}
\nopagebreak
\begin{adjustwidth}{2em}{0pt}

Human-readable conformance report.



\end{adjustwidth}
\end{minipage}
\par\medskip

\par\medskip\noindent\begin{minipage}{\linewidth}
\hypertarget{8861452329391555419}{\texttt{SmallFloats.decode!}}  -- {Method.}
\nopagebreak
\begin{adjustwidth}{2em}{0pt}


\begin{minted}{julia}
decode!(dest::AbstractArray{Float32}, A::AbstractArray{<:Binary}) -> dest
decode!(dest::AbstractArray{Float64}, A::AbstractArray{<:Binary}) -> dest
\end{minted}

\textbf{Exact} bulk ωDecode into an external float array.

Exactness is the whole contract here — it is what distinguishes \texttt{decode!} from \texttt{Float32.(A)}, which rounds like any other Julia conversion — so the destination element type is \textbf{gated on the format}, not assumed. \texttt{datumsexact(X, F)} decides it, and a format whose datums do not all fit \texttt{X} raises rather than silently rounding a bulk array. Every K ≤ 8 format passes both gates, so this is a new refusal only for formats that did not exist before the K ≤ 16 extension.

Where the gate passes and \texttt{K ≤ KSPLIT}, the implementation is unchanged: a Shape-A gather of the (narrowed) constant decode table.



\end{adjustwidth}
\end{minipage}
\par\medskip

\par\medskip\noindent\begin{minipage}{\linewidth}
\hypertarget{2086756398805737224}{\texttt{SmallFloats.decode}}  -- {Method.}
\nopagebreak
\begin{adjustwidth}{2em}{0pt}


\begin{minted}{julia}
decode(v::Binary) -> datumcarrier(typeof(v))
\end{minted}

ωDecode (draft §4.7.2). \textbf{Exact for every datum of every format}, which is why the return type is the format{\textquotesingle}s carrier rather than a fixed one: \texttt{Float64} where the bias allows (432 of the 504 formats), \texttt{Float128} to B = 8192, \texttt{BigFloat} above. \texttt{Float64(v)} is the way to ask for a \texttt{Float64} and accept the rounding.

Two implementations, selected by \texttt{decodepolicy} — dispatch on the representation, not a branch on K (invariant 9, §2 R-F):

\begin{itemize}
\item \texttt{TableDecode} (K ≤ \texttt{KSPLIT}): a \texttt{2{\textasciicircum}K}-entry constant-tuple lookup (bitops plan K2), generated once from \texttt{\_\allowbreak{}decode\_\allowbreak{}compute}, so the two are correct by construction and asserted equivalent exhaustively; constant inputs fold.


\item \texttt{ComputeDecode} (K > \texttt{KSPLIT}): \texttt{\_\allowbreak{}decode\_\allowbreak{}compute} directly. No table — 65 536 entries is what invariant 10 exists to forbid.

\end{itemize}


\end{adjustwidth}
\end{minipage}
\par\medskip

\par\medskip\noindent\begin{minipage}{\linewidth}
\hypertarget{11642396758074844399}{\texttt{SmallFloats.empty\_\allowbreak{}tables!}}  -- {Method.}
\nopagebreak
\begin{adjustwidth}{2em}{0pt}

Drop every cached table and adaptive counter (they rebuild lazily on next use).



\end{adjustwidth}
\end{minipage}
\par\medskip

\par\medskip\noindent\begin{minipage}{\linewidth}
\hypertarget{16181201817687362672}{\texttt{SmallFloats.expbias}}  -- {Method.}
\nopagebreak
\begin{adjustwidth}{2em}{0pt}

Exponent bias: 2{\textasciicircum}(K−P−1) signed, 2{\textasciicircum}(K−P) unsigned.



\end{adjustwidth}
\end{minipage}
\par\medskip

\par\medskip\noindent\begin{minipage}{\linewidth}
\hypertarget{4256747445410702613}{\texttt{SmallFloats.expbitwidth}}  -- {Method.}
\nopagebreak
\begin{adjustwidth}{2em}{0pt}

Width of the exponent field in bits: (K − signbit) − (P − 1).



\end{adjustwidth}
\end{minipage}
\par\medskip

\par\medskip\noindent\begin{minipage}{\linewidth}
\hypertarget{13855860993665101274}{\texttt{SmallFloats.f32\_\allowbreak{}exact}}  -- {Method.}
\nopagebreak
\begin{adjustwidth}{2em}{0pt}


\begin{minted}{julia}
f32_exact(op, f1, f2) -> Bool
\end{minted}

True iff \texttt{op ∈ (\allowbreak{}:Add, :Subtract, :Multiply)} on Float32-decoded operands is exact for every finite pair of (f1, f2) datums — checked once by exhaustive enumeration against a 300-bit BigFloat oracle, then cached. When true, a Float32 intermediate is an exact carrier under every projection mode.



\end{adjustwidth}
\end{minipage}
\par\medskip

\par\medskip\noindent\begin{minipage}{\linewidth}
\hypertarget{12747529899584332009}{\texttt{SmallFloats.f32\_\allowbreak{}impl}}  -- {Method.}
\nopagebreak
\begin{adjustwidth}{2em}{0pt}


\begin{minted}{julia}
f32_impl(op, fr, ρ) -> fn
\end{minted}

The decode32 → guarded-op32 → project kernel for \texttt{op⟨… → fr, ρ⟩}, shaped for \texttt{register\_\allowbreak{}approx!}. Nearest-ties ρ only: directed modes after a nearest-rounded Float32 compute measure κ ≥ 1 (and κ = NaN at saturation boundaries), and stochastic ρ is rejected by κ measurement itself.



\end{adjustwidth}
\end{minipage}
\par\medskip

\par\medskip\noindent\begin{minipage}{\linewidth}
\hypertarget{3988641001870183445}{\texttt{SmallFloats.format}}  -- {Method.}
\nopagebreak
\begin{adjustwidth}{2em}{0pt}


\begin{minted}{julia}
format(K, P, Σ, Δ) -> Type
\end{minted}

The concrete representation type for a format, from \textbf{runtime} parameters — \texttt{format(\allowbreak{}8, 4, true, true) === Binary8p4se}. The programmatic route to a type that can be an array element or a constructor target, and the supported replacement for spelling \texttt{Binary\{K,P,Σ,Δ\}} where a concrete type is meant.

Deliberately a \texttt{Dict} lookup: this path is reflection, it is not performance-sensitive, and it must not pretend to be. Where the parameters are static, the aliases and \texttt{reptype} are the zero-cost routes.



\end{adjustwidth}
\end{minipage}
\par\medskip

\par\medskip\noindent\begin{minipage}{\linewidth}
\hypertarget{10160445188737908012}{\texttt{SmallFloats.formatname}}  -- {Method.}
\nopagebreak
\begin{adjustwidth}{2em}{0pt}

\texttt{formatname(T)} — the draft §3.2 name of a format type.



\end{adjustwidth}
\end{minipage}
\par\medskip

\par\medskip\noindent\begin{minipage}{\linewidth}
\hypertarget{14294533040464309059}{\texttt{SmallFloats.ftz\_\allowbreak{}variant}}  -- {Method.}
\nopagebreak
\begin{adjustwidth}{2em}{0pt}


\begin{minted}{julia}
ftz_variant(op, fr, f1, ρ) -> fn
\end{minted}

Build the Annex-style approximate unary implementation: compute the defined result, then flush subnormal results to zero or \texttt{MinNormalOf(fr)}, whichever is nearer, ties toward zero (magnitude-symmetric for signed formats). Returned \texttt{fn} is suitable for \texttt{register\_\allowbreak{}approx!}; for \texttt{fr} with precision P its κ is 2{\textasciicircum}(P-2) when P ≥ 2 (largest flushed-to-zero subnormal) and 0 when P = 1 (no subnormals exist).



\end{adjustwidth}
\end{minipage}
\par\medskip

\par\medskip\noindent\begin{minipage}{\linewidth}
\hypertarget{16681316778131404680}{\texttt{SmallFloats.isextended}}  -- {Method.}
\nopagebreak
\begin{adjustwidth}{2em}{0pt}

Whether the format{\textquotesingle}s domain is Extended (datum set includes infinities).



\end{adjustwidth}
\end{minipage}
\par\medskip

\par\medskip\noindent\begin{minipage}{\linewidth}
\hypertarget{3128287096853498551}{\texttt{SmallFloats.issigned}}  -- {Method.}
\nopagebreak
\begin{adjustwidth}{2em}{0pt}

Whether the format is Signed (has a sign bit and negative datums).



\end{adjustwidth}
\end{minipage}
\par\medskip

\par\medskip\noindent\begin{minipage}{\linewidth}
\hypertarget{6068508418180240810}{\texttt{SmallFloats.list\_\allowbreak{}approx}}  -- {Method.}
\nopagebreak
\begin{adjustwidth}{2em}{0pt}

Sorted names of all registered κ-approximate implementations.



\end{adjustwidth}
\end{minipage}
\par\medskip

\par\medskip\noindent\begin{minipage}{\linewidth}
\hypertarget{1880978810980444415}{\texttt{SmallFloats.measure\_\allowbreak{}kappa}}  -- {Method.}
\nopagebreak
\begin{adjustwidth}{2em}{0pt}


\begin{minted}{julia}
measure_kappa(fn, op, fr, argformats, ρ; max_exhaustive=2^22, rng, samples=2^20)
    -> (κ::Float64, exhaustive::Bool)
\end{minted}

Measure κ of \texttt{fn(\allowbreak{}x::argformats[1], …)::fr} against the defined results of draft operation \texttt{op} under ρ. Enumerates the full input cross-product when it has at most \texttt{max\_\allowbreak{}exhaustive} points (always true for arity ≤ 2); otherwise measures on \texttt{samples} uniform draws and reports \texttt{exhaustive = false}.



\end{adjustwidth}
\end{minipage}
\par\medskip

\par\medskip\noindent\begin{minipage}{\linewidth}
\hypertarget{2330527590949240408}{\texttt{SmallFloats.nrandbits}}  -- {Method.}
\nopagebreak
\begin{adjustwidth}{2em}{0pt}

Number of random bits N consumed per projected element; 0 for pure modes.



\end{adjustwidth}
\end{minipage}
\par\medskip

\par\medskip\noindent\begin{minipage}{\linewidth}
\hypertarget{11630420834157853592}{\texttt{SmallFloats.projmode}}  -- {Method.}
\nopagebreak
\begin{adjustwidth}{2em}{0pt}

Translate a \texttt{Base.RoundingMode} to its \texttt{RoundingMode3109} counterpart (identity on the latter). Used only at API boundaries. Internals always carry RoundingMode3109.



\end{adjustwidth}
\end{minipage}
\par\medskip

\par\medskip\noindent\begin{minipage}{\linewidth}
\hypertarget{8685776062322191570}{\texttt{SmallFloats.rawvalue}}  -- {Method.}
\nopagebreak
\begin{adjustwidth}{2em}{0pt}

Unsafe raw constructor used by kernels after invariants are established.

Accepts an abstract format and normalizes through \texttt{reptype}; the type is passed \emph{as an argument} to \texttt{\_\allowbreak{}rawvalue} rather than assigned to a local, so it becomes a type parameter of the callee instead of a value the compiler has to fold (technique T6). Kernel call sites pass concrete types and pay nothing.



\end{adjustwidth}
\end{minipage}
\par\medskip

\par\medskip\noindent\begin{minipage}{\linewidth}
\hypertarget{852197704347416372}{\texttt{SmallFloats.register\_\allowbreak{}approx!}}  -- {Method.}
\nopagebreak
\begin{adjustwidth}{2em}{0pt}


\begin{minted}{julia}
register_approx!(name, op, fr, argformats, ρ, fn; κ=nothing, kwargs...) -> ApproxImpl
\end{minted}

Register \texttt{fn} as a named κ-approximate implementation of \texttt{op⟨argformats → fr, ρ⟩}. κ is measured by enumeration (see \texttt{measure\_\allowbreak{}kappa}); a declared \texttt{κ} smaller than the measured value is \textbf{rejected}. Omitting \texttt{κ} declares the measured value. NaN-κ implementations (non-finite mismatches) may be registered only by declaring \texttt{κ=NaN}.



\end{adjustwidth}
\end{minipage}
\par\medskip

\par\medskip\noindent\begin{minipage}{\linewidth}
\hypertarget{15760383384433593536}{\texttt{SmallFloats.register\_\allowbreak{}f32!}}  -- {Method.}
\nopagebreak
\begin{adjustwidth}{2em}{0pt}


\begin{minted}{julia}
register_f32!(op, fr, argformats, ρ; κ=nothing) -> ApproxImpl
\end{minted}

Build, measure, and register the Float32 kernel for \texttt{op⟨argformats → fr, ρ⟩} under the canonical name \texttt{f32/op/args→fr/rm\_\allowbreak{}sm}. Registration measures κ by exhaustive enumeration and fails loudly on understatement — that is the point. Measured baseline (2026-07-24): κ = 0 for all 120 same-format signatures × \{Add, Subtract, Multiply, Divide, Sqrt\} × RNE\_\allowbreak{}\{SatNone, SatFinite, SatPropagate\}.



\end{adjustwidth}
\end{minipage}
\par\medskip

\par\medskip\noindent\begin{minipage}{\linewidth}
\hypertarget{12298054483771101579}{\texttt{SmallFloats.reptype}}  -- {Method.}
\nopagebreak
\begin{adjustwidth}{2em}{0pt}


\begin{minted}{julia}
reptype(F) -> Type
\end{minted}

The concrete representation of a format. The identity on a concrete input, and the \texttt{Binary\{K,P,S,E\} → Code8|Code16} map on an abstract one.

Julia will \textbf{not} exploit {\textquotedbl}this abstract type has exactly one subtype{\textquotedbl}: there are no sealed types, inference cannot narrow an abstract element type to its unique subtype, and nothing stops a third party declaring another \texttt{<: Binary}. So \texttt{reptype} carries real weight rather than being belt-and-braces — it is how a method that holds only the format parameters reaches a constructible type.

The \texttt{Val} barrier is deliberate (invariant 9). The obvious spelling, \texttt{K <= KSPLIT ? Code8\{…\} : Code16\{…\}}, returns a \texttt{Type}, so its inferred return type is a small \texttt{Union} unless the compiler folds the comparison. It usually will; {\textquotedbl}usually{\textquotedbl} cannot underwrite a correctness-critical selection that also gates a zero-allocation guarantee. Behind a \texttt{Val}, each \texttt{\_\allowbreak{}rep} method has exactly one concrete return type and there is nothing to fold. Gate G9 checks it per format.



\end{adjustwidth}
\end{minipage}
\par\medskip

\par\medskip\noindent\begin{minipage}{\linewidth}
\hypertarget{11502893033482809346}{\texttt{SmallFloats.table\_\allowbreak{}bytes}}  -- {Method.}
\nopagebreak
\begin{adjustwidth}{2em}{0pt}

Total bytes currently held by every table cache.



\end{adjustwidth}
\end{minipage}
\par\medskip

\par\medskip\noindent\begin{minipage}{\linewidth}
\hypertarget{2349963207154028089}{\texttt{SmallFloats.table\_\allowbreak{}count}}  -- {Method.}
\nopagebreak
\begin{adjustwidth}{2em}{0pt}

Number of cached unary/binary specializations, both code units.



\end{adjustwidth}
\end{minipage}
\par\medskip

\par\medskip\noindent\begin{minipage}{\linewidth}
\hypertarget{9973158012858070886}{\texttt{SmallFloats.table\_\allowbreak{}policy}}  -- {Method.}
\nopagebreak
\begin{adjustwidth}{2em}{0pt}


\begin{minted}{julia}
table_policy(op, fr, f1[, f2[, f3]], ρ) -> (; shape, entries, bytes, reason)
\end{minted}

Which shape an array call on this signature will take, and why — \texttt{:A} for the table gather, \texttt{:B} for the per-element scalar path.

§4 Stage 5 item 6: {\textquotedbl}this signature is running the compute kernel{\textquotedbl} should be \emph{discoverable}, not inferred from a benchmark that came out slower than expected. Reads the same predicates the kernels do, so it cannot drift from them.



\end{adjustwidth}
\end{minipage}
\par\medskip

\par\medskip\noindent\begin{minipage}{\linewidth}
\hypertarget{13413040251500402145}{\texttt{SmallFloats.ternary\_\allowbreak{}count}}  -- {Method.}
\nopagebreak
\begin{adjustwidth}{2em}{0pt}

Number of cached ternary specializations, both code units.



\end{adjustwidth}
\end{minipage}
\par\medskip

\par\medskip\noindent\begin{minipage}{\linewidth}
\hypertarget{7146146650466703616}{\texttt{SmallFloats.trailingsigbits}}  -- {Method.}
\nopagebreak
\begin{adjustwidth}{2em}{0pt}

Trailing-significand width P − 1 (the stored fraction bits).



\end{adjustwidth}
\end{minipage}
\par\medskip

\par\medskip\noindent\begin{minipage}{\linewidth}
\hypertarget{7419061690047571527}{\texttt{SmallFloats.unregister\_\allowbreak{}approx!}}  -- {Method.}
\nopagebreak
\begin{adjustwidth}{2em}{0pt}

Remove a registered κ-approximate implementation (no-op if absent).



\end{adjustwidth}
\end{minipage}
\par\medskip

\par\medskip\noindent\begin{minipage}{\linewidth}
\hypertarget{5240195886209104356}{\texttt{SmallFloats.vmap!}}  -- {Function.}
\nopagebreak
\begin{adjustwidth}{2em}{0pt}


\begin{minted}{julia}
vmap!(dest, Val(op), fr, ρ, A[, B[, C]]; rng) -> dest
vmap(op, fr, ρ, A...; rng)
\end{minted}

Elementwise draft operation over arrays of \texttt{Binary} values, projecting into \texttt{fr} under ρ. Pure ρ runs the Shape-A gather whenever a table exists or the ternary bitwidth policy grants one; stochastic specializations (and untabled ternary signatures) run the scalar path per element (each element drawing its own R).



\end{adjustwidth}
\end{minipage}
\par\medskip

\par\medskip\noindent\begin{minipage}{\linewidth}
\hypertarget{7247813055970228015}{\texttt{SmallFloats.vmap}}  -- {Method.}
\nopagebreak
\begin{adjustwidth}{2em}{0pt}


\begin{minted}{julia}
vmap(op, fr, ρ, pv::PackedVector; rng=nothing) -> Vector{fr}
\end{minted}

Elementwise \texttt{op} over packed storage: unpack a \texttt{256}-element tile into a byte scratch, run the ordinary byte kernel (\texttt{vmap!}), emit, repeat. Never computes on packed words directly — the deliberate compute-unpacked/store-packed boundary.



\end{adjustwidth}
\end{minipage}
\par\medskip

\par\medskip\noindent\begin{minipage}{\linewidth}
\hypertarget{11729885955705173533}{\texttt{SmallFloats.with\_\allowbreak{}default\_\allowbreak{}projection}}  -- {Function.}
\nopagebreak
\begin{adjustwidth}{2em}{0pt}


\begin{minted}{julia}
with_default_type(f, args...)          -> f(DefaultType(), args...)
with_default_returntype(f, args...)    -> f(DefaultReturnType(), args...)
with_default_projection(f, args...)    -> f(DefaultProjection(), args...)
\end{minted}

Call \texttt{f} with the named session default as its first argument — the supported way to \emph{consume} a default. While the default still holds its initial value (the overwhelmingly common case) the call is statically compiled against that constant: no dynamic dispatch, \texttt{f} fully specialized. After the default is changed, the call crosses a function barrier instead: one dynamic dispatch, everything inside fully specialized.

Allocation contract: the combinator{\textquotesingle}s return type unions both paths. When \texttt{f}{\textquotesingle}s result type does not depend on the default — e.g. \texttt{with\_\allowbreak{}default\_\allowbreak{}projection} with the formats fixed by the caller — the union is concrete and the call is \textbf{zero-allocation with a concretely inferred result}. When the result{\textquotesingle}s type \emph{is} the default (\texttt{with\_\allowbreak{}default\_\allowbreak{}type(\allowbreak{}(\allowbreak{}T, x) -> T(\allowbreak{}x), …)}), the value is computed on the specialized path but boxes once at escape — the irreducible cost of a runtime-chosen type.


\begin{minted}{jlcon}
julia> with_default_type((T, x) -> T(x), 1.5)
Binary8p2se(1.5 ≡ 0x41)

julia> with_default_projection((ρ, x, y) -> Add(Binary8p4se, ρ, x, y),
                               Binary8p4se(1.5), Binary8p4se(0.25))
Binary8p4se(1.75 ≡ 0x46)
\end{minted}



\end{adjustwidth}
\end{minipage}
\par\medskip

\par\medskip\noindent\begin{minipage}{\linewidth}
\hypertarget{5478462901931591645}{\texttt{SmallFloats.with\_\allowbreak{}default\_\allowbreak{}returntype}}  -- {Function.}
\nopagebreak
\begin{adjustwidth}{2em}{0pt}


\begin{minted}{julia}
with_default_type(f, args...)          -> f(DefaultType(), args...)
with_default_returntype(f, args...)    -> f(DefaultReturnType(), args...)
with_default_projection(f, args...)    -> f(DefaultProjection(), args...)
\end{minted}

Call \texttt{f} with the named session default as its first argument — the supported way to \emph{consume} a default. While the default still holds its initial value (the overwhelmingly common case) the call is statically compiled against that constant: no dynamic dispatch, \texttt{f} fully specialized. After the default is changed, the call crosses a function barrier instead: one dynamic dispatch, everything inside fully specialized.

Allocation contract: the combinator{\textquotesingle}s return type unions both paths. When \texttt{f}{\textquotesingle}s result type does not depend on the default — e.g. \texttt{with\_\allowbreak{}default\_\allowbreak{}projection} with the formats fixed by the caller — the union is concrete and the call is \textbf{zero-allocation with a concretely inferred result}. When the result{\textquotesingle}s type \emph{is} the default (\texttt{with\_\allowbreak{}default\_\allowbreak{}type(\allowbreak{}(\allowbreak{}T, x) -> T(\allowbreak{}x), …)}), the value is computed on the specialized path but boxes once at escape — the irreducible cost of a runtime-chosen type.


\begin{minted}{jlcon}
julia> with_default_type((T, x) -> T(x), 1.5)
Binary8p2se(1.5 ≡ 0x41)

julia> with_default_projection((ρ, x, y) -> Add(Binary8p4se, ρ, x, y),
                               Binary8p4se(1.5), Binary8p4se(0.25))
Binary8p4se(1.75 ≡ 0x46)
\end{minted}



\end{adjustwidth}
\end{minipage}
\par\medskip

\par\medskip\noindent\begin{minipage}{\linewidth}
\hypertarget{8589336690729975238}{\texttt{SmallFloats.with\_\allowbreak{}default\_\allowbreak{}type}}  -- {Method.}
\nopagebreak
\begin{adjustwidth}{2em}{0pt}


\begin{minted}{julia}
with_default_type(f, args...)          -> f(DefaultType(), args...)
with_default_returntype(f, args...)    -> f(DefaultReturnType(), args...)
with_default_projection(f, args...)    -> f(DefaultProjection(), args...)
\end{minted}

Call \texttt{f} with the named session default as its first argument — the supported way to \emph{consume} a default. While the default still holds its initial value (the overwhelmingly common case) the call is statically compiled against that constant: no dynamic dispatch, \texttt{f} fully specialized. After the default is changed, the call crosses a function barrier instead: one dynamic dispatch, everything inside fully specialized.

Allocation contract: the combinator{\textquotesingle}s return type unions both paths. When \texttt{f}{\textquotesingle}s result type does not depend on the default — e.g. \texttt{with\_\allowbreak{}default\_\allowbreak{}projection} with the formats fixed by the caller — the union is concrete and the call is \textbf{zero-allocation with a concretely inferred result}. When the result{\textquotesingle}s type \emph{is} the default (\texttt{with\_\allowbreak{}default\_\allowbreak{}type(\allowbreak{}(\allowbreak{}T, x) -> T(\allowbreak{}x), …)}), the value is computed on the specialized path but boxes once at escape — the irreducible cost of a runtime-chosen type.


\begin{minted}{jlcon}
julia> with_default_type((T, x) -> T(x), 1.5)
Binary8p2se(1.5 ≡ 0x41)

julia> with_default_projection((ρ, x, y) -> Add(Binary8p4se, ρ, x, y),
                               Binary8p4se(1.5), Binary8p4se(0.25))
Binary8p4se(1.75 ≡ 0x46)
\end{minted}



\end{adjustwidth}
\end{minipage}
\par\medskip


\section{Internal Reference}



\label{2804742620609579687}{}


This page serves the Insights track: it documents unexported machinery behind the Mental Model, Introducing P3109, and the Technical Guide.



Unexported machinery documented in-source and referenced by the Technical Guide and Technical Examples. Public exported names are in the External Reference.


\par\medskip\noindent\begin{minipage}{\linewidth}
\hypertarget{854920489278715424}{\texttt{SmallFloats.DEFAULT\_\allowbreak{}RBITS}}  -- {Constant.}
\nopagebreak
\begin{adjustwidth}{2em}{0pt}

Default random-bit budget used by no-argument stochastic projection constructors.



\end{adjustwidth}
\end{minipage}
\par\medskip

\par\medskip\noindent\begin{minipage}{\linewidth}
\hypertarget{889067273513895166}{\texttt{SmallFloats.ExactExternalFloat}}  -- {Type.}
\nopagebreak
\begin{adjustwidth}{2em}{0pt}

External float element types whose widening to the Float64 carrier is exact.



\end{adjustwidth}
\end{minipage}
\par\medskip

\par\medskip\noindent\begin{minipage}{\linewidth}
\hypertarget{5954863807747789738}{\texttt{SmallFloats.KMAX}}  -- {Constant.}
\nopagebreak
\begin{adjustwidth}{2em}{0pt}

Highest bitwidth this package implements.



\end{adjustwidth}
\end{minipage}
\par\medskip

\par\medskip\noindent\begin{minipage}{\linewidth}
\hypertarget{9537035174923712895}{\texttt{SmallFloats.KMIN}}  -- {Constant.}
\nopagebreak
\begin{adjustwidth}{2em}{0pt}

Lowest bitwidth the draft defines (\texttt{K = 3}: sign, one exponent bit, one significand bit).



\end{adjustwidth}
\end{minipage}
\par\medskip

\par\medskip\noindent\begin{minipage}{\linewidth}
\hypertarget{6819318616108702912}{\texttt{SmallFloats.KSPLIT}}  -- {Constant.}
\nopagebreak
\begin{adjustwidth}{2em}{0pt}

Largest \texttt{K} whose code point fits a \texttt{UInt8} — the \texttt{Code8}/\allowbreak\texttt{Code16} seam.



\end{adjustwidth}
\end{minipage}
\par\medskip

\par\medskip\noindent\begin{minipage}{\linewidth}
\hypertarget{6899126343738372503}{\texttt{SmallFloats.MaybeR}}  -- {Type.}
\nopagebreak
\begin{adjustwidth}{2em}{0pt}

An explicit stochastic draw \texttt{R}, or \texttt{nothing} to draw one from the rng.



\end{adjustwidth}
\end{minipage}
\par\medskip

\par\medskip\noindent\begin{minipage}{\linewidth}
\hypertarget{9750648865235355525}{\texttt{SmallFloats.TABLE\_\allowbreak{}EAGER\_\allowbreak{}BITS}}  -- {Constant.}
\nopagebreak
\begin{adjustwidth}{2em}{0pt}

Largest table the package will \emph{build}, as \texttt{log2(entries)}. 16 = 65 536 entries.

A bound on \textbf{time}, not memory: a table costs one scalar trip per entry, at a measured {\textasciitilde}1 µs/entry, so this band caps first-call latency at {\textasciitilde}0.1 s while 2{\textasciicircum}24 entries would be tens of seconds — regardless of how comfortably 32 MiB fits in RAM.

16 is not a guess — it is exactly the largest table the K ≤ 8 grid already builds (a binary 8×8 table, 65 536 entries), so \textbf{no K ≤ 8 decision changes} and G5 stays byte-identical by construction. It also admits the case that matters most at wide K: a unary table for a K = 16 format is 2{\textasciicircum}16 entries, the same count, now 128 KiB instead of 64 KiB because the entries are \texttt{UInt16}.



\end{adjustwidth}
\end{minipage}
\par\medskip

\par\medskip\noindent\begin{minipage}{\linewidth}
\hypertarget{13361470731934248224}{\texttt{SmallFloats.TABLE\_\allowbreak{}MAX\_\allowbreak{}BITS}}  -- {Constant.}
\nopagebreak
\begin{adjustwidth}{2em}{0pt}

Hard ceiling on a single table, as \texttt{log2(bytes)} — the byte budget invariant 10 requires. 25 = 32 MiB.

Compared \textbf{as bits}. Computing \texttt{1 << ΣK} in order to decide whether \texttt{1 << ΣK} is too large is exactly the bug the budget exists to prevent: at ΣK = 64 the shift wraps to zero and the impossible allocation looks free.



\end{adjustwidth}
\end{minipage}
\par\medskip

\par\medskip\noindent\begin{minipage}{\linewidth}
\hypertarget{12494916835769691716}{\texttt{SmallFloats.TERNARY\_\allowbreak{}ADAPTIVE\_\allowbreak{}BITS}}  -- {Constant.}
\nopagebreak
\begin{adjustwidth}{2em}{0pt}

K1+K2+K3 up to which a ternary table may be built adaptively once the signature has processed \texttt{TERNARY\_\allowbreak{}BUILD\_\allowbreak{}ELEMS[]} elements (default 21 bits = 2 MiB — the K=7 band). Above this, ternary ops always run the compute kernel.



\end{adjustwidth}
\end{minipage}
\par\medskip

\par\medskip\noindent\begin{minipage}{\linewidth}
\hypertarget{8297071350468471162}{\texttt{SmallFloats.TERNARY\_\allowbreak{}BUILD\_\allowbreak{}ELEMS}}  -- {Constant.}
\nopagebreak
\begin{adjustwidth}{2em}{0pt}

Cumulative element count at which an adaptive-band signature earns its table.



\end{adjustwidth}
\end{minipage}
\par\medskip

\par\medskip\noindent\begin{minipage}{\linewidth}
\hypertarget{6187612789825565941}{\texttt{SmallFloats.TERNARY\_\allowbreak{}CACHE\_\allowbreak{}BYTES}}  -- {Constant.}
\nopagebreak
\begin{adjustwidth}{2em}{0pt}

Byte budget for the ternary cache; least-recently-used tables evict first.



\end{adjustwidth}
\end{minipage}
\par\medskip

\par\medskip\noindent\begin{minipage}{\linewidth}
\hypertarget{17354237254338795853}{\texttt{SmallFloats.TERNARY\_\allowbreak{}EAGER\_\allowbreak{}BITS}}  -- {Constant.}
\nopagebreak
\begin{adjustwidth}{2em}{0pt}

K1+K2+K3 up to which a ternary table is built eagerly on first array call (default 18 bits = 256 KiB — covers every all-K≤6 signature).



\end{adjustwidth}
\end{minipage}
\par\medskip

\par\medskip\noindent\begin{minipage}{\linewidth}
\hypertarget{15945855848429109616}{\texttt{SmallFloats.THREADED\_\allowbreak{}KERNELS}}  -- {Constant.}
\nopagebreak
\begin{adjustwidth}{2em}{0pt}

Master switch for threaded ternary compute loops.



\end{adjustwidth}
\end{minipage}
\par\medskip

\par\medskip\noindent\begin{minipage}{\linewidth}
\hypertarget{8854710295716658150}{\texttt{SmallFloats.THREAD\_\allowbreak{}MIN\_\allowbreak{}ELEMS}}  -- {Constant.}
\nopagebreak
\begin{adjustwidth}{2em}{0pt}

Minimum element count before the pure-ρ ternary compute loop threads.



\end{adjustwidth}
\end{minipage}
\par\medskip

\par\medskip\noindent\begin{minipage}{\linewidth}
\hypertarget{4674484589655001634}{\texttt{Core.Type}}  -- {Method.}
\nopagebreak
\begin{adjustwidth}{2em}{0pt}


\begin{minted}{julia}
(T::Type{<:Binary})(c::Unsigned) -> T
\end{minted}

Construct a value from its \textbf{code point} (validated: \texttt{c} must satisfy the representation invariant for the format{\textquotesingle}s K, or an \texttt{ArgumentError} is thrown; the check folds away when K equals the storage width and for constant codes). \texttt{Binary5p3sf(\allowbreak{}0x08) == Binary5p3sf(\allowbreak{}1.0)}.

\texttt{Unsigned} is the argument-type \emph{class} with code-point semantics, at every width — every other \texttt{Real} (including signed \texttt{Integer}s) constructs by projecting the numeric value. \texttt{Binary8p4se(0x02)} is code point 2 (the second-smallest subnormal), while \texttt{Binary8p4se(2)} is the number 2.0. Widening an \texttt{Unsigned} is lossless, so the only way a wrong-width argument can be wrong is by being out of code range — which throws. \texttt{rawvalue(T, c)} remains the unchecked kernel-internal route.



\end{adjustwidth}
\end{minipage}
\par\medskip

\par\medskip\noindent\begin{minipage}{\linewidth}
\hypertarget{7842307923234431821}{\texttt{SmallFloats.BigExactF}}  -- {Type.}
\nopagebreak
\begin{adjustwidth}{2em}{0pt}

Deferred exact result: \texttt{f()} returns the value as an exact \texttt{BigFloat}.



\end{adjustwidth}
\end{minipage}
\par\medskip

\par\medskip\noindent\begin{minipage}{\linewidth}
\hypertarget{632943239165416513}{\texttt{SmallFloats.Code16}}  -- {Type.}
\nopagebreak
\begin{adjustwidth}{2em}{0pt}


\begin{minted}{julia}
Code16{K,P,SGN,EXT} <: Binary{K,P,SGN,EXT}
\end{minted}

The two-byte representation — every format with \texttt{KSPLIT < K ≤ KMAX}. Differs from its \texttt{Code8} sibling only in the width of the word carrying the code point; all semantics live in the \texttt{Binary} parameters.



\end{adjustwidth}
\end{minipage}
\par\medskip

\par\medskip\noindent\begin{minipage}{\linewidth}
\hypertarget{10616179233481181492}{\texttt{SmallFloats.Code8}}  -- {Type.}
\nopagebreak
\begin{adjustwidth}{2em}{0pt}


\begin{minted}{julia}
Code8{K,P,SGN,EXT} <: Binary{K,P,SGN,EXT}
\end{minted}

The one-byte representation of a P3109 format — every format with \texttt{K ≤ 8}. Differs from its \texttt{Code16} sibling only in the width of the word carrying the code point; all semantics live in the \texttt{Binary} parameters.



\end{adjustwidth}
\end{minipage}
\par\medskip

\par\medskip\noindent\begin{minipage}{\linewidth}
\hypertarget{3168615570275869760}{\texttt{SmallFloats.ComputeDecode}}  -- {Type.}
\nopagebreak
\begin{adjustwidth}{2em}{0pt}

\texttt{2{\textasciicircum}K > 256}: compute from the code point; a 65 536-element tuple literal is not a table, it is a compile-time hazard (invariant 10).



\end{adjustwidth}
\end{minipage}
\par\medskip

\par\medskip\noindent\begin{minipage}{\linewidth}
\hypertarget{5477313120957585094}{\texttt{SmallFloats.DecodePolicy}}  -- {Type.}
\nopagebreak
\begin{adjustwidth}{2em}{0pt}

Decode strategy for a format: a \texttt{2{\textasciicircum}K}-entry constant tuple, or computed.



\end{adjustwidth}
\end{minipage}
\par\medskip

\par\medskip\noindent\begin{minipage}{\linewidth}
\hypertarget{1455587366773904952}{\texttt{SmallFloats.Enclose128F}}  -- {Type.}
\nopagebreak
\begin{adjustwidth}{2em}{0pt}

Correctly-rounded Float128 bracket (plan Site C, Class R): IEEE mandates correct rounding for Float128 \texttt{/}, \texttt{sqrt}, \texttt{fma}, so a nearest-CR result of a value known inexact brackets the truth in the \emph{open} interval \texttt{(\allowbreak{}lo, hi) = (\allowbreak{}prevfloat(\allowbreak{}q), nextfloat(\allowbreak{}q))} with no envelope assumption. \texttt{f} is the MPFR fallback for (near-impossible at P ≤ 8) grid-straddling brackets.



\end{adjustwidth}
\end{minipage}
\par\medskip

\par\medskip\noindent\begin{minipage}{\linewidth}
\hypertarget{16356965793446534219}{\texttt{SmallFloats.EncloseF}}  -- {Type.}
\nopagebreak
\begin{adjustwidth}{2em}{0pt}

Deferred enclosure, resolved by \texttt{\_\allowbreak{}finish} through up to three stages, cheapest first:

\begin{itemize}
\item[1. ] \texttt{yd} — an \emph{eager} Float64 estimate (NaN ⇒ absent) whose libm evaluation is faithful (≤ 1 ulp); the true value is taken to lie within \texttt{|yd|·2{\textasciicircum}-45} (\texttt{\_\allowbreak{}F64\_\allowbreak{}RELEXP}), ≥ 2{\textasciicircum}7 slacker than any faithful-libm error.


\item[2. ] \texttt{fq} — a zero-argument Float128 estimator (plan Site D, Class E; \texttt{nothing} ⇒ absent): true value within \texttt{|y|·2{\textasciicircum}-90} of \texttt{y = fq()} (\texttt{\_\allowbreak{}F128\_\allowbreak{}RELEXP}), a bound ≥ 2{\textasciicircum}18 slacker than any published libquadmath transcendental error, discharged empirically by the differential-build tests.


\item[3. ] \texttt{f(prec)} — the rigorous MPFR ladder: directed \texttt{(lo, hi)::NTuple\{2,BigFloat\}} strictly bracketing the true value.

\end{itemize}
Each estimate stage returns only when the two-sided sticky projection of its envelope agrees on a single code point; any disagreement (or an absent/degenerate estimate) falls through to the next stage. Stage 3 always decides.



\end{adjustwidth}
\end{minipage}
\par\medskip

\par\medskip\noindent\begin{minipage}{\linewidth}
\hypertarget{6912697913028347033}{\texttt{SmallFloats.Head}}  -- {Type.}
\nopagebreak
\begin{adjustwidth}{2em}{0pt}


\begin{minted}{julia}
Head
\end{minted}

Evaluation-carrier selector: a singleton tag, so carrier choice is a method lookup rather than a value the compiler must fold. Rungs are ordered \texttt{HeadF64 < HeadF128 < HeadExact} by carrier width.



\end{adjustwidth}
\end{minipage}
\par\medskip

\par\medskip\noindent\begin{minipage}{\linewidth}
\hypertarget{12343000088784592265}{\texttt{SmallFloats.HeadExact}}  -- {Type.}
\nopagebreak
\begin{adjustwidth}{2em}{0pt}

Rung 3 — unbounded exponent. \texttt{BigFloat} until Stage 7, \texttt{Dyadic} after; the tag is the seam that makes that swap a one-line change (gate G7 pins that the two agree).



\end{adjustwidth}
\end{minipage}
\par\medskip

\par\medskip\noindent\begin{minipage}{\linewidth}
\hypertarget{3528180919873185467}{\texttt{SmallFloats.HeadF128}}  -- {Type.}
\nopagebreak
\begin{adjustwidth}{2em}{0pt}

Rung 2 — \texttt{Float128} arithmetic (the existing escalation machinery, promoted to entry point). Allocation-free: \texttt{Float128} is a bitstype.



\end{adjustwidth}
\end{minipage}
\par\medskip

\par\medskip\noindent\begin{minipage}{\linewidth}
\hypertarget{9123640062122562911}{\texttt{SmallFloats.HeadF64}}  -- {Type.}
\nopagebreak
\begin{adjustwidth}{2em}{0pt}

Rung 1 — \texttt{Float64} arithmetic. Every K ≤ 8 format and 432 of the 504.



\end{adjustwidth}
\end{minipage}
\par\medskip

\par\medskip\noindent\begin{minipage}{\linewidth}
\hypertarget{3305030421944754103}{\texttt{SmallFloats.Rounded}}  -- {Type.}
\nopagebreak
\begin{adjustwidth}{2em}{0pt}

Result of ωRoundToPrecision in canonical integer form: a kind tag plus \texttt{sign · S · 2{\textasciicircum}Q} for finite results (design §5.2). The saturation stage consumes this without ever re-touching the carrier value.



\end{adjustwidth}
\end{minipage}
\par\medskip

\par\medskip\noindent\begin{minipage}{\linewidth}
\hypertarget{844750228189688156}{\texttt{SmallFloats.StickyF}}  -- {Type.}
\nopagebreak
\begin{adjustwidth}{2em}{0pt}

Sticky-head exact result (wide-spread tail, Float128 revision follow-on): the true value is \texttt{v + sgn·ε} for an infinitesimal \texttt{ε > 0} — \texttt{v} carries every bit the projection can consume and \texttt{sgn ∈ \{-1,+1\}} the direction of the neglected tail. Sound because \texttt{v} is either exactly on a rounding threshold of the target grid (where \texttt{sticky} decides, for every mode including stochastic sub-grids) or strictly farther from the nearest threshold than the tail magnitude; the emitting sites in oracle.jl discharge that bound (operand significands ≤ 17 bits, target grids ≤ P−1+N ≤ 67 fractional bits, spreads > the \emph{DE}* thresholds). Non-allocating: replaces the BigFloat escalation for FMA/FAA.



\end{adjustwidth}
\end{minipage}
\par\medskip

\par\medskip\noindent\begin{minipage}{\linewidth}
\hypertarget{12352230725331747513}{\texttt{SmallFloats.TableDecode}}  -- {Type.}
\nopagebreak
\begin{adjustwidth}{2em}{0pt}

\texttt{2{\textasciicircum}K ≤ 256}: the generated constant tuple folds and costs one load.



\end{adjustwidth}
\end{minipage}
\par\medskip

\par\medskip\noindent\begin{minipage}{\linewidth}
\hypertarget{245207806059811915}{\texttt{SmallFloats.TableKey}}  -- {Type.}
\nopagebreak
\begin{adjustwidth}{2em}{0pt}

Value key for the table cache: op + format parameters + ρ parameters (all values, so the cache itself is type-stable and non-allocating to query).



\end{adjustwidth}
\end{minipage}
\par\medskip

\par\medskip\noindent\begin{minipage}{\linewidth}
\hypertarget{9239390415321841063}{\texttt{SmallFloats.TernaryKey}}  -- {Type.}
\nopagebreak
\begin{adjustwidth}{2em}{0pt}

Value key for a ternary table: op + result + three operand formats + ρ.



\end{adjustwidth}
\end{minipage}
\par\medskip

\par\medskip\noindent\begin{minipage}{\linewidth}
\hypertarget{7671132794519427917}{\texttt{Base.nextfloat}}  -- {Method.}
\nopagebreak
\begin{adjustwidth}{2em}{0pt}

NextGreaterThan(v) (draft §4.16): the least datum greater than \texttt{v} in the total order — one step up the code lattice, with NaN → NaN and MaxFinite/+Inf → NaN at the top. \texttt{Base.nextfloat} on \texttt{Binary} is this operation.



\end{adjustwidth}
\end{minipage}
\par\medskip

\par\medskip\noindent\begin{minipage}{\linewidth}
\hypertarget{8297310150874586961}{\texttt{Base.prevfloat}}  -- {Method.}
\nopagebreak
\begin{adjustwidth}{2em}{0pt}

NextLessThan(v) (draft §4.16): the greatest datum less than \texttt{v} in the total order — one step down the code lattice, with NaN → NaN and MinFinite/−Inf → NaN at the bottom. \texttt{Base.prevfloat} on \texttt{Binary} is this operation.



\end{adjustwidth}
\end{minipage}
\par\medskip

\par\medskip\noindent\begin{minipage}{\linewidth}
\hypertarget{17551943659563969876}{\texttt{Base.widen}}  -- {Method.}
\nopagebreak
\begin{adjustwidth}{2em}{0pt}

The public promotion target for \texttt{T} against external numeric types. Always a Julia \texttt{AbstractFloat}, never the internal exact carrier.



\end{adjustwidth}
\end{minipage}
\par\medskip

\par\medskip\noindent\begin{minipage}{\linewidth}
\hypertarget{15942066643973002529}{\texttt{SmallFloats.\_\allowbreak{}bp\_\allowbreak{}element}}  -- {Method.}
\nopagebreak
\begin{adjustwidth}{2em}{0pt}

One element of ωBlockProject: the draft{\textquotesingle}s S-special rows, then ωDivide + ωProject.

Fast-path cascade after the special rows, ordered cheapest-first by result kind: exact Float64 quotient → CR Float128 quotient (exact or one-ulp bracket) → CR-divided Enclose128F bracket → fq-composition envelope (2{\textasciicircum}-89) — each resolved by the two-sided sticky gate; any miss falls through to the rigorous MPFR interval at the bottom.



\end{adjustwidth}
\end{minipage}
\par\medskip

\par\medskip\noindent\begin{minipage}{\linewidth}
\hypertarget{4327771365995947438}{\texttt{SmallFloats.\_\allowbreak{}check\_\allowbreak{}tabulable}}  -- {Method.}
\nopagebreak
\begin{adjustwidth}{2em}{0pt}

Stochastic ρ is a distribution over R, never a table (design §5.4).



\end{adjustwidth}
\end{minipage}
\par\medskip

\par\medskip\noindent\begin{minipage}{\linewidth}
\hypertarget{11960956699752345863}{\texttt{SmallFloats.\_\allowbreak{}comparable}}  -- {Method.}
\nopagebreak
\begin{adjustwidth}{2em}{0pt}

Numeric comparison is defined only when neither operand is NaN — every \texttt{==}/\allowbreak\texttt{<}/\texttt{<=} below returns \texttt{false} otherwise, per IEEE unorderedness.



\end{adjustwidth}
\end{minipage}
\par\medskip

\par\medskip\noindent\begin{minipage}{\linewidth}
\hypertarget{9616758042032316127}{\texttt{SmallFloats.\_\allowbreak{}extexprange}}  -- {Method.}
\nopagebreak
\begin{adjustwidth}{2em}{0pt}

(least subnormal exponent, greatest normal exponent) of an external float.



\end{adjustwidth}
\end{minipage}
\par\medskip

\par\medskip\noindent\begin{minipage}{\linewidth}
\hypertarget{14260418278412798574}{\texttt{SmallFloats.\_\allowbreak{}resolve\_\allowbreak{}rng}}  -- {Method.}
\nopagebreak
\begin{adjustwidth}{2em}{0pt}

Resolve a caller-supplied rng, falling back to the task-local default. Call sites hoist this out of loops so array kernels resolve once per call, not per element.



\end{adjustwidth}
\end{minipage}
\par\medskip

\par\medskip\noindent\begin{minipage}{\linewidth}
\hypertarget{1104655483086699683}{\texttt{SmallFloats.\_\allowbreak{}rng\_\allowbreak{}for}}  -- {Method.}
\nopagebreak
\begin{adjustwidth}{2em}{0pt}

The rng an operation under ρ will actually draw from — \texttt{nothing} for pure ρ, which must never touch RNG state.



\end{adjustwidth}
\end{minipage}
\par\medskip

\par\medskip\noindent\begin{minipage}{\linewidth}
\hypertarget{10743025803044024289}{\texttt{SmallFloats.\_\allowbreak{}rungindex}}  -- {Method.}
\nopagebreak
\begin{adjustwidth}{2em}{0pt}

Per-format rung index from the exponent bias. A pure \texttt{Int} function of the type parameters, so \texttt{Val(\_\allowbreak{}rungindex(F))} is a compile-time constant.

Stated on \texttt{2B} rather than \texttt{B} so it is visibly the \texttt{n = 2} case of \texttt{\_\allowbreak{}rungindex\_\allowbreak{}span} and not an independent constant pair. \texttt{B ≤ 512 ⟺ 2B ≤ 1024} and \texttt{B ≤ 8192 ⟺ 2B ≤ 16384}, so the 432/64/8 partition is unchanged — G9 enumerates it.



\end{adjustwidth}
\end{minipage}
\par\medskip

\par\medskip\noindent\begin{minipage}{\linewidth}
\hypertarget{1496290229282609812}{\texttt{SmallFloats.\_\allowbreak{}rungindex\_\allowbreak{}span}}  -- {Method.}
\nopagebreak
\begin{adjustwidth}{2em}{0pt}

Rung index for a monomial whose factor biases sum to \texttt{ΣB}. This is the primitive form of the boundary rule stated above; the per-format answer is its \texttt{n = 2} same-format case.



\end{adjustwidth}
\end{minipage}
\par\medskip

\par\medskip\noindent\begin{minipage}{\linewidth}
\hypertarget{942381566965426910}{\texttt{SmallFloats.\_\allowbreak{}ternary\_\allowbreak{}table\_\allowbreak{}for}}  -- {Method.}
\nopagebreak
\begin{adjustwidth}{2em}{0pt}


\begin{minted}{julia}
_ternary_table_for(op, fr, f1, f2, f3, ρ, nelems) -> Union{Nothing,Memory{UInt8}}
\end{minted}

The kernel-facing policy gate, called once per array operation with the call{\textquotesingle}s element count. Eager band (Σ Kᵢ ≤ \texttt{TERNARY\_\allowbreak{}EAGER\_\allowbreak{}BITS[]}): fetch/build now. Adaptive band (≤ \texttt{TERNARY\_\allowbreak{}ADAPTIVE\_\allowbreak{}BITS[]}): return the cached table if present, otherwise accumulate \texttt{nelems} against the signature and build only once it has earned \texttt{TERNARY\_\allowbreak{}BUILD\_\allowbreak{}ELEMS[]}. Beyond the adaptive band (the K=8 territory): always \texttt{nothing} — the compute kernel is the right tradeoff there.



\end{adjustwidth}
\end{minipage}
\par\medskip

\par\medskip\noindent\begin{minipage}{\linewidth}
\hypertarget{12144973337474854137}{\texttt{SmallFloats.\_\allowbreak{}unitmask}}  -- {Method.}
\nopagebreak
\begin{adjustwidth}{2em}{0pt}


\begin{minted}{julia}
_unitmask(U, K) -> U
\end{minted}

The low-\texttt{K}-bits mask of storage unit \texttt{U}, built by \textbf{complement} — shift down from \texttt{typemax}, never up from \texttt{one}.

Julia defines a shift at or beyond a type{\textquotesingle}s width as zero, so \texttt{UInt8(1) << 8} and \texttt{UInt16(1) << 16} are both \texttt{0}, and \texttt{2{\textasciicircum}K} computed in a unit of width \texttt{K} is \texttt{0}. The format grid \textbf{always contains a format whose K equals its storage width} (K = 8 on \texttt{UInt8}, K = 16 on \texttt{UInt16}), so that case is structural, not an edge. Here the shift amount is \texttt{width − K ∈ [0, width−1]}, in which the pathological amount is unreachable \emph{by construction}, at every K, for every unit — including a future \texttt{Code32}.



\end{adjustwidth}
\end{minipage}
\par\medskip

\par\medskip\noindent\begin{minipage}{\linewidth}
\hypertarget{1161894050593661291}{\texttt{SmallFloats.\_\allowbreak{}unsupported}}  -- {Method.}
\nopagebreak
\begin{adjustwidth}{2em}{0pt}

Refuse a Base verb with a reason. Absence and refusal must not look alike.



\end{adjustwidth}
\end{minipage}
\par\medskip

\par\medskip\noindent\begin{minipage}{\linewidth}
\hypertarget{14023424897002283292}{\texttt{SmallFloats.apply\_\allowbreak{}op}}  -- {Method.}
\nopagebreak
\begin{adjustwidth}{2em}{0pt}

apply\_\allowbreak{}op(Val(op), fr, ρ, R, xs...) — evaluate \texttt{op}{\textquotesingle}s ω-semantics on decoded operands and project the result into \texttt{fr} under ρ (with random bits R).

The evaluation carrier comes from the \textbf{operands}, not from \texttt{fr}. That is \texttt{carriers.jl}{\textquotesingle}s stated rule — the carrier constrains what \texttt{ωeval} may form, never what may be projected into a format, because \texttt{round\_\allowbreak{}to\_\allowbreak{}precision} works in \texttt{(P, B)} integer space and is exact for any exact input. Stage 3 selected on \texttt{rung(fr)} instead, as a deliberate stand-in while \texttt{decode} still refused K ≥ 9; keeping that would have refused \texttt{Add(\allowbreak{}Binary16p2se, ρ, x8, y8)}, whose operands are ordinary \texttt{Float64} datums and whose exact sum needs no wide carrier at all.

Completeness of the operand-only rule rests on one premise, stated because Stage 3{\textquotesingle}s version of it is what M28 was: the spec register{\textquotesingle}s maximum factor count is \textbf{2}, and the rung boundaries are stated on \texttt{2B} precisely so that a datum on carrier \texttt{C} guarantees its square is also on \texttt{C}. Beyond two factors — \texttt{ScaledMultiply}{\textquotesingle}s four — the monomial bound stops following from the operand types and \texttt{rung(op, Fs...)} must be used with the formats in hand. That is why \texttt{blocks.jl} calls the join and this does not.



\end{adjustwidth}
\end{minipage}
\par\medskip

\par\medskip\noindent\begin{minipage}{\linewidth}
\hypertarget{646269805905669936}{\texttt{SmallFloats.bigprec\_\allowbreak{}prod}}  -- {Method.}
\nopagebreak
\begin{adjustwidth}{2em}{0pt}


\begin{minted}{julia}
bigprec_prod(xs...) -> Int
\end{minted}

\texttt{bigprec} for an expression whose leading term is the \textbf{product} \texttt{x₁·x₂} and whose remaining terms are added to it — the FMA shape. The product{\textquotesingle}s exponent is \texttt{e₁ + e₂}, which lies outside the range of the operands{\textquotesingle} own exponents, so the sum method would understate the spread. Its significand is \texttt{P₁ + P₂} bits wide.



\end{adjustwidth}
\end{minipage}
\par\medskip

\par\medskip\noindent\begin{minipage}{\linewidth}
\hypertarget{14780784727942474206}{\texttt{SmallFloats.blockdecode}}  -- {Method.}
\nopagebreak
\begin{adjustwidth}{2em}{0pt}

ωBlockDecode (draft §5.1.1): per lane ωMultiply(decode(s), decode(xᵢ)), exact on the carrier the block{\textquotesingle}s two formats jointly require.

The carrier is keyed on \texttt{rung(\allowbreak{}Val(\allowbreak{}:Multiply), FS, FE)} — \textbf{not} on either format{\textquotesingle}s own rung, and that distinction is the whole content of this method. A block{\textquotesingle}s scale and element formats are independent, and what has to fit is their \emph{product}: \texttt{carriers.jl}{\textquotesingle}s rule is \texttt{ΣBᵢ ≤ emax}, so two formats that are each rung 1 at \texttt{2B ≤ 1024} can compose into a lane product needing \texttt{B\_\allowbreak{}S + B\_\allowbreak{}E} well past it. Asserting \texttt{carriertype(rung(FE))} would be wrong in the worst way available here — the overflow produces \texttt{±Inf}, a plausible special value, rather than an error.

The assertion stays because it is load-bearing for more than documentation: it is what keeps the returned \texttt{NTuple} concretely typed, and with it the reductions downstream allocation-free. It now names a carrier computed from both formats instead of a constant.



\end{adjustwidth}
\end{minipage}
\par\medskip

\par\medskip\noindent\begin{minipage}{\linewidth}
\hypertarget{10671394415806872186}{\texttt{SmallFloats.blockproject}}  -- {Method.}
\nopagebreak
\begin{adjustwidth}{2em}{0pt}

ωBlockProject (draft §5.1.2) over a tuple of result kinds, scale supplied as a value.



\end{adjustwidth}
\end{minipage}
\par\medskip

\par\medskip\noindent\begin{minipage}{\linewidth}
\hypertarget{6493780492297571350}{\texttt{SmallFloats.codemask}}  -- {Method.}
\nopagebreak
\begin{adjustwidth}{2em}{0pt}

The low-K-bits mask of a format{\textquotesingle}s storage unit (representation invariant 3): the code point occupies the low K bits, the high bits are maintained zero.



\end{adjustwidth}
\end{minipage}
\par\medskip

\par\medskip\noindent\begin{minipage}{\linewidth}
\hypertarget{3840324864146613682}{\texttt{SmallFloats.datumcarrier}}  -- {Method.}
\nopagebreak
\begin{adjustwidth}{2em}{0pt}

The exact carrier for the \emph{datums} of \texttt{F}: what \texttt{decode(v::F)} returns.



\end{adjustwidth}
\end{minipage}
\par\medskip

\par\medskip\noindent\begin{minipage}{\linewidth}
\hypertarget{5824837744343724298}{\texttt{SmallFloats.datumsexact}}  -- {Method.}
\nopagebreak
\begin{adjustwidth}{2em}{0pt}


\begin{minted}{julia}
datumsexact(X, F) -> Bool
\end{minted}

\texttt{true} when \textbf{every} datum of format \texttt{F} is exactly representable in the external float type \texttt{X}.

Three conditions, one per way a conversion can lose information:

\begin{itemize}
\item \texttt{P ≤ precision(X)} — the significand fits;


\item the largest finite datum{\textquotesingle}s exponent is within \texttt{X}{\textquotesingle}s normal range. That exponent is \texttt{B − 1} for a signed format and \texttt{0} for an unsigned one, which is not symmetry-breaking pedantry: an unsigned format spends its whole exponent field below 1, so \texttt{Binary16p1uf}{\textquotesingle}s datums run down to \texttt{2{\textasciicircum}-32769} while its largest is \texttt{0.5};


\item \texttt{2 − P − B ≥} \texttt{X}{\textquotesingle}s least subnormal exponent — and the smallest datum being a \emph{subnormal} of \texttt{X} is fine, because a datum is \texttt{sig · 2{\textasciicircum}(2−P−B)} with integer \texttt{sig}, so once the step is a multiple of \texttt{X}{\textquotesingle}s subnormal step every datum lands on \texttt{X}{\textquotesingle}s grid exactly.

\end{itemize}


\end{adjustwidth}
\end{minipage}
\par\medskip

\par\medskip\noindent\begin{minipage}{\linewidth}
\hypertarget{9427357598855955389}{\texttt{SmallFloats.decodepolicy}}  -- {Method.}
\nopagebreak
\begin{adjustwidth}{2em}{0pt}

Decode strategy: a \texttt{2{\textasciicircum}K ≤ 256} constant tuple, or computed from the code.



\end{adjustwidth}
\end{minipage}
\par\medskip

\par\medskip\noindent\begin{minipage}{\linewidth}
\hypertarget{9355066705305247355}{\texttt{SmallFloats.encode}}  -- {Method.}
\nopagebreak
\begin{adjustwidth}{2em}{0pt}


\begin{minted}{julia}
encode(T, sign, S, Q) -> UInt8   (private; design §3.3)
\end{minted}

ωEncode from canonical integer form: value = sign · S · 2{\textasciicircum}Q, with S ∈ 0:2{\textasciicircum}P (S == 2{\textasciicircum}P is the next-binade carry, draft §4.7.4 NOTE 4). Precondition: the value is in the datum set of \texttt{T} (guaranteed by RoundToPrecision ∘ Saturate).



\end{adjustwidth}
\end{minipage}
\par\medskip

\par\medskip\noindent\begin{minipage}{\linewidth}
\hypertarget{16293177486057421760}{\texttt{SmallFloats.get\_\allowbreak{}table}}  -- {Function.}
\nopagebreak
\begin{adjustwidth}{2em}{0pt}


\begin{minted}{julia}
get_table(op, fr, f1, [f2,] ρ) -> Memory{codeunit_type(fr)}
\end{minted}

Fetch (building and caching on first use) the complete result table for the pure-ρ specialization \texttt{op⟨fr; f1[, f2]⟩ under ρ}: entry \texttt{c + 1} (unary) or \texttt{(c1 << K2) + c2 + 1} (binary) holds the result code point for those operand code points. Throws for stochastic ρ, and throws when the table would exceed the byte budget — its contract is to return a table, so it says so rather than returning \texttt{nothing}. Callers that can fall back want \texttt{table\_\allowbreak{}for}.

\texttt{@noinline} by design: kernels call this once per array operation and index the returned table in their hot loop.



\end{adjustwidth}
\end{minipage}
\par\medskip

\par\medskip\noindent\begin{minipage}{\linewidth}
\hypertarget{17468751084076810785}{\texttt{SmallFloats.get\_\allowbreak{}table}}  -- {Method.}
\nopagebreak
\begin{adjustwidth}{2em}{0pt}


\begin{minted}{julia}
get_table(op, fr, f1, f2, f3, ρ) -> Memory{UInt8}
\end{minted}

Ternary form: entry \texttt{(\allowbreak{}(\allowbreak{}c1 << K2 | c2) << K3) + c3 + 1} holds the result code point. Builds unconditionally (policy lives in \texttt{\_\allowbreak{}ternary\_\allowbreak{}table\_\allowbreak{}for}); pure ρ only.



\end{adjustwidth}
\end{minipage}
\par\medskip

\par\medskip\noindent\begin{minipage}{\linewidth}
\hypertarget{10981503898423870584}{\texttt{SmallFloats.joinhead}}  -- {Method.}
\nopagebreak
\begin{adjustwidth}{2em}{0pt}

The wider of two heads. The carrier join is a max over this order, which is what lets block and scaled operations compose without a second rule.



\end{adjustwidth}
\end{minipage}
\par\medskip

\par\medskip\noindent\begin{minipage}{\linewidth}
\hypertarget{9794079948674987482}{\texttt{SmallFloats.nan\_\allowbreak{}order\_\allowbreak{}key}}  -- {Method.}
\nopagebreak
\begin{adjustwidth}{2em}{0pt}


\begin{minted}{julia}
nan_order_key(F)
\end{minted}

The order key reserved for the single NaN — above every finite key and ±Inf.

Per format, not a constant, and the key type is \texttt{orderkeytype(F)} rather than \texttt{UInt16}. The reason is arithmetic, not tidiness: a format has \texttt{2{\textasciicircum}K} code points mapping to keys \texttt{1 … 2{\textasciicircum}K}, so \texttt{UInt16} runs out at exactly K = 16 — \texttt{UInt16(c) + UInt16(1)} for \texttt{c = 0xffff} wraps \textbf{silently to 0}, putting the largest datum below the smallest. \texttt{orderkeytype} returns \texttt{UInt32} for \texttt{Code16}, which has room for the keys and for a sentinel strictly above all of them.

\emph{(§4 Stage 4 item 3, pulled forward into Stage 3 — §11 M16. The rest of Stage 3 makes wide formats fail loudly; a wraparound in the total order is the one K ≥ 9 defect that would have been silent, so it is closed here rather than carried for a commit.)}



\end{adjustwidth}
\end{minipage}
\par\medskip

\par\medskip\noindent\begin{minipage}{\linewidth}
\hypertarget{3009066085088880286}{\texttt{SmallFloats.orderkeytype}}  -- {Method.}
\nopagebreak
\begin{adjustwidth}{2em}{0pt}

Unsigned type wide enough for a format{\textquotesingle}s \texttt{2{\textasciicircum}K + 1} total-order keys.



\end{adjustwidth}
\end{minipage}
\par\medskip

\par\medskip\noindent\begin{minipage}{\linewidth}
\hypertarget{17503642904345592057}{\texttt{SmallFloats.packing\_\allowbreak{}saves}}  -- {Method.}
\nopagebreak
\begin{adjustwidth}{2em}{0pt}


\begin{minted}{julia}
packing_saves(F) -> Bool
\end{minted}

Whether \texttt{PackedVector\{F\}} is smaller than \texttt{Vector\{F\}} — \texttt{false} exactly when \texttt{bitwidth(F)} equals the storage unit{\textquotesingle}s width, i.e. K = 8 and K = 16.

\textbf{Why this is a query and not a refusal.} §4 Stage 4 item 5 called for \texttt{PackedVector} to refuse loudly at K = 16, on the grounds that packing is the identity there and a silent identity invites benchmark confusion. The grounds are right and the boundary is not: packing is \emph{equally} the identity at K = 8, which has shipped since before the extension and is enumerated by G5{\textquotesingle}s packed section over all 120 formats. A rule that fires on one of two identical cases is not a rule, and the alternative — refusing both — would be a breaking change justified by tidiness.

So the fact is exposed instead of enforced. Benchmarks and the docs read it; \texttt{PackedVector} at those widths stays correct, stays supported, and stays pointless. (§11 M21.)



\end{adjustwidth}
\end{minipage}
\par\medskip

\par\medskip\noindent\begin{minipage}{\linewidth}
\hypertarget{1763536505995871807}{\texttt{SmallFloats.project}}  -- {Method.}
\nopagebreak
\begin{adjustwidth}{2em}{0pt}


\begin{minted}{julia}
project(T, ρ, X; R=0, sticky=0) -> T
\end{minted}

Project a closed-extended-real carrier value into format \texttt{T} under ρ. \texttt{X::Float64} must be \emph{exact} (the Float64 fast path); \texttt{X::BigFloat} likewise (or an interval endpoint with \texttt{sticky ∈ \{-1,+1\}}). \texttt{R} is the stochastic draw.



\end{adjustwidth}
\end{minipage}
\par\medskip

\par\medskip\noindent\begin{minipage}{\linewidth}
\hypertarget{3869604513100043477}{\texttt{SmallFloats.promotecarrier}}  -- {Method.}
\nopagebreak
\begin{adjustwidth}{2em}{0pt}

The public promotion target for \texttt{F} against external numeric types. Always a Julia \texttt{AbstractFloat}, never the internal exact carrier.



\end{adjustwidth}
\end{minipage}
\par\medskip

\par\medskip\noindent\begin{minipage}{\linewidth}
\hypertarget{7234603244356363868}{\texttt{SmallFloats.rung}}  -- {Method.}
\nopagebreak
\begin{adjustwidth}{2em}{0pt}


\begin{minted}{julia}
rung(op::Val, Fs::Type{<:Binary}...) -> Head
\end{minted}

The carrier for evaluating \texttt{op} over operands of formats \texttt{Fs}. A \textbf{join of two lower bounds}, and both are necessary:

\begin{itemize}
\item the \textbf{monomial bound} — \texttt{op}{\textquotesingle}s exact result is a sum of monomials in at most \texttt{opfactors(op)} datum factors, so it needs \texttt{\_\allowbreak{}rungindex\_\allowbreak{}span(\allowbreak{}nf · maxB)} (\texttt{nf · maxB} rather than \texttt{ΣBᵢ} over the actual operands: it dominates every monomial shape, including \texttt{x²} from \texttt{Hypot}, which \texttt{ΣBᵢ} does not, and it coincides with the tight bound in the same-format case that is virtually all traffic);


\item the \textbf{operand bound} — \texttt{decode(v::Fᵢ)} has already produced a value on \texttt{datumcarrier(Fᵢ)}, so the head can never be narrower than any operand{\textquotesingle}s own rung regardless of what the monomial needs.

\end{itemize}
Omitting the second is the subtle error. \texttt{Abs} has one factor, so the monomial bound alone would pick rung 1 for a B = 1024 format — whose datums arrive as \texttt{Float128}, for which \texttt{ωeval(::HeadF64, …)} has no row. The join makes \texttt{lift}{\textquotesingle}s missing narrowing method a consequence rather than a hazard: the head dominates every operand, so nothing ever needs narrowing.

For the same format at two factors this agrees with \texttt{rung(F)} exactly, which is what keeps the 432/64/8 partition and every K ≤ 8 result unchanged.



\end{adjustwidth}
\end{minipage}
\par\medskip

\par\medskip\noindent\begin{minipage}{\linewidth}
\hypertarget{18216053183491810250}{\texttt{SmallFloats.rung}}  -- {Method.}
\nopagebreak
\begin{adjustwidth}{2em}{0pt}


\begin{minted}{julia}
rung(F) -> Head
\end{minted}

The starting carrier for evaluating on datums of format \texttt{F} alone. For an operation over several formats use \texttt{rung(op, Fs...)}, which additionally accounts for the operation{\textquotesingle}s factor count — a \texttt{ScaledMultiply} of two Group A formats can need a wider carrier than either operand does.



\end{adjustwidth}
\end{minipage}
\par\medskip

\par\medskip\noindent\begin{minipage}{\linewidth}
\hypertarget{5631300055872342420}{\texttt{SmallFloats.rungindex}}  -- {Method.}
\nopagebreak
\begin{adjustwidth}{2em}{0pt}

Rung index 1/2/3 of a head — for ordering and for the carrier join.



\end{adjustwidth}
\end{minipage}
\par\medskip

\par\medskip\noindent\begin{minipage}{\linewidth}
\hypertarget{2063626746625773781}{\texttt{SmallFloats.table\_\allowbreak{}for}}  -- {Method.}
\nopagebreak
\begin{adjustwidth}{2em}{0pt}


\begin{minted}{julia}
table_for(op, fr, f1[, f2], ρ) -> Union{Nothing,Memory{codeunit_type(fr)}}
\end{minted}

The kernel-facing gate for unary and binary specializations: the table if policy grants one, \texttt{nothing} if the caller should run the scalar path per element.

Separate from \texttt{get\_\allowbreak{}table} on purpose — \textbf{one name per question.} \texttt{get\_\allowbreak{}table}{\textquotesingle}s contract is {\textquotedbl}return the table{\textquotedbl}, so it throws when it cannot honour that, and that is what the suite, \texttt{conformance} and the precompile workload want. \texttt{table\_\allowbreak{}for} asks the different question {\textquotedbl}should there be a table at all{\textquotedbl}, and answers \texttt{nothing} without prejudice. Its \texttt{Union} return costs nothing because kernels call it once per array operation and branch outside the loop.



\end{adjustwidth}
\end{minipage}
\par\medskip

\par\medskip\noindent\begin{minipage}{\linewidth}
\hypertarget{14316915538520331686}{\texttt{SmallFloats.table\_\allowbreak{}keys}}  -- {Method.}
\nopagebreak
\begin{adjustwidth}{2em}{0pt}

Keys of every cached unary/binary specialization, in a deterministic order.



\end{adjustwidth}
\end{minipage}
\par\medskip

\par\medskip\noindent\begin{minipage}{\linewidth}
\hypertarget{12092491969860390893}{\texttt{SmallFloats.tablebits}}  -- {Method.}
\nopagebreak
\begin{adjustwidth}{2em}{0pt}

\texttt{log2} of the bytes a table over \texttt{Fs} with results in \texttt{fr} would occupy.



\end{adjustwidth}
\end{minipage}
\par\medskip

\par\medskip\noindent\begin{minipage}{\linewidth}
\hypertarget{3948079286591090344}{\texttt{SmallFloats.tablecache}}  -- {Method.}
\nopagebreak
\begin{adjustwidth}{2em}{0pt}

The unary/binary table cache holding results of format \texttt{fr}.



\end{adjustwidth}
\end{minipage}
\par\medskip

\par\medskip\noindent\begin{minipage}{\linewidth}
\hypertarget{17327739502832920628}{\texttt{SmallFloats.ternarycache}}  -- {Method.}
\nopagebreak
\begin{adjustwidth}{2em}{0pt}

The ternary table cache holding results of format \texttt{fr}.



\end{adjustwidth}
\end{minipage}
\par\medskip

\par\medskip\noindent\begin{minipage}{\linewidth}
\hypertarget{2228068531026386954}{\texttt{SmallFloats.unpack\_\allowbreak{}tile!}}  -- {Method.}
\nopagebreak
\begin{adjustwidth}{2em}{0pt}

Unpack \texttt{pv[first:first+len-1]} into \texttt{buf} (byte scratch); tile-granular kernel entry.



\end{adjustwidth}
\end{minipage}
\par\medskip

\par\medskip\noindent\begin{minipage}{\linewidth}
\hypertarget{8374399003204767242}{\texttt{SmallFloats.within\_\allowbreak{}byte\_\allowbreak{}budget}}  -- {Method.}
\nopagebreak
\begin{adjustwidth}{2em}{0pt}

Would this table fit the byte budget? (Memory, not build time.)



\end{adjustwidth}
\end{minipage}
\par\medskip


\part{Insights}


\chapter{Architecture Overview}



\label{12793878815436270562}{}


The internals of SmallFloats.jl organize along two axes, and every page in this Insights track is an elaboration of one or the other.



The first axis is the \textbf{carrier rung} a format{\textquotesingle}s datums decode to, chosen by the format{\textquotesingle}s exponent bias: \texttt{Float64} for the 432 formats whose full dynamic range fits Float64{\textquotesingle}s normal range (rung 1), \texttt{Float128} for formats that need more exponent room but stay in binary floating point (rung 2), and an exact dyadic carrier for the formats — 72 of them — whose range or precision breaks even Float128 (rung 3). The choice is a dispatch decision derived from the exponent bias, never a caller{\textquotesingle}s, and every rung is exact for the datums it carries. The second axis is the \textbf{rigor class} attached to every non-\texttt{Float64} computed result: Class R (unconditional — provable without any accuracy assumption) and Class E (envelope-conditional — faithful-but-not-CR estimates validated by a generous, measured slack). Together the two axes answer, for any operation on any format, {\textquotedbl}what carries the value{\textquotedbl} and {\textquotedbl}what justifies the result.{\textquotedbl}



The pages that follow work through the machinery each axis produces: the encoding and projection engine that every carrier rung feeds into (The Encoding \& Projection Engine), the oracle and the two rigor classes in full (The Oracle \& Rigor Classes), the table and kernel layer that makes the common rung fast (Tables, Kernels \& Sorting), the block layer{\textquotesingle}s exactness argument (Blocks: Exactness without Superaccumulators), and the verification and benchmark doctrine that keeps all of it honest (Verification \& Benchmark Doctrine).



\needspace{36\baselineskip}
\section{Layer map}



\label{9668567837796074949}{}


Source files load in dependency order, each layer speaking only downward:




\begin{table}[h]
\centering
\begin{tabulary}{\linewidth}{L L L}
\toprule
layer & file & provides \\
\toprule
formats & \texttt{formats.jl} & \texttt{Binary\{K,P,SGN,EXT\}}, the 504 named aliases, Group M queries \\
specs & \texttt{projspec.jl} & rounding/saturation singletons, \texttt{ProjSpec\{R,S\}}, the predefined spec grid \\
defaults & \texttt{defaults.jl} & settable session defaults (\texttt{DefaultType}, \texttt{DefaultProjection}, …) behind \texttt{Ref}s; consumed via the speculation guard \\
codec & \texttt{decode\_\allowbreak{}encode.jl} & decode (generated tables + bit-composed compute), encode, order keys, counting sort, \texttt{Class}, Next ops \\
engine & \texttt{project.jl} & \texttt{round\_\allowbreak{}to\_\allowbreak{}precision} (mask-based Float64 core + generic core), \texttt{saturate}, \texttt{project}, \texttt{project\_\allowbreak{}interval} \\
ops & \texttt{ops\_\allowbreak{}scalar.jl} & result-kind protocol, \texttt{apply\_\allowbreak{}op}, the operation registry, the spec register \\
compat & \texttt{juliacompat.jl} & the Base register: declarative op ⇒ Base-function map \\
oracle & \texttt{oracle.jl} & ω-semantics for all 52 operations \\
tables & \texttt{tables.jl} & the pure-ρ result-table cache (unary/binary + bitwidth-gated ternary) \\
kernels & \texttt{kernels.jl} & Shape-A gathers, Shape-B scalar/threaded loops, \texttt{vmap} \\
blocks & \texttt{blocks.jl} & \texttt{Block}, block/scaled ops, exact reductions \\
packed & \texttt{packed.jl} & sub-byte \texttt{PackedVector} storage \\
approx & \texttt{approx.jl} & κ measurement/registry, conformance declaration \\
rand & \texttt{rand.jl} & Random-API hooks: uniform-[0,1) floor projection, clamped normal \\
\bottomrule
\end{tabulary}

\end{table}



\section{see also}



\label{9115667341246479961-dup21}{}


\href{under\_\allowbreak{}engine.md}{The Encoding \& Projection Engine}, \href{under\_\allowbreak{}oracle.md}{The Oracle \& Rigor Classes}, \href{under\_\allowbreak{}tables.md}{Tables, Kernels \& Sorting}, \href{under\_\allowbreak{}blocks.md}{Blocks: Exactness without Superaccumulators}, \href{under\_\allowbreak{}verification.md}{Verification \& Benchmark Doctrine}.



\chapter{The Encoding \& Projection Engine}



\label{10357628553938299882}{}


\section{Encoding and decoding}



\label{8606053261248653135}{}


A value is its code point (\texttt{UInt8}). The bit layout is the draft{\textquotesingle}s: sign (signed formats), biased exponent, trailing significand; one NaN at the negative-zero slot; no −0; ±Inf adjacent to the extremes in extended formats.



\texttt{decode} has \textbf{two shapes, selected by \texttt{decodepolicy(T)}} — a representation decision, not a semantic one.



At \texttt{K ≤ 8} it is a \textbf{\texttt{@generated} constant-tuple lookup}: per format, a \texttt{2{\textasciicircum}K}-tuple of datums built once from the computational decode (so table and computation are correct by construction and asserted equivalent exhaustively). Constant inputs still fold — \texttt{maxfinite\_\allowbreak{}datum(T)} is a compile-time constant — while runtime decode is a single indexed load (≈ 0.7 ns). Above \texttt{K = 8} a \texttt{2{\textasciicircum}16}-tuple is not a constant worth materializing, so the computational decode runs directly.



The computational decode assembles the datum \textbf{by bits} (normalize with \texttt{leading\_\allowbreak{}zeros}, place exponent and mantissa fields, \texttt{reinterpret}) wherever the carrier is \texttt{Float64}.



\textbf{Which carrier that is depends on the format{\textquotesingle}s exponent bias, and is the second of the package{\textquotesingle}s two axes.} \texttt{Float64}{\textquotesingle}s normal range holds every datum of the 432 formats at rung 1 — the property the bit-assembly rests on. It does not hold the other 72: \texttt{Binary16p1uf} reaches \texttt{2{\textasciicircum}32768}. Those decode to \texttt{Float128} (rung 2) or to an exact dyadic carrier (rung 3), chosen by \texttt{datumcarrier(T)}, and every one of them is exact for the datums it accepts. So the invariant is not {\textquotedbl}\texttt{Float64} is the exact carrier for all datums{\textquotedbl} — that was true only while \texttt{K ≤ 8} — but the stronger and still-checkable one:



\begin{quote}
\textbf{every format{\textquotesingle}s datum carrier is exact for that format{\textquotesingle}s datums}, and the choice is a dispatch decision derived from the exponent bias, never a caller{\textquotesingle}s.

\end{quote}


The suite checks it against an independent big-float transliteration of the draft (gate G6, exhaustively over the \texttt{(datum, head)} pairs), and gate G4 checks the consequence that matters: forcing a \emph{wider} carrier never changes a code point, so the choice buys time and never an answer.



Ordering runs on \textbf{integer order keys}: a sign-magnitude fold into an unsigned type wide enough for the format{\textquotesingle}s \texttt{2{\textasciicircum}K + 1} keys (\texttt{UInt16} at K ≤ 8, wider above), monotone with the total order, NaN mapped to \texttt{typemax}. Same-format \texttt{TotalOrder}, \texttt{isless}, and the numeric comparisons are key comparisons ({\textasciitilde}1 ns); since a format has at most \texttt{2{\textasciicircum}K + 1} distinct keys, vectors sort with an \textbf{O(n) counting sort} installed via \texttt{Base.Sort.defalg} (forward and reverse orderings; anything exotic falls back to the stock algorithm).



\section{The projection engine}



\label{8339828346285081570}{}


\texttt{project(\allowbreak{}fr, ρ, X; R, sticky)} is the single write path into a code point:




\nopagebreak[4]
\begin{minted}{text}
RoundToPrecision  →  Saturate  →  Encode
\end{minted}



\textbf{RoundToPrecision} produces a \texttt{Rounded(\allowbreak{}kind, sign, S, Q)} — an exact scaled significand \texttt{S} and exponent \texttt{Q} per the draft{\textquotesingle}s \texttt{Q = max(\allowbreak{}⌊log₂|X|⌋, 1−B) − P + 1}. Two implementations, proven equivalent exhaustively:



\begin{itemize}
\item the \textbf{generic core} (\texttt{\_\allowbreak{}rtp\_\allowbreak{}core}) works on any carrier (\texttt{Float64}, \texttt{Float128}, \texttt{BigFloat}) via exact power-of-two scaling and a fraction ν compared against the mode{\textquotesingle}s decision points;


\item the \textbf{mask-based Float64 core} (\texttt{\_\allowbreak{}rtp\_\allowbreak{}f64}) extracts sign/exponent/significand fields directly, represents ν as a 128-bit fixed-point integer with an OR-mask sticky for bits shifted out, and evaluates every mode — including the stochastic \texttt{⌊ν·2{\textasciicircum}N⌋} comparisons and \texttt{RNITE} ties — as integer field tests. Specials, zeros, and (Convert-only-reachable) subnormal Float64 inputs bail to the generic core.

\end{itemize}


\textbf{Symbolic sticky} is how enclosures talk to the engine: \texttt{sticky ∈ \{−1, 0, +1\}} declares the true value to be the carried value plus an infinitesimal of that sign. The engine folds it into every comparison; the delicate case — true value just \emph{below} a representable dyadic — decrements into the previous binade and sets ν = 1⁻ (encoded in the mask core as an all-ones fixed-point fraction, which reproduces every predicate including the RNITE tie behavior). This is what lets directed modes land exactly on asymptotes: \texttt{tanh → 1⁻} under \texttt{TowardNegative} projects to \texttt{NextLessThan(1)}, automatically.



\textbf{Saturate} classifies against the format{\textquotesingle}s range with two integer comparisons per side, then maps the classification through the draft{\textquotesingle}s twenty saturation rows to \texttt{as-is /\allowbreak{} MaxFinite /\allowbreak{} MinFinite /\allowbreak{} ±Inf /\allowbreak{} NaN}. The comparisons stay cheap because:



\begin{itemize}
\item the extremal magnitude in canonical \texttt{(S, Q)} form is a constant-folded function of the type parameters;


\item the rounded value and the extremum share one lexicographic \texttt{(Q, S)} order (subnormals and the lowest normal binade share the same \texttt{Q});


\item signed formats use sign–magnitude symmetry, and unsigned underflow is just the sign bit;


\item an internal \texttt{HUGEQ} sentinel represents {\textquotedbl}finite but astronomically large{\textquotedbl} (the ν = 1⁻ image of an infinite endpoint), so directed \texttt{SatNone} clamps it to MaxFinite correctly.

\end{itemize}


\textbf{Encode} is pure integer bit assembly, including the significand-carry renormalization and the subnormal/normal field split.



\section{see also}



\label{9115667341246479961-dup22}{}


\href{under\_\allowbreak{}oracle.md}{The Oracle \& Rigor Classes}, \href{under\_\allowbreak{}tables.md}{Tables, Kernels \& Sorting}, \href{under\_\allowbreak{}recipes.md}{Internals Recipes}.



\chapter{The Oracle \& Rigor Classes}



\label{12993964662093357484}{}


Every operation{\textquotesingle}s defined result is computed by \texttt{ωeval}, which returns one of five result kinds; \texttt{apply\_\allowbreak{}op} fast-splits the common one and finishes the rest:




\begin{table}[h]
\centering
\begin{tabulary}{\linewidth}{L L L}
\toprule
kind & meaning & finished by \\
\toprule
\texttt{Float64} & exact (specials; representable arithmetic) & direct \texttt{project} \\
\texttt{Dyadic} & exact at rung 3: an \texttt{Int128} significand and an \texttt{Int64} exponent, \texttt{isbits} & direct \texttt{project} (dyadic carrier) \\
\texttt{Float128} & exact by \textbf{width analysis} & direct \texttt{project} (Float128 carrier) \\
\texttt{StickyF} & wide-spread \texttt{FMA}/\allowbreak\texttt{FAA} tail: head value (\texttt{Float64} or \texttt{Float128}) + tail sign & direct \texttt{project} with \texttt{sticky} set — no allocation \\
\texttt{BigExactF} & exact at a precision \textbf{derived from the operands} — \texttt{bigprec}, not a constant (wide-spread tail for \texttt{Add}; the near-impossible \texttt{FAA} distillation miss) & \texttt{project} on \texttt{BigFloat} \\
\texttt{Enclose128F} & correctly-rounded Float128 \textbf{bracket} & sticky agreement, MPFR fallback \\
\texttt{EncloseF} & MPFR directed enclosure \texttt{f(prec)}, optional Float128 pre-filter \texttt{fq}, optional eager Float64 estimate \texttt{yd} & three-stage: \texttt{yd} → \texttt{fq} → interval protocol \\
\bottomrule
\end{tabulary}

\end{table}



\section{The Dyadic↔Rational bridge}



\label{15909861321242845690}{}


The dyadic carrier converts to and from \texttt{Rational} \textbf{exactly}, and that bridge is the cleanest way to reason about a rung-3 value:




\nopagebreak[4]
\begin{minted}{julia}
using SmallFloats.DyadicNumbers: dyadic_to_rational, rational_to_dyadic, isdyadic

dyadic_to_rational(x)          # Rational{BigInt}, exact — a Dyadic IS S·2^Q
rational_to_dyadic(3 // 4)     # exact, or throws
isdyadic(1 // 3)               # false — no Dyadic represents it
\end{minted}



Three properties are worth knowing. The infinities convert (\texttt{±1//0}, matching \texttt{Rational\{BigInt\}(Inf)}) and NaN does not, because \texttt{Rational} has no NaN slot. A non-dyadic rational is \textbf{refused rather than rounded} — rounding it here would be a rounding outside \texttt{project}. And the round trip is a law about \emph{values}, not fields: \texttt{Dyadic}{\textquotesingle}s significand is deliberately unnormalized, so \texttt{6·2⁻²} and \texttt{3·2⁻¹} are the same value and round-trip to the canonical one.



The full design, with every edge case and the three laws, is in \texttt{other/\allowbreak{}dyadic\_\allowbreak{}rational.md}.



\section{The two rigor classes}



\label{15034539620333151856}{}


Two \textbf{rigor classes} govern every non-\texttt{Float64} path, and their arguments are never mixed:



\textbf{Class R (unconditional)} rests on two facts that hold without any accuracy assumption:



\begin{itemize}
\item \textbf{Width analysis.} Sums of decoded datums are \emph{exactly representable} in \texttt{Float128} whenever operand bits + exponent spread fit 113 bits — checked in advance by integer exponent arithmetic (\texttt{Add} at ΔE ≤ 100, \texttt{FMA} ≤ 92, \texttt{FAA} span ≤ 98). Beyond the threshold, \texttt{Add} escalates to the exact MPFR path — at a precision \textbf{derived from the operands} by \texttt{bigprec}, not the retired 2200-bit constant, which was ample at K ≤ 8 and insufficient for 50 of the 504 formats (gate G2 measures this) — while \texttt{FMA}/\allowbreak\texttt{FAA} take a non-allocating \textbf{sticky-head} shortcut: past that spread the smaller term is provably too small (by construction of the threshold) to affect anything but the tail direction of the larger one, so the result is \texttt{StickyF(head, sign)} — no \texttt{BigFloat} involved. \texttt{FAA}{\textquotesingle}s three-term case runs a bounded Float128 2sum-distillation (Priest-style, ≤ 6 sweeps) to find that head/tail split directly; the residual \texttt{BigExactF} MPFR fallback exists only for the near-impossible case the distillation doesn{\textquotesingle}t converge.


\item \textbf{IEEE correct rounding.} Division and \texttt{sqrt} are correctly rounded at \emph{every} binary width, by mandate. \texttt{Divide}/\allowbreak\texttt{Recip} therefore use the \textbf{Float64} quotient directly: it is within half an ulp of truth, so it serves as \texttt{EncloseF}{\textquotesingle}s eager \texttt{yd} estimate with no Float128 arithmetic on the path at all. \texttt{Sqrt}/\allowbreak\texttt{RSqrt} use the \textbf{Float128} CR result: an inexact nearest-CR value \texttt{q} brackets the truth in the \emph{open} interval \texttt{(\allowbreak{}prevfloat(\allowbreak{}q), nextfloat(\allowbreak{}q))}. Exactness itself is detected by an \texttt{fma} residual test.

\end{itemize}


\textbf{Class E (envelope-conditional).} Faithful-but-not-CR evaluations stand in for an enclosure only inside a generous envelope, resolved by the two-sided sticky gate. Two stages, cheapest first:



\begin{itemize}
\item \emph{Float64 stage:} for operations whose Float64 libm is faithful (≤ 1 ulp ≈ 2⁻⁵² relative — the exp/log families, the hyperbolics, forward trig inside the |x| ≤ 10¹⁵ reduction window, \texttt{atan}, \texttt{asinh}, the softplus composition), the eager estimate \texttt{yd} carries envelope \texttt{E = |yd|·2⁻⁴⁵}, ≥ 2⁷ slacker than any faithful-libm error (measured margin on this machine: ≥ 2⁶·⁶).


\item \emph{Float128 stage:} libquadmath{\textquotesingle}s estimate \texttt{y = fq()} carries \texttt{E = |y|·2⁻⁹⁰}, at least 2¹⁸ slacker than any published libquadmath bound.

\end{itemize}


In both stages, if the sticky-projected endpoints agree, that code point is the answer; if not, the next stage (ultimately the MPFR ladder) decides. An envelope failure can therefore only cost speed, never correctness — unless both endpoints agree \emph{and} the envelope is wrong, which the differential-build tests rule out empirically by building every standard table twice (Float128-first and pure-MPFR) and byte-diffing, and which the exhaustive oracle cross-check re-verifies per operation against the 3072-bit protocol.



\section{The interval protocol}



\label{13833006351776399793}{}


\textbf{The interval protocol} (\texttt{project\_\allowbreak{}interval}) is the termination backbone: evaluate the MPFR enclosure at precision \texttt{p}; if \texttt{lo == hi} the value is exactly representable — project it; otherwise project \texttt{lo} with \texttt{sticky = +1} and \texttt{hi} with \texttt{sticky = −1}; agreement (projection is monotone at fixed \texttt{R}) yields the answer; disagreement doubles \texttt{p}, clamped to the \texttt{maxprec} ceiling (\textbf{derived from the format}, \texttt{max(\allowbreak{}4096, bigprec(\allowbreak{}T) + 64)} — a flat 4096 was ample at K ≤ 8 and could not decide \texttt{ArcCosPi} at \texttt{Binary15p2ue} under a directed or stochastic rule, honored exactly even when it is not a power-of-two multiple of the 256-bit start). Termination requires that the enclosure not be chasing a value the grid can actually hit — which is why the π-scaled operations carry \textbf{Niven peels}: for dyadic arguments, \texttt{tan(πr)} takes rational values only at the quarter-integers (±1) and \texttt{atan2/π} only on the diagonals (±¼, ±¾); those cases are answered exactly before any enclosure is built, and Niven{\textquotesingle}s theorem proves the peel set complete.



\section{see also}



\label{9115667341246479961-dup23}{}


\href{under\_\allowbreak{}engine.md}{The Encoding \& Projection Engine}, \href{under\_\allowbreak{}tables.md}{Tables, Kernels \& Sorting}, \href{howto\_\allowbreak{}add\_\allowbreak{}operation.md}{Add an Operation}.



\chapter{Tables, Kernels \& Sorting}



\label{1013018278037228102}{}


\section{Order keys and counting sort}



\label{9176360583931610048}{}


Ordering runs on \textbf{integer order keys}: a sign-magnitude fold into an unsigned type wide enough for the format{\textquotesingle}s \texttt{2{\textasciicircum}K + 1} keys (\texttt{UInt16} at K ≤ 8, wider above), monotone with the total order, NaN mapped to \texttt{typemax}. Same-format \texttt{TotalOrder}, \texttt{isless}, and the numeric comparisons are key comparisons ({\textasciitilde}1 ns); since a format has at most \texttt{2{\textasciicircum}K + 1} distinct keys, vectors sort with an \textbf{O(n) counting sort} installed via \texttt{Base.Sort.defalg} (forward and reverse orderings; anything exotic falls back to the stock algorithm).



\section{Tables and kernels}



\label{17163678065643276217}{}


For pure specs, unary and binary operations are \textbf{finite functions} — 256 or 65 536 entries — so the kernel layer materializes them once per \texttt{(op, formats, ρ)} into a locked cache (\texttt{Dict\{TableKey, Memory\{UInt8\}\}}, double-checked locking, builds outside the lock) and serves every later array call as a gather: Shape-A, one load per element, measured 0.27 ns/elem unary and 0.5 ns/elem binary. Tables are built \emph{through the scalar path}, so they inherit its bit-exactness; the suite asserts table ≡ scalar over every entry.



Ternary (\texttt{FMA}, \texttt{FAA}, \texttt{Clamp}) is a finite function too — 2{\textasciicircum}(K1+K2+K3) entries — but that count spans four orders of magnitude across the 3–16 bitwidth range (512 B at K=3, 16 MiB at K=8), so one policy doesn{\textquotesingle}t fit the whole range. A separate \texttt{TernaryKey → TernaryEntry} cache (\texttt{\_\allowbreak{}ternary\_\allowbreak{}table\_\allowbreak{}for}, in \texttt{tables.jl}) tiers by Σ bitwidth:



\begin{itemize}
\item \textbf{Eager} (≤ 18 bits, all \texttt{K ≤ 6} combinations, ≤ 256 KiB): builds and caches on the first array call.


\item \textbf{Adaptive} (≤ 21 bits, the \texttt{K = 7} band, up to 2 MiB): accumulates a per-signature element count across calls and builds only once a signature has processed enough elements to amortize the build; a byte-bounded LRU eviction (\texttt{TERNARY\_\allowbreak{}CACHE\_\allowbreak{}BYTES}) guards against many hot signatures coexisting.


\item \textbf{Compute} (\texttt{K = 8}, 16+ MiB — a table is never worth it): \texttt{vmap!} runs Shape-B, the fully specialized scalar pipeline per element, optionally split across \texttt{Threads.@threads} for long enough arrays (each ternary draw is independent under a fixed, non-stochastic ρ, so lanes cannot interact).

\end{itemize}


Every ternary table entry, eager or adaptive, is still built \emph{through the scalar path} — the tiering changes when/whether the cache exists, never what it contains. Stochastic calls of any arity always take Shape-B, with the RNG resolved once per array rather than per element.



\section{see also}



\label{9115667341246479961-dup24}{}


\href{under\_\allowbreak{}engine.md}{The Encoding \& Projection Engine}, \href{under\_\allowbreak{}oracle.md}{The Oracle \& Rigor Classes}, \href{under\_\allowbreak{}recipes.md}{Internals Recipes}.



\chapter{Blocks: Exactness without Superaccumulators}



\label{15688307216380555725}{}


\texttt{blockdecode} produces each lane{\textquotesingle}s \texttt{scale × element} exactly (≤ 17-bit significands in Float64). Reductions then apply \textbf{span filters}: one integer pass over lane exponents decides whether the whole sum (or dot product, with ≤ 33-bit lane products) is exactly representable in \texttt{Float128}; if so, a plain \texttt{Float128} accumulation \emph{is} the exact answer; if not, an exact big-float accumulation takes over. Either way there is exactly one projection, at the end. \texttt{ωBlockProject} follows the draft{\textquotesingle}s special rows for the scale (NaN, 0, ±Inf) and divides each element result by the scale through its own cheapest-first CR-bracket / enclosure cascade (exact Float64 quotient → CR Float128 → bracket/pre-filter → MPFR interval), mirroring the scalar quotient group{\textquotesingle}s rigor arguments.



\section{see also}



\label{9115667341246479961-dup25}{}


\href{under\_\allowbreak{}oracle.md}{The Oracle \& Rigor Classes}, \href{under\_\allowbreak{}tables.md}{Tables, Kernels \& Sorting}, \href{howto\_\allowbreak{}add\_\allowbreak{}operation.md}{Add an Operation}.



\chapter{Verification \& Benchmark Doctrine}



\label{5698011395799222450}{}


\section{Verification doctrine}



\label{12691031809522550080}{}


The value sets are small enough that sampling is never necessary, so the suite enumerates — ≈ 8.9 M assertions in all:



\begin{itemize}
\item formats against an independent draft transliteration (14 679);


\item ordering over all 2.5 M same-format pairs plus Next-op edge tables (7.6 M);


\item every unary operation on every input against a 3072-bit protocol run; divide and the ternaries exhaustively;


\item stochastic R-sweeps with directed-asymptote pins;


\item table ≡ scalar over every entry, including the ternary tiers (eager and adaptive) against the scalar path;


\item the sticky-head \texttt{FMA}/\allowbreak\texttt{FAA} escalation against the MPFR reference across every rounding-mode family and adversarial cancellation cases;


\item blocks against a from-scratch reference composition;


\item Float128 carrier ≡ Float64 carrier, and Float128-first ≡ MPFR differential builds;


\item the mask rounding core ≡ the generic core over datums and reachable sums/products at boundary stochastic budgets N ∈ \{45, 60\};


\item packed round-trips, and κ/conformance behavior.

\end{itemize}


Deterministic \textbf{specialization regressions} (concrete inferred return types at the public entry points, zero warm-path allocation) stand in for timing assertions.



\section{Benchmark doctrine}



\label{14987590650541321188}{}


Recorded after two measurement post-mortems: a benchmark closure over any non-\texttt{const} global measures Julia{\textquotesingle}s dispatch machinery, not the code under test — and it distorts \emph{ratios}, not just absolutes, because dispatch cost varies with call shape (a dynamic keyword call costs {\textasciitilde}1 µs; a dynamic positional call far less; six unresolved interior sites cost six times one).



The rules the shipped \texttt{benchmarking/\allowbreak{}benchmarking.jl} enforces structurally:



\begin{itemize}
\item format types enter as type parameters, never as globals;


\item operands come from untimed setup;


\item functions retrieved reflectively pass through argument barriers to specialize;


\item a preflight aborts the run if warm scalar paths allocate — including the wide-spread \texttt{FMA}/\allowbreak\texttt{FAA} sticky-head path.

\end{itemize}


The table-build section reports both the cold build (cache evicted per sample) and the steady-state warm cache hit, since callers amortize the former through the latter. Vary one binding per variant; verify specialization before believing a number.



\section{Deliberate limitations}



\label{16229184572553167566}{}


No implicit cross-format arithmetic (promotion is to \texttt{Float64}, explicitly). No in-place packed arithmetic. Threading is opt-in and narrow: only the untabled ternary (\texttt{K = 8}) compute kernel threads, and only above a size cutoff and when \texttt{Threads.nthreads() > 1}; every other kernel is single-threaded. \texttt{Irrational}/\allowbreak\texttt{Rational} inputs to \texttt{Convert} are rejected rather than double-rounded silently. The \texttt{Float128} machinery never changes results — disabling it (\texttt{SmallFloats\_\allowbreak{}Float128=disable}) is a tested no-op semantically.



\section{see also}



\label{9115667341246479961-dup26}{}


\href{under\_\allowbreak{}tables.md}{Tables, Kernels \& Sorting}, \href{howto\_\allowbreak{}benchmark.md}{Benchmark Correctly}, \href{howto\_\allowbreak{}verify\_\allowbreak{}custom.md}{Verify Custom Code}.



\chapter{Add an Operation}



\label{9434017197269608002}{}


How to add a new operation to SmallFloats.jl so that it inherits everything the existing catalog has: bit-exact defined results, the full projection-mode vocabulary, generated scalar/array/block surfaces, table caching, exhaustive test coverage, benchmark pickup, and a truthful conformance declaration. Read the Architecture Overview and The Oracle \& Rigor Classes first — this page assumes their vocabulary (the result-kind protocol, the two rigor classes, the interval protocol, the \texttt{yd} envelope stage).



\begin{tcolorbox}[toptitle=-1mm,bottomtitle=1mm,colback=admonition-note!50!white,colframe=admonition-note,title=\textbf{What{\textquotesingle}s generated vs what you write}]
\textbf{Generated for you}, from one registry line: the spec-register method \texttt{Op(\allowbreak{}fr, ρ, operands…; rng, R)}, the same-format convenience method \texttt{Op(x…)}, array methods (Shape-A table gathers or Shape-B loops), the \texttt{Block}/\allowbreak\texttt{Scaled} variants, the export, table cacheability, the conformance listing, and test-harness pickup.

\textbf{You must write}: the ω-mathematics (\texttt{ωeval} in \texttt{src/oracle.jl}) and the six universal duties below — special rows, the result-kind protocol, termination (the Niven duty), envelope justification, ladder validity, and running the gates.

\end{tcolorbox}


The architecture is registry-driven: for the supported arities, \textbf{most of the surface is generated}. Adding an operation is mostly (1) one registry line and (2) one rigorous \texttt{ωeval} method — plus the \emph{duties} that keep the package{\textquotesingle}s guarantees true. Every subsection below states the duties explicitly; they are not optional.



\section{Universal duties (every new operation)}



\label{15376602000632183864}{}


\begin{itemize}
\item[1. ] \textbf{Special rows are spec, not optimization.} Write out the NaN/±Inf/zero rows explicitly in \texttt{ωeval}, in draft style, before any general evaluation. Mark any behavior the draft text under-determines with \texttt{[interp]}.


\item[2. ] \textbf{Result-kind protocol.} \texttt{ωeval} must return one of \texttt{Float64 | Float128 | StickyF | BigExactF | EncloseF | Enclose128F}, with \texttt{Float64}/\allowbreak\texttt{Float128} reserved for \emph{provably exact} results and \texttt{StickyF} for exact-head-plus-tail-direction results whose soundness bound is discharged at the emitting site (see the \texttt{FMA}/\allowbreak\texttt{FAA} wide-spread escalations). Keep the return union narrow (ideally \texttt{\{Float64, EncloseF\}}) — a wide union defeats inference and costs allocations (measured: a 4-type union turned a 50 ns op into 240 ns with 2 allocations).


\item[3. ] \textbf{Termination (the Niven duty).} \texttt{project\_\allowbreak{}interval} terminates only if the enclosure is not chasing a value the projection grid can actually hit. Any input whose \emph{exact} result is a representable dyadic must be \textbf{peeled or detected} before an enclosure is built. For transcendental ops this is usually a finite set of special rows (prove completeness, as the π-family does with Niven{\textquotesingle}s theorem); for algebraic ops it is an exactness test.


\item[4. ] \textbf{Envelope justification.} If you supply an eager \texttt{yd} estimate, you owe an error bound: the true value must lie within \texttt{|yd|·2⁻⁴⁵} (\texttt{\_\allowbreak{}F64\_\allowbreak{}RELEXP}). Acceptable bases, strongest first: IEEE-CR hardware ops (÷, sqrt: ≤ ½ ulp, unconditional); short compositions of CR ops (≤ {\textasciitilde}1.5 ulp); faithful libm (≤ 1 ulp, an empirical-but-published claim); short faithful compositions with condition number ≤ 1 (≤ {\textasciitilde}3 ulp). \textbf{Measure it}: sample your estimator against ≥ 256-bit MPFR truth over the op{\textquotesingle}s domain, including the treacherous regions (near zeros of the result, domain boundaries), and demand ≥ 2⁶ slack under 2⁻⁴⁵. The same duty applies to a Float128 \texttt{fq} estimator against its 2⁻⁹⁰ envelope. Beware \textbf{argument-error amplification}: an estimate computed as \texttt{f(g(x))} where \texttt{g} rounds (e.g. \texttt{sin(π·r)} in Float64) can be catastrophically worse near zeros of \texttt{f} than a native result-domain evaluation (\texttt{sinpi(r)}) — measure the composition you actually ship.


\item[5. ] \textbf{Ladder validity.} The MPFR closure must return a genuine directed enclosure. Whole-expression \texttt{setrounding} is valid only when the composition is monotone in every intermediate (state the argument in a comment, as \texttt{Softplus} does); otherwise round each step in the correct direction explicitly (as \texttt{\_\allowbreak{}mpfr\_\allowbreak{}rsqrt} does for the anti-monotone \texttt{1/√x}).


\item[6. ] \textbf{Gates.} Run the full suite. The unary exhaustive ×hp gate and the coverage check pick up registry ops automatically; extend the hand-curated lists where applicable (monotonicity ops, binary exhaustive sweeps). Then regenerate the benchmark report — \texttt{bench\_\allowbreak{}scalar\_\allowbreak{}ops} iterates the registry lists, so your op appears in all four operand-class tables automatically; \textbf{add a \texttt{\_\allowbreak{}SAFE\_\allowbreak{}DOMAINS} entry} in \texttt{benchmarking/\allowbreak{}benchmarking.jl} if the op is argument-restricted, or the safe-args table will use the oracle-derived fallback pool.

\end{itemize}


\section{Scalar operations}



\label{11509067801894334317}{}


The registry lives in \texttt{src/ops\_\allowbreak{}scalar.jl}. Adding a name to \texttt{\_\allowbreak{}UNARY\_\allowbreak{}OPS} / \texttt{\_\allowbreak{}BINARY\_\allowbreak{}OPS} / \texttt{\_\allowbreak{}TERNARY\_\allowbreak{}OPS} causes, with no further work:



\begin{itemize}
\item registration in \texttt{OP\_\allowbreak{}REGISTRY} (group \texttt{:B}/\allowbreak\texttt{:C}/\texttt{:A} per the loop defaults — override with an explicit \texttt{register\_\allowbreak{}op!} if the default group is wrong);


\item the generated spec-register method \texttt{Op(\allowbreak{}fr, ρ, operands…; rng, R)} and the same-format convenience \texttt{Op(x…)};


\item generated array methods routing through \texttt{vmap} (Shape-A tables for pure ρ — always at arity ≤ 2, bitwidth-gated at arity 3 — Shape-B scalar/threaded loops otherwise);


\item generated \texttt{BlockOp} and \texttt{ScaledOp} variants (\texttt{src/blocks.jl});


\item the export (\texttt{src/SmallFloats.jl} loops the registry);


\item table cacheability, conformance listing, and test-harness coverage.

\end{itemize}


What is \emph{not} generated is the mathematics: \texttt{ωeval} in \texttt{src/oracle.jl}.



\subsection{Adding a unary operation}



\label{8519510493258928519}{}


Worked example: \textbf{\texttt{Cbrt}} (cube root), chosen because it exercises the hardest duty — exactness detection for an algebraic op.



\textbf{Step 1 — registry} (\texttt{src/ops\_\allowbreak{}scalar.jl}): append \texttt{:Cbrt} to \texttt{\_\allowbreak{}UNARY\_\allowbreak{}OPS}. The loop registers it as group \texttt{:B} with arity 1.



\textbf{Step 2 — ω-semantics} (\texttt{src/oracle.jl}), with every duty discharged.



\begin{tcolorbox}[toptitle=-1mm,bottomtitle=1mm,colback=admonition-warning!50!white,colframe=admonition-warning,title=\textbf{Type the row `::AbstractFloat`, not `::Float64`}]
An enclosure-ladder row is reached at \textbf{every rung}. \texttt{\_\allowbreak{}LADDER\_\allowbreak{}OPS} is derived by subtraction in \texttt{oracle.jl} — whatever is neither an exact selection, nor exact arithmetic, nor \texttt{Convert} — so a new registry operation automatically gets \texttt{\_\allowbreak{}ωf128}, \texttt{\_\allowbreak{}ωexact} and \texttt{\_\allowbreak{}ωdyadic} rows that forward to \emph{this} method with \texttt{Float128}, \texttt{BigFloat} or converted-from-\texttt{Dyadic} operands.

A row typed \texttt{::Float64} therefore works for the 432 formats at rung 1 and raises a \texttt{MethodError} for the other 72. The package types 22 of its own ladder rows \texttt{::AbstractFloat} for exactly this reason; the 6 that are \texttt{::Float64} are the exact-arithmetic rows, which have per-head siblings written separately.

Gate G10 catches the omission — it calls every registry operation at every rung and asserts only that the call returns — and G4 asserts that \texttt{\_\allowbreak{}EXACT\_\allowbreak{}SELECTION}, \texttt{\_\allowbreak{}EXACT\_\allowbreak{}ARITH}, \texttt{\_\allowbreak{}LADDER\_\allowbreak{}OPS} and \texttt{Convert} partition \texttt{OP\_\allowbreak{}REGISTRY} exactly, so an operation cannot be classified into none of them and silently lose its wide-rung rows.

\end{tcolorbox}



\begin{minted}{julia}
# Exact-cube detection (termination duty): x = ±m·2^(3k) with m an odd perfect
# cube ⇒ cbrt(x) is a Float64-exact dyadic that project_interval must never
# chase. Oracle-path cost is irrelevant; correctness of the peel is everything.
function _exact_cbrt(x::Float64)
    fr, e = frexp(abs(x))               # |x| = fr·2^e, fr ∈ [0.5, 1)
    m = Int64(ldexp(fr, 53)); e -= 53   # |x| = m·2^e exactly
    tz = trailing_zeros(m); m >>= tz; e += tz
    e % 3 == 0 || return nothing
    r = round(Int64, cbrt(Float64(m)))
    for c in max(r - 1, 1):(r + 1)      # cbrt is faithful ⇒ true root ∈ r ± 1
        Int128(c)^3 == Int128(m) && return flipsign(ldexp(Float64(c), e ÷ 3), x)
    end
    nothing
end

function ωeval(::Val{:Cbrt}, x::AbstractFloat)     # ::AbstractFloat, NOT ::Float64
    isnan(x) && return NaN                      # special rows first: they are spec
    isinf(x) && return x                        # cbrt(±∞) = ±∞
    iszero(x) && return 0.0
    c = _exact_cbrt(x)
    c === nothing || return c                   # exact dyadic result — peeled
    # yd: Float64 cbrt is faithful (≤ 1 ulp ≈ 2⁻⁵²; measure it) ⇒ ≥ 2⁷ slack
    # under the 2⁻⁴⁵ envelope. fq: libquadmath cbrtq under the 2⁻⁹⁰ envelope.
    # Ladder: cbrt is monotone ⇒ whole-expression directed rounding encloses.
    _encl1(cbrt, x; fq=() -> cbrt(Float128(x)), yd=cbrt(x))
end
\end{minted}



\textbf{Step 3 — optional surface}: a Base veneer (\texttt{Base.cbrt(\allowbreak{}x::Binary) = Cbrt(\allowbreak{}x)} next to the others in \texttt{ops\_\allowbreak{}scalar.jl}), a docstring, and — since cbrt is monotone — its symbol in the test suite{\textquotesingle}s monotonicity list.



\textbf{Step 4 — validate}: envelope measurement for \texttt{cbrt} (include tiny and huge \texttt{x}; cbrt has no dangerous regions — the result is never near zero except at zero, which is peeled), then the full suite: the exhaustive ×hp gate now checks \texttt{Cbrt} on every code point of the reference format under four modes against the 3072-bit protocol, automatically.



\subsection{Adding a binary operation}



\label{7660721820420346737}{}


Worked example: \textbf{\texttt{LogAddExp}} — \texttt{log(eˣ + eʸ)} — chosen because it exercises composed estimators and the ladder-monotonicity argument.



\textbf{Step 1 — registry}: append \texttt{:LogAddExp} to \texttt{\_\allowbreak{}BINARY\_\allowbreak{}OPS} (default group \texttt{:C}; the special five arithmetic names get \texttt{:A}).



\textbf{Step 2 — ω-semantics}. One generic evaluator serves all three precisions — Float64 (\texttt{yd}), Float128 (\texttt{fq}), and BigFloat (the ladder):




\nopagebreak[4]
\begin{minted}{julia}
# Stable form: h + log1p(exp(l − h)), h = max, l = min. Every step is monotone
# nondecreasing in (x, y), so whole-expression directed rounding in _mpfr2 is a
# valid enclosure (the same argument Softplus records).
_logaddexp(a, b) = (h = max(a, b); l = min(a, b); h + log1p(exp(l - h)))

function ωeval(::Val{:LogAddExp}, x::AbstractFloat, y::AbstractFloat)
    (isnan(x) | isnan(y)) && return NaN
    (x == Inf || y == Inf) && return Inf        # ∞ dominates
    x == -Inf && return y                       # log(0 + eʸ) = y, exactly
    y == -Inf && return x
    # yd: two faithful libm calls plus additions, condition ≤ 1 throughout
    # (result ≥ max(x, y)) ⇒ ≤ ~3 ulp ≈ 2⁻⁵⁰ — 2⁵ slack; measure it, including
    # x ≈ y (where log1p(exp(0)) dominates) and widely separated operands.
    _encl2(_logaddexp, x, y;
           fq=() -> _logaddexp(Float128(x), Float128(y)),
           yd=_logaddexp(x, y))
end
\end{minted}



\textbf{Termination note}: can \texttt{logaddexp(x, y)} be exactly dyadic for dyadic operands with the special rows peeled? \texttt{log(eˣ + eʸ) = d} requires \texttt{eˣ + eʸ = e{\textasciicircum}d} — by Lindemann–Weierstrass this forces relations that dyadic arguments cannot satisfy except on the peeled rows, so no further peel is needed; \textbf{record that argument in a comment}. When you cannot prove such a statement for your op, hunt for exact cases the way \texttt{\_\allowbreak{}exact\_\allowbreak{}cbrt} does.



\textbf{Step 3 — tests}: binary ops are not in the unary exhaustive gate; add a divide-style exhaustive sweep (256×256 operand pairs against a ladder-only reference, a handful of modes) to the suite, and an envelope-sanity loop entry if you introduced a new libm dependence.



\subsection{Adding a ternary operation}



\label{10687193682775408327}{}


Worked example: \textbf{\texttt{FMS}} — fused multiply-subtract, \texttt{x·y − z} — chosen because the cheapest correct implementation is \emph{delegation}, and delegation is a first-class technique here.



\textbf{Step 1 — registry}: append \texttt{:FMS} to \texttt{\_\allowbreak{}TERNARY\_\allowbreak{}OPS} (group \texttt{:A} by the loop). The generated array method routes through the same ternary bitwidth policy every other ternary op uses (\texttt{tables.jl}{\textquotesingle}s \texttt{\_\allowbreak{}ternary\_\allowbreak{}table\_\allowbreak{}for}): small operand formats table (eagerly or adaptively), \texttt{K = 8} runs the — optionally threaded — Shape-B scalar loop, and stochastic ρ always draws per element. No op-specific work needed; the policy keys on formats and ρ, not on which op it is.



\textbf{Step 2 — ω-semantics} by delegation to \texttt{FMA}{\textquotesingle}s already-verified analysis (width thresholds, \texttt{\_\allowbreak{}twosum} exactness, wide-spread fallback):




\nopagebreak[4]
\begin{minted}{julia}
# x·y − z ≡ FMA(x, y, −z). Negation of z follows the Subtract convention:
# preserve NaN, and map −0 to +0 so the single-zero datum model is respected.
ωeval(::Val{:FMS}, x::Float64, y::Float64, z::Float64) =
    ωeval(Val(:FMA), x, y, isnan(z) ? z : (iszero(z) ? 0.0 : -z))
\end{minted}



Delegation inherits FMA{\textquotesingle}s result kinds, its exactness thresholds, \emph{and} its test pedigree — but verify the sign algebra of your reduction on the special rows (here: \texttt{FMS(∞, 1, ∞)} must be NaN, which the delegation produces because \texttt{FMA(∞, 1, −∞)} is \texttt{∞ − ∞}). Add the delegated op to the ternary exhaustive sweep in the suite.



\subsection{Adding a quaternary operation}



\label{1028380857650662652}{}


Arity 4 is \textbf{not yet plumbed} — the generation loops, kernels, and block layer handle arities 1–3. Adding a quaternary op is therefore two tasks: extend the plumbing (once), then add the op. Worked example: \textbf{\texttt{FMMA}} — \texttt{x·y + z·w}, a two-term dot product.



\textbf{Plumbing (one-time), four sites:}



\begin{itemize}
\item[1. ] \texttt{src/ops\_\allowbreak{}scalar.jl} — register manually after the arity loops (\texttt{register\_\allowbreak{}op!(\allowbreak{}:FMMA, 4, :A)}), and add an \texttt{op.arity == 4} branch to the spec-register generation loop, mirroring the ternary branch with a fourth operand:


\nopagebreak[4]
\begin{minted}{julia}
# …continuing the existing `if op.arity == …` cascade:
if op.arity == 4
    @eval begin
        @inline function $name(fr::Type{<:Binary}, ρ::ProjSpec,
                               x::Binary, y::Binary, z::Binary, w::Binary;
                               rng::MaybeRNG=nothing, R::Union{Nothing,Int}=nothing)
            apply_op($V(), fr, ρ, _drawR(ρ, rng, R),
                     decode(x), decode(y), decode(z), decode(w))
        end
        @inline $name(x::T, y::T, z::T, w::T; kw...) where {T<:Binary} =
            $name(T, default_projspec(T), x, y, z, w; kw...)
    end
end
\end{minted}

(\texttt{apply\_\allowbreak{}op} and \texttt{ωeval} are already variadic — no change needed there.)


\item[2. ] \texttt{src/kernels.jl} — a four-array \texttt{vmap!} method. The plain Shape-B loop (copy the \emph{body} of the ternary method{\textquotesingle}s compute branch, add operand \texttt{D}) is enough to get correct results; the ternary method{\textquotesingle}s table-tiering and threading are a bitwidth-driven optimization layered on top (\texttt{tables.jl}{\textquotesingle}s \texttt{\_\allowbreak{}ternary\_\allowbreak{}table\_\allowbreak{}for} and its LRU cache), not required plumbing — extend to arity 4 only if a table win at that arity is worth the 2{\textasciicircum}(4K) growth. Also add an arity-4 branch in the registry-generated array-surface loop and the matching rng-threading overload (mirroring the three-array one) so stochastic ρ dispatches correctly.


\item[3. ] \texttt{src/blocks.jl} — either an arity-4 branch in the Block/Scaled generation loop (four operand blocks, four \texttt{blockdecode}s) or an explicit skip (\texttt{op.arity > 3 \&\& continue}) if a block form is not wanted; be deliberate.


\item[4. ] \texttt{src/approx.jl} — \texttt{conformance\_\allowbreak{}report} prints arities with \texttt{for a in 1:3}; widen to \texttt{1:4}.

\end{itemize}


\textbf{The op itself} follows the \texttt{FAA}/\allowbreak\texttt{FMA} width-analysis playbook exactly:




\nopagebreak[4]
\begin{minted}{julia}
const _DE_FMMA = 92     # two exact ≤17-bit products: 18 + ΔE ≤ 113, margin 3

function ωeval(::Val{:FMMA}, x::Float64, y::Float64, z::Float64, w::Float64)
    (isnan(x) | isnan(y) | isnan(z) | isnan(w)) && return NaN
    (((iszero(x) && isinf(y)) || (isinf(x) && iszero(y))) ||
     ((iszero(z) && isinf(w)) || (isinf(z) && iszero(w)))) && return NaN
    p = x * y; q = z * w                          # exact: ≤8-bit significands
    if isinf(p) || isinf(q)
        (isinf(p) && isinf(q) && p != q) && return NaN
        return isinf(p) ? p : q
    end
    s, e = _twosum(p, q)
    e == 0.0 && return iszero(s) ? 0.0 : s        # exact in Float64
    (_f128() && !iszero(p) && !iszero(q) && _expdiff(p, q) <= _DE_FMMA) &&
        return Float128(p) + Float128(q)          # exact by width
    BigExactF(() -> setprecision(() ->
        BigFloat(x) * BigFloat(y) + BigFloat(z) * BigFloat(w),
        BigFloat, bigprec_prod(x, y, z, w)))     # DERIVED, never a constant
end
\end{minted}



\begin{tcolorbox}[toptitle=-1mm,bottomtitle=1mm,colback=admonition-warning!50!white,colframe=admonition-warning,title=\textbf{Derive the precision from the operands — never a constant}]
This escalation used to read \texttt{\_\allowbreak{}BIGP}, a 2200-bit constant. It was ample for every format that existed at K ≤ 8 and \textbf{insufficient for 50 of the 504}, and it did not announce that: it returned a plausible wrong number, which gate G2 was written red to record (\texttt{1//1} where the exact answer is \texttt{2//1}, the operand returned unchanged because a residual 30 000 binades below the rounding position fell off the accumulator).

\texttt{bigprec(F…)} derives the precision from the formats; \texttt{bigprec(x…)} and \texttt{bigprec\_\allowbreak{}prod(x…)} derive it from the operands in flight. Gate G2 asserts the derivation is sufficient across the whole grid, so a new operation that uses it inherits that proof. A new constant would need its own.

\end{tcolorbox}


\textbf{Testing duty is heavier at arity 4}: the full cross-product is 2³² tuples, so the suite entry must be \emph{sampled} (seeded, ladder-referenced, ≥ 10⁵ tuples across several modes) plus hand-picked exhaustion of the special-row algebra and the width-threshold boundary (\texttt{ΔE ∈ \{91, 92, 93\}} constructions). Document the sampling in the test, since it breaks the suite{\textquotesingle}s {\textquotedbl}enumerate, never sample{\textquotedbl} norm — that exception must be visible, not silent.



\section{Block operations}



\label{6513894782918662055}{}


The block layer (\texttt{src/blocks.jl}) implements the draft{\textquotesingle}s elementwise schema \textbf{once} — \texttt{blockdecode → ω-op lanewise → blockproject} — and generates every registry op{\textquotesingle}s \texttt{BlockOp}/\allowbreak\texttt{ScaledOp} from it. The reductions are hand-written on a second shared pattern. Which subsection you need depends on whether your operation fits one of those two patterns.



\subsection{Adding a scaled block operation}



\label{10512613773019865025}{}


If the scalar operation is (or can be) a \textbf{registry op, the scaled form is free}: registering \texttt{:Cbrt} above already produced \texttt{ScaledCbrt(fr, ρ, s, x)} and \texttt{BlockCbrt(fr, ρ, b, sr)} with the draft{\textquotesingle}s §5.4 semantics — decode \texttt{scale × element} exactly per lane, apply the ω-op, \texttt{BlockProject} against the result scale through the division cascade. \textbf{Prefer this route}; it is the package{\textquotesingle}s non-divergence mechanism.



Write a scaled op by hand only when it is \emph{not} an elementwise image of a registry op. The primitive to compose with is \texttt{\_\allowbreak{}bp\_\allowbreak{}element(\allowbreak{}fr, ρ, R, res, Sdat)} — it owns the draft{\textquotesingle}s scale-special rows (NaN/0/±Inf scale) and the entire division-by-scale rigor cascade. Example, a scaled affine \texttt{s₁·x + s₂·y} projected against a \emph{unit} result scale:




\nopagebreak[4]
\begin{minted}{julia}
function ScaledAffine(fr::Type{<:Binary}, ρ::ProjSpec,
                      s1::Binary, x::Binary, s2::Binary, y::Binary;
                      rng::MaybeRNG=nothing)
    # Each (scale, element) pair is a TWO-FACTOR MONOMIAL and gets its own head
    # from the two formats — the same rule `blockdecode` applies to a block.
    h1 = rung(Val(:Multiply), typeof(s1), typeof(x))
    h2 = rung(Val(:Multiply), typeof(s2), typeof(y))
    a = ωeval(h1, Val(:Multiply), lift(h1, decode(s1)), lift(h1, decode(x)))::carriertype(h1)
    b = ωeval(h2, Val(:Multiply), lift(h2, decode(s2)), lift(h2, decode(y)))::carriertype(h2)
    # The two products may land on DIFFERENT carriers, so the add runs at their
    # join with each operand lifted into it.
    h = _joinheads(a, b)
    res = ωeval(h, Val(:Add), lift(h, a), lift(h, b))
    _bp_element(fr, ρ, _drawR(ρ, rng, nothing), res, 1.0)
end
\end{minted}



\textbf{The assertion is \texttt{::carriertype(h)}, and writing \texttt{::Float64} there is a real defect rather than a stylistic choice.} It is true for the 432 formats at rung 1 and an outright \texttt{TypeError} for the other 72, whose lanes decode to \texttt{Float128} or to the exact dyadic carrier. The generated \texttt{Scaled*} surface in \texttt{blocks.jl} carried exactly that mistake and it is recorded as §11 M44 in the execution log — the guidance here used to teach it.



What the assertion \emph{is} load-bearing for survives the correction: \texttt{Multiply} is closed on the carrier at every rung (gate G4 pins this), so the row cannot escalate to a lazy \texttt{BigExactF} here, and saying so is what lets \texttt{\_\allowbreak{}joinheads} below see carrier values only. JET finds the missing statement before any test can.



The \texttt{Add} result may be any kind in the protocol — \texttt{\_\allowbreak{}bp\_\allowbreak{}element} finishes all of them. Dividing by a scale of \texttt{1.0} is the identity fast path.



\subsection{Adding an elementwise block operation}



\label{15527859041760837550}{}


Same schema, whole-block granularity. The three obligations: decode operand blocks with \texttt{blockdecode} (exact), evaluate lanewise with \texttt{ωeval} (never with ad-hoc Float64 math), and finish with \texttt{blockproject} against the supplied result scale (never with per-lane \texttt{project} — the scale-division rows belong to \texttt{blockproject}). Example — \texttt{BlockRelu}, an elementwise \texttt{max(x, 0)} that reuses the verified \texttt{MaximumNumber} semantics:




\nopagebreak[4]
\begin{minted}{julia}
"""BlockRelu(fr, ρ, b, sr): lanewise max(xᵢ, 0) — NaN lanes propagate NaN per
MaximumNumber's number-preferring rows only when both operands are NaN, i.e.
never here; state your lane semantics this explicitly for any new op."""
function BlockRelu(fr::Type{<:Binary}, ρ::ProjSpec, b::Block{B}, sr::Binary;
                   rng::MaybeRNG=nothing) where {B}
    X = blockdecode(b)
    Z = ntuple(i -> ωeval(Val(:MaximumNumber), X[i], 0.0), Val(B))
    blockproject(fr, ρ, sr, Z; rng)
end
\end{minted}



Add the name to the block-surface export list and to \texttt{conformance()}{\textquotesingle}s \texttt{blocknames} extension vector so the declaration stays truthful. Test against a from-scratch reference composition (the suite{\textquotesingle}s existing block tests show the pattern) — element format × scale format × B, exhaustively where feasible.



\subsection{Adding a reductive block operation}



\label{9686810033048460081}{}


Reductions do \textbf{not} fit the elementwise schema; they follow the second pattern, visible in \texttt{BlockReduceAdd}/\allowbreak\texttt{BlockDotProduct}: (1) resolve the ∞/NaN \textbf{fold algebra} on Float64 classifications first; (2) an integer \textbf{span filter} decides whether the whole reduction is exactly representable in \texttt{Float128} — if so, plain \texttt{Float128} accumulation \emph{is} the exact answer; (3) otherwise a \texttt{BigExactF} exact accumulator at provably sufficient precision; (4) exactly \textbf{one} projection, at the very end, via \texttt{\_\allowbreak{}finish}. Never round intermediates.



Worked example — \textbf{\texttt{BlockSumOfSquares}}, \texttt{Σ xᵢ²}:




\nopagebreak[4]
\begin{minted}{julia}
"""BlockSumOfSquares(fr, ρ, b): project(Σ decode-laneᵢ²) — one rounding total.
Fold algebra: any NaN lane ⇒ NaN; else any ±Inf lane ⇒ +Inf (squares cannot
cancel — no opposite-infinity NaN case exists, unlike ReduceAdd); else exact."""
function BlockSumOfSquares(fr::Type{<:Binary}, ρ::ProjSpec, b::Block{B};
                           rng::MaybeRNG=nothing, R::Union{Nothing,Int}=nothing) where {B}
    X = blockdecode(b)
    res = if any(isnan, X)
        NaN
    elseif any(isinf, X)
        Inf
    elseif all(iszero, X)
        0.0
    else
        # span filter: lanes carry ≤17-bit significands ⇒ squares ≤34 bits and
        # a square's exponent is 2·(lane exponent); the sum is exact in Float128
        # when 34 + 2·span + ⌈log₂B⌉ + 1 ≤ 113  ⇒  2·span + ⌈log₂B⌉ ≤ 78.
        if _f128() && 2 * _expspan(X) + _log2ceil(B) <= 78
            acc = Float128(0)
            for v in X
                q = Float128(v)
                acc += q * q                      # exact squares, exact sum
            end
            acc
        else
            BigExactF(() -> setprecision(BigFloat, _REDPREC) do
                acc = BigFloat(0)
                for v in X
                    acc += BigFloat(v)^2          # exact at _REDPREC
                end
                acc
            end)
        end
    end
    _finish(fr, ρ, _drawR(ρ, rng, R), res)
end
\end{minted}



The two lines that demand your care in any new reduction are the \textbf{fold algebra} (derive it from the draft{\textquotesingle}s reduce definition on the extended reals — here the absence of an opposite-infinities NaN case is a \emph{theorem about squares}, stated in the docstring) and the \textbf{span-filter inequality} (derive it from operand widths, write the derivation in the comment, and add boundary-condition tests at the threshold, exactly as the suite does for the existing \texttt{\_\allowbreak{}DE\_\allowbreak{}*} constants). Everything else — the exact accumulator, the single \texttt{\_\allowbreak{}finish}, stochastic \texttt{R} plumbing — is pattern.



\section{After any addition: the closing checklist}



\label{16308398306918412518}{}


\begin{itemize}
\item Full test suite green, including your new sweep entries.


\item Envelope measurements recorded for any new \texttt{yd}/\allowbreak\texttt{fq} estimator.


\item Benchmark report regenerated; \texttt{\_\allowbreak{}SAFE\_\allowbreak{}DOMAINS} entry added if the op is argument-restricted; sanity-check the op{\textquotesingle}s row in the safe-args table.


\item Docstring on the public name; \texttt{[interp]} markers where you interpreted.


\item \texttt{conformance\_\allowbreak{}report()} shows the op (automatic for registry ops; manual block additions extend \texttt{blocknames}).

\end{itemize}


\section{see also}



\label{9115667341246479961-dup27}{}


\href{under\_\allowbreak{}oracle.md}{The Oracle \& Rigor Classes}, \href{under\_\allowbreak{}blocks.md}{Blocks: Exactness without Superaccumulators}, \href{howto\_\allowbreak{}verify\_\allowbreak{}custom.md}{Verify Custom Code}.



\chapter{Verify Custom Code}



\label{11797808841028591813}{}


Trust, then verify — these patterns turn {\textquotedbl}is my code exact?{\textquotedbl} into an enumeration. Formats are small enough that {\textquotedbl}test on a few points{\textquotedbl} is never necessary; the recipes below enumerate the input space (or, where the space is too large to enumerate, sweep it exhaustively along the dimension that matters) and compare against an independently computed reference. Several examples use unexported internals (\texttt{SmallFloats.round\_\allowbreak{}to\_\allowbreak{}precision}, \texttt{SmallFloats.project}, …); those are stable enough to learn from but are not covered by semantic-versioning guarantees.



\section{Verifying a quantizer exhaustively against an independent reference}



\label{15080892262823953367}{}


Here every in-range datum of \texttt{Binary8p3se} is projected into \texttt{Binary8p4se} and checked against a brute-force nearest search under 256-bit arithmetic (distance agreement; the draft{\textquotesingle}s tie rule is the projection engine{\textquotesingle}s job and is pinned elsewhere in the suite):




\nopagebreak[4]
\begin{minipage}{\linewidth}
\begin{minted}{julia}
using SmallFloats

function ref_nearest_distance(::Type{T}, x) where {T}
    fins = [v for v in (rawvalue(T, UInt8(c)) for c in 0:255) if isfinite(decode(v))]
    minimum(abs(setprecision(() -> BigFloat(decode(v)) - BigFloat(x), BigFloat, 256))
            for v in fins)
end

function verify_quantizer()
    ok = true
    for c in 0x00:0xff
        x = decode(rawvalue(Binary8p3se, c))
        (isfinite(x) && abs(x) <= decode(MaxFiniteOf(Binary8p4se))) || continue
        got = Convert(Binary8p4se, RNE_SN, x)
        ok &= abs(decode(got) - x) == Float64(ref_nearest_distance(Binary8p4se, x))
    end
    ok
end
verify_quantizer()
\end{minted}




\begin{minted}{text}
true
\end{minted}
\end{minipage}




This pattern — enumerate the inputs, compare against an independently written reference — is how the entire package is tested, and it is available to \emph{your} quantization code at trivial cost.



\section{Proving your fused kernel exact: BlockDotProduct vs 512-bit truth}



\label{13954472669056568806}{}


The block layer promises \emph{one} projection with exact lane arithmetic. Trust, then verify — against big-float truth, over random blocks with mixed scales and an honest special-value mix:




\nopagebreak[4]
\begin{minipage}{\linewidth}
\begin{minted}{julia}
using SmallFloats, Random

rng = Xoshiro(5)
mk() = Block(Binary8p1uf(2.0^rand(rng, -3:3)),
             ntuple(_ -> rawvalue(Binary8p4se, UInt8(rand(rng, 0:255))), 32))

function verify_dots(trials)
    for _ in 1:trials
        bx, by = mk(), mk()
        lx = [decode(bx.s) * decode(v) for v in bx.x]
        ly = [decode(by.s) * decode(v) for v in by.x]
        (any(!isfinite, lx) || any(!isfinite, ly)) && continue
        truth = setprecision(() -> sum(BigFloat(lx[i]) * BigFloat(ly[i]) for i in 1:32),
                             BigFloat, 512)
        got = BlockDotProduct(Binary8p4se, RNE_SN, bx, by)
        ref = setprecision(() -> SmallFloats.project(Binary8p4se, RNE_SN, truth), BigFloat, 512)
        codepoint(got) == codepoint(ref) || return false
    end
    true
end
verify_dots(200)
\end{minted}




\begin{minted}{text}
true
\end{minted}
\end{minipage}




The same shape verifies any custom fused kernel you write: compose the truth in \texttt{BigFloat}, project once, compare code points.



\section{Stochastic rounding, audited: the full-R sweep}



\label{6013932738708787830}{}


Every stochastic projection is a deterministic function of the draw \texttt{R}, so its \emph{distribution} is checkable exactly. For \texttt{x = 2 + 3/64} in \texttt{Binary8p4se} the fraction is ν = 3/16 of an ulp, so \texttt{StochasticA\{4\}} must round up for exactly 3 of the 16 draws:




\nopagebreak[4]
\begin{minipage}{\linewidth}
\begin{minted}{julia}
σ4 = RSA_SN(4)                  # ≡ ProjSpec(StochasticA{4}(), SatNone())
x = 2.0 + 3/64
count(decode(SmallFloats.project(Binary8p4se, σ4, x; R)) == 2.25 for R in 0:15)
\end{minted}




\begin{minted}{text}
3
\end{minted}
\end{minipage}




Sweeping \texttt{R} like this turns {\textquotedbl}is my stochastic pipeline unbiased?{\textquotedbl} from a statistical question into an exhaustive one — the pattern the shipped test suite uses.



\section{Ranking safety: an exhaustive monotonicity audit of score conversion}



\label{8732453457535267162}{}


Decision and ranking systems consume scores through \emph{order}. If scores are produced in one format and compared in another, the conversion must be monotone — and formats are small enough to prove it by enumeration rather than trust it. Walk every \texttt{Binary8p3se} datum in total order and check that its \texttt{Binary8p4se} image never goes backward on the integer order keys:




\nopagebreak[4]
\begin{minipage}{\linewidth}
\begin{minted}{julia}
using SmallFloats
using SmallFloats: order_key

function monotone_conversion(::Type{From}, ::Type{To}) where {From,To}
    prev = nothing
    for v in sort(From.(0x00:UInt8(2^bitwidth(From) - 1)))
        isnan(decode(v)) && continue
        g = Convert(To, RNE_SN, v)
        prev !== nothing && order_key(g) < order_key(prev) && return false
        prev = g
    end
    true
end
monotone_conversion(Binary8p3se, Binary8p4se)
\end{minted}




\begin{minted}{text}
true
\end{minted}
\end{minipage}




Run this for every \texttt{(From, To, ρ)} triple your system actually uses; it is a few hundred integer comparisons per triple. A ranking pipeline whose conversions all pass this audit cannot invert a preference by changing formats — a guarantee no amount of spot-testing provides.



\section{The κ-margin decision-rule audit}



\label{3192217673468291530}{}


Search under a compute budget often wants a cheap, approximate evaluation function. The κ registry turns {\textquotedbl}how approximate?{\textquotedbl} into a \emph{measured} code-point bound — and code-point bounds compose into a decision rule: \textbf{if two defined evaluations differ by more than 2κ code points along the total order, the approximate evaluator cannot invert their comparison.} Verify the rule exhaustively for a κ = 2 evaluator:




\nopagebreak[4]
\begin{minipage}{\linewidth}
\begin{minted}{julia}
using SmallFloats: order_key

fast(x) = (r = Exp(Binary8p4se, RNE_SN, x);
           isfinite(decode(r)) ? NextGreaterThan(NextGreaterThan(r)) : r)
κ, exhaustive = measure_kappa(fast, :Exp, Binary8p4se, (Binary8p4se,), RNE_SN)

function margin_audit(fast, κ)
    codes = [rawvalue(Binary8p4se, UInt8(c)) for c in 0:255]
    safe = violations = 0
    for a in codes, b in codes
        da, db = Exp(Binary8p4se, RNE_SN, a), Exp(Binary8p4se, RNE_SN, b)
        (isnan(decode(da)) || isnan(decode(db))) && continue
        codedistance(da, db) > 2κ || continue
        safe += 1
        (order_key(fast(a)) < order_key(fast(b))) ==
            (order_key(da) < order_key(db)) || (violations += 1)
    end
    (safe, violations)
end
(κ, exhaustive, margin_audit(fast, κ)...)
\end{minted}




\begin{minted}{text}
(2.0, true, 49522, 0)
\end{minted}
\end{minipage}




49,522 operand pairs clear the 2κ margin, and the approximate evaluator agrees with the defined ordering on every one of them — zero violations, exhaustively. Inside the margin, comparisons are genuinely undecidable at this κ; a search can treat sub-margin comparisons as ties to expand, or escalate those few nodes to the exact evaluator. Either way the pruning is \emph{provably} sound, with κ measured at registration rather than promised.



\section{see also}



\label{9115667341246479961-dup28}{}


\href{howto\_\allowbreak{}add\_\allowbreak{}operation.md}{Add an Operation}, \href{howto\_\allowbreak{}benchmark.md}{Benchmark Correctly}, \href{under\_\allowbreak{}recipes.md}{Internals Recipes}.



\chapter{Benchmark Correctly}



\label{1754891814411159897}{}


Recorded after two measurement post-mortems: a benchmark closure over any non-\texttt{const} global measures Julia{\textquotesingle}s dispatch machinery, not the code under test — and it distorts \emph{ratios}, not just absolutes, because dispatch cost varies with call shape (a dynamic keyword call costs {\textasciitilde}1 µs; a dynamic positional call far less; six unresolved interior sites cost six times one).



\section{The rules}



\label{5346814332251914868}{}


The rules the shipped \texttt{benchmarking/\allowbreak{}benchmarking.jl} enforces structurally:



\begin{itemize}
\item format types enter as type parameters, never as globals;


\item operands come from untimed setup;


\item functions retrieved reflectively pass through argument barriers to specialize;


\item a preflight aborts the run if warm scalar paths allocate — including the wide-spread \texttt{FMA}/\allowbreak\texttt{FAA} sticky-head path.

\end{itemize}


The table-build section reports both the cold build (cache evicted per sample) and the steady-state warm cache hit, since callers amortize the former through the latter. Vary one binding per variant; verify specialization before believing a number.



\section{Benchmarking without measuring the dispatcher}



\label{585007679975857547}{}


The package{\textquotesingle}s benchmark doctrine, in one snippet (needs the \texttt{benchmarking/} environment for Chairmarks). Format types enter as \textbf{type parameters}; operands come from Chairmarks{\textquotesingle} \emph{untimed} \texttt{setup}; and if you retrieve functions reflectively (\texttt{getfield(\allowbreak{}SmallFloats, op)}), pass them through an argument barrier so they specialize:




\nopagebreak[4]
\begin{minipage}{\linewidth}
\begin{minted}{julia}
using Chairmarks, SmallFloats, Random
using Statistics: median

function bench_add(::Type{T}) where {T<:Binary}          # T: type parameter, not a global
    pool = [rawvalue(T, rand(UInt8)) for _ in 1:4096]
    @be (rand(pool), rand(pool)) (t -> Add(T, RNE_SN, t[1], t[2]))(_)
end

b = bench_add(Binary8p4se)
(round(median(b).time * 1e9; digits=1), median(b).allocs)
\end{minted}




\begin{minted}{text}
(16.3, 0.0)          # ns per full scalar Add, zero allocations
\end{minted}
\end{minipage}




\begin{tcolorbox}[toptitle=-1mm,bottomtitle=1mm,colback=admonition-warning!50!white,colframe=admonition-warning,title=\textbf{The 60× benchmarking slowdown}]
The same call with \texttt{T} read from a non-\texttt{const} global measures {\textasciitilde}1 µs — Julia{\textquotesingle}s dynamic keyword dispatch, not this package. Two project post-mortems trace to exactly this mistake; the shipped \texttt{benchmarking/\allowbreak{}benchmarking.jl} asserts specialization (zero warm-path allocation) before it believes any number, and so should yours.

\end{tcolorbox}


\section{see also}



\label{9115667341246479961-dup29}{}


\href{under\_\allowbreak{}verification.md}{Verification \& Benchmark Doctrine}, \href{howto\_\allowbreak{}verify\_\allowbreak{}custom.md}{Verify Custom Code}, \href{under\_\allowbreak{}recipes.md}{Internals Recipes}.



\chapter{Internals Recipes}



\label{3785606410798990400}{}


Outputs below are captured from real sessions. These examples use unexported internals (\texttt{SmallFloats.round\_\allowbreak{}to\_\allowbreak{}precision}, \texttt{SmallFloats.project}, \texttt{SmallFloats.get\_\allowbreak{}table}, \texttt{SmallFloats.order\_\allowbreak{}key}, …); those are stable enough to learn from but are not covered by semantic-versioning guarantees.



\section{Watching RoundToPrecision work}



\label{2755073532632396567}{}


\texttt{round\_\allowbreak{}to\_\allowbreak{}precision} returns the draft{\textquotesingle}s exact \texttt{(sign, S, Q)} decomposition before saturation ever sees it:




\nopagebreak[4]
\begin{minipage}{\linewidth}
\begin{minted}{julia}
using SmallFloats
using SmallFloats: round_to_precision

r = round_to_precision(4, 8, NearestTiesToEven(), 2.30078125, 0, 0)   # P=4, bias=8
(r.sign, r.S, r.Q, r.S * 2.0^r.Q)
\end{minted}




\begin{minted}{text}
(1, 9, -2, 2.25)          # S·2^Q = 9/4: the nearest P=4 value
\end{minted}
\end{minipage}




\section{Symbolic sticky: projecting {\textquotedbl}just below{\textquotedbl} a number}



\label{1121493031502271068-dup1}{}


The engine accepts \texttt{sticky ∈ \{-1, 0, +1\}} meaning {\textquotedbl}the true value is the carrier plus an infinitesimal of this sign{\textquotedbl}. This is how enclosure endpoints and asymptotes (\texttt{tanh → 1⁻}) are projected exactly:




\nopagebreak[4]
\begin{minipage}{\linewidth}
\begin{minted}{julia}
using SmallFloats: project
project(Binary8p4se, ProjSpec(TowardNegative(), SatNone()), 1.0; sticky=-1)
\end{minted}




\begin{minted}{text}
Binary8p4se(0.9375 ≡ 0x3f)     # == NextLessThan(1.0): the engine crossed the binade
\end{minted}
\end{minipage}




\section{Tables are the scalar path, memoized}



\label{17860062526638989794}{}



\begin{minipage}{\linewidth}
\begin{minted}{julia}
using SmallFloats: get_table
empty_tables!()
tbl = get_table(:Exp, Binary8p4se, Binary8p4se, RNE_SN)
(tbl[Int(0x45) + 1],
 codepoint(Exp(Binary8p4se, RNE_SN, rawvalue(Binary8p4se, 0x45))),
 table_bytes())
\end{minted}




\begin{minted}{text}
(0x52, 0x52, 256)              # table entry ≡ scalar result; one 256-byte table cached
\end{minted}
\end{minipage}




\section{Order keys}



\label{11135872524320346846}{}


Ordering is integer arithmetic: a sign-magnitude fold, monotone with the total order, NaN at the top. This is what comparisons and the O(n) counting sort run on:




\nopagebreak[4]
\begin{minipage}{\linewidth}
\begin{minted}{julia}
using SmallFloats: order_key
[(v, order_key(v)) for v in (Binary8p4se(-1.0), Binary8p4se(0.0),
                             Binary8p4se(1.0), Binary8p4se(NaN))]
\end{minted}




\begin{minted}{text}
-1.0 → 64      0.0 → 129      1.0 → 193      NaN → 65535
\end{minted}
\end{minipage}




\section{see also}



\label{9115667341246479961-dup30}{}


\href{under\_\allowbreak{}engine.md}{The Encoding \& Projection Engine}, \href{under\_\allowbreak{}tables.md}{Tables, Kernels \& Sorting}, \href{howto\_\allowbreak{}verify\_\allowbreak{}custom.md}{Verify Custom Code}.



\part{Support}


\chapter{Troubleshooting}



\label{5264174653950164674}{}


Symptom-first reference. Find what you are seeing, not what you did wrong — each entry names the symptom, the cause, the fix, and where to read the full explanation.



\section{Values \& Construction}



\label{13839706789210099024}{}


\textbf{Symptom: \texttt{T(0x02)} gives a different value than \texttt{T(2)}, and it looks like a bug.} Cause: an \texttt{Unsigned} argument is a \emph{code point}, at every bitwidth — \texttt{T(0x02)} constructs code point \texttt{0x02}, while \texttt{T(2)} projects the numeric value two. Signedness of the integer type is what distinguishes them, so this holds for \texttt{UInt8}, \texttt{UInt16}, and any other \texttt{Unsigned}. Fix: use \texttt{T(Int(c))} when you mean the number, \texttt{T(c::Unsigned)} when you mean the code point, or call \texttt{Convert} when intent should be unmistakable. Full explanation: Cheat Sheet, Values, Code Points \& Conversion.



\textbf{Symptom: \texttt{Binary8p4se(0x02)} and \texttt{convert(\allowbreak{}Binary8p4se, 0x02)} disagree.} Cause: \texttt{convert} is the deliberate exception to the Unsigned-is-a-code-point rule and is value-preserving — \texttt{Binary8p4se(0x02)} is code point 2, while \texttt{convert(\allowbreak{}Binary8p4se, 0x02)} is the value 2.0. Promotion, \texttt{similar}, and \texttt{fill} depend on \texttt{convert} behaving this way. Fix: know which one you are calling; prefer \texttt{T(x)} or \texttt{Convert(T, ρ, x)} in your own code and reserve bare \texttt{convert} for places Base itself calls it. Full explanation: Cheat Sheet, Values, Code Points \& Conversion.



\textbf{Symptom: constructing from a \texttt{Rational} throws.} Cause: \texttt{Rational} is not one of the supported construction carriers. Fix: convert explicitly to an exact supported carrier (a \texttt{Float64}, \texttt{Dyadic}, or similar) or knowingly to \texttt{Float64} first, then construct. Full explanation: Cheat Sheet, Values and code points.



\textbf{Symptom: \texttt{T(-0.0) === zero(T)} surprises code that expects signed zero.} Cause: every P3109 format has exactly one zero; there is no negative zero to distinguish. Fix: stop testing for negative zero; if you need {\textquotedbl}approached from below,{\textquotedbl} use the projection engine{\textquotesingle}s \texttt{sticky} argument instead of relying on a stored sign. Full explanation: Mental Model, The Encoding \& Projection Engine.



\textbf{Symptom: a value built with \texttt{rawvalue} decodes to garbage or a non-existent datum.} Cause: \texttt{rawvalue(T, c)} is unchecked code-point construction — it skips the range validation that \texttt{T(c::Unsigned)} performs, so an out-of-range integer produces a bit pattern with no corresponding intent-checked datum. Fix: prefer validated \texttt{T(c::Unsigned)} outside performance kernels; reserve \texttt{rawvalue} for hot loops where the caller has already proven \texttt{c} is in range. Full explanation: Cheat Sheet, Values and code points.



\textbf{Symptom: dividing by zero gives \texttt{NaN}, not \texttt{Inf} as IEEE-754 code expects.} Cause: P3109 defines \texttt{x / 0 → NaN} (and \texttt{Recip(±Inf) = 0}), which is not IEEE division-by-zero semantics. Fix: do not port IEEE-754 zero-division assumptions across; check for zero divisors explicitly if your algorithm depends on signed infinities from division. Full explanation: Cheat Sheet, Values and code points; Introducing P3109.



\section{Arithmetic \& Semantics}



\label{6618291483251033039}{}


\textbf{Symptom: \texttt{Multiply}/\allowbreak\texttt{Add} overflowed and you expected a clamp, but got \texttt{Inf} (or vice versa).} Cause: \texttt{SatNone} does not mean {\textquotedbl}no saturation{\textquotedbl} in the clamping sense — it applies the draft{\textquotesingle}s domain-, signedness-, and direction-dependent rows, which can still produce \texttt{Inf}, \texttt{MaxFinite}, or \texttt{NaN} depending on the case. Only \texttt{SatFinite} guarantees clamping to the finite range. Fix: choose \texttt{SatFinite} whenever clamping is actually required; do not assume \texttt{SatNone} behaves like \texttt{clamp}. Full explanation: Cheat Sheet, Projection specifications.



\textbf{Symptom: mixing two \texttt{Binary} formats in one expression throws a promotion error instead of silently widening.} Cause: this is deliberate — the package refuses silent cross-format promotion so that no rounding step happens somewhere nobody wrote. Fix: convert explicitly, \texttt{Convert(T, ρ, x)}, to the format you actually want before combining. Full explanation: Cheat Sheet, Values and code points.



\textbf{Symptom: \texttt{rem}, \texttt{mod}, or a similar Base numeric function refuses to run, or errors in a way that looks like a missing method.} Cause: some generic Base numeric functions are defined in terms of operations or promotion behavior this package deliberately does not supply by default (for example, functions that would require an implicit cross-format promotion or an IEEE-754-only special case). This is a refusal, not an omission bug. Fix: express the computation with the package{\textquotesingle}s explicit operation catalog (\texttt{Add}, \texttt{Subtract}, \texttt{Multiply}, \texttt{Divide}, \texttt{FMA}, …) and \texttt{Convert} calls instead of relying on the generic fallback. Full explanation: Cheat Sheet, Scalar operation catalog.



\section{Arrays \& Performance}



\label{15207003463717874679}{}


\textbf{Symptom: an array of \texttt{Binary} values is slow, or every element looks boxed.} Cause: the element type is the \emph{abstract} format \texttt{Binary\{K,P,Σ,Δ\}}, not a concrete alias — the abstract format is not \texttt{isbits}, so a generic-element array boxes every entry. Fix: use a concrete element type, \texttt{Vector\{Binary8p4se\}} or \texttt{Vector\{format(K,P,Σ,Δ)\}}. Full explanation: Cheat Sheet, Common pitfalls (below); Architecture Overview.



\textbf{Symptom: \texttt{Vector\{Binary\{K,P,Σ,Δ\}\}(undef, n)} silently boxes every element, even though you thought you had a concrete type.} Cause: \texttt{Binary\{K,P,Σ,Δ\}} with type-parameter variables is abstract, and \texttt{Vector\{…\}(undef, n)} cannot be intercepted or redirected to a concrete representation the way constructors can. Fix: use \texttt{Vector\{Binary8p4se\}} (a concrete alias) or \texttt{Vector\{format(K,P,Σ,Δ)\}} (concrete once \texttt{K,P,Σ,Δ} are literal values); \texttt{similar} normalizes correctly and is safe. Full explanation: Cheat Sheet, Common pitfalls.



\textbf{Symptom: benchmarks report {\textasciitilde}1 microsecond per call for an operation you know is sub-nanosecond warm.} Cause: the format type or projection spec was read from a non-\texttt{const} global, so every call pays Julia{\textquotesingle}s dynamic dispatch instead of running the specialized method. This is a benchmarking-methodology bug, not a package performance regression — it has been traced to this exact mistake in real post-mortems. Fix: keep format types and projection specs in \texttt{const} bindings, function arguments, or type parameters; put runtime-selected-format code behind a function barrier; verify zero warm-path allocation before trusting a number. Full explanation: Technical Examples, Benchmarking without measuring the dispatcher; Verification \& Benchmark Doctrine.



\textbf{Symptom: the first array call for a given format/operation is much slower than every call after it.} Cause: deterministic unary/binary operations build a cached result table on first use for that specialization; you measured the build, not the steady state. Fix: warm the specialization before benchmarking, and benchmark warm calls separately from the first call. Use \texttt{table\_\allowbreak{}bytes()} to inspect the cache footprint and \texttt{empty\_\allowbreak{}tables!()} to reset it. Full explanation: Cheat Sheet, Performance checklist (pointer); Tables, Kernels \& Sorting.



\section{Random \& Stochastic}



\label{16404060453799710041}{}


\textbf{Symptom: a stochastic-rounding pipeline gives different results every run, even though you want to compare against a fixed baseline.} Cause: stochastic projections (\texttt{RSA}, \texttt{RSB}, \texttt{RSC}) consume random bits from an RNG (or an explicit draw), and without a fixed source that draw changes every run. Fix: supply a seeded \texttt{rng} (for example \texttt{Xoshiro(1)}), or pass an explicit \texttt{R} in \texttt{0:(2{\textasciicircum}N - 1)} for exact, reproducible draws — ideal in tests. Full explanation: Cheat Sheet, Stochastic projections; Reproducible Stochastic Rounding.



\textbf{Symptom: you passed a seed but two calls in the same pipeline still diverge.} Cause: the uniform draw and the stochastic rounding draw must come from the \emph{same} rng stream to stay reproducible together; if you construct a second \texttt{Xoshiro} mid-pipeline instead of threading the same \texttt{rng} through, the streams desynchronize. Fix: thread one \texttt{rng} object through every call in the pipeline that needs to agree. Full explanation: Cheat Sheet, Stochastic projections.



\section{see also}



\label{9115667341246479961-dup31}{}


\href{cheatsheet.md}{Cheat Sheet}, \href{glossary.md}{Glossary}, \href{explain\_\allowbreak{}performance.md}{Performance Model}, \href{under\_\allowbreak{}verification.md}{Verification \& Benchmark Doctrine}.



\chapter{Glossary}



\label{5283940344096028924}{}


Alphabetical, one or two lines per term, with a plain-text pointer to the page where the term is defined in full.




\begin{table}[h]
\centering
\nopagebreak[3]
\begin{tabulary}{\linewidth}{L L L}
\toprule
Term & Definition & See \\
\toprule
Carrier & The concrete numeric type an exact intermediate or datum is held on during computation — \texttt{Float64} at rung 1, \texttt{Float128} at rung 2, or an exact dyadic carrier at rung 3, chosen by the format{\textquotesingle}s exponent bias, never by the caller. & Architecture Overview; The Encoding \& Projection Engine \\
Code point & The \texttt{2{\textasciicircum}K}-valued integer identity of a value within its format — what a \texttt{Binary} value actually stores, in bijection with the format{\textquotesingle}s datum set. & Mental Model \\
Code unit & The storage word a format{\textquotesingle}s code points are packed into — \texttt{UInt8} at \texttt{K ≤ 8}, \texttt{UInt16} above — returned by \texttt{codeunit\_\allowbreak{}type(T)} and by \texttt{codepoint(x)}. & Cheat Sheet, Values and code points \\
Datum & The exact mathematical value a code point denotes — a member of the closed extended reals (finite dyadic values, ±Inf, NaN). & Mental Model \\
Datum carrier & The specific carrier type that is exact for a given format{\textquotesingle}s datums, named by \texttt{datumcarrier(T)}; not to be confused with the promotion carrier used for mixed-type Base arithmetic. & The Encoding \& Projection Engine \\
Dyadic & An exact rational-like carrier, \texttt{S · 2{\textasciicircum}Q} with an \texttt{Int128} significand and \texttt{Int64} exponent, \texttt{isbits}, used at rung 3 where even \texttt{Float128} cannot hold a format{\textquotesingle}s range or precision; converts to and from Julia{\textquotesingle}s \texttt{Rational} exactly. & The Oracle \& Rigor Classes; Design Rationales (dyadic\_\allowbreak{}rational.md) \\
κ (kappa) & The maximum code-point distance, measured along the total order, between an approximate implementation{\textquotesingle}s result and the defined result — measured by exhaustive enumeration at registration time, never declared unchecked. & The Oracle \& Rigor Classes; Conformance and approximations \\
Niven peel & The termination argument for π-scaled trigonometric operations: a finite set of exactly-representable cases split off ({\textquotedbl}peeled{\textquotedbl}) before interval enclosure begins, with Niven{\textquotesingle}s theorem proving the peel set is complete. & The Oracle \& Rigor Classes \\
Omega-operation (ω) & The auxiliary, mathematically exact operation (\texttt{ωeval}) that a defined operation{\textquotesingle}s result is the projection of — the thing that is computed \emph{before} rounding and saturation ever see it. & The Oracle \& Rigor Classes \\
Order key & An integer, sign-magnitude-folded encoding of a value used for comparisons and sorting — monotone with the total order, NaN mapped to the maximum key, enabling O(n) counting sort. & The Encoding \& Projection Engine \\
Projection & The single write path that produces every code point in the package: exact result → RoundToPrecision → Saturate → Encode. There is no second rounding routine anywhere in the package. & The Encoding \& Projection Engine \\
Projection specification (ProjSpec) & The pair \texttt{(\allowbreak{}rounding mode, saturation mode)} that parameterizes a projection, constructed as \texttt{ProjSpec(\allowbreak{}rounding, saturation)} or via a named constant like \texttt{RNE\_\allowbreak{}SN}. & Cheat Sheet, Projection specifications \\
Promotion carrier & The ordinary Julia float type (\texttt{Float64} at rung 1, \texttt{BigFloat} at rung 2/3) that a \texttt{Binary} value promotes to when combined with an ordinary number — named by \texttt{promotecarrier(T)}; distinct from the datum carrier used internally. & Design Rationales (promoterules.md); Julia Compatibility \\
Result kind & One of the five categories \texttt{ωeval} returns a computed result as (for example exact-\texttt{Float64}, exact-\texttt{Dyadic}, or an enclosure needing further work), which \texttt{apply\_\allowbreak{}op} dispatches on to finish the projection. & The Oracle \& Rigor Classes \\
Rigor class (R/E) & The two grades of justification for a non-\texttt{Float64} computed result: Class R (unconditional, provable without any accuracy assumption) and Class E (envelope-conditional, a faithful-but-not-correctly-rounded estimate validated by a measured slack). & Architecture Overview; The Oracle \& Rigor Classes \\
Rounding mode & One of the six deterministic modes (\texttt{NearestTiesToEven}, \texttt{NearestTiesToAway}, \texttt{TowardPositive}, \texttt{TowardNegative}, \texttt{TowardZero}, \texttt{ToOdd}) or three stochastic variants (\texttt{RSA}, \texttt{RSB}, \texttt{RSC}) that make up half of a projection specification. & Cheat Sheet, Projection specifications \\
Rung & One of three carrier tiers a format is assigned to by its exponent bias: rung 1 (\texttt{Float64}, 432 formats), rung 2 (\texttt{Float128}, formats whose range needs more exponent room), rung 3 (exact dyadic carrier, formats whose range or precision breaks \texttt{Float128}). & Architecture Overview \\
Saturation mode & One of three policies for out-of-range projection results: \texttt{SatFinite} (clamp to the finite range), \texttt{SatPropagate} (preserve representable infinities, clamp other overflow), \texttt{SatNone} (apply the draft{\textquotesingle}s domain-, signedness-, and direction-dependent rows). & Cheat Sheet, Projection specifications \\
Saturation rows & The draft{\textquotesingle}s twenty-row table mapping a rounded value{\textquotesingle}s range classification to its final disposition (as-is, MaxFinite, MinFinite, ±Inf, NaN), consulted by the Saturate stage of projection. & The Encoding \& Projection Engine \\
Span filter & An integer pre-pass over lane exponents in a block reduction that proves a narrower carrier already covers every value the reduction can produce, licensing a faster accumulation path without loss of exactness. & Blocks: Exactness without Superaccumulators \\
Sticky (symbolic) & A \texttt{sticky ∈ \{-1, 0, +1\}} argument to the projection engine meaning {\textquotedbl}the true value is the carrier value plus an infinitesimal of this sign{\textquotedbl} — how enclosure endpoints and asymptotes (\texttt{tanh → 1⁻}) project exactly without inventing a fake carrier value. & The Encoding \& Projection Engine \\
Table gather & The array-kernel strategy of computing a result by indexing a precomputed, cached lookup table (built once through the scalar path) rather than recomputing per element. & Tables, Kernels \& Sorting \\
Total order & The draft{\textquotesingle}s total ordering over a format{\textquotesingle}s value set, with a single NaN placed at one end deterministically; the basis for \texttt{TotalOrder}, \texttt{isless}, comparisons, and counting sort. & The Encoding \& Projection Engine \\
\bottomrule
\end{tabulary}

\end{table}



\section{see also}



\label{9115667341246479961-dup32}{}


\href{cheatsheet.md}{Cheat Sheet}, \href{troubleshooting.md}{Troubleshooting}, \href{under\_\allowbreak{}architecture.md}{Architecture Overview}.



\chapter{Examples from AI}



\label{11561040251548532114}{}


A map of the worked examples. Each is identified in boldface and briefly described. The source for each example is linkded.



\section{General AI}



\label{12995498770613901834}{}


\textbf{Log-odds belief updating.} Accumulating log-likelihood ratios in an 8-bit format: the running belief can drift far from the exact value while the threshold decision it drives stays correct — order survives quantization even when magnitude does not. Source: \href{https://github.com/JeffreySarnoff/SmallFloats.jl/blob/main/docs/oldsrc/user_examples.md\#L119}{User Examples — General AI}.



\textbf{Fuzzy inference in unsigned formats.} Membership degrees on \texttt{[0, 1]} evaluated with an unsigned finite format; the classic fuzzy connectives (Gödel t-norm, s-norm, product t-norm) are bit-exact registry operations, so a fuzzy rule base reproduces to the code point across machines. Source: \href{https://github.com/JeffreySarnoff/SmallFloats.jl/blob/main/docs/oldsrc/user_examples.md\#L156}{User Examples — General AI}.



\textbf{Search-heuristic ordering and a monotonicity audit.} Quantizing 200 heuristic scores to 8 bits collapses many distinct scores onto shared code points, yet argmax and top-k survive; a companion exhaustive walk proves a format conversion never inverts the total order. Sources: \href{https://github.com/JeffreySarnoff/SmallFloats.jl/blob/main/docs/oldsrc/user_examples.md\#L183}{User Examples — General AI}, \href{https://github.com/JeffreySarnoff/SmallFloats.jl/blob/main/docs/oldsrc/technical_examples.md\#L75}{Technical Examples — General AI}.



\textbf{κ-safe decision margins for approximate evaluators.} Turning a measured κ bound into a provable rule — two evaluations more than 2κ code points apart cannot have their comparison inverted by a κ-bounded approximate evaluator — and verifying the rule exhaustively. Source: \href{https://github.com/JeffreySarnoff/SmallFloats.jl/blob/main/docs/oldsrc/technical_examples.md\#L112}{Technical Examples — General AI}.



\section{Machine Learning}



\label{11381515430689351323}{}


\textbf{Quantizing a weight tensor and measuring the damage.} Projecting a Gaussian weight tensor into an 8-bit format and reporting MSE, max error, and overflow fraction as the three-number check for any quantization. Source: \href{https://github.com/JeffreySarnoff/SmallFloats.jl/blob/main/docs/oldsrc/user_examples.md\#L211}{User Examples — Machine Learning}.



\textbf{Picking a format: precision vs range.} A table of RMSE and MaxFinite across four \texttt{K = 8} formats trading significand bits for exponent range on the same data. Source: \href{https://github.com/JeffreySarnoff/SmallFloats.jl/blob/main/docs/oldsrc/user_examples.md\#L231}{User Examples — Machine Learning}.



\textbf{Stochastic unbiasedness.} Nearest rounding is systematically biased on a fixed input; stochastic rounding is unbiased on average — the property that matters for accumulating many small contributions in low precision. Source: \href{https://github.com/JeffreySarnoff/SmallFloats.jl/blob/main/docs/oldsrc/user_examples.md\#L255}{User Examples — Machine Learning}.



\textbf{Exhaustive quantizer verification.} Checking every in-range code point of a quantizer against an independent 256-bit reference computation, rather than spot-checking. Source: \href{https://github.com/JeffreySarnoff/SmallFloats.jl/blob/main/docs/oldsrc/technical_examples.md\#L157}{Technical Examples — Machine Learning}.



\textbf{κ workflow with rejection.} Registering an approximate kernel, measuring its κ by exhaustive enumeration, and watching the registry refuse an understated κ declaration. Source: \href{https://github.com/JeffreySarnoff/SmallFloats.jl/blob/main/docs/oldsrc/technical_examples.md\#L197}{Technical Examples — Machine Learning}.



\textbf{Conformance export.} Producing a machine-readable conformance dictionary suitable for attaching to experiment artifacts. Source: \href{https://github.com/JeffreySarnoff/SmallFloats.jl/blob/main/docs/oldsrc/technical_examples.md\#L221}{Technical Examples — Machine Learning}.



\section{Deep Learning}



\label{15209394077164581735}{}


\textbf{MX block quantization and the staging pitfall.} Per-row block scales tracking dynamic range a bare element format cannot, plus the specific case of staging through the narrow element format before block-quantizing — which silently discards the benefit. Source: \href{https://github.com/JeffreySarnoff/SmallFloats.jl/blob/main/docs/oldsrc/user_examples.md\#L277}{User Examples — Deep Learning}.



\textbf{One-rounding dot products.} \texttt{BlockDotProduct} computing every lane product and the accumulation exactly, then rounding exactly once. Source: \href{https://github.com/JeffreySarnoff/SmallFloats.jl/blob/main/docs/oldsrc/user_examples.md\#L322}{User Examples — Deep Learning}.



\textbf{Swamping and stall demo.} Accumulating many small gradient steps under nearest rounding stalls completely once the accumulator dwarfs the increment; stochastic rounding keeps absorbing the increments in expectation. Source: \href{https://github.com/JeffreySarnoff/SmallFloats.jl/blob/main/docs/oldsrc/user_examples.md\#L340}{User Examples — Deep Learning}.



\textbf{Activation LUT audit.} An 8-bit unary activation is a 256-byte lookup table; auditing saturation fraction and distinct output codes is cheap enough to run for every candidate nonlinearity. Source: \href{https://github.com/JeffreySarnoff/SmallFloats.jl/blob/main/docs/oldsrc/user_examples.md\#L369}{User Examples — Deep Learning}.



\textbf{Hard-tanh κ measurement.} Measuring the worst-case code-point deviation of a piecewise hardware-style activation against the defined activation, exhaustively, and registering it under the measured bound. Source: \href{https://github.com/JeffreySarnoff/SmallFloats.jl/blob/main/docs/oldsrc/technical_examples.md\#L291}{Technical Examples — Deep Learning}.



\textbf{Packing a quantized model.} Storing a million-element quantized model in a \texttt{PackedVector} at the format{\textquotesingle}s true bitwidth instead of a full byte per element. Source: \href{https://github.com/JeffreySarnoff/SmallFloats.jl/blob/main/docs/oldsrc/user_examples.md\#L415}{User Examples — Deep Learning}.



\textbf{BlockDotProduct 512-bit verification.} Proving a fused block dot-product kernel exact against 512-bit big-float truth over random blocks with mixed scales and a special-value mix — the pattern to reuse for any custom fused kernel. Source: \href{https://github.com/JeffreySarnoff/SmallFloats.jl/blob/main/docs/oldsrc/technical_examples.md\#L234}{Technical Examples — Deep Learning}.



\section{see also}



\label{9115667341246479961-dup33}{}


\href{howto\_\allowbreak{}choose\_\allowbreak{}format.md}{How-To Guides}, \href{under\_\allowbreak{}architecture.md}{Architecture Overview}, \href{cheatsheet.md}{Cheat Sheet}.



\chapter{Cheat Sheet}



\label{17671780375889550826}{}


A one-page reference for the SmallFloats.jl operations you are most likely to write. For explanations and semantics, see the User Guide; for implementation details, see the Technical Guide.



Single-topic cards: Format Card · Ops Card · Blocks \& Packed Card (appendix pages below).



\section{Start here}



\label{11995297219376789591}{}



\begin{minted}{julia}
using SmallFloats
using Random: Xoshiro

T = Binary8p4se
ρ = RNE_SN

x = T(1.6)
y = T(0.25)
z = Add(T, ρ, x, y)
\end{minted}



The explicit operation shape is always:




\nopagebreak[4]
\begin{minted}{julia}
Op(result_format, projection_spec, operands...; rng, R)
\end{minted}



For same-format operands under the default projection, ordinary Julia syntax is available:




\nopagebreak[4]
\begin{minted}{julia}
z == x + y
exp(x)
fma(x, y, z)
\end{minted}



\needspace{26\baselineskip}
\section{Format names}



\label{4610047011167023553}{}



\begin{minted}{text}
Binary K p P (s|u) (e|f)
       │   │  │     └─ extended (Inf) or finite domain
       │   │  └────── signed or unsigned
       │   └───────── significand precision, including the implicit bit
       └───────────── total bitwidth, 3 through 16
\end{minted}



Examples:




\begin{table}[h]
\centering
\begin{tabulary}{\linewidth}{L L}
\toprule
Name & Meaning \\
\toprule
\texttt{Binary8p4se} & 8-bit, precision 4, signed, extended \\
\texttt{Binary8p4sf} & 8-bit, precision 4, signed, finite \\
\texttt{Binary6p3ue} & 6-bit, precision 3, unsigned, extended \\
\texttt{Binary5p5uf} & 5-bit, precision 5, unsigned, finite \\
\texttt{Binary16p6se} & 16-bit, precision 6, signed, extended (opt-in export) \\
\texttt{Binary16p1uf} & 16-bit, precision 1, unsigned, finite — the MX scale shape \\
\bottomrule
\end{tabulary}

\end{table}



The alias is the \emph{representation} of the parametric format, related by \texttt{<:} and \textbf{not} by \texttt{===}:




\nopagebreak[4]
\begin{minted}{julia}
Binary8p4se <: Binary{8, 4, true, true}    # true
Binary8p4se === Binary{8, 4, true, true}   # FALSE — `Binary` is abstract
format(16, 6, true, true)                  # the alias, from runtime parameters
\end{minted}



\texttt{using SmallFloats} exports the 120 aliases at K ≤ 8. For the other 384: \texttt{using SmallFloats.Formats}, or \texttt{SmallFloats.Binary16p6se}, or \texttt{format(...)}.



Useful format queries:




\nopagebreak[4]
\begin{minted}{julia}
bitwidth(T)       # K
precision(T)      # P
issigned(T)
isextended(T)
expbias(T)
expbitwidth(T)
trailingsigbits(T)

MaxFiniteOf(T)
MinFiniteOf(T)
MinPositiveOf(T)
MinNormalOf(T)
MaxSubnormalOf(T)
\end{minted}



Every query above also takes a \textbf{value} instead of a type — \texttt{bitwidth(x)} ≡ \texttt{bitwidth(typeof(x))} — and folds to the same constant:




\nopagebreak[4]
\begin{minted}{julia}
x = Binary8p4se(1.6)
bitwidth(x)       # 8
MaxFiniteOf(x)    # Binary8p4se(224.0 ≡ 0x7e)
\end{minted}



\section{Values and code points}



\label{1762024146591479678}{}



\begin{minted}{julia}
x = T(1.6)                    # numeric value: project into T
x = Convert(T, ρ, 1.6)       # explicit projection

c = codepoint(x)              # the code point, in the format's storage unit
x = T(c)                      # an Unsigned means validated CODE POINT
x = rawvalue(T, c)            # unchecked code-point construction

d = decode(x)                 # the exact datum, on the format's own carrier
\end{minted}



\begin{tcolorbox}[toptitle=-1mm,bottomtitle=1mm,colback=admonition-warning!50!white,colframe=admonition-warning,title=\textbf{`Unsigned` means code point — at every width}]
\texttt{T(0x02)} constructs code point \texttt{0x02}; \texttt{T(2)} projects the numeric value two. Signedness of the integer type is what distinguishes them, so the rule holds for \texttt{UInt8}, \texttt{UInt16} and any other \texttt{Unsigned}: \texttt{Binary16p6se(0x0002)} is code point 2, not the value 2.0. Use \texttt{Convert} when intent should be unmistakable.

\texttt{codepoint} returns the format{\textquotesingle}s \emph{storage unit} — \texttt{UInt8} at K ≤ 8, \texttt{UInt16} above — and that is \texttt{codeunit\_\allowbreak{}type(T)}. A code-point argument is range-checked against \texttt{2{\textasciicircum}K} whatever its width, so passing a wider \texttt{Unsigned} than necessary is lossless and safe.

\texttt{convert} is the deliberate exception and is value-preserving: \texttt{Binary8p4se(0x02)} is code point 2, while \texttt{convert(\allowbreak{}Binary8p4se, 0x02)} is the value 2.0. Promotion, \texttt{similar} and \texttt{fill} depend on that.

\end{tcolorbox}


Random values (uniform [0,1) floor-projected; normal round-to-nearest, tails clamped to MaxFinite; \texttt{randn} signed formats only):




\nopagebreak[4]
\begin{minted}{julia}
rand(T); rand(T, dims); rand!(A)      # uniform on [0, 1)
randn(T); randn(T, dims); randn!(A)   # standard normal, always finite
rand(Xoshiro(1), T)                   # explicit rng as usual

rand(rng, T;  projection=RTP_SN)   # scalar forms: land under any ProjSpec
randn(rng, T; projection=RTZ_SF) # (stochastic ρ draws R from same rng)
\end{minted}



Classification and stepping:




\nopagebreak[4]
\begin{minted}{julia}
isnan(x); isinf(x); isfinite(x); iszero(x)
signbit(x); issubnormal(x)
Class(x)
NextGreaterThan(x); NextLessThan(x)
nextfloat(x); prevfloat(x)
TotalOrder(x, y)
\end{minted}



\section{Projection specifications}



\label{714366806163910976}{}


A projection specification is \texttt{(\allowbreak{}rounding mode, saturation mode)}:




\nopagebreak[4]
\begin{minted}{julia}
ρ = ProjSpec(TowardPositive(), SatFinite())
roundingmode(ρ)
saturationmode(ρ)
\end{minted}



\needspace{38\baselineskip}
\subsection{Deterministic projections}



\label{18424161794924729217}{}



\begin{table}[h]
\centering
\begin{tabulary}{\linewidth}{L L L L}
\toprule
Rounding & \texttt{SatFinite} & \texttt{SatPropagate} & \texttt{SatNone} \\
\toprule
nearest, ties even & \texttt{RNE\_\allowbreak{}SF} & \texttt{RNE\_\allowbreak{}SP} & \texttt{RNE\_\allowbreak{}SN} \\
nearest, ties away & \texttt{RNA\_\allowbreak{}SF} & \texttt{RNA\_\allowbreak{}SP} & \texttt{RNA\_\allowbreak{}SN} \\
toward +∞ & \texttt{RTP\_\allowbreak{}SF} & \texttt{RTP\_\allowbreak{}SP} & \texttt{RTP\_\allowbreak{}SN} \\
toward −∞ & \texttt{RTN\_\allowbreak{}SF} & \texttt{RTN\_\allowbreak{}SP} & \texttt{RTN\_\allowbreak{}SN} \\
toward zero & \texttt{RTZ\_\allowbreak{}SF} & \texttt{RTZ\_\allowbreak{}SP} & \texttt{RTZ\_\allowbreak{}SN} \\
round to odd & \texttt{RTO\_\allowbreak{}SF} & \texttt{RTO\_\allowbreak{}SP} & \texttt{RTO\_\allowbreak{}SN} \\
\bottomrule
\end{tabulary}

\end{table}



\texttt{RNE\_\allowbreak{}SN} is the package-wide default. Saturation modes mean:




\begin{table}[h]
\centering
\begin{tabulary}{\linewidth}{L L}
\toprule
Mode & Out-of-range behavior \\
\toprule
\texttt{SatFinite} & clamp everything to the finite range \\
\texttt{SatPropagate} & preserve representable infinities; clamp other overflow \\
\texttt{SatNone} & apply the draft{\textquotesingle}s domain-, signedness-, and direction-dependent rows \\
\bottomrule
\end{tabulary}

\end{table}



\needspace{19\baselineskip}
\subsection{Stochastic projections}



\label{13387950688615589261}{}


Constructors are grouped by stochastic variant:




\begin{table}[h]
\centering
\begin{tabulary}{\linewidth}{L L L L}
\toprule
Variant & \texttt{SatFinite} & \texttt{SatPropagate} & \texttt{SatNone} \\
\toprule
A & \texttt{RSA\_\allowbreak{}SF(N)} & \texttt{RSA\_\allowbreak{}SP(N)} & \texttt{RSA\_\allowbreak{}SN(N)} \\
B & \texttt{RSB\_\allowbreak{}SF(N)} & \texttt{RSB\_\allowbreak{}SP(N)} & \texttt{RSB\_\allowbreak{}SN(N)} \\
C & \texttt{RSC\_\allowbreak{}SF(N)} & \texttt{RSC\_\allowbreak{}SP(N)} & \texttt{RSC\_\allowbreak{}SN(N)} \\
\bottomrule
\end{tabulary}

\end{table}



\texttt{N} is the random-bit budget, \texttt{1 ≤ N ≤ 60}. Omitting it uses \texttt{N = 8}.




\begin{minted}{julia}
σ = RSA_SN(8)

Add(T, σ, x, y; rng=Xoshiro(1))  # reproducible stream
Add(T, σ, x, y; R=17)            # exact draw, ideal for tests

isstochastic(σ)                   # true
nrandbits(σ)                      # 8
\end{minted}



For an \texttt{N}-bit mode, explicit \texttt{R} must be in \texttt{0:(2{\textasciicircum}N - 1)}.



\subsection{Session defaults}



\label{3970345369238056532}{}


Read with \texttt{DefaultX()}, set with \texttt{DefaultX!(v)}:




\nopagebreak[4]
\begin{minted}{julia}
DefaultType()            # Binary8p2se     DefaultType!(Binary8p4se)
DefaultReturnType()      # Binary8p2se     DefaultReturnType!(Binary8p3se)
DefaultRoundingMode()    # NearestTiesToEven()
DefaultSaturationMode()  # SatNone()
DefaultProjection()      # RNE_SN
DefaultRNG()             # Xoshiro         DefaultRNG!(Xoshiro(42))
DefaultRbits()           # 8               DefaultRbits!(16)
\end{minted}



Setting a rounding/saturation component rebuilds \texttt{DefaultProjection}; setting \texttt{DefaultProjection!} directly updates both components. Always: \texttt{DefaultProjection(\allowbreak{}) === ProjSpec(\allowbreak{}DefaultRoundingMode(\allowbreak{}), DefaultSaturationMode(\allowbreak{}))}.



The convenience methods (\texttt{a + b}, \texttt{Exp(x)}, \texttt{T(2.1)}) \textbf{do} consult \texttt{DefaultProjection}, so changing it changes their results globally. They read it through a speculation guard: allocation-free while the default holds its initial value, one dispatch per call after you change it. Explicit forms (\texttt{Add(T, ρ, x, y)}) are unaffected by the session default.



Consume a default via the combinators — never by computing on a bare \texttt{DefaultX()} read. No dispatch while the default is unchanged, one barrier dispatch after a change; zero-alloc when \texttt{f}{\textquotesingle}s result type doesn{\textquotesingle}t depend on the default (a default-typed result boxes once at escape):




\nopagebreak[4]
\begin{minted}{julia}
with_default_type((T, x) -> T(x), 1.5)          # Binary8p2se(1.5 ≡ 0x41)
with_default_projection((ρ, x, y) -> Add(T, ρ, x, y), x, y)
with_default_returntype(f, args...)              # f(DefaultReturnType(), args...)
\end{minted}



\needspace{34\baselineskip}
\section{Scalar operation catalog}



\label{9067118353150552449}{}



\begin{minted}{julia}
# explicit result format and projection
Add(T, ρ, x, y)
FMA(T, ρ, x, y, z)
Convert(T, ρ, external_value)

# same-format default-projection convenience
Add(x, y)
FMA(x, y, z)
\end{minted}




\begin{table}[h]
\centering
\begin{tabulary}{\linewidth}{L L}
\toprule
Arity & Operations \\
\toprule
Unary & \texttt{Abs}, \texttt{Negate}, \texttt{Sqrt}, \texttt{RSqrt}, \texttt{Recip}, \texttt{Exp}, \texttt{Exp2}, \texttt{ExpMinusOne}, \texttt{Log}, \texttt{Log2}, \texttt{LogOnePlus}, \texttt{Softplus}, \texttt{Sin}, \texttt{Cos}, \texttt{Tan}, \texttt{ArcSin}, \texttt{ArcCos}, \texttt{ArcTan}, \texttt{Sinh}, \texttt{Cosh}, \texttt{Tanh}, \texttt{ArcSinh}, \texttt{ArcCosh}, \texttt{ArcTanh}, \texttt{SinPi}, \texttt{CosPi}, \texttt{TanPi}, \texttt{ArcSinPi}, \texttt{ArcCosPi}, \texttt{ArcTanPi} \\
Binary & \texttt{CopySign}, \texttt{Add}, \texttt{Subtract}, \texttt{Multiply}, \texttt{Divide}, \texttt{Hypot}, \texttt{ArcTan2}, \texttt{ArcTan2Pi}, \texttt{Maximum}, \texttt{Minimum}, \texttt{MaximumNumber}, \texttt{MinimumNumber}, \texttt{MaximumMagnitude}, \texttt{MinimumMagnitude}, \texttt{MaximumMagnitudeNumber}, \texttt{MinimumMagnitudeNumber}, \texttt{MinimumFinite}, \texttt{MaximumFinite} \\
Ternary & \texttt{FMA}, \texttt{FAA}, \texttt{Clamp} \\
Conversion & \texttt{Convert} \\
\bottomrule
\end{tabulary}

\end{table}



Common Base spellings under \texttt{RNE\_\allowbreak{}SN}:




\nopagebreak[4]
\begin{minted}{julia}
x + y; x - y; x * y; x / y
-x; abs(x); inv(x); sqrt(x)
exp(x); exp2(x); expm1(x)
log(x); log2(x); log1p(x)
sin(x); cos(x); tan(x)
asin(x); acos(x); atan(x); atan(y, x)
sinh(x); cosh(x); tanh(x)
min(x, y); max(x, y); clamp(x, y, z)
fma(x, y, z); muladd(x, y, z)
\end{minted}



\section{Arrays}



\label{13119822296912292823}{}


Every registered operation has elementwise array methods:




\nopagebreak[4]
\begin{minted}{julia}
A = T.([1.0, 1.5, 2.0])
B = T.([0.25, 0.5, 0.75])

C = Add(T, ρ, A, B)
E = Exp(T, ρ, A)
F = FMA(T, ρ, A, B, C)

C = vmap(:Add, T, ρ, A, B)
vmap!(C, Val(:Add), T, ρ, A, B)
\end{minted}



Operands and destination must have matching axes. Deterministic unary/binary operations use cached result tables; affordable ternary signatures may also use tables. Stochastic operations compute each element and consume one draw per projection.




\begin{minted}{julia}
table_bytes()
empty_tables!()
\end{minted}



\section{Blocks and scaled operations}



\label{17754733778172498420-dup1}{}



\begin{minted}{julia}
FS = Binary8p1uf
FE = Binary8p4se

b = Block(FS(4.0), FE(1.5), FE(-0.75), FE(2.0), FE(0.5))

blocksize(b)
scaleformat(b)
elemformat(b)

ConvertFromBlock(T, ρ, b)
BlockReduceAdd(T, ρ, b)
BlockReduceMultiply(T, ρ, b)
BlockDotProduct(T, ρ, b, b)
\end{minted}



Every scalar operation except \texttt{Convert} has generated block and scaled forms:




\nopagebreak[4]
\begin{minted}{julia}
BlockAdd(T, ρ, b1, b2, result_scale)
BlockExp(T, ρ, b, result_scale)

ScaledAdd(T, ρ, scale1, x1, scale2, x2)
ScaledExp(T, ρ, scale, x)
\end{minted}



Quantize against a supplied or automatically selected scale:




\nopagebreak[4]
\begin{minted}{julia}
ConvertToBlock(FS, FE, ρ, values_tuple, scale)
ConvertToBlockMaxAbsFinite(FS, FE, scale_ρ, element_ρ, values_tuple)
\end{minted}



\texttt{BlockVector} stores many equal-shape blocks in a structure-of-arrays layout.



\section{Packed storage}



\label{8567984038512272605}{}



\begin{minted}{julia}
v = Binary5p2se.([0.5, 1.0, 1.5, 2.0])
pv = PackedVector(v)

pv[2]
pv[2] = Binary5p2se(0.75)
collect(pv)

out = vmap(:Exp, Binary5p2se, ρ, pv)
\end{minted}



\texttt{PackedVector} stores each code point in exactly \texttt{bitwidth(F)} bits. Computation unpacks tiles internally; packed arithmetic is deliberately not in-place.



\section{Conformance and approximations}



\label{9146403741965371237}{}



\begin{minted}{julia}
conformance()
conformance_dict()
conformance_report()
draft_revision()

κ, exhaustive = measure_kappa(fn, :Exp, T, (T,), ρ)
register_approx!(:fast_exp, :Exp, T, (T,), ρ, fn; κ)
impl = approx(:fast_exp)
kappa(impl)
kappa_measured(impl)
list_approx()
unregister_approx!(:fast_exp)
\end{minted}



Approximate implementations are never substituted into the default API.



\section{Common pitfalls}



\label{8987423192913392676}{}


See \href{troubleshooting.md}{Troubleshooting} for the full symptom-first list.



\section{Performance checklist}



\label{15543690777485812783}{}


See \href{under\_\allowbreak{}verification.md}{Verification \& Benchmark Doctrine} and the \href{explain\_\allowbreak{}performance.md}{Performance Model} for the full checklist.



For worked applications, continue to \href{examples\_\allowbreak{}index.md}{Examples from AI}. For correctness and benchmark methodology, continue to \href{under\_\allowbreak{}architecture.md}{Insights}.



\chapter{SmallFloats.jl benchmark report}



\label{918403416514326915}{}


Generated: 2026-07-30 12:32 UTC  ·  Julia 1.12.6  ·  alderlake (32 logical CPUs, 1 Julia thread)  ·  Float128 paths: enabled  ·  Chairmarks 1.3.1



Reference format for per-operation tables: \texttt{Binary8p4se} under \texttt{(\allowbreak{}NearestTiesToEven, SatNone)}. Every table names its operand class: the scalar-operation tables appear in four variants — all code points (NaN and ±Inf sampled), NaN excluded, finite-only, and per-operation in-domain — and every other sampled table uses the all-code-points pool, identified in its note. Times are per call; minima first, medians alongside. Methodology per the recorded benchmark doctrine: type-parameterized barriers, untimed setup, specialization preflight.



\needspace{22\baselineskip}
\section{Core primitives}



\label{10766655646123147915}{}


The decode/compare/step/classify layer plus the projection engine. Operands: all code points — NaN and ±Inf sampled.




\begin{table}[h]
\centering
\begin{tabulary}{\linewidth}{L L L L}
\toprule
operation & min & median & allocs \\
\toprule
\texttt{decode} & 1.6 ns & 1.6 ns & 0 \\
\texttt{order\_\allowbreak{}key} & 1.6 ns & 1.8 ns & 0 \\
\texttt{x < y} & 1.8 ns & 1.8 ns & 0 \\
\texttt{TotalOrder} & 2.1 ns & 2.1 ns & 0 \\
\texttt{Class} & 1.6 ns & 2.5 ns & 0 \\
\texttt{NextGreaterThan} & 1.8 ns & 2.1 ns & 0 \\
\texttt{project (RNE·SatNone)} & 2.2 ns & 6.7 ns & 0 \\
\texttt{project (\allowbreak{}StochasticA[8], R drawn)} & 2.1 ns & 7.1 ns & 0 \\
\bottomrule
\end{tabulary}

\end{table}



\section{Scalar operations}



\label{11509067801894334317-dup1}{}


Each arity is measured under four operand classes — safe args, no NaN/Inf, no NaN, and all code points — so operand-class effects (NaN fast rows, ±Inf special rows) are separable.




\clearpage



\subsection{Unary (30)}



\label{2812026322430809105-dup1}{}


\needspace{47\baselineskip}
\subsubsection{Safe args}



\label{15024080724600605730}{}


Finite operands within each operation{\textquotesingle}s safe domain — the fully unmasked per-operation scalar cost. The argument-restricted ops (Sqrt, RSqrt, Log, Log2, LogOnePlus, Recip, Divide, ArcSin, ArcCos, ArcCosh, ArcTanh) draw from explicit per-argument safe-domain predicates; every other op uses finite operand tuples whose defined result is not NaN (oracle-derived). Sorted by median.




\begin{table}[h]
\centering
\begin{tabulary}{\linewidth}{L L L L}
\toprule
operation & min & median & allocs \\
\toprule
\texttt{Abs} & 4.9 ns & 6.7 ns & 0 \\
\texttt{Negate} & 4.8 ns & 6.9 ns & 0 \\
\texttt{Recip} & 7.9 ns & 18.0 ns & 0 \\
\texttt{Sqrt} & 5.7 ns & 18.7 ns & 0 \\
\texttt{RSqrt} & 9.6 ns & 20.7 ns & 0 \\
\texttt{Cosh} & 7.3 ns & 21.4 ns & 0 \\
\texttt{Sin} & 6.2 ns & 21.7 ns & 0 \\
\texttt{Cos} & 7.3 ns & 21.7 ns & 0 \\
\texttt{Exp2} & 7.3 ns & 22.4 ns & 0 \\
\texttt{Exp} & 7.4 ns & 22.5 ns & 0 \\
\texttt{ArcSin} & 5.7 ns & 23.2 ns & 0 \\
\texttt{ExpMinusOne} & 5.3 ns & 23.2 ns & 0 \\
\texttt{Tanh} & 5.8 ns & 23.2 ns & 0 \\
\texttt{ArcCos} & 5.7 ns & 24.3 ns & 0 \\
\texttt{ArcSinPi} & 5.5 ns & 25.5 ns & 0 \\
\texttt{Sinh} & 5.7 ns & 25.6 ns & 0 \\
\texttt{Log} & 5.5 ns & 25.7 ns & 0 \\
\texttt{Log2} & 5.9 ns & 26.6 ns & 0 \\
\texttt{Tan} & 5.3 ns & 26.6 ns & 0 \\
\texttt{LogOnePlus} & 5.4 ns & 26.7 ns & 0 \\
\texttt{ArcCosPi} & 6.1 ns & 26.8 ns & 0 \\
\texttt{CosPi} & 6.9 ns & 27.8 ns & 0 \\
\texttt{SinPi} & 5.6 ns & 28.1 ns & 0 \\
\texttt{ArcTan} & 5.2 ns & 28.4 ns & 0 \\
\texttt{ArcTanPi} & 5.1 ns & 31.0 ns & 0 \\
\texttt{ArcTanh} & 6.2 ns & 31.1 ns & 0 \\
\texttt{TanPi} & 6.2 ns & 31.1 ns & 0 \\
\texttt{ArcCosh} & 5.5 ns & 33.3 ns & 0 \\
\texttt{Softplus} & 13.8 ns & 33.3 ns & 0 \\
\texttt{ArcSinh} & 6.2 ns & 33.4 ns & 0 \\
\bottomrule
\end{tabulary}

\end{table}



\needspace{44\baselineskip}
\subsubsection{No NaN, Inf args}



\label{16036960291186107490}{}


Operands exclude NaN and ±Inf; finite datums only (zeros and subnormals kept). Domain-restricted ops still take NaN fast rows on out-of-domain finite operands. Sorted by median.




\begin{table}[h]
\centering
\begin{tabulary}{\linewidth}{L L L L}
\toprule
operation & min & median & allocs \\
\toprule
\texttt{Sqrt} & 1.9 ns & 2.1 ns & 0 \\
\texttt{ArcCosh} & 2.1 ns & 2.1 ns & 0 \\
\texttt{ArcSinPi} & 1.9 ns & 2.2 ns & 0 \\
\texttt{Abs} & 4.3 ns & 6.7 ns & 0 \\
\texttt{Negate} & 4.8 ns & 6.9 ns & 0 \\
\texttt{Log2} & 2.1 ns & 8.6 ns & 0 \\
\texttt{RSqrt} & 2.1 ns & 9.5 ns & 0 \\
\texttt{Recip} & 1.9 ns & 17.8 ns & 0 \\
\texttt{Cos} & 7.1 ns & 21.2 ns & 0 \\
\texttt{Sin} & 5.4 ns & 21.5 ns & 0 \\
\texttt{Cosh} & 7.0 ns & 21.6 ns & 0 \\
\texttt{Exp2} & 7.3 ns & 22.3 ns & 0 \\
\texttt{Exp} & 7.0 ns & 22.4 ns & 0 \\
\texttt{ExpMinusOne} & 4.5 ns & 22.8 ns & 0 \\
\texttt{ArcSin} & 2.3 ns & 22.8 ns & 0 \\
\texttt{Tanh} & 5.7 ns & 23.1 ns & 0 \\
\texttt{ArcCos} & 2.5 ns & 23.8 ns & 0 \\
\texttt{Log} & 2.3 ns & 25.3 ns & 0 \\
\texttt{Sinh} & 5.1 ns & 25.6 ns & 0 \\
\texttt{LogOnePlus} & 1.9 ns & 26.3 ns & 0 \\
\texttt{ArcCosPi} & 2.1 ns & 26.5 ns & 0 \\
\texttt{Tan} & 6.1 ns & 26.6 ns & 0 \\
\texttt{CosPi} & 6.4 ns & 27.4 ns & 0 \\
\texttt{SinPi} & 5.2 ns & 27.7 ns & 0 \\
\texttt{ArcTan} & 5.5 ns & 28.3 ns & 0 \\
\texttt{ArcTanh} & 2.1 ns & 28.7 ns & 0 \\
\texttt{TanPi} & 5.3 ns & 30.8 ns & 0 \\
\texttt{ArcTanPi} & 4.8 ns & 30.8 ns & 0 \\
\texttt{Softplus} & 14.3 ns & 33.5 ns & 0 \\
\texttt{ArcSinh} & 5.0 ns & 33.6 ns & 0 \\
\bottomrule
\end{tabulary}

\end{table}



\needspace{43\baselineskip}
\subsubsection{No NaN args}



\label{11020842090045916616}{}


Operands exclude the NaN code point; ±Inf and every finite datum are sampled. Sorted by median.




\begin{table}[h]
\centering
\begin{tabulary}{\linewidth}{L L L L}
\toprule
operation & min & median & allocs \\
\toprule
\texttt{ArcCosh} & 2.1 ns & 2.1 ns & 0 \\
\texttt{RSqrt} & 2.1 ns & 2.2 ns & 0 \\
\texttt{Log} & 2.3 ns & 2.5 ns & 0 \\
\texttt{ArcSin} & 2.5 ns & 2.6 ns & 0 \\
\texttt{Sqrt} & 1.9 ns & 5.1 ns & 0 \\
\texttt{Abs} & 4.4 ns & 6.7 ns & 0 \\
\texttt{Negate} & 4.4 ns & 6.8 ns & 0 \\
\texttt{ArcSinPi} & 2.0 ns & 7.3 ns & 0 \\
\texttt{Log2} & 2.1 ns & 8.6 ns & 0 \\
\texttt{Recip} & 1.9 ns & 17.8 ns & 0 \\
\texttt{Cosh} & 5.0 ns & 21.2 ns & 0 \\
\texttt{Sin} & 2.1 ns & 21.4 ns & 0 \\
\texttt{Cos} & 1.9 ns & 21.5 ns & 0 \\
\texttt{Exp2} & 5.5 ns & 22.3 ns & 0 \\
\texttt{Exp} & 5.0 ns & 22.3 ns & 0 \\
\texttt{Tanh} & 5.7 ns & 23.0 ns & 0 \\
\texttt{ExpMinusOne} & 4.6 ns & 23.1 ns & 0 \\
\texttt{ArcCos} & 2.6 ns & 23.8 ns & 0 \\
\texttt{Sinh} & 5.1 ns & 25.6 ns & 0 \\
\texttt{LogOnePlus} & 1.9 ns & 26.2 ns & 0 \\
\texttt{ArcCosPi} & 2.1 ns & 26.6 ns & 0 \\
\texttt{Tan} & 2.0 ns & 27.2 ns & 0 \\
\texttt{CosPi} & 2.5 ns & 27.5 ns & 0 \\
\texttt{SinPi} & 2.1 ns & 27.7 ns & 0 \\
\texttt{ArcTan} & 5.7 ns & 28.3 ns & 0 \\
\texttt{ArcTanh} & 2.1 ns & 28.6 ns & 0 \\
\texttt{ArcTanPi} & 4.7 ns & 30.8 ns & 0 \\
\texttt{TanPi} & 2.1 ns & 30.8 ns & 0 \\
\texttt{Softplus} & 4.8 ns & 33.5 ns & 0 \\
\texttt{ArcSinh} & 4.5 ns & 33.6 ns & 0 \\
\bottomrule
\end{tabulary}

\end{table}



\needspace{46\baselineskip}
\subsubsection{All code points}



\label{8134844262983830456}{}


Operands drawn uniformly over ALL code points — NaN and ±Inf are sampled; medians of domain-restricted ops are diluted by instant NaN rows. Sorted by median. Transcendental rows mix special-row fast returns with enclosure-path evaluations, so these are \emph{scalar-path} costs; bulk unary work routes through 256-byte tables (see Array kernels).




\begin{table}[h]
\centering
\begin{tabulary}{\linewidth}{L L L L}
\toprule
operation & min & median & allocs \\
\toprule
\texttt{Sqrt} & 1.9 ns & 2.0 ns & 0 \\
\texttt{ArcCosh} & 2.1 ns & 2.1 ns & 0 \\
\texttt{ArcSinPi} & 1.9 ns & 2.1 ns & 0 \\
\texttt{RSqrt} & 2.1 ns & 2.1 ns & 0 \\
\texttt{Log} & 2.3 ns & 3.8 ns & 0 \\
\texttt{Log2} & 2.1 ns & 5.5 ns & 0 \\
\texttt{ArcSin} & 2.1 ns & 5.8 ns & 0 \\
\texttt{ArcTanh} & 2.1 ns & 6.6 ns & 0 \\
\texttt{Abs} & 2.0 ns & 6.7 ns & 0 \\
\texttt{Negate} & 2.1 ns & 6.8 ns & 0 \\
\texttt{Recip} & 1.8 ns & 17.8 ns & 0 \\
\texttt{Cosh} & 2.3 ns & 21.0 ns & 0 \\
\texttt{Cos} & 1.9 ns & 21.3 ns & 0 \\
\texttt{Sin} & 2.1 ns & 21.5 ns & 0 \\
\texttt{Exp2} & 2.1 ns & 22.3 ns & 0 \\
\texttt{Exp} & 2.0 ns & 22.3 ns & 0 \\
\texttt{ExpMinusOne} & 2.0 ns & 22.7 ns & 0 \\
\texttt{Tanh} & 2.3 ns & 23.1 ns & 0 \\
\texttt{ArcCos} & 2.3 ns & 23.7 ns & 0 \\
\texttt{Sinh} & 2.3 ns & 25.5 ns & 0 \\
\texttt{LogOnePlus} & 1.8 ns & 26.4 ns & 0 \\
\texttt{Tan} & 2.0 ns & 26.5 ns & 0 \\
\texttt{ArcCosPi} & 1.9 ns & 26.6 ns & 0 \\
\texttt{CosPi} & 2.3 ns & 27.5 ns & 0 \\
\texttt{SinPi} & 2.0 ns & 27.7 ns & 0 \\
\texttt{ArcTan} & 2.2 ns & 28.2 ns & 0 \\
\texttt{ArcTanPi} & 2.1 ns & 30.8 ns & 0 \\
\texttt{TanPi} & 1.9 ns & 30.8 ns & 0 \\
\texttt{Softplus} & 2.0 ns & 33.5 ns & 0 \\
\texttt{ArcSinh} & 2.1 ns & 33.6 ns & 0 \\
\bottomrule
\end{tabulary}

\end{table}




\clearpage



\subsection{Binary (18)}



\label{9493958918153706001-dup1}{}


\needspace{35\baselineskip}
\subsubsection{Safe args}



\label{1931488528371625443}{}


Finite operands within each operation{\textquotesingle}s safe domain — the fully unmasked per-operation scalar cost. The argument-restricted ops (Sqrt, RSqrt, Log, Log2, LogOnePlus, Recip, Divide, ArcSin, ArcCos, ArcCosh, ArcTanh) draw from explicit per-argument safe-domain predicates; every other op uses finite operand tuples whose defined result is not NaN (oracle-derived). Sorted by median.




\begin{table}[h]
\centering
\begin{tabulary}{\linewidth}{L L L L}
\toprule
operation & min & median & allocs \\
\toprule
\texttt{MinimumMagnitude} & 4.9 ns & 7.1 ns & 0 \\
\texttt{MaximumMagnitude} & 6.9 ns & 7.1 ns & 0 \\
\texttt{CopySign} & 4.8 ns & 7.2 ns & 0 \\
\texttt{Maximum} & 4.5 ns & 7.2 ns & 0 \\
\texttt{Minimum} & 4.6 ns & 7.2 ns & 0 \\
\texttt{MaximumFinite} & 4.8 ns & 7.2 ns & 0 \\
\texttt{MinimumFinite} & 4.9 ns & 7.3 ns & 0 \\
\texttt{MinimumNumber} & 4.8 ns & 7.3 ns & 0 \\
\texttt{MaximumNumber} & 5.0 ns & 7.3 ns & 0 \\
\texttt{MaximumMagnitudeNumber} & 7.1 ns & 7.3 ns & 0 \\
\texttt{MinimumMagnitudeNumber} & 4.8 ns & 7.3 ns & 0 \\
\texttt{Multiply} & 4.9 ns & 7.4 ns & 0 \\
\texttt{Add} & 7.1 ns & 9.0 ns & 0 \\
\texttt{Subtract} & 7.4 ns & 9.2 ns & 0 \\
\texttt{Divide} & 6.4 ns & 18.3 ns & 0 \\
\texttt{Hypot} & 7.8 ns & 28.0 ns & 0 \\
\texttt{ArcTan2} & 8.0 ns & 32.7 ns & 0 \\
\texttt{ArcTan2Pi} & 7.2 ns & 35.3 ns & 0 \\
\bottomrule
\end{tabulary}

\end{table}



\needspace{32\baselineskip}
\subsubsection{No NaN, Inf args}



\label{9842352165407780921}{}


Operands exclude NaN and ±Inf; finite datums only (zeros and subnormals kept). Domain-restricted ops still take NaN fast rows on out-of-domain finite operands. Sorted by median.




\begin{table}[h]
\centering
\begin{tabulary}{\linewidth}{L L L L}
\toprule
operation & min & median & allocs \\
\toprule
\texttt{MinimumMagnitude} & 4.8 ns & 7.1 ns & 0 \\
\texttt{MaximumMagnitude} & 6.9 ns & 7.1 ns & 0 \\
\texttt{Minimum} & 4.7 ns & 7.2 ns & 0 \\
\texttt{Maximum} & 4.6 ns & 7.2 ns & 0 \\
\texttt{CopySign} & 4.7 ns & 7.2 ns & 0 \\
\texttt{MinimumFinite} & 5.0 ns & 7.2 ns & 0 \\
\texttt{MaximumFinite} & 5.1 ns & 7.2 ns & 0 \\
\texttt{MinimumNumber} & 6.1 ns & 7.3 ns & 0 \\
\texttt{MaximumNumber} & 4.8 ns & 7.3 ns & 0 \\
\texttt{MinimumMagnitudeNumber} & 4.8 ns & 7.3 ns & 0 \\
\texttt{MaximumMagnitudeNumber} & 7.2 ns & 7.3 ns & 0 \\
\texttt{Multiply} & 4.9 ns & 7.4 ns & 0 \\
\texttt{Add} & 7.1 ns & 9.0 ns & 0 \\
\texttt{Subtract} & 7.9 ns & 9.5 ns & 0 \\
\texttt{Divide} & 2.6 ns & 18.3 ns & 0 \\
\texttt{Hypot} & 8.2 ns & 28.0 ns & 0 \\
\texttt{ArcTan2} & 9.0 ns & 32.7 ns & 0 \\
\texttt{ArcTan2Pi} & 7.1 ns & 35.4 ns & 0 \\
\bottomrule
\end{tabulary}

\end{table}



\needspace{31\baselineskip}
\subsubsection{No NaN args}



\label{5394694367249928840}{}


Operands exclude the NaN code point; ±Inf and every finite datum are sampled. Sorted by median.




\begin{table}[h]
\centering
\begin{tabulary}{\linewidth}{L L L L}
\toprule
operation & min & median & allocs \\
\toprule
\texttt{MinimumMagnitude} & 4.6 ns & 7.1 ns & 0 \\
\texttt{MaximumMagnitude} & 5.0 ns & 7.1 ns & 0 \\
\texttt{CopySign} & 5.0 ns & 7.2 ns & 0 \\
\texttt{Minimum} & 4.6 ns & 7.2 ns & 0 \\
\texttt{Maximum} & 4.6 ns & 7.2 ns & 0 \\
\texttt{MaximumFinite} & 5.0 ns & 7.2 ns & 0 \\
\texttt{MinimumFinite} & 5.0 ns & 7.2 ns & 0 \\
\texttt{MinimumNumber} & 5.1 ns & 7.3 ns & 0 \\
\texttt{MaximumNumber} & 4.6 ns & 7.3 ns & 0 \\
\texttt{MinimumMagnitudeNumber} & 4.6 ns & 7.3 ns & 0 \\
\texttt{MaximumMagnitudeNumber} & 5.1 ns & 7.3 ns & 0 \\
\texttt{Multiply} & 5.0 ns & 7.4 ns & 0 \\
\texttt{Add} & 6.3 ns & 9.0 ns & 0 \\
\texttt{Subtract} & 7.2 ns & 9.5 ns & 0 \\
\texttt{Divide} & 2.6 ns & 18.3 ns & 0 \\
\texttt{Hypot} & 6.1 ns & 28.0 ns & 0 \\
\texttt{ArcTan2} & 7.1 ns & 32.7 ns & 0 \\
\texttt{ArcTan2Pi} & 6.7 ns & 35.4 ns & 0 \\
\bottomrule
\end{tabulary}

\end{table}



\needspace{32\baselineskip}
\subsubsection{All code points}



\label{11499910499181378231}{}


Operands drawn uniformly over ALL code points — NaN and ±Inf are sampled; medians of domain-restricted ops are diluted by instant NaN rows. Sorted by median.




\begin{table}[h]
\centering
\begin{tabulary}{\linewidth}{L L L L}
\toprule
operation & min & median & allocs \\
\toprule
\texttt{MinimumMagnitude} & 2.1 ns & 7.1 ns & 0 \\
\texttt{MaximumMagnitude} & 1.9 ns & 7.1 ns & 0 \\
\texttt{CopySign} & 2.1 ns & 7.2 ns & 0 \\
\texttt{Minimum} & 2.1 ns & 7.2 ns & 0 \\
\texttt{Maximum} & 1.8 ns & 7.2 ns & 0 \\
\texttt{MaximumFinite} & 5.6 ns & 7.2 ns & 0 \\
\texttt{MinimumFinite} & 5.0 ns & 7.2 ns & 0 \\
\texttt{MinimumNumber} & 5.1 ns & 7.3 ns & 0 \\
\texttt{MaximumNumber} & 4.6 ns & 7.3 ns & 0 \\
\texttt{MinimumMagnitudeNumber} & 4.6 ns & 7.3 ns & 0 \\
\texttt{MaximumMagnitudeNumber} & 5.0 ns & 7.4 ns & 0 \\
\texttt{Multiply} & 2.1 ns & 7.4 ns & 0 \\
\texttt{Add} & 3.7 ns & 9.0 ns & 0 \\
\texttt{Subtract} & 4.6 ns & 9.5 ns & 0 \\
\texttt{Divide} & 2.5 ns & 18.3 ns & 0 \\
\texttt{Hypot} & 3.0 ns & 28.0 ns & 0 \\
\texttt{ArcTan2} & 4.6 ns & 32.7 ns & 0 \\
\texttt{ArcTan2Pi} & 3.8 ns & 35.3 ns & 0 \\
\bottomrule
\end{tabulary}

\end{table}



\subsection{Ternary (3)}



\label{1836852574804533888-dup1}{}


\needspace{20\baselineskip}
\subsubsection{Safe args}



\label{9973129846987395370}{}


Finite operands within each operation{\textquotesingle}s safe domain — the fully unmasked per-operation scalar cost. The argument-restricted ops (Sqrt, RSqrt, Log, Log2, LogOnePlus, Recip, Divide, ArcSin, ArcCos, ArcCosh, ArcTanh) draw from explicit per-argument safe-domain predicates; every other op uses finite operand tuples whose defined result is not NaN (oracle-derived). Sorted by median.




\begin{table}[h]
\centering
\begin{tabulary}{\linewidth}{L L L L}
\toprule
operation & min & median & allocs \\
\toprule
\texttt{Clamp} & 5.1 ns & 8.0 ns & 0 \\
\texttt{FMA} & 8.1 ns & 9.6 ns & 0 \\
\texttt{FAA} & 8.3 ns & 9.9 ns & 0 \\
\bottomrule
\end{tabulary}

\end{table}



\needspace{17\baselineskip}
\subsubsection{No NaN, Inf args}



\label{15876135725619624151}{}


Operands exclude NaN and ±Inf; finite datums only (zeros and subnormals kept). Domain-restricted ops still take NaN fast rows on out-of-domain finite operands. Sorted by median.




\begin{table}[h]
\centering
\begin{tabulary}{\linewidth}{L L L L}
\toprule
operation & min & median & allocs \\
\toprule
\texttt{Clamp} & 5.5 ns & 8.0 ns & 0 \\
\texttt{FMA} & 8.3 ns & 9.7 ns & 0 \\
\texttt{FAA} & 8.1 ns & 9.7 ns & 0 \\
\bottomrule
\end{tabulary}

\end{table}



\needspace{16\baselineskip}
\subsubsection{No NaN args}



\label{14279209625200149733}{}


Operands exclude the NaN code point; ±Inf and every finite datum are sampled. Sorted by median.




\begin{table}[h]
\centering
\begin{tabulary}{\linewidth}{L L L L}
\toprule
operation & min & median & allocs \\
\toprule
\texttt{Clamp} & 5.2 ns & 8.0 ns & 0 \\
\texttt{FMA} & 6.9 ns & 9.7 ns & 0 \\
\texttt{FAA} & 7.1 ns & 9.7 ns & 0 \\
\bottomrule
\end{tabulary}

\end{table}



\needspace{17\baselineskip}
\subsubsection{All code points}



\label{689699942421014554}{}


Operands drawn uniformly over ALL code points — NaN and ±Inf are sampled; medians of domain-restricted ops are diluted by instant NaN rows. Sorted by median.




\begin{table}[h]
\centering
\begin{tabulary}{\linewidth}{L L L L}
\toprule
operation & min & median & allocs \\
\toprule
\texttt{Clamp} & 2.1 ns & 8.0 ns & 0 \\
\texttt{FMA} & 3.7 ns & 9.7 ns & 0 \\
\texttt{FAA} & 4.1 ns & 9.7 ns & 0 \\
\bottomrule
\end{tabulary}

\end{table}



\section{Sensitivity studies}



\label{14065249436593140616}{}


How the scalar costs move with the format and with the projection specification.



\needspace{30\baselineskip}
\subsection{Format sensitivity}



\label{18234318524667261557}{}


Same three binary ops across formats; \texttt{Binary8p1uf} exercises the wide-exponent-spread escalations, small-K formats the tiny-table regime. Operands: all code points — NaN and ±Inf sampled.




\begin{table}[h]
\centering
\begin{tabulary}{\linewidth}{L L L L}
\toprule
operation & min & median & allocs \\
\toprule
\texttt{Add⟨Binary8p4se⟩} & 3.7 ns & 9.0 ns & 0 \\
\texttt{Divide⟨Binary8p4se⟩} & 2.5 ns & 18.3 ns & 0 \\
\texttt{Multiply⟨Binary8p4se⟩} & 2.1 ns & 7.4 ns & 0 \\
\texttt{Add⟨Binary8p3sf⟩} & 3.9 ns & 8.7 ns & 0 \\
\texttt{Divide⟨Binary8p3sf⟩} & 2.6 ns & 16.4 ns & 0 \\
\texttt{Multiply⟨Binary8p3sf⟩} & 2.1 ns & 7.2 ns & 0 \\
\texttt{Add⟨Binary8p1uf⟩} & 4.7 ns & 74.0 ns & 0 \\
\texttt{Divide⟨Binary8p1uf⟩} & 2.6 ns & 8.4 ns & 0 \\
\texttt{Multiply⟨Binary8p1uf⟩} & 2.1 ns & 6.9 ns & 0 \\
\texttt{Add⟨Binary5p2se⟩} & 4.0 ns & 9.1 ns & 0 \\
\texttt{Divide⟨Binary5p2se⟩} & 2.5 ns & 9.0 ns & 0 \\
\texttt{Multiply⟨Binary5p2se⟩} & 1.9 ns & 7.4 ns & 0 \\
\texttt{Add⟨Binary3p1se⟩} & 4.0 ns & 7.6 ns & 0 \\
\texttt{Divide⟨Binary3p1se⟩} & 2.6 ns & 6.3 ns & 0 \\
\texttt{Multiply⟨Binary3p1se⟩} & 2.1 ns & 5.9 ns & 0 \\
\bottomrule
\end{tabulary}

\end{table}



\needspace{25\baselineskip}
\subsection{Projection by rounding/saturation mode}



\label{15662701981835679547}{}


\texttt{project(\allowbreak{}Binary8p4se, ρ, x)} over the mode vocabulary (stochastic budgets N = 8). Operands: all code points — NaN and ±Inf sampled.




\begin{table}[h]
\centering
\begin{tabulary}{\linewidth}{L L L L}
\toprule
operation & min & median & allocs \\
\toprule
\texttt{NearestTiesToEven} & 2.0 ns & 6.7 ns & 0 \\
\texttt{NearestTiesToAway} & 2.8 ns & 7.7 ns & 0 \\
\texttt{TowardPositive} & 3.0 ns & 6.0 ns & 0 \\
\texttt{TowardNegative} & 2.6 ns & 5.9 ns & 0 \\
\texttt{TowardZero} & 2.5 ns & 5.3 ns & 0 \\
\texttt{ToOdd} & 2.0 ns & 6.6 ns & 0 \\
\texttt{StochasticA[8]} & 2.0 ns & 7.0 ns & 0 \\
\texttt{StochasticB[8]} & 2.5 ns & 7.1 ns & 0 \\
\texttt{StochasticC[8]} & 2.3 ns & 7.1 ns & 0 \\
\texttt{RNE · SatFinite} & 2.1 ns & 7.6 ns & 0 \\
\texttt{RNE · SatPropagate} & 2.1 ns & 7.2 ns & 0 \\
\bottomrule
\end{tabulary}

\end{table}



\section{Kernels and storage}



\label{9956393270762268485}{}


Array-shaped work: the table-gather and compute kernels, the ternary bitwidth policy, sorting, table builds, blocks, and packed/converted storage.



\needspace{26\baselineskip}
\subsection{Array kernels (vmap)}



\label{14819160431636425252}{}


Warm caches: table specializations prebuilt, so table rows measure the gather; scalar-loop rows measure the full compute pipeline per element. The ternary row here is \texttt{Binary8p4se} (K=8, always the compute path); see the next section for how the ternary bitwidth policy behaves across K. Operands: all code points — NaN and ±Inf sampled.




\begin{table}[h]
\centering
\begin{tabulary}{\linewidth}{L L L L L}
\toprule
operation & min & median & allocs & per element \\
\toprule
\texttt{vmap unary (\allowbreak{}table gather), n=65536} & 8.56 μs & 8.71 μs & 0 & 0.13 ns/elem — 7.52 Gelem/s \\
\texttt{vmap binary (\allowbreak{}table gather), n=65536} & 16.45 μs & 16.81 μs & 0 & 0.26 ns/elem — 3.9 Gelem/s \\
\texttt{vmap ternary (\allowbreak{}scalar loop), n=65536} & 1.01 ms & 1.02 ms & 0 & 15.54 ns/elem — 0.06 Gelem/s \\
\texttt{vmap binary stochastic (\allowbreak{}scalar loop), n=65536} & 646.57 μs & 655.78 μs & 0 & 10.01 ns/elem — 0.1 Gelem/s \\
\texttt{vmap unary through PackedVector, n=65536} & 76.97 μs & 90.29 μs & 773 & 1.38 ns/elem — 0.73 Gelem/s \\
\bottomrule
\end{tabulary}

\end{table}



\needspace{34\baselineskip}
\subsection{Ternary bitwidth tiers (FMA/FAA)}



\label{14443973703327008650}{}


\texttt{FMA}/\allowbreak\texttt{FAA}/\texttt{Clamp} are total functions on \texttt{2{\textasciicircum}(K1+K2+K3)} code points, but that count spans 512 B (K=3) to 16 MiB (K=8), so the array kernel tables small operand formats eagerly, tables mid-size ones adaptively (after enough elements amortize the build; not shown here — see the adaptive-cache gate in \texttt{test/ternary\_\allowbreak{}opt.jl}), and always runs the scalar compute kernel at K=8, threaded above a size cutoff when \texttt{Threads.nthreads() > 1}. Each tier{\textquotesingle}s optimized row is paired with a scalar-loop baseline (policy Refs forced off around the measurement, restored after) so the win is visible per tier; this process has 1 Julia thread. No threaded/sequential comparison below — rerun with \texttt{julia -t N} (N > 1) to see it. Same reference format, ρ, and operand pool discipline as Array kernels above.




\begin{table}[h]
\centering
\begin{tabulary}{\linewidth}{L L L L L}
\toprule
operation & min & median & allocs & per element \\
\toprule
\texttt{FMA K=4 (\allowbreak{}eager table), n=65536} & 19.58 μs & 20.46 μs & 0 & 0.31 ns/elem — 3.2 Gelem/s \\
\texttt{FMA K=4 (\allowbreak{}eager table), scalar-loop baseline, n=65536} & 967.88 μs & 974.03 μs & 0 & 14.86 ns/elem — 0.07 Gelem/s \\
\texttt{FMA K=6 (\allowbreak{}eager table), n=65536} & 24.03 μs & 24.34 μs & 0 & 0.37 ns/elem — 2.69 Gelem/s \\
\texttt{FMA K=6 (\allowbreak{}eager table), scalar-loop baseline, n=65536} & 1.0 ms & 1.01 ms & 0 & 15.46 ns/elem — 0.06 Gelem/s \\
\texttt{FMA K=8 (\allowbreak{}compute), n=65536} & 990.69 μs & 1.0 ms & 0 & 15.31 ns/elem — 0.07 Gelem/s \\
\texttt{FMA K=8 (\allowbreak{}compute), scalar-loop baseline, n=65536} & 990.11 μs & 1.0 ms & 0 & 15.27 ns/elem — 0.07 Gelem/s \\
\bottomrule
\end{tabulary}

\end{table}



\needspace{17\baselineskip}
\subsection{Sorting (64 K values)}



\label{10624969220324633770}{}


Counting sort is installed as the default algorithm for \texttt{Binary} vectors. Operands: all code points — NaN and ±Inf sampled.




\begin{table}[h]
\centering
\begin{tabulary}{\linewidth}{L L L L}
\toprule
operation & min & median & allocs \\
\toprule
\texttt{sort! (\allowbreak{}counting sort via defalg), n=65536} & 155.36 μs & 157.4 μs & 2 \\
\texttt{sort! (\allowbreak{}stock comparison sort), n=65536} & 1.36 ms & 1.37 ms & 3 \\
\texttt{sort! rev=true (\allowbreak{}counting sort), n=65536} & 155.48 μs & 157.95 μs & 2 \\
\bottomrule
\end{tabulary}

\end{table}



\needspace{25\baselineskip}
\subsection{Table builds (oracle + projection, Float128-first)}



\label{17249747280115456051}{}


Cold cache per sample (\texttt{empty\_\allowbreak{}tables!} in untimed setup); JIT pre-warmed. The warm-hit column is the steady-state cost of \texttt{get\_\allowbreak{}table} when the specialization is already cached (min / median). Table entries enumerate every code point by construction (NaN and ±Inf included).




\begin{table}[h]
\centering
\begin{tabulary}{\linewidth}{L L L L L}
\toprule
operation & min & median & allocs & warm hit \\
\toprule
\texttt{Exp⟨8p4se⟩ (256 entries)} & 30.31 μs & 31.16 μs & 257 & 65.8 ns / 67.4 ns \\
\texttt{Tanh⟨8p4se⟩ (\allowbreak{}256 entries)} & 29.49 μs & 30.93 μs & 257 & 65.4 ns / 67.5 ns \\
\texttt{Add⟨8p4se×8p4se⟩ (\allowbreak{}64 K entries)} & 13.77 ms & 13.93 ms & 458754 & 65.0 ns / 66.6 ns \\
\texttt{Divide⟨8p4se×8p4se⟩ (\allowbreak{}64 K)} & 14.02 ms & 14.14 ms & 458754 & 65.1 ns / 66.6 ns \\
\texttt{Add⟨8p1uf×8p1uf⟩ (\allowbreak{}64 K, wide-spread)} & 27.08 ms & 29.5 ms & 995304 & 65.8 ns / 67.8 ns \\
\bottomrule
\end{tabulary}

\end{table}



\needspace{20\baselineskip}
\subsection{Block and scaled operations}



\label{17033350743052450193}{}


Elements \texttt{Binary8p4se}, scales \texttt{Binary8p1uf}, B = 32. Operands: all code points — NaN and ±Inf sampled.




\begin{table}[h]
\centering
\begin{tabulary}{\linewidth}{L L L L L}
\toprule
operation & min & median & allocs & per lane \\
\toprule
\texttt{BlockAdd (B=32)} & 963.7 ns & 1.17 μs & 35 & 36.65 ns/lane \\
\texttt{BlockDotProduct → 8p4se (\allowbreak{}B=32)} & 211.3 ns & 265.3 ns & 3 & 8.29 ns/lane \\
\texttt{BlockReduceAdd → 8p4se (\allowbreak{}B=32)} & 43.5 ns & 489.4 ns & 0 & 15.29 ns/lane \\
\texttt{ConvertToBlockMaxAbsFinite (\allowbreak{}B=32)} & 3.79 μs & 3.95 μs & 111 & 123.3 ns/lane \\
\bottomrule
\end{tabulary}

\end{table}



\needspace{22\baselineskip}
\subsection{Conversions and packed storage}



\label{4461369721169500046}{}


Operands: all code points — NaN and ±Inf sampled.




\begin{table}[h]
\centering
\begin{tabulary}{\linewidth}{L L L L L}
\toprule
operation & min & median & allocs & per element \\
\toprule
\texttt{T(\allowbreak{}::UInt8) code-point constructor} & 1.8 ns & 1.8 ns & 0 &  \\
\texttt{rawvalue (\allowbreak{}unchecked kernel route)} & 1.8 ns & 1.8 ns & 0 &  \\
\texttt{T(\allowbreak{}::Float64) numeric constructor (\allowbreak{}projects)} & 3.2 ns & 7.7 ns & 0 &  \\
\texttt{Convert 8p4se → 8p3se (\allowbreak{}scalar)} & 2.1 ns & 6.7 ns & 0 &  \\
\texttt{Float64 → 8p4se (\allowbreak{}project)} & 2.0 ns & 6.7 ns & 0 &  \\
\texttt{PackedVector pack, n=65536} & 33.47 μs & 35.43 μs & 3 & 0.54 ns/elem \\
\texttt{PackedVector unpack (\allowbreak{}collect), n=65536} & 28.86 μs & 29.05 μs & 3 & 0.44 ns/elem \\
\bottomrule
\end{tabulary}

\end{table}



{\rule{\textwidth}{1pt}}


\emph{All numbers from this machine/run; absolute values vary by host. Regenerate with \texttt{julia --project=benchmarking benchmarking/\allowbreak{}benchmarking.jl benchmarking/\allowbreak{}benchmark\_\allowbreak{}report.md}.}



\part{Standard \& Design Papers}


\chapter{The Standard \& Design Papers}



\label{887308718192388530}{}


An index to the curated \texttt{docs/other/} collection: the standard{\textquotesingle}s concept map, the design rationales behind specific implementation choices, and the working papers that planned the K > 8 extension. Each entry carries a status label — read it before treating a document as current.



\section{The standard}



\label{15833458192391556191}{}


\textbf{P3109 Draft Notes} — \href{https://github.com/JeffreySarnoff/SmallFloats.jl/blob/main/docs/other/IEEE_D1_concepts.md}{\texttt{IEEE\_\allowbreak{}D1\_\allowbreak{}concepts.md}}. A concept map of the whole IEEE P3109/D1 draft (July 2026): format-defining parameters, value space, operation machinery, and projection, with every node linked back to the draft{\textquotesingle}s line numbers. Status: normative-description.



The full draft transliteration, \texttt{IEEE\_\allowbreak{}D1.md}, is repo-only and is not part of this curated set — it is linked out from the concept map and from the working papers below (see Decision D1). Find it at \href{https://github.com/JeffreySarnoff/SmallFloats.jl/blob/main/docs/other/IEEE_D1.md}{\texttt{IEEE\_\allowbreak{}D1.md}} in the repository.



\section{Design Rationales}



\label{15541647678802253679}{}


\textbf{dyadic\_\allowbreak{}rational.md} — \href{https://github.com/JeffreySarnoff/SmallFloats.jl/blob/main/docs/other/dyadic_rational.md}{\texttt{dyadic\_\allowbreak{}rational.md}}. The \texttt{Dyadic} ↔ \texttt{Rational} bridge: how the rung-3 exact carrier converts to and from Julia{\textquotesingle}s \texttt{Rational} exactly, and the three laws that conversion must obey (round-trip on every finite dyadic, ±∞ agreement with Base{\textquotesingle}s \texttt{1//0}, and rejection only at the true NaN slot). Status: design-rationale.



\textbf{promoterules.md} — \href{https://github.com/JeffreySarnoff/SmallFloats.jl/blob/main/docs/other/promoterules.md}{\texttt{promoterules.md}}. Works out why \texttt{promote\_\allowbreak{}rule} between two P3109 formats cannot mean {\textquotedbl}widen to a bigger format{\textquotedbl} the way \texttt{Int}/\allowbreak\texttt{Float64} promotion does, and what the promotion carrier is instead. Status: design-rationale, partially superseded by the promotion-carrier section in the Cheat Sheet and Julia Compatibility guide — those pages reflect the shipped \texttt{promotecarrier(T)} behavior, and this document should be read as the reasoning behind it rather than as the current API description.



\textbf{Float32inside.md} — \href{https://github.com/JeffreySarnoff/SmallFloats.jl/blob/main/docs/other/Float32inside.md}{\texttt{Float32inside.md}}. A design study asking whether \texttt{Float32} can replace \texttt{Float64} as the internal carrier; verdict is no as a universal replacement, yes as a surface type and as a gated compute carrier for specific kernels. Status: exploratory.



\textbf{Float32more.md} — \href{https://github.com/JeffreySarnoff/SmallFloats.jl/blob/main/docs/other/Float32more.md}{\texttt{Float32more.md}}. Companion to Float32inside.md: a detailed, empirically validated implementation plan for the Float32/BFloat16 surface, the exactness-gated compute path, and κ-registered Float32 kernels, with measured pass/fail counts per milestone. Status: exploratory.



\section{Extension Working Papers}



\label{17585610141340629196}{}


Four documents plan and re-plan the extension of the format grid from \texttt{K ≤ 8} (120 at K ≤ 8, retained as the default exports) to \texttt{3 ≤ K ≤ 16} (504 formats). Read in this order; each says explicitly where it overrides its predecessor.



\textbf{extendingK.md} — \href{https://github.com/JeffreySarnoff/SmallFloats.jl/blob/main/docs/other/extendingK.md}{\texttt{extendingK.md}}. The first pass: what the draft permits, the two axes (storage unit and evaluation carrier/rung) that govern the extension, and format counts by starting rung. Status: exploratory/historical.



\textbf{doingtheextensions.md} — \href{https://github.com/JeffreySarnoff/SmallFloats.jl/blob/main/docs/other/doingtheextensions.md}{\texttt{doingtheextensions.md}}. Turns the plan into an architecture — the representation lattice, the carrier dispatch lattice, and generated per-format trait methods — staged so every step compiles and passes independently. Status: exploratory/historical.



\textbf{implementextensions.md} — \href{https://github.com/JeffreySarnoff/SmallFloats.jl/blob/main/docs/other/implementextensions.md}{\texttt{implementextensions.md}}. The execution plan: a line-by-line audit of the source tree that found ten findings (errors or gaps) against the inherited plan, followed by a ten-stage, single-commit-per-stage implementation schedule. Status: exploratory/historical.



\textbf{morejulian.md} — \href{https://github.com/JeffreySarnoff/SmallFloats.jl/blob/main/docs/other/morejulian.md}{\texttt{morejulian.md}}. A separate audit of where the package{\textquotesingle}s Julia interface departs from idiomatic Julia — ranked by user impact rather than by extension stage — covering broken \texttt{AbstractFloat} contract methods and deliberate departures alike. Status: exploratory/historical.



\section{see also}



\label{9115667341246479961-dup34}{}


\href{cheatsheet.md}{Cheat Sheet}, \href{glossary.md}{Glossary}, \href{under\_\allowbreak{}architecture.md}{Architecture Overview}.



\end{document}
